\documentclass[11pt]{article}
\usepackage[margin=1.1in]{geometry}
\usepackage{amsmath,amssymb,amsthm,mathtools,booktabs,array,tabularx,longtable}
\usepackage[dvipsnames]{xcolor}
\usepackage[colorlinks=true,linkcolor=MidnightBlue,citecolor=MidnightBlue,urlcolor=MidnightBlue]{hyperref}
\usepackage{fancyhdr}
\pagestyle{fancy}\fancyhf{}\renewcommand{\headrulewidth}{0pt}
\fancyfoot[R]{\thepage}
\fancyfoot[L]{\scriptsize\color{gray}Generated by the Transport Lab, not the version of record. Official version: \href{https://doi.org/10.5281/zenodo.22721790}{doi:10.5281/zenodo.22721790}}
\fancypagestyle{plain}{\fancyhf{}\fancyfoot[R]{\thepage}\fancyfoot[L]{\scriptsize\color{gray}Generated by the Transport Lab, not the version of record. Official version: \href{https://doi.org/10.5281/zenodo.22721790}{doi:10.5281/zenodo.22721790}}\renewcommand{\headrulewidth}{0pt}}
\newtheorem{theorem}{Theorem}[section]
\newtheorem{proposition}[theorem]{Proposition}
\newtheorem{lemma}[theorem]{Lemma}
\newtheorem{corollary}[theorem]{Corollary}
\newtheorem{conjecture}[theorem]{Conjecture}
\newtheorem{refutation}[theorem]{Refutation}
\newtheorem{prediction}[theorem]{Prediction}
\newtheorem{principle}[theorem]{Principle}
\newtheorem{observation}[theorem]{Observation}
\newtheorem{axiom}[theorem]{Axiom}
\theoremstyle{definition}
\newtheorem{definition}[theorem]{Definition}
\newtheorem{fact}[theorem]{Fact}
\newtheorem{assembled}[theorem]{Assembly}
\newtheorem{postulate}[theorem]{Postulate}
\newtheorem{premise}[theorem]{Premise}
\newtheorem{hypothesis}[theorem]{Hypothesis}
\newtheorem{remark}[theorem]{Remark}
\newtheorem{example}[theorem]{Example}
\theoremstyle{plain}
\setlength{\marginparwidth}{0.8in}
\newcommand{\changedmark}[2]{\marginpar{\raggedright\scriptsize\color{#1}#2}}
\providecommand{\Aa}{}\renewcommand{\Aa}{\mathcal{A}}
\providecommand{\Ad}{}\renewcommand{\Ad}{\widetilde{\mathrm{Ad}}}
\providecommand{\Aev}{}\renewcommand{\Aev}{\mathcal{A}^{+}}
\providecommand{\Aod}{}\renewcommand{\Aod}{\mathcal{A}^{-}}
\providecommand{\B}{}\renewcommand{\B}{\mathrm{B}}
\providecommand{\C}{}\renewcommand{\C}{\mathbb{C}}
\providecommand{\CAR}{}\renewcommand{\CAR}{\mathrm{CAR}}
\providecommand{\CS}{}\renewcommand{\CS}{{\mathbb{S}^{+}}}
\providecommand{\CSm}{}\renewcommand{\CSm}{{\mathbb{S}^{-}}}
\providecommand{\Cl}{}\renewcommand{\Cl}{\mathrm{Cl}}
\providecommand{\Con}{}\renewcommand{\Con}{\mathcal{C}}
\providecommand{\Dl}{}\renewcommand{\Dl}{\Delta}
\providecommand{\EN}{}\renewcommand{\EN}{E_{N}}
\providecommand{\End}{}\renewcommand{\End}{\mathrm{End}}
\providecommand{\F}{}\renewcommand{\F}{\mathbb{F}}
\providecommand{\G}{}\renewcommand{\G}{\Gamma}
\providecommand{\HH}{}\renewcommand{\HH}{\mathbb{H}}
\providecommand{\IS}{}\renewcommand{\IS}{I\sigma_3}
\providecommand{\J}{}\renewcommand{\J}{\mathbf{J}}
\providecommand{\LH}{}\renewcommand{\LH}{L_{H}}
\providecommand{\Mat}{}\renewcommand{\Mat}{\mathrm{M}}
\providecommand{\Mo}{}\renewcommand{\Mo}{M^{\circ}}
\providecommand{\Mp}{}\renewcommand{\Mp}{M^{+}}
\providecommand{\Msun}{}\renewcommand{\Msun}{M_{\odot}}
\providecommand{\Mt}{}\renewcommand{\Mt}{\widetilde{M}}
\providecommand{\Ncone}{}\renewcommand{\Ncone}{\mathcal{N}}
\providecommand{\Neff}{}\renewcommand{\Neff}{N_{\mathrm{eff}}}
\providecommand{\PCS}{}\renewcommand{\PCS}{\mathbb{P}{\mathbb{S}^{+}}}
\providecommand{\PT}{}\renewcommand{\PT}{\mathbb{PT}}
\providecommand{\Pin}{}\renewcommand{\Pin}{\mathrm{Pin}}
\providecommand{\Q}{}\renewcommand{\Q}{\mathbb{Q}}
\providecommand{\R}{}\renewcommand{\R}{\mathbb{R}}
\providecommand{\Ric}{}\renewcommand{\Ric}{\mathrm{Ric}}
\providecommand{\SBH}{}\renewcommand{\SBH}{S_{\mathrm{BH}}}
\providecommand{\Sc}{}\renewcommand{\Sc}{\mathbb{S}}
\providecommand{\Sch}{}\renewcommand{\Sch}{\mathrm{P}}
\providecommand{\Sig}{}\renewcommand{\Sig}{\Sigma}
\providecommand{\Som}{}\renewcommand{\Som}{S_\omega}
\providecommand{\Sp}{}\renewcommand{\Sp}{\mathbb{S}}
\providecommand{\Sph}{}\renewcommand{\Sph}{\mathrm{S}}
\providecommand{\Spi}{}\renewcommand{\Spi}{S_\pi}
\providecommand{\Spin}{}\renewcommand{\Spin}{\mathrm{Spin}}
\providecommand{\Sth}{}\renewcommand{\Sth}{S_{\mathrm{th}}}
\providecommand{\T}{}\renewcommand{\T}{\mathbb{T}}
\providecommand{\TH}{}\renewcommand{\TH}{T_{\mathrm{H}}}
\providecommand{\Tr}{}\renewcommand{\Tr}{\operatorname{Tr}}
\providecommand{\UHF}{}\renewcommand{\UHF}{\mathrm{UHF}}
\providecommand{\Z}{}\renewcommand{\Z}{\mathbb{Z}}
\providecommand{\Zc}{}\renewcommand{\Zc}{\mathrm{Z}}
\providecommand{\cmg}{}\renewcommand{\cmg}{\mathrm{cm^{2}\,g^{-1}}}
\providecommand{\dS}{}\renewcommand{\dS}{\mathrm{dS}}
\providecommand{\dsc}{}\renewcommand{\dsc}{\mathrm{dsc}}
\providecommand{\dtPG}{}\renewcommand{\dtPG}{dt}
\providecommand{\dtS}{}\renewcommand{\dtS}{dt_{\mathrm{S}}}
\providecommand{\ellv}{}\renewcommand{\ellv}{\ell}
\providecommand{\er}{}\renewcommand{\er}{e_r}
\providecommand{\et}{}\renewcommand{\et}{e_t}
\providecommand{\ev}{}\renewcommand{\ev}{\mathrm{ev}}
\providecommand{\gl}{}\renewcommand{\gl}{\mathfrak{gl}}
\providecommand{\hatV}{}\renewcommand{\hatV}{\widehat{V}}
\providecommand{\hb}{}\renewcommand{\hb}{\bar{h}}
\providecommand{\hf}{}\renewcommand{\hf}{\hat{f}}
\providecommand{\id}{}\renewcommand{\id}{\mathrm{id}}
\providecommand{\kpc}{}\renewcommand{\kpc}{\mathrm{kpc}}
\providecommand{\lP}{}\renewcommand{\lP}{\ell_{P}}
\providecommand{\mD}{}\renewcommand{\mD}{\mathbf{m}_{D}}
\providecommand{\mb}{}\renewcommand{\mb}{\bar m}
\providecommand{\mnu}{}\renewcommand{\mnu}{\Sigma m_{\nu}}
\providecommand{\nInf}{}\renewcommand{\nInf}{n_\infty}
\providecommand{\nO}{}\renewcommand{\nO}{n_0}
\providecommand{\ninf}{}\renewcommand{\ninf}{n_{\infty}}
\providecommand{\nz}{}\renewcommand{\nz}{n_{0}}
\providecommand{\proj}{}\renewcommand{\proj}{\operatorname{proj}}
\providecommand{\pv}{}\renewcommand{\pv}{\pi}
\providecommand{\rad}{}\renewcommand{\rad}{\mathrm{rad}}
\providecommand{\sgn}{}\renewcommand{\sgn}{\operatorname{sgn}}
\providecommand{\shA}{}\renewcommand{\shA}{\mathcal{A}}
\providecommand{\shC}{}\renewcommand{\shC}{\mathcal{C}}
\providecommand{\shE}{}\renewcommand{\shE}{\mathcal{E}}
\providecommand{\shJ}{}\renewcommand{\shJ}{\mathcal{J}}
\providecommand{\shK}{}\renewcommand{\shK}{\mathcal{K}}
\providecommand{\shL}{}\renewcommand{\shL}{L}
\providecommand{\shW}{}\renewcommand{\shW}{\mathcal{W}}
\providecommand{\sll}{}\renewcommand{\sll}{\mathfrak{sl}}
\providecommand{\so}{}\renewcommand{\so}{\mathfrak{so}}
\providecommand{\su}{}\renewcommand{\su}{\mathfrak{su}}
\providecommand{\supp}{}\renewcommand{\supp}{\operatorname{supp}}
\providecommand{\tr}{}\renewcommand{\tr}{\operatorname{tr}}
\providecommand{\uu}{}\renewcommand{\uu}{\mathfrak{u}}
\providecommand{\br}{}\renewcommand{\br}[1]{\langle #1\rangle}
\providecommand{\bra}{}\renewcommand{\bra}[1]{\langle #1|}
\providecommand{\file}{}\renewcommand{\file}[1]{\texttt{#1}}
\providecommand{\hp}{}\renewcommand{\hp}[2]{( #1 \,|\, #2 )}
\providecommand{\ip}{}\renewcommand{\ip}[2]{\langle #1, #2\rangle}
\providecommand{\ket}{}\renewcommand{\ket}[1]{|#1\rangle}
\providecommand{\rel}{}\renewcommand{\rel}[2]{D\bigl(#1 \,\|\, #2\bigr)}
\providecommand{\rev}{}\renewcommand{\rev}[1]{\widetilde{#1}}
\providecommand{\slashed}{}\renewcommand{\slashed}[1]{{#1}\mkern-10mu/}
\title{\bfseries Division by Zero at the Gravitational Singularity\\[4pt]\large\itshape The Transport Principle Applied to the Black Hole, and the Information Paradox as a Category Error}
\author{Leon Butler\\ \normalsize\href{mailto:lbutler.math@gmail.com}{\texttt{lbutler.math@gmail.com}}\\ \normalsize ORCID: \href{https://orcid.org/0009-0006-4449-5617}{0009-0006-4449-5617}}
\date{Generated by the Transport Lab, 27 October 2026 \textperiodcentered{} Transport Papers, T5\\[3pt]\small Official version: \href{https://doi.org/10.5281/zenodo.22721790}{doi:10.5281/zenodo.22721790}}
\begin{document}
\maketitle
\begin{abstract}
\normalsize
In gauge theory gravity the Schwarzschild solution is carried by an invertibility-demanding linear map, and the singularity is where the theory divides by zero. The series' transport algorithm is run on that division. The gauge field is a unipotent ray at every radius, so the divergence lies in the connection bivector; at the centre the light cones fold onto a null direction, a type-N obstruction whose forced repair is the Witt completion $\Cl(1,3) \to \Cl(2,4)$. No classical value repairs the dying directions, by two no-go theorems; the completion installs canonical anticommutation relations, one fermionic mode per boundary point, and the curvature factorises exactly into $-M/r^3$ times a fixed shape operator, selecting the boost plane. Descent to the metric level is, mode by mode, the fermion parity average, so every physical state has thermal single-mode marginals, and one named premise, KMS at the Euclidean section's period, sets their temperature at Hawking's and leaves the joint state free. Bosonisation at s-wave scope gives the Planck sector from a named quasi-freeness premise. The descent's multi-mode cost is a relative entropy carried by inter-quanta correlations: on the pointwise-descent reading, information is lost in $\Cl(1,3)$ and preserved in $\Cl(2,4)$, a category error between levels. In a rate model with boundary-mode locality the radiation entropy is exactly the window's thermal content, giving the Page envelope without typicality. Inside, the modular flow is the $t$-line holonomy of the collar, at Hubble-normalised rate $1/3$; the premised locality's pairing part is the tangent of the vacuum orbit.
\end{abstract}
\medskip
\noindent\textbf{MSC 2020.} Primary 83C57, 81T05; Secondary 15A66, 46L60, 94A17.\\
\textbf{Keywords.} black hole information paradox; Hawking radiation; Clifford algebra; conformal group; modular theory; thermofield double; Page curve; bosonisation.
\newpage
\noindent \textbf{Authorship and disclosure. }The founding idea of the series is the author's: that division by zero in a Clifford algebra should be read not as a prohibition but as a transport into a higher-dimensional algebra, and that the transported division need not be accompanied by an inverse operation. The direction of investigation and the critical review of the results are likewise the author's. The mathematical development and the drafting of the text were carried out in extended collaboration with Claude Fable 5 (Anthropic), an AI system.

\noindent\fbox{\parbox{\dimexpr\linewidth-2\fboxsep-2\fboxrule}{\textbf{This is not the version of record.} The official version of this paper is the one deposited on Zenodo, \href{https://doi.org/10.5281/zenodo.22721790}{doi:10.5281/zenodo.22721790}. This PDF was generated by the Transport Lab on 27 October 2026 and assembled from approved parts of its store: the results, proofs and text of the paper as the store holds them, each result with its premises and machine checks. It is regenerated when a result it contains is corrected, so it may differ from the deposited paper. Since the deposited version 1.0: 6 results corrected, 10 changed, 19 reworded (wording, notation or citations only), 7 results added and 33 cited results added. Each revised or added result is marked in the margin (\textcolor{BrickRed}{corrected}, \textcolor{BurntOrange}{changed}, \textcolor{gray}{reworded}), and the last section, \emph{Changes since the deposited version}, lists them with the reason for each.}}
\newpage

\section{Introduction}

\phantomsection\label{c:note.singularity.setting-introduction}
The information paradox is usually posed as a conflict between two theories. This paper poses it as a comparison of entropies computed in two algebras, and computes them.

The setting is gauge theory gravity (GTG), the formulation of general relativity in the spacetime algebra $\Cl(1,3)$ due to Lasenby, Doran and Gull \cite{LDG1998}, in which gravity is a gauge field $\hb$, a position-dependent linear map on vectors, and every metrical statement is built from $\hb^{-1}$: the theory's geometric content exists exactly where $\hb$ is invertible. the section ``Step 1: locate the failure'' sets the theory out exactly, field equations, connection bivector and gauge laws included, because the chain's first step is a statement about those equations and must be made in them. For Schwarzschild in the Newtonian gauge, $\hb$ is invertible in closed form at every $r > 0$, horizon included, and is unimodular there, so nothing determinant-like ever degenerates: the divergence is carried entirely by the connection bivector, and at $r = 0$ the description fails in every gauge, through a null kernel, in the exact rescaled-limit sense of Proposition \ref{c:singularity.degeneration-exactly}. The singularity is a division-by-zero locus in the literal, algebraic sense, and the transport principle of this series, the root paper T1 \cite{ButlerT1}, is a general theory of exactly such loci: it classifies the obstruction, forces a minimal completion, and quantifies what any return to the original level must lose. Its operational form is the six-step transport algorithm of \cite[\S 8.1, Table 1]{ButlerT1}, and this paper is that algorithm run end to end on the black hole; the table accompanying \ref{c:singularity.imported-equivalence-einstein-equations} gives the step-by-step correspondence.

The paper runs the resulting chain from the classical field to the Page curve, numbered by the transport algorithm of the root paper with the physics expanded where it lies (Steps 1 to 3, 4a to 4e, 5, 6a and 6b), each step carrying its declared status (proved, or proved from named premises). The steps are summarised in the table accompanying \ref{c:singularity.imported-equivalence-einstein-equations} and constitute the paper's outline. Steps 1 through 5 are theorems, at their stated scopes: the location and type of the failure, the forced ascent to $\Cl(2,4)$, the installation of canonical anticommutation relations where classical repair is barred, on the approach by Fact \ref{c:singularity.imported-no-c-number-repair} and at the boundary itself by Fact \ref{c:singularity.imported-no-equivariant-value-wall}, the selection of the boost plane and its rate, the identity of boost and number operator, single-mode thermality at exactly $\TH$ with the state's marginals fixed by univalence and one regularity premise rather than assumed, the bosonic Planck sector by bosonisation in the s-wave-reduced sector under the named joint quasi-freeness premise (Q), and the lossy descent with its exact multi-mode cost (Theorem \ref{c:singularity.deficit-theorem-cost-descent-exactly}). Steps 6a and 6b remain conditional in substance, but every condition is a \emph{named postulate}: the boundary-mode capacity is a theorem from three axioms (curvature cutoff, cell density, counted epoch, the causal-epoch integral computed exactly in Lemma \ref{c:singularity.causal-epoch-computed-exactly}), the radiation entropy of the rate model is an exact identity from localisation, purity, capacity, injection and boundary-mode locality, with the Page envelope its corollary and exact saturation confined to the zero-frequency idealisation, of which the linear weights supply exactly one mode, the sea level, so no window of two or more members ever saturates, the no-attic condition is derived rather than assumed, and the envelope with its endpoint needs no typicality and no assumption about where correlations are placed, boundary-mode locality being the one dynamical premise. The chain's conclusion on the paradox is a dissolution rather than a resolution: information is lost at the $\Cl(1,3)$ level and preserved at the $\Cl(2,4)$ level, these are the same physics read at two levels, and the discrepancy between the levels is not destruction but a fixed algebraic projection whose thermal output is the Hawking flux itself and whose size is an equation: the relative entropy of Theorem \ref{c:singularity.deficit-theorem-cost-descent-exactly}, carried entirely by the correlations of the quanta (Corollary \ref{c:singularity.deficit-carried-quanta-correlations}).

The paper also contains an exactly solvable finite-mode model (the section ``The finite-mode model: the dissolution step by step'') in which the dissolution can be followed step by step, at both levels. In the model, every single radiated quantum is \emph{provably, exactly} thermal, by a parity-protection lemma stating that no encoding by phase scrambling can be visible to any single-quantum measurement; the joint radiation state is pure; the information is recoverable from inter-quanta correlations; the mutual-information computation per emission, together with the Page-shaped entropy sequence, is computed in exact arithmetic and machine-verified; and the multi-mode descent is applied to the radiation outright (the section ``The descent, applied: the two levels in numbers''), so that the two-level conclusion is exhibited in exact numbers (the table accompanying \ref{c:singularity.purity-correlation-transfer-recovery} and Proposition \ref{c:singularity.two-levels-finite-model}) rather than argued. The model is a toy by construction and is offered as one: its value is that every claim in it is a computation, so the chain's conceptual endpoint (thermality as the record of what the projection cannot see; correlations as its carrier) is exhibited rather than argued.

A word on the standard of proof, binding throughout. The environments Theorem, Proposition and Lemma, with a Proof, are reserved for results proved \emph{in this paper}, at the standard of T1 \cite{ButlerT1}. Results proved elsewhere that the chain needs are stated as \textbf{Facts (imported)}, with their precise provenance; the typography tracks the epistemic content, so a reader can see at a glance which links of the chain this paper carries and which it inherits. Four links not previously proved are proved here outright: the finite-dimensional modular theory behind Step 4d; the state selection that completes it, the identification of the descent with the fermion parity average and the rigidity theorem that follows (Theorems \ref{c:singularity.descent-parity-average}, \ref{c:singularity.ascended-marginals-are-diagonal}, \ref{c:singularity.rate-rigidity-temperature-not-algebra-intrinsic}); the curvature behind Step 4b, derived in the gauge formalism in the text (Theorem \ref{c:singularity.curvature-gauge-field-exact-shape-decoupling}), with the appendix as an independent metric-picture cross-check; and the multi-mode descent behind Steps 5 and 6b, with its deficit theorem (Theorem \ref{c:singularity.deficit-theorem-cost-descent-exactly}).

Throughout, the physics the chain touches is kept in the text deliberately: $\Cl(2,4)$ as the conformal algebra and the twistor connection (Step 3), Israel's thermofield double and the Bisognano-Wichmann-Sewell mechanism (Step 4d), the island dictionary and the Hayden-Preskill mirror (Step 6b). These points of contact frame what kind of object the chain produces at each stage; the derivations do not depend on them.

 \begin{tabular}{c>{ \arraybackslash}p{4.7cm}>{ \arraybackslash}p{4.6cm}l}
\toprule
Step & Content & Key formula & Status \\
\midrule
1 & \emph{locate the failure}: unipotent ray; cones fold & $\hb = 1 - \beta P$, $P^2 = 0$ & proved \\
2 & \emph{diagnose the type}: N, and type C excluded & $p - q \equiv 6 \bmod 8$ & proved, given Facts \ref{c:sheaf.schwarzschild-degeneration-boundary-tame}, \ref{c:singularity.imported-one-sided-completion-singularity} \\
3 & \emph{apply the rule}: the forced ascent; conformal/twistor algebra & $\Cl(1,3) \to \Cl(2,4)$: $vv^* {+}\, v^*v = 1$ & proved, given Facts \ref{c:sheaf.schwarzschild-degeneration-boundary-tame}, \ref{c:singularity.imported-one-sided-completion-singularity} \\
4a & \emph{compute in the completion}: the transport instantiated; operators forced & $\iota(s,x,t) = (2t, x, \tfrac12 s)$ & proved, given Facts \ref{c:sheaf.schwarzschild-degeneration-boundary-tame}, \ref{c:singularity.imported-one-sided-completion-singularity} \\
4b & \emph{compute in the completion}: curvature selects the plane; the rate & $R(B) = -\tfrac{M}{r^3}S(B)$; $\kappa = 1/4M$ & proved (Schw.) \\
4c & \emph{compute in the completion}: boost $=$ number operator $=$ descent projectors & $\sigma = e_0 e_1 = 2N - 1$ & proved, given Facts \ref{c:sheaf.schwarzschild-degeneration-boundary-tame}, \ref{c:singularity.imported-one-sided-completion-singularity} \\
4d & \emph{compute in the completion}: thermality; the marginals fixed, the joint state free & $\tan\theta = e^{-\omega/2T}$; $\beta = 2\pi/\kappa$ & proved from (P), (R) \\
4e & \emph{compute in the completion}: photons by bosonisation; Bose-Einstein at $\TH$ & $[a_m, a_{-m}] = m$ & proved (s-wave; the Planck factor from (Q)) \\
5 & \emph{descend, and state what descent costs}: deficit $=$ correlations & $S(E\rho) - S(\rho) = \rel{\rho}{E\rho}$ & proved; the reading by Prop.~\ref{c:singularity.p-reduced-metric-level-reading-many} \\
6a & \emph{iterate, or stop}: the clock stops; boundary-mode capacity; area law & $dN_{\mathrm{6a}}/dt = \nu/\TH$; $C = N_{\mathrm{6a}}\ln 2$ & proved from (C1)-(C3) \\
6b & \emph{iterate, or stop}: two levels; Page envelope; eviction & $S_R = \sum_{k \in W} h_k \leq \min(\Sth, C)$ & proved from (L), (U), (C), (I), (L$_E$) \\
\midrule
-- & \textbf{new}: finite-mode correlation model, both levels & exact & proved (model) \\
\bottomrule
\end{tabular}

\emph{The transport chain, with statuses. The step names in italics are those of the transport algorithm of \cite[\S 8.1]{ButlerT1}, verbatim, and the section titles carry the same names; the table accompanying \ref{c:singularity.imported-equivalence-einstein-equations} gives the bijection. Each status is argued where the step is proved, and every premise is a named one stated in the text: univalence (P) and horizon regularity (R), the KMS condition at the period the Euclidean section fixes, imposed mode by mode, in Step 4d; joint quasi-freeness (Q) in Step 4e; the metric-level reading $\EN$, licensed in Step 5 by Proposition \ref{c:singularity.p-reduced-metric-level-reading-many} up to its definitional antecedent; the capacity postulates (C1)-(C3) in Step 6a; localisation (L), purity (U), capacity (C) and injection (I), collected as Premise \ref{c:singularity.rate-models-kinematic-inputs-l-u} in Step 6b, with boundary-mode locality (L$_E$) stated as Premise \ref{c:singularity.l-boundary-mode-locality} in Step 4d and used in Step 6b, of which (L)'s exhaustiveness clause is the one carrying conjectural weight, and (I)'s freshness clause, like (L$_E$)'s never-return clause, is the rate model's own assumption rather than an import. Given them the statements are theorems; Step 6b uses no typicality and claims the envelope, the exact rising branch, the quantified falling branch and the endpoint; exact saturation of the falling branch holds only in the zero-frequency idealisation, of which the linear weights supply exactly one mode.}

\subsection*{The chain as the transport algorithm}
\phantomsection\label{c:note.singularity.setting-chain-transport-algorithm}
The transport principle's operational form is the six-step transport algorithm of the root paper, \cite[\S 8.1, Table 1]{ButlerT1}, every step of which is licensed there by a theorem; and the worked-instance remark of that paper, \cite[Rem.~8.3]{ButlerT1}, promises that the present paper mirrors the algorithm row for row. the table accompanying \ref{c:singularity.imported-equivalence-einstein-equations} makes the promise explicit, and the chain is \emph{numbered by the algorithm}: Steps 1 to 3 coincide with the algorithm's; the algorithm's Step 4 (compute in the completion) expands into Steps 4a to 4e, because the computation in the completion is where the temperature, the statistics and the Planck sector are produced; Step 5 (descend, and state what descent costs) is the descent computation, quantified exactly by the deficit theorem; and Step 6 (iterate, or stop) is discharged by the mod-8 clock stopping at $d \equiv 6$, with Steps 6a and 6b counting what the boundary mode space at which it stops holds. The reader who wants each step's licence should read the table's middle column as a citation list into T1 \cite{ButlerT1}. The dictionary is: its Steps 1 to 3 are unchanged; its Steps 4 to 8 are here 4a, 4b, 4c, 4d, 4e; its Steps 9 and 10 are here 6a and 6b; and the quantification of the descent, folded into its Step 4, stands here as Step 5.

 \begin{tabular}{>{ \arraybackslash}p{2.9cm}>{ \arraybackslash}p{3.2cm}>{ \arraybackslash}p{6.6cm}}
\toprule
Algorithm step \cite[\S 8.1]{ButlerT1} & Warrant there & Instantiated here \\
\midrule
1. locate the failure & Lem.~2.1 & Step 1 (the section ``Step 1: locate the failure''): from the field equations; nowhere at $r > 0$, unavoidably at $r = 0$; in $\Omega$, not in $\det \hb$ \\
2. diagnose the type & Prop.~3.3; Thms 3.6, 3.9 & Step 2 (the section ``Step 2: diagnose the type''): type N (null kernel direction, Prop.~\ref{c:singularity.degeneration-exactly}); type C excluded, $p - q \equiv 6 \bmod 8$ \\
3. apply the rule & Thms 4.10, 4.13, 4.19 & Step 3 (the section ``Step 3: apply the rule''): Rule 1, $\Cl(1,3) \to \Cl(2,4)$, the CAR pair installed \\
4. compute in the completion & Thm 4.10; Prop.~4.21 & Steps 4a to 4e (the section ``Step 4a: compute in the completion; the transport instantiated, and the operators forced''-the section ``Step 4e: compute in the completion; from fermions to photons''): the transport instantiated, the plane and the rate, modular flow, $\TH$, Fermi-Dirac, Planck sector \\
5. descend, and state what descent costs & Lem.~4.23; Prop.~7.5; Thm 7.6 & Step 5 (the section ``Step 5: descend, and state what descent costs''): $\EN$, the deficit $\rel{\rho}{\EN\rho}$, the quanta correlations; applied in Steps 4d, 6b and the model (the section ``The descent, applied: the two levels in numbers'' in numbers) \\
6. iterate, or stop & Thm 7.1; Prop.~7.4 & Steps 6a and 6b (the section ``Step 6a: iterate, or stop; the clock stops, and the boundary-mode capacity''-the section ``Step 6b: iterate, or stop; two levels, the Page envelope, and eviction''): the tower stops at $d \equiv 6$ by T3 \cite{ButlerT3}, and the boundary mode space at which it stops is counted \\
\bottomrule
\end{tabular}

\emph{The transport algorithm of T1 \cite{ButlerT1}, instantiated. The middle column names the theorems of the root paper licensing each step; the right column names where this paper executes it.}

\subsection*{Prior art}
\phantomsection\label{c:note.singularity.setting-prior-art}
The physics targets are classical: Hawking's flux and temperature \cite{Hawking1975}; Israel's thermofield-double reading of the eternal black hole \cite{Israel1976}; the identification of modular flow with the boost for wedge and horizon algebras by Bisognano-Wichmann and Sewell \cite{BW1976, Sewell1982}; Page's curve and typicality theorem \cite{Page1993}; the quantum-extremal-surface derivations of the Page curve \cite{Penington2020, AEMM2019}; the Hayden-Preskill mirror \cite{HP2007}; the interior-volume growth \cite{CR2015}. The algebraic inputs are this series': the no-go and minimal completion, the taxonomy, the pinching and one-way theorems, and the transport algorithm this paper instantiates, from T1 \cite{ButlerT1}; the tower and CAR limit, from T3 \cite{ButlerT3}; and the localisation of the ladder on the degenerate locus, from T4 \cite{ButlerT4}. GTG is \cite{LDG1998, DoranLasenby2003}; bosonisation and the boson-fermion correspondence are \cite{PressleySegal1986}, met from the physics side by the near-horizon two-dimensional reduction \cite{RW2005}; twistors are \cite{Penrose1967}. Klein's inequality, the single classical input of the deficit theorem, is standard \cite{NielsenChuang}. This paper is the chain's self-contained form, with the statuses argued inline, the imports replaced by citations to the series' proved papers, and the finite-mode model new.

\subsection*{Contribution}
\phantomsection\label{c:note.singularity.setting-contribution}
At the standard of T1 \cite{ButlerT1}, with statuses declared per step.
(1) The chain, end to end, in one self-contained paper, as an explicit instance of the transport algorithm of \cite[\S 8.1]{ButlerT1}: the section ``Step 1: locate the failure'' through the section ``Step 6b: iterate, or stop; two levels, the Page envelope, and eviction'', mapped in the table accompanying \ref{c:singularity.imported-equivalence-einstein-equations}. (2) The gauge theory, exactly: the field equations stated, the connection bivector of the Newtonian-gauge solution computed in closed form with its Newtonian-field identity and divergence rate, the radial curvature derived from the gauge field alone, the degeneration made exact as a rescaled limit, and non-removability proved by a curvature invariant (the section ``Step 1: locate the failure''); the curvature's exact shape decoupling $R = -(M/r^3)S$ proved in the formalism, with the modular bivector of Step 4c matched to the curvature-selected plane by equation (Theorem \ref{c:singularity.curvature-gauge-field-exact-shape-decoupling}, Corollary \ref{c:singularity.matching-steps-4b-4c-join-equation}). (3) The dissolution stated exactly: the two-level theorem-schema of the section ``Step 6b: iterate, or stop; two levels, the Page envelope, and eviction'', with the capacity arm of the envelope proved from the Schmidt inequality and the thermal arm from boundary-mode locality, the eviction theorem proved in the rate model from the named inputs (L), (U), (C), (I) and (L$_E$), the radiation entropy being exactly the window's thermal content, with the no-attic condition derived rather than assumed (Proposition \ref{c:singularity.no-attic-derived-model}) and the falling branch quantified exactly, its saturation confined to the zero-frequency idealisation, the linear weights supplying exactly one such mode so that no window of two or more saturates (the claim made without typicality is the envelope, the rising branch, the quantified falling branch and the endpoint, not the saturated curve), and the level discrepancy itself now an equation: the multi-mode descent $\EN$, its deficit theorem $S(E\rho) - S(\rho) = \rel{\rho}{E\rho}$, and the identification of the deficit with the inter-quanta correlations (Step 5, the section ``Step 5: descend, and state what descent costs''). (4) The state selected, not assumed (the section ``What selects the state''): the descent identified with the fermion parity average, so that the descended form is state independent (Theorem \ref{c:singularity.descended-form-state-independent}) and per-mode thermality follows from univalence rather than from a symmetry hypothesis (Theorem \ref{c:singularity.ascended-marginals-are-diagonal}); the rigidity theorem showing that the temperature is not algebra intrinsic and that the calibration enters exactly once (Theorem \ref{c:singularity.rate-rigidity-temperature-not-algebra-intrinsic}); the rate fixed by one premise, the KMS condition at the geometric period, with the period computed from the smoothness of the Euclidean section (Proposition \ref{c:singularity.euclidean-section-period}, Theorem \ref{c:singularity.rate-discharge-hypothesis-iii}), the sign of the statistics located in the ungraded KMS identity (Proposition \ref{c:singularity.two-kms-functionals-sign}), and the premise's location shown to lie outside the framework's reach (Proposition \ref{c:singularity.why-calibration-cannot-be-internal}); global purity derived as the canonical purification (Theorem \ref{c:singularity.purity-output-partner-commutant}); and the over-determination theorem showing that no equilibrium principle can supply the correlations Steps 5 and 6b need (Theorem \ref{c:singularity.exact-equilibrium-over-determines-boundary-modes}). (5) The finite-mode model of the section ``The finite-mode model: the dissolution step by step'': the parity-protection lemma (single-quantum marginals are exactly thermal under arbitrary phase scrambling), global purity, correlation transfer, the exact mutual-information table, and the descent applied outright, with both levels exhibited in exact numbers (Proposition \ref{c:singularity.two-levels-finite-model}). (6) The modular flow followed inside (the section ``The modular flow of the boundary modes inside, and the entangling half of (L$_E$)''): identified with the $t$-line holonomy of the interior collar imported from T6 \cite{ButlerT6}, with rate $M/r^2$ per Killing time and Hubble-normalised limit $1/3$, so that a KMS state of that flow on the collar, if one exists, has the universal inverse temperature $6\pi$ (Proposition \ref{c:singularity.modular-flow-boundary-modes-holonomy-universal}); and the premised half of (L$_E$) read in $\mathfrak{so}(2\nu)$, its pairing part the tangent of the vacuum orbit and its differential part in the stabiliser (Proposition \ref{c:singularity.l-stabiliser-tangent}). (7) Machine verification of every computational identity in the paper, chain, gauge sector, state selection and model alike, in exact rational arithmetic \cite{ButlerVerifyT05}.

\section*{What we assume}
The results below are taken as given in this paper; each result that rests on one says so.
\begin{postulate}[C1: curvature cutoff]\label{c:singularity.c1-curvature-cutoff}
\changedmark{BurntOrange}{changed since the deposit}
The q-number boundary mode space of the singularity is supported on the stretched singularity, the surface on which the curvature reaches the Planck scale: $2M/r^3 = 1$ in Planck units, i.e.\ $r_c = (2M)^{1/3}$.

\emph{Note (form and calibration).} The form of this premise is forced, not chosen: an algebraic (pointwise) curvature criterion at a mass-independent value selects a member of the family $r_c^{3} = 2M/\lambda$, of which this premise is $\lambda = 1$, and, with an isometry-invariant density for Premise \ref{c:singularity.c2-cell-density}, the leading count sees the threshold $\lambda$ and the density $\nu$ only through the one number $\kappa_{s} = \nu\lambda^{-1/2}$ (Theorem \emph{The forms of (C1) and (C2) are forced by an algebraic curvature criterion and isometry invariance, leaving one number}). At leading order this premise together with (C2) is equivalent to (C2$'$), the modes carried by the exceptional set itself with $\nu' = \tfrac23\kappa_{s}$ modes per Planck area of its Kasner measure $\mu_{E} = 12\pi M$ per unit Killing time; the two differ only at relative order $(2M)^{-2/3}$ (Theorem \emph{At leading order the slice of (C1) can be dropped, the modes counted on the exceptional set with its Kasner measure}). The premise is kept as stated.

\emph{Scope (Schwarzschild calibration).} This premise, with Premise \ref{c:singularity.c2-cell-density} and the intrinsic-clock premise (IC) (Premise \emph{Intrinsic clock (IC), flagged premise} [T8, Principle 4.1]), is read as a calibration of the count on the Schwarzschild family only. Read literally on other holes the count fails the area law: on a Schwarzschild-de Sitter hole, with the counted epoch read as the hole's modular period, it is $u/(3u - 2)$ times $A_{b}/4$, $u = 2M/r_{b}$, for every $\Lambda > 0$, and no horizon's period gives $A_{b}/4$ (Theorem \emph{On a Schwarzschild-de Sitter hole the cutoff count over the hole's modular period is $u/(3u - 2)$ times $A_{b}/4$, $u = 2M/r_{b}$, an entropy per unit horizon area $1/(4(1 - 2\Lambda/3w_{h}))$, $w_{h} = 2M/r_{b}^{3}$, unbounded towards the Nariai limit; so the count per unit area is not curvature-independent and clause (Loc) cannot be obtained from (C1), (C2) and (IC) read with a horizon's modular period under one mode set}); for a spinning hole with $a \ge 2(2M)^{2/3}$ the cutoff locus is a timelike tube beyond the inner horizon and the literal count misses the area law by a factor of order $M^{-1/3}$ (Theorem \emph{For Kerr with a >= 2(2M)\^{}{2/3} the curvature-cutoff locus is a timelike tube beyond the inner horizon, and the counting rule read literally misses the area law by a factor of order M\^{}{-1/3}}). Away from the Schwarzschild family the count is read through the horizon placement, the modes on the horizon with a density fixed by flat tangent data, as clause (Loc) of Premise \emph{(C$_{\mathrm{loc}}$): on every local Rindler horizon the entropy of a cut is the integral of a curvature-independent area density (Loc) and, before matter crosses, the modes on the light rays are jointly thermal for the local boost (R$_{\mathrm{loc}}$)} states.
\end{postulate}

\begin{postulate}[C2: cell density]\label{c:singularity.c2-cell-density}
\changedmark{BrickRed}{corrected since the deposit}
The boundary mode space carries $\nu$ fermionic modes per Planck cell of the stretched singularity, with $\nu$ a constant that this premise leaves free (the entropy paper's area-law condition with the counted epoch at the modular period requires $\nu = 1/(16\pi\ln 2) \approx 0.029$ in the capacity convention, $1/16\pi \approx 0.020$ in the 1-nat convention, Prediction \emph{One boundary mode per Bekenstein-Mukhanov area quantum} [T8, T17 stance 6], so not of order one; its corank factor is fixed at one: the forced weighted completion gives exactly one ladder pair per point of the blown-up singular boundary, Remark \ref{c:note.singularity.corank-limit-weighted-answer} and \cite[Thm.~8.7, \S 8.9]{ButlerT4}; what remains in $\nu$ is the Planck-cell calibration).

\emph{Note (form and calibration).} The form of this premise is forced, not chosen: a density invariant under the isometries of the slice, which act transitively, is a constant multiple $\nu$ of the induced volume, and with the slice family $r_c^{3} = 2M/\lambda$ that an algebraic, mass-independent curvature criterion forces for Premise \ref{c:singularity.c1-curvature-cutoff}, the leading count sees $\nu$ and the threshold $\lambda$ only through the one number $\kappa_{s} = \nu\lambda^{-1/2}$ (Theorem \emph{The forms of (C1) and (C2) are forced by an algebraic curvature criterion and isometry invariance, leaving one number}); the registered density is a value of $\kappa_{s}$, equal to $\nu$ at $\lambda = 1$. At leading order this premise together with (C1) is equivalent to (C2$'$), the modes carried by the exceptional set itself with $\nu' = \tfrac23\kappa_{s}$ modes per Planck area of its Kasner measure $\mu_{E} = 12\pi M$ per unit Killing time; the two differ only at relative order $(2M)^{-2/3}$ (Theorem \emph{At leading order the slice of (C1) can be dropped, the modes counted on the exceptional set with its Kasner measure}). The premise is kept as stated.

\emph{Scope (Schwarzschild calibration).} This premise, with Premise \ref{c:singularity.c1-curvature-cutoff} and the intrinsic-clock premise (IC) (Premise \emph{Intrinsic clock (IC), flagged premise} [T8, Principle 4.1]), is read as a calibration of the count on the Schwarzschild family only: the density it calibrates there does not carry over to other holes through the cutoff placement. Read literally on a Schwarzschild-de Sitter hole, with the counted epoch read as the hole's modular period and the density fixed at $\Lambda = 0$, the count is $u/(3u - 2)$ times $A_{b}/4$, $u = 2M/r_{b}$, for every $\Lambda > 0$, and no horizon's period gives $A_{b}/4$ (Theorem \emph{On a Schwarzschild-de Sitter hole the cutoff count over the hole's modular period is $u/(3u - 2)$ times $A_{b}/4$, $u = 2M/r_{b}$, an entropy per unit horizon area $1/(4(1 - 2\Lambda/3w_{h}))$, $w_{h} = 2M/r_{b}^{3}$, unbounded towards the Nariai limit; so the count per unit area is not curvature-independent and clause (Loc) cannot be obtained from (C1), (C2) and (IC) read with a horizon's modular period under one mode set}); for a spinning hole with $a \ge 2(2M)^{2/3}$ the literal count misses the area law by a factor of order $M^{-1/3}$ (Theorem \emph{For Kerr with a >= 2(2M)\^{}{2/3} the curvature-cutoff locus is a timelike tube beyond the inner horizon, and the counting rule read literally misses the area law by a factor of order M\^{}{-1/3}}). Away from the Schwarzschild family the count is read through the horizon placement, the modes on the horizon with a density fixed by flat tangent data, as clause (Loc) of Premise \emph{(C$_{\mathrm{loc}}$): on every local Rindler horizon the entropy of a cut is the integral of a curvature-independent area density (Loc) and, before matter crosses, the modes on the light rays are jointly thermal for the local boost (R$_{\mathrm{loc}}$)} states.
\end{postulate}

\begin{postulate}[C3: counted epoch]\label{c:singularity.c3-counted-epoch}
The span of Killing time counted at any one moment of the exterior hole is the causal epoch: the portion of the singularity in the causal future of a single horizon-crossing, of order one Schwarzschild time by Lemma \ref{c:singularity.causal-epoch-computed-exactly}.
\end{postulate}

\begin{postulate}[(L$_E$): boundary-mode locality, as descent covariance on finite windows]\label{c:singularity.l-boundary-mode-locality}
\changedmark{BurntOrange}{changed since the deposit}
\emph{Premise (L$_{E}$), boundary-mode locality.} On every finite window of the boundary modes, $v_{1}, \dots, v_{n}$ satisfying the CAR, the internal dynamics $\alpha_{t}$ of the boundary modes commutes with the descent at all times:
\[
\alpha_{t}\circ\mathcal{E} \;=\; \mathcal{E}\circ\alpha_{t} \quad \text{for all } t, \qquad \mathcal{E} = E_{1}\circ\cdots\circ E_{n}, \quad E_{k}(x) = N_{k}xN_{k} + (1 - N_{k})x(1 - N_{k}),
\]
$\mathcal{E}$ the descent of the window (Proposition \ref{c:singularity.descent-what-it-erases}).

\emph{The version 1 wording, equivalent.} By Theorem \emph{On finite windows, boundary-mode locality is equivalent to the dynamics commuting with the descent for all times}, on every finite window this holds if and only if the dynamics is generated by elements of the abelian algebra $\mathcal{D}$ that the $\sigma_{k} = 2N_{k} - 1$ generate, diagonal in the joint number basis: the internal dynamics is generated by the geometric data the framework supplies on $E$, the boost bivector field $x \mapsto \sigma_{r}(x)$ and its products. Cross-point products such as $\sigma_{j}\sigma_{k}$ correlate phases without moving occupation; hopping $v_{j}^{*}v_{k} + v_{k}^{*}v_{j}$, pairing $v_{j}^{*}v_{k}^{*} + v_{k}v_{j}$ and every other generator with a component outside $\mathcal{D}$ are excluded by the one clause. Within the class of generators the framework supplies, the exclusion of hopping is a theorem (Theorem \ref{c:singularity.class-relative-no-go-no-supplied}); what the premise adds is the exclusion of unsupplied generators, pairing terms and the possible differential dynamics, which Proposition \ref{c:singularity.l-stabiliser-tangent} places in the tangent of the vacuum orbit and in its stabiliser.

\emph{What it forbids.} By Theorem \emph{Number-diagonal dynamics carries operator support only to directly coupled modes, with no exponential growth, so no dynamics (L\_E) admits realises the kappa\^{}{-1} ln N spreading law in real time}, every dynamics the premise admits carries an evolved mode operator only to the modes directly coupled to it, and every local commutator and out-of-time-order correlator is almost periodic and bounded: the premise forbids exponential real-time commutator growth. The real-time spreading law $t_{*} = \kappa^{-1}\ln N$ of Proposition \emph{The spreading law, and the chaos bound met, not saturated} [T8, Asm 6.1] is therefore not realised under it, and the scrambling time is a modular-time statement: the growth $e^{\kappa t}$ is a rate of the rotor flow, the KMS-strip continuation of the phase flow that is the windows' time evolution (Theorem \emph{The phase flow is the time evolution; the rotor flow is its KMS continuation}). The premise is kept over fast real-time scrambling; the two cannot both hold.
\end{postulate}

\begin{postulate}[(P): univalence]\label{c:singularity.p-univalence}
Physical states of the boundary modes are even: $\omega \circ \alpha = \omega$, equivalently $\omega$ annihilates the odd sector $\Aod$.
\end{postulate}

\begin{postulate}[(Q): joint quasi-freeness, with the s-wave weights at the half-integer offset]\label{c:singularity.q-joint-quasi-freeness-s-wave}
\changedmark{BurntOrange}{changed since the deposit}
In the stationary idealisation of the boundary modes, and for this step alone, the joint state of the mode array is taken to be the gauge-invariant quasi-free state with the marginal data that (P) and (R) (Premise \ref{c:singularity.r-regularity-fixed-locus}) fix unconditionally (Theorem \ref{c:singularity.rate-discharge-hypothesis-iii}, Corollary \ref{c:singularity.chains-state-class-stability}): writing $\psi_k$ for the fermion modes of the array of Theorem \ref{c:singularity.bosonisation-bridge-s-wave-sector}, its two-point function is $\langle \psi^\dagger_k \psi_l \rangle = \delta_{kl}\,n_k$, and every higher correlation function factorises by Wick's theorem, so that all connected correlators vanish and the joint state is KMS at $\beta = 2\pi/\kappa$ for the joint modular flow. The mode weights are evenly spaced with the array's sea level midway between two sites,
\[
\omega_k = \kappa\bigl(k + \tfrac12\bigr), \qquad k \in \Z ,
\]
so that no mode has zero weight, the weights are symmetric under the particle-hole relabelling $k \mapsto -1-k$, and the Killing flow at one modular period $2\pi/\kappa$ is the parity automorphism $\psi \mapsto -\psi$.

\emph{Why this offset.} The weights $\omega_k$ of Definition \ref{c:obj.singularity.geometric-modular-condition-g} are fixed nowhere else in the framework: given (P), (R) and the quasi-free clause, an affine law $\omega_k = ck + d$ is what the currents' being normal modes of the boost flow requires, particle-hole symmetry leaves exactly two offset classes, $d = 0$ and $d = c/2$, and the slope is the value $\beta c = 2\pi$ (Theorem \emph{Premise (Q) splits into the maximum-entropy member of the (P),(R) class and an affine weight law with two free data}(ii), (iv)). The class $d = c/2$ is adopted as a consequence of one reading, not as a bare premise: the array's modes are read as the discretised modes of the field theory's spinor fields on the horizon's light rays, $\psi(u) \mapsto e^{\kappa t/2}\psi(e^{\kappa t}u)$, of boost weight $\tfrac12$ (the weight at which the boost preserves the anticommutator), in the state that Premise \emph{(N$_{+}$): on the continuum of a horizon's null generators the state of the modes is a gauge-invariant quasi-free ground state for the affine null translations along the generators} (N$_{+}$) selects, a ground state for the affine null translations of the generators (Theorem \emph{A ground state for the affine null translations of a horizon's generators restricts on either side of the edge to the continuum I-state at rapidity period $2\pi$, so the continuum restatements of (R) for the eternal hole, (R$_{\Lambda}$), (R$_{\mathrm{loc}}$) and the joint KMS part of (Q) follow from one premise with no temperature in it, posable only on the continuum}). A field of boost weight $\tfrac12$ changes sign under a full turn of the regular Euclidean section, and $\kappa(k + \tfrac12)$ is the offset at which the array's Killing flow at one modular period carries that sign (Theorem \emph{With the sea level between sites the s-wave array carries exactly half of ln 2 at the Hawking temperature}). The reading is a decision, not a derivation: the Euclidean-period sign $R_{2\pi} = -1$ (Proposition \emph{Where the $+1$ of Fermi-Dirac comes from} [T5, Cor 5.4]) holds in a KMS state of the array at either offset, so identifying the array's real-time recurrence at one period with it is part of the reading.

\emph{On the continuum.} The joint KMS part of this premise is not a separate assumption on the continuum mode space of a horizon's generators: there it follows from (N$_{+}$), whose null ground state restricts on either side of the edge to the continuum I-state at $\beta = 2\pi/\kappa$, jointly KMS and even (Theorem \emph{A ground state for the affine null translations of a horizon's generators restricts on either side of the edge to the continuum I-state at rapidity period $2\pi$, so the continuum restatements of (R) for the eternal hole, (R$_{\Lambda}$), (R$_{\mathrm{loc}}$) and the joint KMS part of (Q) follow from one premise with no temperature in it, posable only on the continuum}(ii), (iii)). The stationary array statement here is its point-spectrum discretisation, on which (N$_{+}$) itself cannot be posed (part (iv) there). The quasi-free clause is a hypothesis of (N$_{+}$) as well, and (N$_{+}$) does not give the weight clause.

\emph{Consequences of the offset.} At these weights one array carries exactly $\tfrac12\ln 2$ at the Hawking temperature, given the classical theta-function identities, so under the entropy-matching hypothesis two arrays carry the capacity $\ln 2$ of one registered boundary mode with no zero-weight mode (Theorem \emph{With the sea level between sites the s-wave array carries exactly half of ln 2 at the Hawking temperature}(i), (iii)); the current sector of one array is one oscillator tower, one helicity channel, at either offset (Theorem \emph{Under (Q) an array's ln 2 is its zero-mode charge doublet, which no current with m != 0 connects, and its current sector is one oscillator tower}(ii), (iv)), so one registered mode then corresponds to two arrays and two helicity channels. Since every mode has nonzero weight, the gluing equation of Theorem \ref{c:singularity.rate-discharge-hypothesis-iii} fixes $\beta$ at every mode of the array, and the zero-weight scopes of Theorems \ref{c:singularity.descended-form-state-independent}, \ref{c:singularity.universality-temperature-under-g} and \ref{c:singularity.rate-discharge-hypothesis-iii} do not arise on it.

\emph{Version 2.} The source states the weight clause as $\omega_k = \kappa k$ with the sea level at $k = 0$, one mode of zero weight at which the gluing equation fixes $n_0 = \tfrac12$ but not $\beta$. Results computed under that offset refer to version 2.
\end{postulate}

\begin{postulate}[(R): regularity at the fixed locus]\label{c:singularity.r-regularity-fixed-locus}
\changedmark{BurntOrange}{changed since the deposit}
For every mode $k$ and every pair $a, b$ of elements of its ladder $\{v_k, v_k^{*}\}$, the correlation function $F_{a,b}(t) := \omega\bigl(a\,\alpha_t(b)\bigr)$ of the Killing flow $\alpha_t = \operatorname{Ad}(e^{it\omega_k K_k})$, an entire function of $t$, continues in imaginary Killing time to a single-valued function on the regular Euclidean section of the Schwarzschild gravitational field on flat spacetime (Proposition \ref{c:singularity.euclidean-section-period}; Definition \emph{Gauge Theory Gravity: the gravitational field on flat spacetime}): its two orderings are the two sides of the cut, and they meet after one period,
\[
F_{a,b}\bigl(t + i\,p\bigr) \;=\; \omega\bigl(\alpha_t(b)\,a\bigr), \qquad p = \frac{2\pi}{\kappa} .
\]
In words: each single-mode marginal is KMS for the Killing flow at exactly the period the Euclidean section's smoothness at its tip dictates. The regularity is local, posed at the tip of the Euclidean section (a polar chart about the tip of the Wick-rotated flat background); for a state that is KMS for the Killing flow, this is also the period at which it is regular across the horizon locus in the Painlev\'e-Gullstrand chart.

\emph{Note (continuum restatement).} Posed on the continuum of the horizon's null generators (Definition \emph{Continuum mode space} [T7, Def 4.1]), the content of (R) for the eternal (Kruskal) hole follows from Premise \emph{(N$_{+}$): on the continuum of a horizon's null generators the state of the modes is a gauge-invariant quasi-free ground state for the affine null translations along the generators}: by Theorem \emph{A ground state for the affine null translations of a horizon's generators restricts on either side of the edge to the continuum I-state at rapidity period $2\pi$, so the continuum restatements of (R) for the eternal hole, (R$_{\Lambda}$), (R$_{\mathrm{loc}}$) and the joint KMS part of (Q) follow from one premise with no temperature in it, posable only on the continuum}(iii) the generators' modes are then jointly KMS for the Killing flow at $\beta = 8\pi M$. The wording above, on the ladders $\{v_{k}, v_{k}^{*}\}$ with a phase-acting flow, is a discretisation of that statement, recovered only for modes of vanishing Mellin width; (N$_{+}$) itself cannot be posed on point spectrum. A hole formed by collapse is not covered by (N$_{+}$), and (R) remains a premise. (The whole generator of the eternal horizon, which that case of (N$_{+}$) uses, lies beyond one flat chart; see that premise.)
\end{postulate}

\begin{postulate}[The rate model's kinematic inputs: (L), (U), (C), (I)]\label{c:singularity.rate-models-kinematic-inputs-l-u}
\leavevmode
\begin{itemize}
\item[(L)] \emph{Localisation, derived in part.} By \cite[Thm.~4.4, Cor.~4.5, Thms 8.7, 8.12]{ButlerT4} and Step 4a, the completion installs q-number degrees of freedom exactly where division fails, which licenses the window mode space $L(t)$ inside the closure of the degenerate locus; the radiation $R(t)$ is a q-number sector on the model's own grounds, its quanta being the emitted modes carried off $E$ by the emission rather than outputs of the localisation theorems; and that these two exhaust the sectors is the model's exhaustiveness postulate. The total state space at each stage is $H_{L(t)} \otimes H_{R(t)}$.
\item[(U)] \emph{Purity.} The ascended description is globally pure, as in Fact \ref{c:singularity.assembled-page-curve-up-typicality}(i), each step of the evolution unitary.
\item[(C)] \emph{Capacity.} $\dim H_{L(t)} = 2^{|W(t)|}$ exactly, the number $|W(t)|$ of window modes filling to the mode count $N_{\mathrm{6a}}$ of Step 6a and non-increasing thereafter; the Page reading $C(t) \approx S_{\mathrm{BH}}(M(t))$ is for the full window.
\item[(I)] \emph{Injection.} The $k$-th emitted quantum enters in a thermofield pair $c_k|00\rangle + s_k|11\rangle$ with a fresh window mode, the pair form and the marginals from Theorem \ref{c:singularity.modular-descent} and Corollary \ref{c:singularity.chains-state-class-stability}; its \emph{freshness}, the clause that the pair arrives uncorrelated with everything already present, comes from neither cited result (Corollary \ref{c:singularity.chains-state-class-stability}(3) leaves exactly such correlations free) and is an independent assumption of the rate model, marked as such precisely as (L$_E$)'s never-return clause is. The quantum carries the marginal entropy $h_k := H(c_k^{2}) \leq \ln 2$, with $H(p) := -p\ln p - (1-p)\ln(1-p)$ the binary entropy in nats and equality exactly at the maximally entangled pair $c_k = s_k$; the thermal accumulation is $\Sth(t) := \sum_{k \leq t} h_k$.
\end{itemize}
\end{postulate}

\section{Step 1: locate the failure}

\phantomsection\label{c:note.singularity.setting-step-1-locate-failure}
Conventions: spacetime algebra $\Cl(1,3)$ with $\gamma_0^2 = +1$, $\gamma_i^2 = -1$; in the $(t,r)$ plane write $\et := \gamma_0$, $\er := \gamma_1$, so $\et^2 = 1$, $\er^2 = -1$, $\et \cdot \er = 0$. Units $G = c = \hbar = 1$; the Einstein coupling in the field equations is written $8\pi$, the symbol $\kappa$ being reserved throughout for the surface gravity of Step 4b.

\subsection*{The gauge theory, exactly}
\phantomsection\label{c:note.singularity.setting-gauge-theory-exactly}
In GTG, gravity is a pair of gauge fields on flat spacetime: the position-dependent linear map $\hb(a)$ on vectors, gauging local frame rescalings, and the connection bivector $\Omega(a)$, gauging local Lorentz rotations \cite{LDG1998} (Definition \emph{Gauge Theory Gravity: the gravitational field on flat spacetime}). The covariant derivative and the field strength are
\begin{equation}
D_a M \;=\; a\cdot\nabla M + \Omega(a) \times M, \qquad
R(a \wedge b) \;=\; a\cdot\nabla\,\Omega(b) - b\cdot\nabla\,\Omega(a) + \Omega(a) \times \Omega(b),
\end{equation}
with $A \times B := \tfrac12(AB - BA)$ the commutator product, so that $[D_a, D_b]M = R(a \wedge b) \times M$. The curvature operator on physical bivectors is $\mathcal{R}(B) = -R(\hb(B))$, in the store's curvature sign: Lasenby, Doran and Gull \cite{LDG1998} and \cite[Ch.~14]{DoranLasenby2003} use $R(\hb(B))$, and this single overall sign is the only difference. Contracting once and twice with the vector derivative gives the Ricci and scalar curvatures, $\mathcal{R}(b) = \partial_a \cdot \mathcal{R}(a \wedge b)$ and $\mathcal{R} = \partial_a \cdot \mathcal{R}(a)$, and the field equations of the theory are
\begin{equation}
\mathcal{R}(a) - \tfrac12\, a\, \mathcal{R} \;=\; -8\pi\, T(a), \qquad
D \wedge \hb^{T}(a) \;=\; \hb^{T}(\nabla \wedge a) ,
\end{equation}
where $\hb^{T}$ is the adjoint, $a \cdot \hb^{T}(b) = \hb(a) \cdot b$. The first is Einstein's equation in bivector-contracted vector form, with $T(a)$ the matter stress-energy (in the sign of \cite{LDG1998} it reads $+8\pi T(a)$); the second is the vanishing of torsion, which determines $\Omega$ algebraically from $\hb$ and its derivatives, so that $\hb$ carries the theory's entire dynamical content. The metric every observer measures is
\begin{equation}
g(a,b) \;=\; \hb^{-1}(a) \cdot \hb^{-1}(b),
\end{equation}
and the inverse in that formula is where the theory can fail: the geometric reading of the theory exists exactly where $\hb$ is invertible. The theory's two gauge freedoms act by
\begin{equation}
\hb(a) \;\mapsto\; \hb(\tilde U a U), \quad\text{equivalently}\quad \hb^{T}(a) \;\mapsto\; U\,\hb^{T}(a)\,\tilde U, \qquad
\Omega(a) \;\mapsto\; U\,\Omega(a)\,\tilde U - 2\,(a\cdot\nabla U)\,\tilde U
\end{equation}
for a rotor field $U$ (rotation gauge), together with the position-gauge displacements, which act by pullback \cite{LDG1998}; the curvature transforms covariantly under both, $\mathcal{R}(B) \mapsto U\,\mathcal{R}(\tilde U B U)\,\tilde U$ under the rotation gauge, so every scalar built from it is gauge-invariant. Every operation in these displays is interior to $\Cl(1,3)$: nothing invokes a curved manifold beyond the flat background.

\begin{fact}[Imported: equivalence with the Einstein equations]\label{c:singularity.imported-equivalence-einstein-equations}
On any region where $\hb$ is invertible, the GTG field equations $\mathcal{R}(a) - \tfrac12\, a\, \mathcal{R} \;=\; 8\pi\, T(a), \qquad D \wedge \hb(a) \;=\; \hb(\nabla \wedge a)$ for $(\hb, \Omega)$ are equivalent to the Einstein field equations for the metric $g(a,b) \;=\; \hb^{-1}(a) \cdot \hb^{-1}(b)$; in vacuum, to Ricci-flatness of $g$. Moreover the curvatures correspond, not only the equations: on the invertible region the GTG curvature bivector $\mathcal{R}(B)$, evaluated on the coordinate planes $g_\mu \wedge g_\nu$ of the frame $g_\mu = \hb^{-1}(e_\mu)$ and contracted with the reciprocal planes, reproduces the Riemann tensor $R_{\mu\nu}{}^{\rho\sigma}$ of $g$, so the curvature operator on bivectors, its spectrum, its eigenplanes and its scalar invariants are the same objects in the two pictures. \emph{Provenance:} Lasenby, Doran and Gull \cite[\S 4.3 and App.~C]{LDG1998}; textbook form, including the dictionary between $\mathcal{R}(B)$ and the Riemann tensor, \cite[Ch.~13-14]{DoranLasenby2003}. It is the second clause that licenses the assembly of Theorem \ref{c:singularity.curvature-gauge-field-exact-shape-decoupling}'s spectrum from the two formalisms.
\emph{Source:} A. Lasenby, C. Doran and S. Gull, Gravity, gauge theories and geometric algebra, Phil. Trans. R. Soc. Lond. A 356 (1998), 487--582., \S 4 (\S 4.3, covariant forms of the field equations), App.~C; C. Doran and A. Lasenby, Geometric Algebra for Physicists, Cambridge University Press (2003)., Ch.~13-14.
\end{fact}

\subsection*{The Schwarzschild gauge field, its connection, and where the divergence lies}
\phantomsection\label{c:note.singularity.setting-schwarzschild-gauge-field-connection-where}
For Schwarzschild in the Newtonian (Painlev\'e-Gullstrand) gauge the gauge field is
\begin{equation}
\hb(a) \;=\; a - \beta\,(a \cdot \et)\,\er \;=\; (1 - \beta P)(a), \qquad \beta(r) = \sqrt{2M/r}, \qquad P(a) := (a \cdot \et)\,\er .
\end{equation}

\begin{proposition}[The unipotent ray]\label{c:singularity.unipotent-ray}
$P^2 = 0$ exactly, at every radius; consequently $\hb$ is invertible at every $r > 0$, in closed form,
\[
\hb^{-1} \;=\; 1 + \beta P ,
\]
horizon included. The Schwarzschild gauge trajectory $\{1 - \beta P : \beta \geq 0\}$ is a ray of unipotent maps, and the horizon ($\beta = 1$) is algebraically unremarkable: division does not fail there.
\end{proposition}
\begin{proof}
$P^2(a) = P\bigl((a\cdot\et)\er\bigr) = (a\cdot\et)(\er\cdot\et)\,\er = 0$ since $\er \cdot \et = 0$. Then $(1 - \beta P)(1 + \beta P) = 1 - \beta^2 P^2 = 1$.

\end{proof}

\phantomsection\label{c:note.singularity.setting-schwarzschild-gauge-field-connection-where-2}
The metric $g(a,b) \;=\; \hb^{-1}(a) \cdot \hb^{-1}(b)$ that $\hb(a) \;=\; a - \beta\,(a \cdot \et)\,\er \;=\; (1 - \beta P)(a), \qquad \beta(r) = \sqrt{2M/r}, \qquad P(a) := (a \cdot \et)\,\er$ defines is computed blade-exactly in the section ``The metric of the gauge field, and its static form'': it is the Painlev\'e-Gullstrand line element $ds^2 = dt^2 - (dr + \beta\,dt)^2 - r^2 d\Omega^2$, with orthonormal coframe
\begin{equation}
\theta^0 = dt, \qquad \theta^1 = dr + \beta\,dt, \qquad \theta^2 = r\,d\vartheta, \qquad \theta^3 = r\sin\vartheta\,d\varphi .
\end{equation}
The next proposition computes the second gauge field, which turns out to be the carrier of the divergence.

\begin{proposition}[The connection bivector, computed]\label{c:singularity.connection-bivector-computed}
\changedmark{gray}{reworded}
For the Painlev\'e-Gullstrand gauge field $\hb = 1 - \beta P$ of Fact \emph{The gauge field and its cometric} [T4, Fact 8.1] on flat spacetime (Definition \emph{Gauge Theory Gravity: the gravitational field on flat spacetime}), whose physical-vector map in components is the coframe $\theta^0 = dt, \qquad \theta^1 = dr + \beta\,dt, \qquad \theta^2 = r\,d\vartheta, \qquad \theta^3 = r\sin\vartheta\,d\varphi$, $\theta^a(v) = \gamma^a \cdot \hb^{-1}(v)$, write $\sigma_r$ for the unit radial boost bivector of that coframe (the sign convention is fixed once, in Remark \ref{c:note.singularity.frame-made-explicit}; the verification script fixes it by computation). Then:
\begin{enumerate}
\item all connection data is algebraic in $(r, \beta)$: $\beta' = -\beta/2r$ and $\beta\beta' = -M/r^2$ exactly, so every coefficient below lies in the field $\mathbb{Q}(r)[\beta]$ with $\beta^2 = 2M/r$;
\item the torsion condition of $\mathcal{R}(a) - \tfrac12\, a\, \mathcal{R} \;=\; 8\pi\, T(a), \qquad D \wedge \hb(a) \;=\; \hb(\nabla \wedge a)$ is solved by the rotation-gauge field $\Omega$ whose components in this coframe are the connection one-forms
\[
\omega^1{}_0 = -\beta'\,\theta^1, \qquad
\omega^2{}_0 = -\tfrac{\beta}{r}\,\theta^2, \qquad
\omega^3{}_0 = -\tfrac{\beta}{r}\,\theta^3, \qquad
\omega^2{}_1 = \tfrac{1}{r}\,\theta^2, \qquad
\omega^3{}_1 = \tfrac{1}{r}\,\theta^3, \qquad
\omega^3{}_2 = \tfrac{\cot\vartheta}{r}\,\theta^3,
\]
verified by substitution in the first structure equations: $d\theta^0 = 0$ and $-\omega^0{}_1 \wedge \theta^1 = \beta'\,\theta^1 \wedge \theta^1 = 0$; $d\theta^1 = \beta'\,dr\wedge dt = \beta'\,\theta^1\wedge\theta^0 = -\,\omega^1{}_0\wedge\theta^0$; and $d\theta^2 = \tfrac1r\,dr\wedge\theta^2 = \tfrac1r(\theta^1 - \beta\theta^0)\wedge\theta^2 = -\,\omega^2{}_0\wedge\theta^0 - \omega^2{}_1\wedge\theta^1$, with the $\theta^3$ leg the round sphere's standard connection;
\item evaluated on the background directions $\partial_t$, $\partial_r$ of the $(t,r)$ plane, the connection bivector is
\[
\Omega(\partial_t) \;=\; -\beta\beta'\,\sigma_r \;=\; \frac{M}{r^2}\,\sigma_r ,
\qquad
\Omega(\partial_r) \;=\; -\beta'\,\sigma_r \;=\; \frac{\beta}{2r}\,\sigma_r \;=\; \sqrt{\frac{M}{2r^3}}\;\sigma_r :
\]
the time-direction connection is the Newtonian field $M/r^2$, exactly and at every radius, and the radial-direction connection diverges as $r \to 0$ at the rate $r^{-3/2}$, the square root of the field-strength scale $M/r^3$.
\end{enumerate}
\end{proposition}
\begin{proof}
The structure equations below are the torsion condition written in components of the coframe $\theta^a(v) = \gamma^a \cdot \hb^{-1}(v)$, and the one-forms $\omega^a{}_b$ are the components of the rotation-gauge field, so that evaluating them on a background direction gives $\Omega(a)$ (Fact \emph{Imported: the dictionary between Gauge Theory Gravity and curved-spacetime language}(i) and (ii)). (1) $\beta = (2M)^{1/2} r^{-1/2}$ gives $\beta' = -\tfrac12 (2M)^{1/2} r^{-3/2} = -\beta/2r$ and $\beta\beta' = -\beta^2/2r = -M/r^2$. (2) The displayed substitutions, each one line; the coefficients of $\omega^2{}_0$ and $\omega^2{}_1$ are read off the decomposition $dr = \theta^1 - \beta\theta^0$, which is where the gauge field's off-diagonal character enters the connection. The script verifies the $\theta^0$, $\theta^1$ and $\theta^2$ structure equations as exact identities in $\mathbb{Q}(r)[\beta]$. (3) The frame duals of $\theta^0 = dt, \qquad \theta^1 = dr + \beta\,dt, \qquad \theta^2 = r\,d\vartheta, \qquad \theta^3 = r\sin\vartheta\,d\varphi$ satisfy $\partial_t = \hat e_0 + \beta\,\hat e_1$ and $\partial_r = \hat e_1$, so $\omega^1{}_0(\partial_t) = -\beta'\,\theta^1(\partial_t) = -\beta'\beta$ and $\omega^1{}_0(\partial_r) = -\beta'$; assembling into the bivector-valued $\Omega$ gives the display, and the numerical identities are (1).

\end{proof}

\begin{proposition}[The solution, and where the divergence lies]\label{c:singularity.solution-where-divergence-lies}
The pair $(\hb, \Omega)$ of $\hb(a) \;=\; a - \beta\,(a \cdot \et)\,\er \;=\; (1 - \beta P)(a), \qquad \beta(r) = \sqrt{2M/r}, \qquad P(a) := (a \cdot \et)\,\er$ and Proposition \ref{c:singularity.connection-bivector-computed} solves the vacuum field equations $\mathcal{R}(a) - \tfrac12\, a\, \mathcal{R} \;=\; 8\pi\, T(a), \qquad D \wedge \hb(a) \;=\; \hb(\nabla \wedge a)$ at every $r > 0$, and the divergence of the classical description resides in $\Omega$ and in the curvature, not in $\hb$:
\begin{enumerate}
\item (The radial curvature, from the gauge field alone.) On the radial plane both connection values are multiples of the single bivector $\sigma_r$, so the commutator term of the field strength the relation displayed with \ref{c:singularity.imported-equivalence-einstein-equations} vanishes identically and nothing depends on $t$; hence, in two lines,
\[
\mathcal{R}(\partial_t \wedge \partial_r)
\;=\; \partial_t\,\Omega(\partial_r) - \partial_r\,\Omega(\partial_t) + \Omega(\partial_t)\times\Omega(\partial_r)
\;=\; -\,\partial_r\Bigl(\frac{M}{r^2}\Bigr)\sigma_r
\;=\; \frac{2M}{r^3}\,\sigma_r :
\]
the radial plane is an eigenplane of the curvature with eigenvalue of magnitude $2M/r^3$, derived without leaving the gauge formalism.
\item (Vacuum.) The full curvature operator is diagonal on the six frame planes with eigenvalues of magnitudes $(2M/r^3, M/r^3, M/r^3, M/r^3, M/r^3, 2M/r^3)$ and vanishing Ricci contraction: the mixed and angular planes are computed in the section ``The curvature operator, by the Cartan structure equations'' in the static chart, where the computation is shortest, and the two pictures agree on the invertible region by Fact \ref{c:singularity.imported-equivalence-einstein-equations}. The shape this spectrum assembles into is Theorem \ref{c:singularity.curvature-gauge-field-exact-shape-decoupling}'s task.
\item (Unimodularity.) $\det \hb = 1$ exactly, at every radius, horizon included: $P$ is rank-one nilpotent, so $1 - \beta P$ is unipotent as a matrix. Likewise the $(t,r)$ block of the metric $g(a,b) \;=\; \hb^{-1}(a) \cdot \hb^{-1}(b)$ has determinant $-1$ at every radius. Nothing determinant-like degenerates at any $r > 0$: at finite radius the classical field is a volume-preserving shear, and the entire divergence of the approach is carried by the connection (Proposition \ref{c:singularity.connection-bivector-computed}(3)) and by the curvature, both of which diverge, and by the metric components themselves, which diverge without a rescaled limit unless one is taken exactly as in Proposition \ref{c:singularity.degeneration-exactly}.
\end{enumerate}
\end{proposition}
\begin{proof}
(1) The display, using Proposition \ref{c:singularity.connection-bivector-computed}(3) and $-\,\partial_r(Mr^{-2}) = 2M/r^3$; the vanishing of the commutator is $\sigma_r \times \sigma_r = 0$. The script checks the identity $(\beta\beta')' = 2M/r^3$ in $\mathbb{Q}(r)[\beta]$, the vanishing of the angular wedge contribution, and the two-form computation itself. (2) the section ``The curvature operator, by the Cartan structure equations'', whose eigenvalues are verified as rational-function identities; vacuum is the vanishing of the Ricci contraction of that spectrum, equivalently Fact \ref{c:singularity.imported-equivalence-einstein-equations} against the Ricci-flatness of Schwarzschild. (3) $\det(1 - \beta P) = 1$ for nilpotent $P$ (unipotent matrices have unit determinant; the script evaluates the determinant exactly at sample $\beta$ including the horizon), and the PG $(t,r)$ block $\bigl(\begin{smallmatrix} 1-\beta^2 & -\beta \\ -\beta & -1 \end{smallmatrix}\bigr)$ has determinant $-(1-\beta^2) - \beta^2 = -1$.

\end{proof}

\begin{proposition}[The cones fold]\label{c:singularity.cones-fold}
Track the null directions of the $(t,r)$ plane in the metric $g(a,b) \;=\; \hb^{-1}(a) \cdot \hb^{-1}(b)$, $g(a,a) = \hb^{-1}(a)^{2}$ with $\hb^{-1} = 1 + \beta P$: for $a = \et + \lambda\,\er$,
\[
\hb^{-1}(a) = \et + (\lambda + \beta)\,\er, \qquad g(a,a) = 1 - (\lambda + \beta)^2 ,
\]
so $a$ is $g$-null exactly when $\lambda = -\beta \pm 1$: the physical cone's two null directions are $\et + (-\beta \pm 1)\er$, the radial null curves $dr/dt = -\beta \pm 1$ of the Painlev\'e-Gullstrand line element of the section ``The metric of the gauge field, and its static form''. At the horizon ($\beta = 1$) they are $\lambda = 0$, the hovering ray, and $\lambda = -2$, an ingoing one. Normalising each by its $\er$-coefficient and letting $r \to 0$ (so $\beta \to \infty$),
\[
\tfrac{1}{-\beta \pm 1}\,\et + \er \;\longrightarrow\; \er :
\]
both branches of the light cone collapse onto the single direction $\er$, both tipping towards $-\er$, inward, on the way. Simultaneously $g(\et, \et) = 1 - \beta^2 \to -\infty$: no rescaling of the classical field stays finite and nondegenerate. The degeneration is through a null kernel direction in the exact sense of Proposition \ref{c:singularity.degeneration-exactly}, and it is not removable by gauge choice (Proposition \ref{c:singularity.non-removability-invariant-argument}).
\end{proposition}
\begin{proof}
The displayed computations, using $\et \cdot \er = 0$ to remove the cross term in $\hb^{-1}(a)^{2}$ (orthogonal vectors anticommute, \ref{c:std-algebra.orthogonal-anticommute}) (note that it is $\hb^{-1}$, not $\hb$, that enters the metric: the $g$-null vectors are the $\hb$-images of the flat null vectors, $\hb(\et \pm \er) = \et + (\pm 1 - \beta)\er$); the limit is elementary. The two claims of the final sentence are the two propositions cited (\ref{c:singularity.degeneration-exactly} and \ref{c:singularity.non-removability-invariant-argument}), which are the exact forms this proposition's limits point at.

\end{proof}

\subsection*{Where division fails}
\phantomsection\label{c:note.singularity.setting-where-division-fails}
The folded direction deserves a caution before it is used: the display above is projective (each normalised vector is exactly $g$-null at finite $\beta$, but its limit $\er$ has $g(\er, \er) = -1$ at every radius), and the metric itself does not converge at all. The correct statement requires a rescaling, and the next proposition supplies it.

\begin{proposition}[The degeneration, exactly]\label{c:singularity.degeneration-exactly}
As $r \to 0$, exactly and with every identity holding at finite $\beta$:
\begin{enumerate}
\item (The map.) $\beta^{-1}\hb = \beta^{-1}\,\mathrm{id} - P$, so $\beta^{-1}\hb \to -P$: a rank-one map with image $\R\,\er$ and kernel $\et^{\perp} = \operatorname{span}(\er, e_\vartheta, e_\varphi)$.
\item (The metric.) $\beta^{-2} g(a,b) = -(a\cdot\et)(b\cdot\et) + \beta^{-1}\bigl[(a\cdot\et)(b\cdot\er) + (b\cdot\et)(a\cdot\er)\bigr] + \beta^{-2}(a \cdot b)$, so $\beta^{-2}g \to \hat g$ with
\[
\hat g(a,b) \;=\; -(a\cdot\et)(b\cdot\et) :
\]
a form of rank one and signature $(0,1,3)$, whose radical is $\et^{\perp}$, the kernel of the limit map.
\item (The null direction, located.) The folded direction $\er$ of Proposition \ref{c:singularity.cones-fold} lies in the radical of $\hat g$ and is $\hat g$-null, while $g(\er,\er) = -1$ at every finite radius: the nullity of the degenerate direction is a property of the limit form, not of the flat background, and this is the exact content of the phrase ``degeneration through a null kernel direction''. It is the datum Step 2 uses.
\end{enumerate}
\end{proposition}
\begin{proof}
(1) and (2) are the displayed algebraic identities, which hold at every finite $\beta$ (expand $g(a,b) \;=\; \hb^{-1}(a) \cdot \hb^{-1}(b)$ with $\hb^{-1} = 1 + \beta P$ and divide); the limits are then immediate, and the rank, kernel, image and signature statements are finite linear algebra, checked exactly in the script (the signature by exact congruence diagonalisation). (3) $\hat g(\er, \er) = 0$ since $\er \cdot \et = 0$, and $\er$ lies in the radical computed in (2); $g(\er,\er) = \er^2 = -1$.

\end{proof}

\begin{proposition}[Non-removability: the invariant argument]\label{c:singularity.non-removability-invariant-argument}
The curvature scalar
\[
K \;:=\; R_{abcd}R^{abcd} \;=\; 4\Bigl(2\,\bigl(\tfrac{2M}{r^3}\bigr)^2 + 4\,\bigl(\tfrac{M}{r^3}\bigr)^2\Bigr) \;=\; \frac{48\,M^2}{r^6}
\;\longrightarrow\; \infty \quad (r \to 0),
\]
and $K$ is invariant under every gauge transformation of the theory, $r$-dependent transformations of arbitrary size included. Hence no gauge choice removes the degeneration at $r = 0$: the division-by-zero locus is a property of the solution, not of its presentation.
\end{proposition}
\begin{proof}
The middle equality evaluates $K$ from the eigenvalue spectrum of Proposition \ref{c:singularity.solution-where-divergence-lies}(2) (each plane contributes four times its eigenvalue squared), and is verified as an exact rational-function identity in the script. Invariance: under the rotation-gauge law $\hb(a) \;\mapsto\; R\,\hb(a)\,\tilde R, \qquad \Omega(a) \;\mapsto\; R\,\Omega(a)\,\tilde R - 2\,(a\cdot\nabla R)\,\tilde R$ the curvature transforms by conjugation, $\mathcal{R}(B) \mapsto R\,\mathcal{R}(B)\tilde R$, which preserves every scalar contraction; under position-gauge displacements it transforms by pullback, which relocates but cannot alter the function's unbounded range on any neighbourhood of the singular locus. A divergent invariant cannot be transformed to a finite one.

\end{proof}

\begin{remark}[The corank of the limit, and the weighted answer]\label{c:note.singularity.corank-limit-weighted-answer}
Proposition \ref{c:singularity.degeneration-exactly} used the uniform rescaling weight $\beta^{-2}$, and under it the limit form has corank three. The choice of weight is a genuine choice, and it is not an open one: the Schwarzschild composite of the sheaf paper settles it. There, after the real blow-up of the worldline (which absorbs the two angular rates as the exceptional divisor's own tautological rescaling), the rescaling weights are \emph{forced}, per the geometrically named dying slot, and equal the Kasner exponents $(-\tfrac13, \tfrac23, \tfrac23)$ \cite[Thm.~8.16]{ButlerT4}; the weighted coframe module carrying them is unique in the rain-bounded class \cite[Thm.~8.18]{ButlerT4}; and in that module the degeneration is boundary-tame of corank exactly \emph{one}, with radical generated by $k = \sqrt r\,dr + \sqrt{2M}\,dt_{\mathrm{PG}}$, conormal to the infalling congruence, whose $G$-length is $-r$ exactly at every radius, the Killing-time covector $dt_s$ spanning the radical line only on $E$ itself \cite[Thm.~8.7]{ButlerT4}. The corank-three limit of Proposition \ref{c:singularity.degeneration-exactly} is the unweighted, uniform-rescaling shadow of that structure (the two forms rescale different objects, the metric on vectors here, the cometric on the weighted coframes there, so there is no contradiction; the forced answer is the weighted one). Where the answer bites is Step 6a: Axiom \ref{c:singularity.c2-cell-density}'s cell density carries a corank factor, and that factor is fixed: one ladder pair per point of the blown-up singular boundary \cite[\S 8.9]{ButlerT4}, so the corank factor in $\nu$ is one, and what remains open in $\nu$ is the cell calibration and the epoch, assigned to the counting paper (T8). Nothing in Steps 1 through 4e depends on any of this.
\end{remark}

\phantomsection\label{c:note.singularity.setting-where-division-fails-2}
a figure of the original displays the folding. The section's propositions locate the division exactly: nowhere at $r > 0$ (Propositions \ref{c:singularity.unipotent-ray} and \ref{c:singularity.solution-where-divergence-lies}(3): the field solves the equations, invertibly and unimodularly), unavoidably at $r = 0$ (Propositions \ref{c:singularity.degeneration-exactly} and \ref{c:singularity.non-removability-invariant-argument}), and through a direction that is null for the limit form (Proposition \ref{c:singularity.degeneration-exactly}(3)). That last fact is the input the series' taxonomy needs, and this is the algorithm's Step 1 discharged \cite[\S 8.1, Step 1]{ButlerT1}.

\section{Step 2: diagnose the type}

\begin{proposition}[Type N, and only type N]\label{c:singularity.type-n-only-type-n}
The obstruction at $r = 0$ is type N in the taxonomy of \cite[\S 3]{ButlerT1}: an isotropic direction without a metric partner, in the degenerate sense in which the root paper hands such cases to the sheaf paper. Type C cannot occur: $\Cl(1,3)$ has $p - q = -2 \equiv 6 \pmod 8$, hence is central simple ($\cong \Mat_2(\HH)$), has no central idempotents, and offers ideal fusion nothing to fuse. Type P, which is present in $\Cl(1,3)$ as in every algebra of its size, demands no transport: its rule is the rule of no rule \cite[\S 6, Prop.~6.1]{ButlerT1}, so it neither licenses nor obstructs the ascent. The operative transport rule is Rule 1, the Witt completion, and by \cite[Thm.~3.6]{ButlerT3} no floor of the resulting tower ever develops a type-C obstruction either. This is the algorithm's Step 2 \cite[\S 8.1, Step 2]{ButlerT1}, with both diagnostic questions answered by computation.
\end{proposition}
\begin{proof}
The degenerating direction is null for the limit form (Proposition \ref{c:singularity.degeneration-exactly}(3)). This is the type-N datum of \cite[Def.~3.1, Prop.~3.3]{ButlerT1} (\ref{c:transport.existence-locus-type-n}) in a degenerate setting that the root paper's taxonomy does not itself cover: there, type N is a null vector of a nondegenerate form, which acquires its annihilating idempotent through a hyperbolic partner, whereas here $\er$ lies in the radical of $\hat g$ and has no partner at all. The root paper is explicit that its algorithm assumes nondegeneracy and that the degenerate case is the sheaf paper's \cite[Rem.~4.15, Rem.~8.4]{ButlerT1}; the warrant for treating the radical direction as type N and applying Rule 1 fibrewise is therefore the localisation theorem of \cite[Thm.~4.4, Cor.~4.5]{ButlerT4} (\ref{c:sheaf.one-sided-completion}, \ref{c:sheaf.localised-transport}) in the Schwarzschild form \cite[Thms 8.7, 8.12]{ButlerT4} (\ref{c:sheaf.schwarzschild-degeneration-boundary-tame}, \ref{c:sheaf.one-sided-completion-singularity}), which Step 3 invokes, and which supplies the hyperbolic partner by completion exactly on $E$. The absence of type C is the classification: the mod-8 arithmetic and the classification $\Cl(1,3) \cong \Mat_2(\HH)$ are classical \cite{Porteous1995} (\ref{c:primer.classification-real-clifford-algebras}, \ref{c:std-algebra.sta-matrix}); the existence criterion for type C ($p - q \equiv 1 \bmod 4$) is \cite[Thm.~3.6]{ButlerT1} (\ref{c:transport.existence-classification-type-c}), and $6 \not\equiv 1 \bmod 4$; the root paper's own worked signature for $\Cl(1,3)$ is \cite[Ex.~3.12]{ButlerT1}. That type P imposes nothing is \cite[Prop.~6.1]{ButlerT1} (\ref{c:transport.type-p-self-repairing}): the identity assignment repairs every idempotent at cost zero. The tower statement is the frozen clock of T3 \cite{ButlerT3} (\ref{c:tower.rigidity-clock-frozen}). The script checks the mod-8 datum and the centre; the type-N diagnosis itself is a citation, not a computation.

\end{proof}

\section{Step 3: apply the rule}

\phantomsection\label{c:note.singularity.setting-step-3-apply-rule}
Rule 1 adjoins one hyperbolic plane, minimally and uniquely \cite[Thms.~4.10, 4.13]{ButlerT1}, with minimality holding even against arbitrary unital embeddings \cite[Thm.~4.19]{ButlerT1}: $\Cl(1,3) \hookrightarrow \Cl(2,4)$, with adjoined null pair $(v, v^*)$ satisfying
\begin{equation}
v^2 = (v^*)^2 = 0, \qquad v v^* + v^* v = 1 . \tag{CAR}
\end{equation}
At the geometric level, the localisation theorem of the sheaf paper \cite[Thm.~4.4, Cor.~4.5]{ButlerT4}, extended to exact Schwarzschild by its composite \cite[Thms 8.7, 8.12]{ButlerT4}, performs this ascent fibrewise over the degenerate locus and only there: the completed pair exists exactly where the division fails. Those theorems are stated for tame interior walls, whose hypotheses exact Schwarzschild does not satisfy (the singular locus is a worldline of codimension three, and the degeneration is anisotropic); the import at the black hole is nevertheless not by analogy, because the Schwarzschild composite of \cite[\S 8]{ButlerT4} extends the localisation to this very case: after the real blow-up of the worldline and in the forced weighted coframe module, the degeneration is boundary-tame of corank one with the Killing-time covector as its radical \cite[Thm.~8.7]{ButlerT4}, the completion is performed with the ladder pair supported exactly on the blown-up singularity $E \cong \R \times S^2$, and no shadow wall forms \cite[Thm.~8.12]{ButlerT4}. The ascent this paper runs is therefore theorem-backed at the black hole itself. This is the algorithm's Step 3 \cite[\S 8.1, Step 3]{ButlerT1}.

\begin{lemma}[Pseudo-inverse identities, and the hyperbolic atom]\label{c:singularity.pseudo-inverse-identities-hyperbolic-atom}
\changedmark{gray}{reworded}
From $v^2 = (v^*)^2 = 0, \qquad v v^* + v^* v = 1$ alone: $v v^* v = v$ and $v^* v v^* = v^*$, so each member of the pair is a von Neumann pseudo-inverse of the other, with no genuine inverse existing ($v$ is a zero divisor). The pair is realisable inside a single $(1,1)$ block: in $\Cl(1,1)$ with $e_0^2 = 1$, $e_1^2 = -1$, the vectors
\[
v = \tfrac12(e_0 + e_1), \qquad v^* = \tfrac12(e_0 - e_1)
\]
satisfy $v^2 = (v^*)^2 = 0, \qquad v v^* + v^* v = 1$ exactly. This block, the \emph{hyperbolic atom}, is the unique unit of transport that, at the level of elements, needs no ascent of its own \cite[Rem.~4.5]{ButlerT1}.
\end{lemma}
\begin{proof}
$v v^* v = v(1 - v v^*)^{\vphantom{*}} = v - v^2 v^* = v$ using $v^2 = (v^*)^2 = 0, \qquad v v^* + v^* v = 1$ and $v^2 = 0$; mirror for $v^*$. In the model: $v^2 = \tfrac14(e_0^2 + e_1^2) = \tfrac14(1 - 1) = 0$ (the cross terms cancel by anticommutation, \ref{c:std-algebra.orthogonal-anticommute}), likewise $(v^*)^2 = 0$, and, by \ref{c:std-algebra.vector-squares}, $v v^* + v^* v = 2\,(v \cdot v^*) = \tfrac12\bigl(e_0^2 - e_1^2\bigr) = \tfrac12(1 + 1) = 1$.

\end{proof}

\phantomsection\label{c:note.singularity.setting-step-3-apply-rule-2}
\emph{Where the ascent lands.} $\Cl(2,4)$ is not an anonymous algebra. It is the conformal geometric algebra of Minkowski space, obtained by adjoining the null pair $(\nO, \nInf)$ to the spacetime algebra; it is the algebra of the conformal group $\mathrm{SO}(2,4)$ of compactified Minkowski space; the spinor module of its even part is $\C^4$, Penrose's twistor space \cite{Penrose1967}; and $\mathrm{SO}(2,4)$ is the isometry group underlying $\mathrm{AdS}_5$, hence the symmetry algebra of the best-studied holographic correspondence. The mechanical rule, applied to the unique locus in classical physics where division by zero is unavoidable, selects the algebra Penrose has argued for five decades is the correct language for precisely this regime. There is independent physical evidence pointing the same way: near a generic singularity every mass scale becomes negligible against the curvature scale and the dynamics is asymptotically conformally invariant, the essential observation of BKL analysis and of conformal cyclic cosmology. The metric degenerates at $r = 0$; the light cones survive; and $\Cl(2,4)$ is the light cones' algebra.

\section{Step 4a: compute in the completion; the transport instantiated, and the operators forced}

\subsection*{The point-level transport, instantiated}
\phantomsection\label{c:note.singularity.setting-point-level-transport-instantiated}
The algorithm's Step 4 is executed on points, not elements \cite[\S 8.1, Step 4]{ButlerT1}. This subsection carries it out for the Schwarzschild fibre, so that the operator content of the rest of this section rests on a computed instance rather than on the imported theorem alone.

\begin{definition}[The fibrewise conformal lift]\label{c:obj.singularity.fibrewise-conformal-lift}
Let $p$ lie in the invertible region, so that the tangent space $(T_pM, g_p)$ is nondegenerate of signature $(1,3)$, and write $Q_p(x) := g_p(x,x)$.
Adjoin to $T_pM$ the hyperbolic plane $H = \operatorname{span}(\nO, \nInf)$ of Rule 1, with $\nO^2 = \nInf^2 = 0$ and $\nO \cdot \nInf = -1$ (the normalisation of T1 \cite{ButlerT1}, under which the lift below is null: $X_p(x)^{2} = Q_p(x)\,(1 + \nO\cdot\nInf) = 0$), giving the completion $\Cl(2,4)$ of the section ``Step 3: apply the rule''. The \emph{lift} of $x \in T_pM$ is the null ray
\begin{equation}
X_p(x) \;:=\; \bigl[\, \nO + x + \tfrac12 Q_p(x)\, \nInf \,\bigr] \;\in\; \mathcal{Q}_p ,
\end{equation}
and in the coordinates $N = s\,\nO + x + t\,\nInf$, written $(s,x,t)$, the quadric is cut out by the incidence condition $Q_p(x) = 2st$, the transported reciprocal is the linear involution $\iota(s,x,t) = (2t,\, x,\, \tfrac12 s)$, and the read-out is $\dsc(s,x,t) = x/s$, undefined where $s$ vanishes \cite[Thm.~4.10, Prop.~4.21, Lem.~4.23]{ButlerT1}. For a covector $\xi$ at an invertible point, the lift is that of its metric dual, $X_p(\xi) := X_p(\xi^{\sharp})$, computed directly from the dual norm since $Q_p(\xi^{\sharp}) = Q^{*}_p(\xi)$; clause (4) of Proposition \ref{c:singularity.transport-schwarzschild-fibre} applies it in this form.
\end{definition}

\begin{proposition}[The transport, on the Schwarzschild fibre]\label{c:singularity.transport-schwarzschild-fibre}
\changedmark{gray}{reworded}
Fix $p$ in the invertible region of the Schwarzschild gravitational field $(\hb, \Omega)$ on flat spacetime (Definition \emph{Gauge Theory Gravity: the gravitational field on flat spacetime}), take $T_pM$ to be the background vectors at $p$, and write $Q_p(x) := Q(\hb^{-1}(x))$ for the square of the physical vector over $x$ and $Q^*_p(\xi) := Q(\hb^{T}(\xi))$ for the dual norm of a covector. Vectors and covectors are read in the physical space: the gauge field carries the flat background to the physical vectors, a vector $x$ to $\hb^{-1}(x)$ and a covector $\xi$ to $\hb^{T}(\xi)$, the lift $X_p(x)$ carries $\hb^{-1}(x)$ in its vector slot, and the values in (2) to (4) are taken through this identification ($x/Q_p(x)$ is the physical vector $\hb^{-1}(x)/Q_p(x)$, $x + \tfrac12 \nInf$ is $\hb^{-1}(x) + \tfrac12 \nInf$, and $\xi/r$ is $\hb^{T}(\xi)/r$). In the static gauge the Gram of $Q_p$ is the derived metric $g_p = \operatorname{diag}(f, -f^{-1}, -r^2, -r^2\sin^2\theta)$ with $f = 1 - 2M/r$. Then:
\begin{enumerate}
\item[(1)] $X_p(x)$ satisfies the incidence condition for every $x \in T_pM$, and $\iota$ preserves it and is an involution.
\item[(2)] \emph{(Invertible directions: the transport returns to the original space.)} If $Q_p(x) \neq 0$ then $\dsc\bigl(\iota(X_p(x))\bigr) = x/Q_p(x)$, the classical reciprocal.
\item[(3)] \emph{(Null directions: the transport does not return.)} If $Q_p(x) = 0$, that is if $x$ is physically null, in particular for the radial null direction $x = \et + f\,\er$, then $\iota(X_p(x)) = [\,x + \tfrac12 \nInf\,]$, an ideal point, and $\dsc$ is undefined on it.
\item[(4)] \emph{(The dying direction, and where the descent fails.)} Let $\xi := k = \sqrt r\,dr + \sqrt{2M}\,dt_{\mathrm{PG}}$ be the generator of the radical line of the degeneration in the weighted Painlev\'e-Gullstrand coframe of \cite[Thm.~8.7]{ButlerT4}, in the static chart $\xi = \sqrt{2M}\,dt_s + (\sqrt r/f)\,dr$. Its dual norm is $Q^*_p(\xi) = -r$ exactly, at every radius, the horizon included: nonzero for every $r > 0$ and tending to $0$ as $r \to 0^+$. Consequently $\iota(X_p(\xi))$ descends at every point $r > 0$, to $\dsc = -\xi/r$, and converges, as $r \to 0^+$, to an ideal point at which $\dsc$ is undefined. The read-out fails exactly on $E$ and nowhere else in the physical closure.
\end{enumerate}
\end{proposition}
\begin{proof}
Here $Q_p(x) = Q(\hb^{-1}(x))$ and $Q^*_p(\xi) = Q(\hb^{T}(\xi))$, the derived metric and cometric of the gravitational field at $p$, whose static-gauge Grams are those displayed (Fact \emph{Imported: the dictionary between Gauge Theory Gravity and curved-spacetime language}(i)). (1) With $X_p(x) = (1, x, \tfrac12 Q_p(x))$ (the triples of \ref{c:transport.swap-rescale-normal-form}) the incidence condition reads $Q_p(x) = 2 \cdot 1 \cdot \tfrac12 Q_p(x)$, an identity. Applying $\iota$ gives $(Q_p(x), x, \tfrac12)$, whose incidence check is $2 \cdot Q_p(x) \cdot \tfrac12 = Q_p(x)$, again an identity; and $\iota^2(s,x,t) = \iota(2t, x, \tfrac12 s) = (s, x, t)$.

(2) $\iota(X_p(x)) = (Q_p(x), x, \tfrac12)$, so $\dsc$ returns $x/Q_p(x)$, which is the classical reciprocal $\rho(x)$ of \cite[Lem.~4.23(2)]{ButlerT1} (\ref{c:transport.descent-read-out-map}).

(3) For the stated $x$, $Q_p(x) = f \cdot 1 - f^{-1} \cdot f^2 = f - f = 0$, so $X_p(x) = (1, x, 0)$ and $\iota(X_p(x)) = (0, x, \tfrac12)$: the first slot vanishes, which is the defining condition of an ideal point, and $\dsc$ requires division by it. The value is $[\,x + \tfrac12 \nInf\,]$, exactly the ideal value of \cite[Thm.~4.10(3)]{ButlerT1} (\ref{c:transport.totality}).

(4) With $t_{\mathrm{PG}} = t_s + \int \beta f^{-1}dr$ and $\beta^{2} = 2M/r$, $dt_{\mathrm{PG}} = dt_s + (\beta/f)dr$, so $k = \sqrt{2M}\,dt_s + \sqrt r\,(1 + 2M/rf)\,dr = \sqrt{2M}\,dt_s + (\sqrt r/f)\,dr$, using $f + 2M/r = 1$. In the static cometric $(g_p^{-1})^{tt} = 1/f$, $(g_p^{-1})^{rr} = -f$, so $Q^*_p(\xi) = 2M/f - f\,(r/f^{2}) = (2M - r)/f = -r$, an identity of rational functions in $r$ with no pole at the horizon; in the Painlev\'e-Gullstrand frame of \cite[Thm.~8.7]{ButlerT4} (\ref{c:sheaf.schwarzschild-degeneration-boundary-tame}) the same value is read off the polynomial Gram matrix as $g_k = (m^{2} - r) - 2m^{2} + m^{2} = -r$ with $m = \sqrt{2M}$, which is that theorem's clause (3). Hence $\iota(X_p(\xi)) = (-r, \xi, \tfrac12)$ has nonvanishing first slot at every $r > 0$, where $\dsc$ returns $-\xi/r$, and first slot tending to $0$ as $r \to 0^+$, which is the ideal stratum. (The Killing-time covector $dt_s$ itself, with $Q^*_p(dt_s) = 1/f$, spans the radical line only on $E$ \cite[Thm.~8.7(4)]{ButlerT4}, and its static norm has a pole at the horizon where nothing degenerates; $k$ is the regular generator.) One further line closes the junction with Step 1, computed in the Painlev\'e-Gullstrand chart where Propositions \ref{c:singularity.cones-fold} and \ref{c:singularity.degeneration-exactly} fold the cones: the PG cometric has $g^{tt} = 1$, $g^{tr} = -\beta$, $g^{rr} = \beta^{2} - 1$, so for $k = \sqrt{2M}\,dt_{\mathrm{PG}} + \sqrt r\,dr$ the dual's $t$-component is $\sqrt{2M} - \beta\sqrt r = 0$ identically and its $r$-component is $-\beta\sqrt{2M} + (\beta^{2} - 1)\sqrt r = -\sqrt r$: $k^{\sharp} = -\sqrt r\,\er$ exactly, at every radius and with no horizon pole, in the same frame as Step 1, so the corank-one radical of the completion and Step 1's dying direction $\er$ are the same direction outright, not merely in the limit. That the failure locus is exactly $E$, and that no shadow wall forms elsewhere, is \cite[Thm.~8.12]{ButlerT4} (\ref{c:sheaf.one-sided-completion-singularity}). All identities in (1) to (4) are machine-verified in exact symbolic arithmetic (the section ``Machine verification'', parts L1 to L14), and the nullity of the lift, $X_p(x)^{2} = 0$ in $\Cl(2,4)$ with the stated normalisation, blade-exactly.

\end{proof}

\begin{remark}[What this instance shows, and what it does not]\label{c:note.singularity.what-instance-shows-what-it-does}
The computation is the algorithm's own diagnosis, made concrete: in the completion the division never fails, and what fails is the return to the original space. The answer to ``divide by the dying direction'' exists as a point of $\mathcal{Q}_p$ and is not a vector of spacetime, so no assignment of a classical partner value can be read off it. That is the point-level counterpart of the element-level no-go recorded next as Fact \ref{c:singularity.imported-no-c-number-repair}, and the two together are what force the repair to be delivered as operators: the value has left the space, and the algebra keeps only the ladder pair $v^2 = (v^*)^2 = 0, \qquad v v^* + v^* v = 1$ that records it.
\end{remark}

\begin{assembled}[Imported: no c-number repair]\label{c:singularity.imported-no-c-number-repair}
No equivariant assignment of classical partner values to the dying directions exists, in the minimal ascent or in any enlargement. Mechanism: the stabiliser of the null direction $[v]$ contains the parabolic null rotations, the candidate partners form an affine space on which those rotations act by translation, and a translation of an affine space has no fixed point; cohomologically, the obstruction is the tautological class of the unipotent stabiliser, nonzero in every mixed-signature dimension above two. \emph{Provenance:} the element-level no-go is \cite[Thm.~4.2]{ButlerT1}; the totality of the transported division is \cite[Thm.~4.10]{ButlerT1}; the cohomological form, with the torsor made exact modulo the left ideal and the class identified as $\theta \in H^1(N_v; U) = \operatorname{Hom}(N_v, U)$, is \cite[Thm.~4.8]{ButlerT1}. The dimension-two exception, where the class vanishes, is exactly Lemma \ref{c:singularity.pseudo-inverse-identities-hyperbolic-atom}'s self-completeness. \emph{Scope:} these theorems are proved where the ambient form is nondegenerate and the dying direction has a hyperbolic partner, which is the situation at every point of the invertible region, whose limits toward $E$ the no-go therefore governs; at $E$ itself the partner is exactly what is missing (Proposition \ref{c:singularity.type-n-only-type-n}), so the import does not apply there verbatim. The boundary case is supplied by the no-value theorem at the wall, \cite[Thm.~10.4]{ButlerT4}, reproduced verbatim as Fact \ref{c:singularity.imported-no-equivariant-value-wall}: no equivariant classical value exists at $E$ either, the completion's own transversal-independence carrying the equivariance. The claim that no classical value can repair the dying directions therefore holds on the approach by \cite[Thm.~4.2]{ButlerT1} and at the boundary by Fact \ref{c:singularity.imported-no-equivariant-value-wall}.
\end{assembled}
\begin{proof}
No equivariant assignment of classical partner values to the dying directions exists, in the minimal ascent or in any enlargement. Mechanism: the stabiliser of the null direction $[v]$ contains the parabolic null rotations, the candidate partners form an affine space on which those rotations act by translation, and a translation of an affine space has no fixed point; cohomologically, the obstruction is the tautological class of the unipotent stabiliser, nonzero in every mixed-signature dimension above two. \emph{Provenance:} the element-level no-go is \cite[Thm.~4.2]{ButlerT1} (\ref{c:transport.no-go-level-elements}); the totality of the transported division is \cite[Thm.~4.10]{ButlerT1}; the cohomological form, with the torsor made exact modulo the left ideal and the class identified as $\theta \in H^1(N_v; U) = \operatorname{Hom}(N_v, U)$, is \cite[Thm.~4.8]{ButlerT1} (\ref{c:transport.cohomological-form-no-go}). The dimension-two exception, where the class vanishes, is exactly Lemma \ref{c:singularity.pseudo-inverse-identities-hyperbolic-atom}'s self-completeness. \emph{Scope:} these theorems are proved where the ambient form is nondegenerate and the dying direction has a hyperbolic partner, which is the situation at every point of the invertible region, whose limits toward $E$ the no-go therefore governs; at $E$ itself the partner is exactly what is missing (Proposition \ref{c:singularity.type-n-only-type-n}), so the import does not apply there verbatim. The boundary case is supplied by the no-value theorem at the wall, \cite[Thm.~10.4]{ButlerT4} (\ref{c:sheaf.no-equivariant-value-wall}), reproduced verbatim as Fact \ref{c:singularity.imported-no-equivariant-value-wall}: no equivariant classical value exists at $E$ either, the completion's own transversal-independence carrying the equivariance. The claim that no classical value can repair the dying directions therefore holds on the approach by \cite[Thm.~4.2]{ButlerT1} and at the boundary by Fact \ref{c:singularity.imported-no-equivariant-value-wall}.

\end{proof}

\phantomsection\label{c:note.singularity.setting-point-level-transport-instantiated-2}
What Rule 1 installs instead of a value is the pair $v^2 = (v^*)^2 = 0, \qquad v v^* + v^* v = 1$ itself: the canonical anticommutation relations of one fermionic mode. Nothing commuting satisfies these relations; objects obeying them are irreducibly operators on a two-state space. By the localisation of \cite[Cor.~4.5]{ButlerT4}, in the form the Schwarzschild composite gives it at the black hole \cite[Thm.~8.12]{ButlerT4} (reproduced verbatim, hypotheses included, as Facts \ref{c:sheaf.schwarzschild-degeneration-boundary-tame} and \ref{c:singularity.imported-one-sided-completion-singularity} in the section ``The imported localisation theorems, verbatim''), exactly one such mode is at each point of the blown-up singular locus $E \cong \R \times S^2$, and nowhere else in the physical closure: the composite's corank-one theorem \cite[Thm.~8.7]{ButlerT4} makes the count one pair per point, not one per dying coordinate direction. Quantisation is demanded locally, exactly where the classical description divides by zero, and is output rather than input: no quantisation procedure is imposed anywhere; the strength of ``demanded'' is uniform: barred on the approach by Fact \ref{c:singularity.imported-no-c-number-repair} and at the boundary itself by Fact \ref{c:singularity.imported-no-equivariant-value-wall}.

\begin{proposition}[The descent, and what it erases]\label{c:singularity.descent-what-it-erases}
Set $P := v^* v$. Then $P^2 = P$, $1 - P = v v^*$, and in the defining $2 \times 2$ representation ($v = E_{12}$, $v^* = E_{21}$),
\[
P = \begin{pmatrix} 0 & 0 \\ 0 & 1 \end{pmatrix}, \qquad
E(x) := P x P + (1{-}P)\,x\,(1{-}P), \qquad
E\begin{pmatrix} a & b \\ c & d \end{pmatrix} = \begin{pmatrix} a & 0 \\ 0 & d \end{pmatrix} .
\]
Descent retains the occupation data and erases the coherences $b, c$.
\end{proposition}
\begin{proof}
$P^2 = v^*(v v^*)v = v^*(1 - v^* v)v = v^* v - (v^*)^2 v^2 = P$, using $v^2 = (v^*)^2 = 0, \qquad v v^* + v^* v = 1$ (the CAR of a Witt pair, \ref{c:std-algebra.witt-basis-car}) twice and $v^2 = (v^*)^2 = 0$; then $1 - P = vv^*$ is $v^2 = (v^*)^2 = 0, \qquad v v^* + v^* v = 1$ itself. In the defining representation (\ref{c:std-algebra.cl11-matrix}) $v = E_{12}$, $v^*= E_{21}$ satisfy $v^2 = (v^*)^2 = 0, \qquad v v^* + v^* v = 1$ by matrix multiplication, $P = E_{22}$, and the displayed action of $E$ is the block computation $PxP + (1{-}P)x(1{-}P) = \operatorname{diag}(a, d)$.

\end{proof}

\begin{definition}[The grading, and the parity automorphism]\label{c:obj.singularity.grading-parity-automorphism}
Write $\Aa$ for the algebra generated by the completed pair, $\Aa \cong \Mat_2(\C)$ as a factor. It carries the $\mathbb{Z}_2$ grading it inherits from the Clifford algebra: $\Aev := \operatorname{span}\{1, \sigma\}$ is the even part and $\Aod := \operatorname{span}\{v, v^{*}\}$ the odd part, where $\sigma := v^{*}v - vv^{*}$ is the bivector of Theorem \ref{c:singularity.modular-structure-completed-pair}. The \emph{parity automorphism} is $\alpha(x) := \sigma x \sigma$, well defined because $\sigma^{2} = 1$; it fixes $\Aev$ pointwise and acts by $-1$ on $\Aod$. On the full boundary mode space the parity automorphism is $\operatorname{Ad}(\Gamma)$ with $\Gamma := \prod_k \sigma_k$, the total fermion parity; the mode-wise refinement, one parity automorphism $\operatorname{Ad}(\sigma_k)$ per point of $E$, is Proposition \ref{c:singularity.p-reduced-metric-level-reading-many}'s task, and on two or more modes the two generate different averages.
\end{definition}

\begin{theorem}[The descent is the parity average]\label{c:singularity.descent-parity-average}
For every $x \in \Aa$,
\[
E(x) \;=\; \tfrac12\bigl(x + \alpha(x)\bigr) \;=\; \tfrac12\bigl(x + \sigma x \sigma\bigr) .
\]
Consequently $E$ is the group average over $\{\id, \alpha\}$, hence the unique conditional expectation of $\Aa$ onto its fixed-point algebra $\Aev$; its range is $\Aev$ and its kernel is exactly $\Aod$, the adjoined ladder sector; and $E$ preserves every state $\omega$ with $\omega \circ \alpha = \omega$.
\end{theorem}
\begin{proof}
In the defining representation $v = E_{12}$, $v^{*} = E_{21}$, so $P = E_{22}$ and $\sigma = 2P - 1 = \operatorname{diag}(-1, +1)$; then
\[
\alpha\begin{pmatrix} a & b \\ c & d \end{pmatrix}
= \begin{pmatrix} -1 & 0 \\ 0 & 1\end{pmatrix}\begin{pmatrix} a & b \\ c & d \end{pmatrix}\begin{pmatrix} -1 & 0 \\ 0 & 1\end{pmatrix}
= \begin{pmatrix} a & -b \\ -c & d \end{pmatrix},
\]
so $\tfrac12(x + \alpha(x)) = \operatorname{diag}(a,d)$, which is the value of $E$ computed in Proposition \ref{c:singularity.descent-what-it-erases}. The average of a finite group acting by automorphisms is a conditional expectation onto the fixed-point algebra. It is moreover the unique conditional expectation of $\Aa$ onto $\Aev$, by bimodularity: if $F$ is any conditional expectation onto $\Aev$, then for odd $x$, since $\sigma \in \Aev$ and $\sigma x \sigma = -x$, bimodularity gives $\sigma F(x) \sigma = F(\sigma x \sigma) = -F(x)$, while $F(x) \in \Aev$ is fixed by conjugation with $\sigma$, so $F(x) = -F(x) = 0$: every conditional expectation onto $\Aev$ annihilates the $(-1)$-eigenspace of $\alpha$, which is $\Aod = \operatorname{span}\{v, v^{*}\}$, and agrees with the identity on $\Aev$, hence equals $E$ on $\Aa = \Aev \oplus \Aod$. State preservation is $\omega(E(x)) = \tfrac12(\omega(x) + \omega(\alpha(x))) = \omega(x)$. The script verifies the identity, idempotence, the conditional-expectation property $E(axb) = aE(x)b$ for even $a, b$, and the range and kernel, on generic rational data.

\end{proof}

\begin{remark}[The metric level is the parity even level]\label{c:note.singularity.metric-level-parity-even-level}
Theorem \ref{c:singularity.descent-parity-average} sharpens the chain's slogan that descent keeps the records and loses the coherences. What descent keeps is the even subalgebra, and the sector the completion installs is exactly the odd one (its ladder pair mixes the physical and adjoined planes: $k^{*} = \nO/2\nu$ against the dying direction $k$ itself, Fact \ref{c:singularity.imported-one-sided-completion-singularity}; this $\nu$ is T04's coupling constant, not Step 6a's cell density): the two-level split is the Clifford algebra's own $\mathbb{Z}_2$ grading, not an extra structure imposed on it. This is a single-mode statement: on the full boundary mode space the metric-level algebra the chain uses is the joint-diagonal one, the fixed points of one parity per point of $E$, strictly smaller than the total-parity-even algebra, and its licence is Proposition \ref{c:singularity.p-reduced-metric-level-reading-many}, not the grading alone. The one-way character recorded next as Fact \ref{c:singularity.imported-descent-one-way-conditional-expectation} is then the statement that a $\mathbb{Z}_2$ average cannot be undone, and the uniqueness clause there is strengthened rather than merely imported: Fact \ref{c:singularity.imported-descent-one-way-conditional-expectation}'s Takesaki clause is uniqueness among trace-preserving conditional expectations, while Theorem \ref{c:singularity.descent-parity-average} makes $E$ the unique conditional expectation onto $\Aev$ outright, trace-preserving or not; and the grading is not a choice. Step 4d turns this identification into the state-selection theorem of the section ``What selects the state''.
\end{remark}

\begin{fact}[Imported: the descent is a one-way conditional expectation (external ingredients)]\label{c:singularity.imported-descent-one-way-conditional-expectation-external}
The results from outside the series on which \ref{c:singularity.imported-descent-one-way-conditional-expectation} draws, as cited there: M. Takesaki, Conditional expectations in von Neumann algebras, J. Funct. Anal. 9 (1972), 306--321.
\emph{Source:} M. Takesaki, Conditional expectations in von Neumann algebras, J. Funct. Anal. 9 (1972), 306--321..
\end{fact}

\begin{assembled}[Imported: the descent is a one-way conditional expectation]\label{c:singularity.imported-descent-one-way-conditional-expectation}
\changedmark{gray}{reworded}
$E$ is a conditional expectation onto the commutant of $P$, its kernel is exactly the adjoined ladder sector, and no unital algebra homomorphism from the completed algebra back to its base exists at all: ascent is two-way for what was lifted from the original space, one-way for what was built in the completion. It is moreover the unique conditional expectation onto that subalgebra preserving the trace. \emph{Provenance:} the conditional expectation property, the kernel and the non-existence of a return homomorphism are \cite[Prop.~7.5, Thm.~7.6]{ButlerT1}; the uniqueness of the trace-preserving conditional expectation onto a von Neumann subalgebra is Takesaki's theorem \cite{Takesaki1972}, and in the present order-two case it is also Theorem \ref{c:singularity.descent-parity-average}.
\end{assembled}
\begin{proof}
$E$ is a conditional expectation onto the commutant of $P$, its kernel is exactly the adjoined ladder sector, and no unital algebra homomorphism from the completed algebra back to its base exists at all: ascent is two-way for what was lifted from the original space, one-way for what was built in the completion. It is moreover the unique conditional expectation onto that subalgebra preserving the trace. \emph{Provenance:} the conditional-expectation property, the kernel and the non-existence of a return homomorphism are \cite[Prop.~7.5, Thm.~7.6]{ButlerT1} (\ref{c:transport.pinching-kernel}, \ref{c:transport.door-one-way}); the uniqueness of the trace-preserving conditional expectation onto a von Neumann subalgebra is Takesaki's theorem \cite{Takesaki1972}, and in the present order-two case it is also Theorem \ref{c:singularity.descent-parity-average}.

This restates \ref{c:singularity.imported-descent-one-way-conditional-expectation-external}.

\end{proof}

\section{Steps 4b and 4c: compute in the completion; the plane, the rate, and the boost as number operator}

\phantomsection\label{c:note.singularity.setting-steps-4b-4c-compute-completion}
Step 4b is stated first in its gauge-theoretic form: the curvature of the gauge field factorises exactly, at every radius, into a divergent scalar times a fixed shape operator, and the shape operator is what selects the plane. The appendix's Cartan computation is an independent cross-check in the metric picture.

\begin{theorem}[The curvature of the gauge field: exact shape decoupling]\label{c:singularity.curvature-gauge-field-exact-shape-decoupling}
\changedmark{gray}{reworded}
Let $\sigma_r$ be the unit bivector of the radial boost plane, and let $\Pi$ denote the projection of the six-dimensional bivector space onto $\operatorname{span}\{\sigma_r, I\sigma_r\}$, with $I$ the pseudoscalar. The radial plane is $\hb$-invariant,
\[
\hb(\et) \wedge \hb(\er) \;=\; (\et - \beta\er) \wedge \er \;=\; \et \wedge \er ,
\]
so the coordinate and physical radial planes coincide and $\sigma_r$ needs no chart to define. Then, exactly and at every radius,
\[
\mathcal{R}(B) \;=\; -\,\frac{M}{r^3}\,S(B), \qquad
S \;:=\; 1 - 3\Pi \;=\; -\tfrac12\bigl(B \mapsto B + 3\,\sigma_r B \sigma_r\bigr),
\]
in the sign convention of Proposition \ref{c:singularity.solution-where-divergence-lies}(1) and of the Cartan computation of the section ``The Schwarzschild curvature and the surface gravity: the metric-picture cross-check'', which agree: $\mathcal{R}(\sigma_r) = +2M/r^{3}\,\sigma_r$ while $S(\sigma_r) = -2\sigma_r$ (only the eigenvalue ratio, the trace, and the identity of the distinguished plane pair enter the chain, and each is convention-free). The shape operator satisfies
\[
S^2 + S = 2 ,
\]
has spectrum $(-2, -2, 1, 1, 1, 1)$ with the doubled eigenvalue attained exactly on $\{\sigma_r, I\sigma_r\}$, and is trace-free on $\Lambda^2$, though that verifies an algebraic identity rather than field equations: $\tr(1 - 3\Pi) = 6 - 6 = 0$ holds automatically for any rank-two $\Pi$, and trace-freeness on $\Lambda^2$ is the vanishing of the scalar curvature, not Ricci flatness, which is carried by the solution of the section ``Step 1: locate the failure'' and not by this operator identity. The entire approach to the singularity is carried by the single scalar $M/r^3$.
\end{theorem}
\begin{proof}
Three parts, each machine-checked. \emph{The algebra of $S$.} For $B$ in the plane pair: $\sigma_r\,\sigma_r\,\sigma_r = \sigma_r$ since $\sigma_r^2 = 1$, and $\sigma_r (I\sigma_r)\sigma_r = I\sigma_r^3 = I\sigma_r$ since $I$ commutes with every even element (\ref{c:std-algebra.pseudoscalar-commutation}); so $-\tfrac12(B + 3\sigma_r B\sigma_r) = -2B$ there. For each of the four mixed planes, $B$ shares exactly one index with the radial plane, hence anticommutes with $\sigma_r$ (\ref{c:std-algebra.orthogonal-anticommute}): $\sigma_r B \sigma_r = -B\sigma_r^2 = -B$, giving $+B$. This is $1 - 3\Pi$; then $S^2 = 1 + 3\Pi$ (idempotency of $\Pi$) gives $S^2 + S = 2$, the spectrum is read off, and the trace is $2\cdot(-2) + 4\cdot 1 = 0$. \emph{The eigenvalues.} On the radial plane the eigenvalue comes from the gauge field alone: Proposition \ref{c:singularity.solution-where-divergence-lies}(1) computed, in two lines and with the commutator term vanishing identically,
\[
\mathcal{R}(\partial_t \wedge \partial_r) \;=\; -\,\partial_r\Bigl(\frac{M}{r^2}\Bigr)\,\sigma_r \;=\; \frac{2M}{r^3}\,\sigma_r .
\]
The mixed and angular planes are computed in the section ``The curvature operator, by the Cartan structure equations'' by the Cartan structure equations in the static chart, giving $-M/r^3$ on the four mixed planes and the doubled value again on the angular dual; curvature is a pointwise invariant, and the two pictures agree on the invertible region by Fact \ref{c:singularity.imported-equivalence-einstein-equations}. The eigenvalue ratio is $-2 : 1$, which is the ratio of $S$'s spectrum. \emph{Assembly.} An operator diagonal on the six planes with eigenvalues proportional to $(-2,-2,1,1,1,1)$ is $\pm(M/r^3)(1 - 3\Pi)$, and the sign is fixed by the radial plane: $\mathcal{R}(\sigma_r) = +2M/r^{3}\,\sigma_r$ against $S(\sigma_r) = -2\sigma_r$ gives the minus sign, and the mixed planes, $-M/r^{3}$ against $S = +1$, confirm it. The $\hb$-invariance display is one line of anticommutation ($\er \wedge \er = 0$). The script checks the conjugation table, the operator identities, the spectrum and trace, the eigenvalue ratio against the appendix's rational-function values, the sign of the identity on the radial and mixed planes, and the invariance display, blade-exactly.

\end{proof}

\phantomsection\label{c:note.singularity.setting-steps-4b-4c-compute-completion-2}
In classical language the spectrum is the Petrov type-D structure of Schwarzschild, $\Psi_2 = -M/r^{3}$ its one nonvanishing Weyl scalar \cite{Chandrasekhar1983}; the theorem's contribution is the derivation inside the gauge formalism and the exact operator identity $S^{2} + S = 2$, not the spectrum itself.

\begin{proposition}[The plane and the rate]\label{c:singularity.plane-rate}
The plane in which the completed pair lies is not chosen: the doubled eigenvalue of Theorem \ref{c:singularity.curvature-gauge-field-exact-shape-decoupling} singles out the pair $\{\sigma_r, I\sigma_r\}$ invariantly at every radius, and its boost member $\sigma_r$, the member with square $+1$ up to the orientation sign, is at the horizon the boost plane of the Killing flow, with surface gravity
\[
\kappa = \frac{1}{4M} .
\]
The geometric normalisation this fixes: an interval of $2\pi$ in rapidity corresponds to $2\pi/\kappa = 8\pi M$ of Schwarzschild time. This single conversion factor is where the numerical value of the temperature enters the chain.
\end{proposition}
\begin{proof}
The selection is Theorem \ref{c:singularity.curvature-gauge-field-exact-shape-decoupling}; $(I\sigma_r)^2 = -1$ while $\sigma_r^2 = +1$ separates the pair's two members invariantly. The surface gravity is computed as the redshifted proper acceleration in the section ``The surface gravity'' (\ref{c:singularity.computed}): $\kappa = \tfrac12 f'(r_h) = 1/4M$. The broader geometric development (the parallel ideal point, the weighted limit, the Kerr tests) remains the task of the series' geometry paper T6 \cite{ButlerT6}; nothing below uses more than what Theorem \ref{c:singularity.curvature-gauge-field-exact-shape-decoupling} and the appendix prove.

\end{proof}

\begin{theorem}[Modular structure of the completed pair]\label{c:singularity.modular-structure-completed-pair}
Let $(v, v^*)$ satisfy $v^2 = (v^*)^2 = 0, \qquad v v^* + v^* v = 1$ and define $\sigma := v^* v - v v^*$. Then:
\begin{enumerate}
\item $\sigma^2 = 1$; and in the hyperbolic atom of Lemma \ref{c:singularity.pseudo-inverse-identities-hyperbolic-atom}, $\sigma = e_0 e_1$ exactly, the unit boost bivector of the plane;
\item the rotor flow $R_\lambda = \exp(\lambda\sigma/2)$ acts by $v \mapsto e^{-\lambda} v$ and $v^* \mapsto e^{+\lambda} v^*$: the pair diagonalises the boost with opposite weights, the structure of outgoing and ingoing horizon modes;
\item the spectral projectors $\tfrac12(1 \pm \sigma)$ of the boost are identically the occupation projectors $v^* v$ and $v v^*$: the descent projectors of Proposition \ref{c:singularity.descent-what-it-erases};
\item in the fermionic representation $\sigma = 2N - 1$: the generator of the geometric boost is, up to normalisation, the number operator, $K := \tfrac{\sigma}{2} = N - \tfrac12$;
\item $\mathrm{Spin}^+(1,1)$ is one-dimensional, so the boost flow is the \emph{only} one-parameter flow implemented by real rotors of the atom. The scope matters and is stated: the mode algebra $\Mat_2(\C)$ admits a three-parameter family of one-parameter unitary groups, and the uniqueness claimed is within the class the geometry supplies, real rotor flows of $\Cl(1,1)$; Proposition \ref{c:singularity.euclidean-rotor-where-it-lies} later leaves that class deliberately, its Euclidean rotor being complex.
\end{enumerate}
One object, three readings: boost generator, number operator, generator of the descent projectors.
\end{theorem}
\begin{proof}
(1) $\sigma^2 = (v^*v)^2 + (vv^*)^2 - v^*v\,vv^* - vv^*\,v^*v = v^*v + vv^* - 0 - 0 = 1$, using idempotency (Proposition \ref{c:singularity.descent-what-it-erases}) and $v^2 = (v^*)^2 = 0$ to remove the cross terms. In the atom (\ref{c:singularity.pseudo-inverse-identities-hyperbolic-atom}): $v v^* = \tfrac14(e_0 + e_1)(e_0 - e_1) = \tfrac14(1 - e_0e_1 + e_1e_0 + 1) = \tfrac12(1 - e_0e_1)$, so $v^*v = 1 - vv^* = \tfrac12(1 + e_0e_1)$ and $\sigma = e_0 e_1$. (2) From $\sigma v = -v$ and $v\sigma = v$ (in the atom: $e_0e_1(e_0 + e_1) = -(e_1 + e_0)$ and $(e_0+e_1)e_0e_1 = e_1 + e_0$; abstractly from $v^2 = (v^*)^2 = 0, \qquad v v^* + v^* v = 1$), and $\sigma^2 = 1$ giving $e^{-\lambda\sigma} = \cosh\lambda - \sigma\sinh\lambda$:
\[
e^{\lambda\sigma/2}\,v\,e^{-\lambda\sigma/2} = v\,e^{-\lambda\sigma} = v\cosh\lambda - v\sigma\,\sinh\lambda = (\cosh\lambda - \sinh\lambda)\,v = e^{-\lambda}v ,
\]
using $v\sigma = v$ once; mirror-symmetrically for $v^*$. (3) is $v^*v = \tfrac12(1 + \sigma)$. (4) $N = v^*v$ counts occupation in the two-state module, and $\sigma = 2v^*v - 1$. (5) The rotor group of $\Cl(1,1)$ is generated by its single unit bivector.

\end{proof}

\begin{corollary}[The matching: Steps 4b and 4c join by equation]\label{c:singularity.matching-steps-4b-4c-join-equation}
\changedmark{BrickRed}{corrected since the deposit}
Realise the hyperbolic atom of Lemma \ref{c:singularity.pseudo-inverse-identities-hyperbolic-atom} on the radial plane, $e_0 = \hat e_t$, $e_1 = \hat e_r$. Where this realisation lies must be stated, because it is available exactly where the ascent is not needed: at every $r > 0$ the physical radial plane has signature $(1,1)$ and carries a CAR pair inside $\Cl(1,3)$ with no ascent at all, but at $r = 0$ the radial plane itself degenerates, its radical the dying direction $k$ (corank one in the forced weighting, Remark \ref{c:note.singularity.corank-limit-weighted-answer} and Fact \ref{c:sheaf.schwarzschild-degeneration-boundary-tame}(2); Proposition \ref{c:singularity.degeneration-exactly}'s uniform rescaling degenerates the whole form), so no CAR pair supported on the radial plane survives there, even though the weighted module still contains hyperbolic planes (Fact \ref{c:sheaf.schwarzschild-degeneration-boundary-tame}(2): $V_1$ carries $\mathrm{diag}(m^2, -1, -1)$). On $E$ the atom is carried by $\operatorname{span}(k, \nO/2\nu)$, one leg the dying direction itself and one adjoined (Fact \ref{c:singularity.imported-one-sided-completion-singularity}, Remark \ref{c:note.singularity.metric-level-parity-even-level}): the odd sector is not a pre-existing grading of $\Cl(1,3)$ precisely because one of its legs lies outside. Nor is the matching an identification of two coexisting copies, Fact \ref{c:singularity.imported-one-sided-completion-singularity}'s pair existing exactly over $E$: the two constructions share the leg $k$. At every $r > 0$ the radial plane contains $k^{\sharp} = -\sqrt r\,\er$ (Proposition \ref{c:singularity.transport-schwarzschild-fibre}(4)), the plane degenerates onto $\R k$ as $r \to 0$, and the completion re-attaches $\nO/2\nu$ to that same $k$. What is shared is the direction and the degeneration onto it, not a null leg of the atom: $k^{\sharp}$ is spacelike at every $r > 0$ while the atom's legs are null, and the $E$-atom's bivector $k \wedge \nO/2\nu$ is not $\sigma_r$. The identification of the modular bivector with the curvature-selected plane is therefore made at $r > 0$, where Theorem \ref{c:singularity.curvature-gauge-field-exact-shape-decoupling} holds, and \emph{transported} to $E$ along the degeneration; that transport is the matching's residual step, stated as such rather than absorbed into the shared direction, and the two-level computations of Steps 4c through 6b concern the abstract atom, whichever realisation carries it. Then the modular bivector of Theorem \ref{c:singularity.modular-structure-completed-pair},
\[
\sigma \;=\; v^*v - vv^* \;=\; e_0 e_1 ,
\]
spans the radial plane and satisfies the pair of equations
\[
S(\sigma) = -2\,\sigma, \qquad \sigma^2 = +1 ,
\]
and this pair characterises $\sigma$ up to sign among the unit frame bivectors: the eigenvalue equation confines $\sigma$ to $\{\sigma_r, I\sigma_r\}$ (Theorem \ref{c:singularity.curvature-gauge-field-exact-shape-decoupling}), and the square excludes the dual, for which $(I\sigma_r)^2 = -1$. The bivector that generates Step 4c's modular flow, the phase flow $\operatorname{Ad}(e^{it\sigma/2})$, whose real rotor flow is $\exp(\lambda\sigma/2)$, is therefore the same object Step 4b's curvature distinguishes, identified by equation rather than by name; the residual sign is the orientation of the boost, the labelling of outgoing against ingoing, which is exactly the $\mathbb{Z}_2$ that Theorem \ref{c:singularity.modular-structure-completed-pair}(2) attaches to the weights of $v$ and $v^*$. And by the localisation theorem \cite[Thm.~4.4, Cor.~4.5]{ButlerT4} the pair, hence the atom, hence $\sigma$, is at the degenerate radial direction, which is where Proposition \ref{c:singularity.degeneration-exactly} located the division: the plane the curvature selects, the plane the completed pair lies in, and the plane through which division fails are one plane. One plane, not one locus: the curvature selects it at every $r > 0$, the calibration reads it at the horizon (Proposition \ref{c:singularity.plane-rate}), the ladder inhabits it on $E$, and Proposition \ref{c:singularity.why-calibration-cannot-be-internal} records that the calibration locus and the singular locus do not meet; what carries the identification across them is the frame-robustness of the plane (Remark \ref{c:note.singularity.frame-made-explicit}) and the radius-independence of its selection, not a claim that the loci coincide.
\end{corollary}
\begin{proof}
$\sigma = e_0 e_1$ is Theorem \ref{c:singularity.modular-structure-completed-pair}(1) in the stated realisation of the hyperbolic atom (\ref{c:singularity.pseudo-inverse-identities-hyperbolic-atom}), and $e_0 e_1$ spans the radial plane by construction. $S(\sigma_r) = -2\sigma_r$, $\sigma_r^2 = +1$ and $(I\sigma_r)^2 = -1$ are entries of Theorem \ref{c:singularity.curvature-gauge-field-exact-shape-decoupling}'s conjugation table; $S$ is linear, so the eigenvalue equation holds for $\pm\sigma_r$ alike, and among the six unit frame bivectors the two displayed equations are satisfied by $\pm\sigma_r$ and nothing else (the mixed planes fail the first, the dual fails the second). The script checks the characterisation exhaustively over the six planes. The location clause rests on \ref{c:sheaf.one-sided-completion} and \ref{c:sheaf.localised-transport}, which place the ladder pair, hence the atom, exactly over the singular locus, and on \ref{c:singularity.degeneration-exactly}, which locates the division at the degenerate radial direction.

\end{proof}

\begin{remark}[The frame, made explicit]\label{c:note.singularity.frame-made-explicit}
Which frame carries these statements deserves one sentence each, since the $\hb$-image frame is exactly what degenerates. The bivector $\sigma_r$ is defined in the orthonormal frame of the solution, and Theorem \ref{c:singularity.curvature-gauge-field-exact-shape-decoupling}'s invariance display shows the radial \emph{plane} is the same whether computed in the flat background frame or the physical one: the plane is frame-robust precisely where the argument needs it, and no choice is introduced unstated. The sign of $\sigma_r$ (equivalently, which null leg of the plane is called outgoing) is an orientation convention; every appearance of $\sigma_r$ in this paper is either quadratic in it or explicitly sign-convention-flagged, and the verification script fixes one convention and checks every identity in it. The conversion of Proposition \ref{c:singularity.plane-rate}, $2\pi$ of rapidity equal to $8\pi M$ of Killing time, is a statement about the Killing flow's boost in this plane; it is the only place in the chain where geometry is converted into a number carrying units, and Theorem \ref{c:singularity.modular-descent} is where that number becomes the temperature. One more frame fact, imported, settles the question whether the s-wave structure might exist only in coordinates that degenerate: the radial null dyad of this plane, read as covectors, is carried by the gauge field onto forms with \emph{constant unit} components in the infalling orthonormal frame, at every radius up to the wall \cite[Prop.~8.23, Rem.~8.24]{ButlerT4}.
The structure Steps 4d and 4e run on is exactly finite in the frame of the observer who actually crosses, even though the individual coordinate legs diverge.
\end{remark}

\section{Step 4d: compute in the completion; thermality, and the state that carries it}

\phantomsection\label{c:note.singularity.setting-step-4d-compute-completion-thermality}
Temperature is a property of a state, and a fixed projector is state-independent: the projection alone provably cannot select a temperature, and this paper does not pretend otherwise. What the descent structure supplies, exactly, is the complete operator structure of thermality: the modular Hamiltonian ($K = N - \tfrac12$, Theorem \ref{c:singularity.modular-structure-completed-pair}), its spectral projectors (the descent projectors), the localisation plane and the rate (Proposition \ref{c:singularity.plane-rate}), and the uniqueness of the flow. The theorem of the section ``The operator structure, and the theorem it supports'' is then that any state meeting three physically standard conditions descends to the Hawking thermal matrix, exactly. the section ``What selects the state'' takes the next question: what selects the state. The answer is that two of the three conditions are theorems, the third is an output, and the whole residue is a single real number whose value is provably not fixable inside the algebra and is fixed by one named premise. One notational warning: from here to the end of Step 6b, $\beta$ is an inverse temperature; the Painlev\'e-Gullstrand velocity $\beta(r)$ of the section ``Step 1: locate the failure'' reappears only in the appendix and in the verification section.

\begin{lemma}[Modular theory of a matrix factor, proved outright]\label{c:singularity.modular-theory-matrix-factor-proved-outright}
Let $H = \C^n \otimes \C^n$, let $M = \Mat_n(\C) \otimes 1$ act on the first factor, and let $\Omega = \sum_i \lambda_i\, |ii\rangle$ with all $\lambda_i > 0$ and $\sum \lambda_i^2 = 1$ (any cyclic separating vector of $M$ has this form in its Schmidt basis). Define the Tomita operator as the closure of $S : a\Omega \mapsto a^*\Omega$ ($a \in M$). Then its polar decomposition $S = J\Delta^{1/2}$ has
\[
\Delta = \rho \otimes \bar\rho^{\,-1}, \qquad \rho = \sum_i \lambda_i^2\,|i\rangle\langle i| , \qquad J\,|ij\rangle = |ji\rangle \ \text{(antilinearly)},
\]
and the modular flow on $M$ is $\sigma_t(a) = \rho^{it} a\, \rho^{-it}$. Moreover the state $\omega = \langle\Omega, \cdot\,\Omega\rangle$ satisfies the KMS identity $\omega(ab) = \omega(b\,\rho\,a\,\rho^{-1})$ for all $a, b \in M$.
\end{lemma}
\begin{proof}
Everything is finite-dimensional, so closures are automatic. On the matrix units $a = |i\rangle\langle j| \otimes 1$: $a\Omega = \lambda_j |ij\rangle$ and $a^*\Omega = \lambda_i |ji\rangle$, so $S$ is the antilinear map $|ij\rangle \mapsto (\lambda_i/\lambda_j)\,|ji\rangle$, and since the $a\Omega$ span $H$ ($\Omega$ is cyclic), this determines $S$. Define $\Delta |ij\rangle := (\lambda_i^2/\lambda_j^2)|ij\rangle$, which is exactly $(\rho \otimes \bar\rho^{\,-1})|ij\rangle$, positive and invertible, and let $J$ be the antilinear extension of $|ij\rangle \mapsto |ji\rangle$, which is antiunitary with $J^2 = 1$. Then $J\Delta^{1/2}|ij\rangle = (\lambda_i/\lambda_j) J|ij\rangle = (\lambda_i/\lambda_j)|ji\rangle = S|ij\rangle$ (the eigenvalues are real, so antilinearity causes no conjugation), and uniqueness of the polar decomposition gives the claim. The flow: $\Delta^{it}(b \otimes 1)\Delta^{-it} = (\rho^{it}\otimes\bar\rho^{-it})(b \otimes 1)(\rho^{-it}\otimes\bar\rho^{\,it}) = \rho^{it} b \rho^{-it} \otimes 1$, interior to $M$. KMS: the reduced density matrix of $\Omega$ on the first factor is $\rho$, so $\omega(x) = \tr(\rho x)$ for $x \in \Mat_n(\C)$, and $\omega(b\,\rho a \rho^{-1}) = \tr(\rho\, b\, \rho a \rho^{-1}) = \tr(\rho\, a b) = \omega(ab)$ by cyclicity of the trace.

\end{proof}

\begin{lemma}[Schmidt forcing, proved outright]\label{c:singularity.schmidt-forcing-proved-outright}
Realise the boost on the doubled system by the one-parameter unitary group $e^{i t G}$, $G = N \otimes 1 - 1 \otimes N$ (the pair carries the flow of Theorem \ref{c:singularity.modular-structure-completed-pair}, the partner the mirrored flow). If a unit vector $\Omega \in \C^2 \otimes \C^2$ is invariant up to phase under this group and is cyclic and separating for the pair factor, then
\[
\Omega \;=\; \cos\theta\,|00\rangle + \sin\theta\,|11\rangle \qquad \text{with } \cos\theta\sin\theta \neq 0 ,
\]
up to a global phase and a phase convention on $|11\rangle$: the Schmidt basis is the number basis.
\end{lemma}
\begin{proof}
Invariance up to phase under a one-parameter unitary group means $\Omega$ is an eigenvector of its generator. $G$ is diagonal in the product number basis with $G|ij\rangle = (i - j)|ij\rangle$: eigenvalue $0$ on $\operatorname{span}\{|00\rangle, |11\rangle\}$, eigenvalue $+1$ on $\C|10\rangle$, eigenvalue $-1$ on $\C|01\rangle$. An eigenvector for $\pm 1$ is a product vector, and a product vector is annihilated by some nonzero $a \in M$ (respectively fails cyclicity), contradicting separating; so $\Omega = c_0|00\rangle + c_1|11\rangle$. If either coefficient vanished, $\Omega$ would again be a product vector; so both are nonzero, and phases can be absorbed into the basis vectors, leaving the stated real form.

\end{proof}

\begin{theorem}[Modular descent]\label{c:singularity.modular-descent}
\changedmark{BurntOrange}{changed since the deposit}
Let the ascended algebra of the horizon pair carry any state that is (i) pure, (ii) invariant under the boost flow generated by $\sigma$, and (iii) entangling across the two ends of the pair (cyclic and separating for the one-sided subalgebra). Then the descended state, the partial trace over the partner, which for such a state coincides with the restriction under the pinching $E$ of Proposition \ref{c:singularity.descent-what-it-erases} (Proposition \ref{c:singularity.pinching-against-partial-trace-scope-coincidence}(i); the coincidence is exactly single-mode, Proposition \ref{c:singularity.pinching-against-partial-trace-scope-coincidence}(ii)), is KMS for the phase flow $\operatorname{Ad}(e^{it\sigma/2})$ of $\sigma$ at some inverse temperature $\beta \in \R$ in rapidity, and the hypotheses leave $\beta$ free, its sign included: nothing in (i) to (iii) forces $\beta > 0$ (the descended form of Theorem \ref{c:singularity.descended-form-state-independent} allows the closed range $[-\infty, +\infty]$, and hypothesis (iii) removes exactly the endpoints), and positivity arrives with (R). If the modular parameter is normalised to the geometric period, one modular period being one period $2\pi/\kappa$ of rapidity-time (the Bisognano-Wichmann normalisation, which the section ``What selects the state'' names as premise (R) and proves irreducible), then $\beta = 2\pi$ in rapidity and, by Proposition \ref{c:singularity.plane-rate}, the physical temperature is
\[
T \;=\; \frac{\kappa}{2\pi} \;=\; \frac{1}{8\pi M} \;=\; \TH ,
\]
exactly the Hawking temperature. For the CAR pair, the reduced density matrix per mode of frequency $\omega$ is then the Fermi-Dirac matrix $\operatorname{diag}(1, e^{-\omega/T})/Z$ with occupancy $1/(e^{\omega/T} + 1)$. The hypotheses are sufficient, not necessary: a state violating (ii) or (iii) need not descend thermally, and the projection cannot make it do so, but some such states do (the proof exhibits one), and Theorem \ref{c:singularity.descended-form-state-independent} explains why no hypothesis on the state can be necessary for the \emph{form}.
\end{theorem}
\begin{proof}
The pair algebra is the $2 \times 2$ factor generated by $(v, v^*)$; its commutant in the doubled system is the partner copy, and hypotheses (i) and (iii) say exactly that the state vector is cyclic and separating for the pair factor. By Lemma \ref{c:singularity.schmidt-forcing-proved-outright}, hypothesis (ii) forces
\[
|\Omega\rangle = \cos\theta\,|00\rangle + \sin\theta\,|11\rangle, \qquad \cos\theta\sin\theta \neq 0 ,
\]
and by Lemma \ref{c:singularity.modular-theory-matrix-factor-proved-outright} (with $n = 2$, $\lambda_0 = \cos\theta$, $\lambda_1 = \sin\theta$) the modular operator of $|\Omega\rangle$ is $\Delta = \rho \otimes \bar\rho^{\,-1}$ with $\rho = \operatorname{diag}(\cos^2\theta, \sin^2\theta)$ and modular flow $\sigma_t = \operatorname{Ad}(\rho^{it})$ on the pair factor, KMS identity included. So far the state is Gibbs for $K$ at an inverse temperature $\beta$ determined by $\tan^{2}\theta = e^{-\beta}$, and nothing in (i) to (iii) constrains $\theta$: this is the residue the section ``What selects the state'' takes up. It remains to normalise, and this one step is an import. Assign the mode its frequency by scaling the boost generator: the physical generator of the flow is $\omega K$, $K = N - \tfrac12$ (Theorem \ref{c:singularity.modular-structure-completed-pair}(4)), and Proposition \ref{c:singularity.plane-rate} converts rapidity to Schwarzschild time at rate $\kappa$. The modular group of Lemma \ref{c:singularity.modular-theory-matrix-factor-proved-outright} runs at parameter $t$ with generator $\log\Delta$; demanding, by premise (R) (\ref{c:singularity.r-regularity-fixed-locus}), that modular parameter $2\pi$ correspond to the geometric flow of $\omega K$ over one period $2\pi/\kappa$ of rapidity-time forces $\log\rho = -(2\pi/\kappa)\,\omega N + \text{const}$, i.e.
\[
\rho \;\propto\; e^{-(2\pi/\kappa)\,\omega N}
\quad\Longleftrightarrow\quad
\tan\theta = e^{-\omega/2T}, \qquad T = \frac{\kappa}{2\pi} ,
\]
since $\rho = \operatorname{diag}(\cos^2\theta, \sin^2\theta)$ and $N = \operatorname{diag}(0,1)$. The partial trace over the partner is then $\rho_{\text{mode}} = \operatorname{diag}(1, e^{-\omega/T})/Z$: Fermi-Dirac at exactly $\TH$, manifestly KMS for the flow of $K$ by the last clause of Lemma \ref{c:singularity.modular-theory-matrix-factor-proved-outright}. What the hypotheses do and do not force: if (ii) fails, Lemma \ref{c:singularity.schmidt-forcing-proved-outright}'s eigenvector argument fails and the Schmidt basis need not be the number basis, so the reduction need not be diagonal in $N$ and then is not a Gibbs state of $K$ (take $\Omega = c\,|{+}\rangle|0\rangle + s\,|{-}\rangle|1\rangle$ with $|{\pm}\rangle = (|0\rangle \pm |1\rangle)/\sqrt2$ and $c \neq s$, cyclic and separating, whose reduction $\tfrac12\begin{psmallmatrix} 1 & c^{2} - s^{2} \\ c^{2} - s^{2} & 1\end{psmallmatrix}$ has nonzero coherence); but the converse fails, since $\Omega = c|01\rangle + s|10\rangle$ with $cs \neq 0$ is pure, cyclic and separating and not a $G$-eigenvector, yet reduces to $\operatorname{diag}(c^{2}, s^{2})$, Fermi-Dirac at exactly $\TH$ when $s^{2}/c^{2} = e^{-\omega/\TH}$. If (iii) fails, one Schmidt coefficient vanishes and the reduction is pure, not thermal at any finite temperature. The script checks both examples. In general, conditions (i)-(iii) are the hypotheses of Tomita-Takesaki theory, and the identification of modular flow with the boost for horizon-localised algebras is the Bisognano-Wichmann-Sewell mechanism \cite{BW1976, Sewell1982}; the chain's contribution is that, within the class of real rotor flows the geometry supplies, the flow is algebraically \emph{forced} (Theorem \ref{c:singularity.modular-structure-completed-pair}(5), with the scope stated there) rather than assumed, and that the single-mode instance is proved outright above rather than quoted.

\end{proof}

\begin{remark}[Israel's reading, recovered]\label{c:note.singularity.israels-reading-recovered}
The natural state satisfying (i)-(iii) is the thermofield double taken across the two ends of the completed pair, which is Israel's 1976 reading of the eternal black hole \cite{Israel1976}, now expressed in the chain's own terms: the global vacuum is entanglement between the element and its adjoined partner, and Hawking radiation is what descent sees of that entanglement. The slogan of Step 4a (the exterior radiation as the shadow of the return map) is hereby made quantitative: the shadow is the partial trace, identified mode by mode with the pinching $E$ by Proposition \ref{c:singularity.pinching-against-partial-trace-scope-coincidence}(i) (the coincidence scoped precisely in Step 5, which this identification anticipates), and its output is thermal at exactly $\TH$. What the identification deliberately does not extend to is the full boundary mode space, where the two maps differ (Proposition \ref{c:singularity.pinching-against-partial-trace-scope-coincidence}(ii)) and the computation is Theorem \ref{c:singularity.deficit-theorem-cost-descent-exactly}'s. The word ``natural'' is made precise in the section ``What selects the state'': by Theorem \ref{c:singularity.purity-output-partner-commutant} the thermofield double is not one state among those satisfying the hypotheses but the canonical purification of the one-sided state, the partner being the commutant produced by the modular conjugation.
\end{remark}

\begin{remark}[Three wrinkles, recorded]\label{c:note.singularity.three-wrinkles-recorded}
First, statistics: the chain's pairs are fermionic, so the directly produced spectrum is Fermi-Dirac, correct for the fermionic sector of Hawking emission; the bosonic Planck sector is Step 4e's task and is derived, not assumed. Second, coherence with the count: at $\omega \sim \TH$ the thermal entropy per fermionic mode is of order $\ln 2$, one bit, exactly the per-mode capacity that converts Step 6a's mode count $N_{\mathrm{6a}}$ into the boundary-mode capacity $C = N_{\mathrm{6a}}\ln 2$ of Proposition \ref{c:singularity.area-law-test}: the thermodynamics and the counting cohere. Third, the role of the hyperbolic atom is now complete: $\Cl(1,1)$, the unique self-complete unit of transport (Lemma \ref{c:singularity.pseudo-inverse-identities-hyperbolic-atom}), is the horizon's boost plane, matched to the curvature-selected plane by Corollary \ref{c:singularity.matching-steps-4b-4c-join-equation}, and its intrinsic structure, the unique rotor flow with its spectral projectors, is the modular structure of black hole thermodynamics. The atom that needs no ascent is the one that carries the temperature.
\end{remark}

\subsection*{What selects the state}
\phantomsection\label{c:note.singularity.setting-what-selects-state}
Theorem \ref{c:singularity.modular-descent} is conditional on three hypotheses, and the status line of the table accompanying \ref{c:singularity.imported-equivalence-einstein-equations} would read \emph{up to a state} because of them. This subsection shows that the first is an output rather than an input, that the second is not needed for the form at all and that its per-mode content follows from a statement about observables, that what remains is exactly one real number per mode, that this number cannot be fixed inside the algebra, and that one named premise fixes it, mode by mode, and with it the third hypothesis. The central step is Theorem \ref{c:singularity.descent-parity-average}: the descent is the parity average, so the question ``which states descend thermally'' is a question about a $\mathbb{Z}_2$ grading, and gradings are easily supplied.

\begin{theorem}[The descended form is state independent]\label{c:singularity.descended-form-state-independent}
Let $\varrho$ be the density matrix of \emph{any} state on the boundary mode space, and $\varrho_k$ its marginal at mode $k$. Then $E_k(\varrho_k)$ is diagonal in that mode's number basis, and therefore
\[
E_k(\varrho_k) \;=\; \frac{e^{-\beta_k \omega_k N_k}}{Z_k}
\]
for a unique $\beta_k \in [-\infty, +\infty]$, determined by $\beta_k\omega_k = \log\bigl(p_k/(1-p_k)\bigr)$ with $p_k$ the descended vacancy probability. (Stated for $\omega_k \neq 0$: at a zero-weight mode the right side is $\tfrac12\,\mathrm{id}$ for every $\beta_k$, so uniqueness always fails there, and existence fails unless $p_k = \tfrac12$.) The descended single-mode state is a Gibbs state of the boost generator $K_k$ whatever the ascended state was: no hypothesis on the state is used.
\end{theorem}
\begin{proof}
By Theorem \ref{c:singularity.descent-parity-average} the range of $E_k$ is the even subalgebra, which in the number basis is the diagonal. A diagonal density matrix on a two-dimensional space is $\operatorname{diag}(p_k, 1 - p_k)$, and $\beta_k \mapsto (1 + e^{\beta_k\omega_k})^{-1}$ is a bijection from $\R$ onto $(0,1)$, so every such matrix with $0 < p_k < 1$ is the stated Gibbs matrix for exactly one $\beta_k$; the two boundary values are the limits $\beta_k = \pm\infty$.

\end{proof}

\begin{remark}[The unconditional reading, and how it sharpens the gap]\label{c:note.singularity.unconditional-reading-how-it-sharpens-gap}
Theorem \ref{c:singularity.descended-form-state-independent} says less than it appears to, and that matters. For a single fermionic mode ``thermal'' is not a discriminating property: after the descent the mode's state space is one-dimensional and every point of it is a Gibbs state. So no part of the physical content of this step lies in the \emph{shape} of the reduced density matrix, and the chain must not claim the Fermi-Dirac form as a prediction: the descent produces that form from any input whatever. All of the content lies in the number $\beta_k$, in the claim that it is the same at every mode, and in the claim that its value is $2\pi/\kappa$. The gap is therefore not ``which class of states''; it is ``which point of a line, and why the same point at every mode''. The rest of this subsection is about that line.
\end{remark}

\begin{fact}[Imported: the univalence superselection rule]\label{c:singularity.imported-univalence-superselection-rule}
In a theory whose field algebra is $\mathbb{Z}_2$ graded by fermion number, coherent superpositions of states of opposite fermion parity are not observable: the observable algebra is the even subalgebra and physical states are even. \emph{Provenance:} Wick, Wightman and Wigner \cite{WWW1952}; the graded structure of CAR algebras and their even subalgebras is \cite[\S 5.2]{BR2}. What the rule superselects is the total parity, a single $\mathbb{Z}_2$; the mode-wise strengthening that the multi-mode descent of Step 5 needs is Proposition \ref{c:singularity.p-reduced-metric-level-reading-many}, stated there, and is not part of this import.
\emph{Source:} G. C. Wick, A. S. Wightman and E. P. Wigner, The intrinsic parity of elementary particles, Phys. Rev. 88 (1952), 101--105.; O. Bratteli and D. W. Robinson, Operator Algebras and Quantum Statistical Mechanics 2: Equilibrium States. Models in Quantum Statistical Mechanics, 2nd ed., Springer (1997)., \S 5.2.
\end{fact}

\phantomsection\label{c:note.singularity.setting-what-selects-state-2}
The chain is entitled to (P) for a reason it supplies itself: the boundary modes form a CAR algebra, one pair per point of $E$ (Fact \ref{c:singularity.imported-one-sided-completion-singularity}; the tower limit \cite[Thm.~5.6]{ButlerT3} is its algebraic model), its grading is the Clifford grading of Definition \ref{c:obj.singularity.grading-parity-automorphism}, and by Theorem \ref{c:singularity.descent-parity-average} its even part is, mode by mode, exactly the metric-level algebra (the multi-mode metric level is the joint-diagonal algebra, whose licence is Proposition \ref{c:singularity.p-reduced-metric-level-reading-many}, one parity per point rather than the product of all of them). What (P) asserts is that the odd sector carries no expectation values, which by that theorem is the same as asserting that the ascended and descended readings agree \emph{mode by mode}. That is a much weaker demand than invariance of the whole state under the boost flow, and it is not a choice of state at all: it is a statement about which elements of the boundary mode space are observables.

\begin{theorem}[Ascended marginals are diagonal]\label{c:singularity.ascended-marginals-are-diagonal}
Let $\omega$ satisfy (P). Then $\omega(v_k) = \omega(v_k^{*}) = 0$ at every mode, so the ascended marginal $\varrho_k$ is diagonal in the number basis, coincides with $E_k(\varrho_k)$, and is by Theorem \ref{c:singularity.descended-form-state-independent} the Gibbs state of $K_k$ at a unique $\beta_k$ for $\omega_k \neq 0$ (at a zero-weight mode diagonality still holds, and the Gibbs form requires $p_k = \tfrac12$, Theorem \ref{c:singularity.descended-form-state-independent}'s parenthesis). Consequently partial trace and pinching coincide on every single-mode marginal, which is Proposition \ref{c:singularity.pinching-against-partial-trace-scope-coincidence}(i) with hypothesis (ii) of Theorem \ref{c:singularity.modular-descent} removed; and the parity protection lemma of the finite-mode model, Lemma \ref{c:singularity.parity-protection}, is the special case of this theorem for the model's states, which are even because every configuration in their support has even total occupation.
\end{theorem}
\begin{proof}
$v_k$ and $v_k^{*}$ are odd, so $\alpha(v_k) = -v_k$ and (P, \ref{c:singularity.p-univalence}) gives $\omega(v_k) = \omega(\alpha(v_k)) = -\omega(v_k)$, hence $\omega(v_k) = 0$, and likewise for $v_k^{*}$. The off-diagonal entries of $\varrho_k$ in the number basis are exactly $\omega(v_k)$ and $\omega(v_k^{*})$, those being the mode's two off-diagonal matrix units. Diagonality is the coincidence with $E_k(\varrho_k)$, and Theorem \ref{c:singularity.descended-form-state-independent} supplies the Gibbs form. The argument uses no finiteness and holds verbatim on the CAR algebra over $E$. For the model: its states are supported on configurations with $R_k = L_k$ at every $k$, whose total occupation $2\sum_k b_k$ is even, and a state supported in a single parity sector is even. The script verifies both clauses, with a parity-mixed control whose marginal coherence is nonzero.

\end{proof}

\begin{remark}[What hypothesis (ii) cost, and what (P) costs]\label{c:note.singularity.what-hypothesis-ii-cost-what-p}
Theorem \ref{c:singularity.ascended-marginals-are-diagonal} replaces a symmetry assumption about the state by a statement about which elements are observable. The replacement is strictly weaker: hypothesis (ii) constrains the whole state, (P) constrains one complex number per mode. It is also more robust, and Theorem \ref{c:singularity.exact-equilibrium-over-determines-boundary-modes} below shows that the robustness is necessary, not merely convenient: hypothesis (ii) is an equilibrium condition, the evaporating hole is not in equilibrium, and the exact equilibrium condition turns out to be too strong for what Steps 5 and 6b need.
\end{remark}

\subsection*{One temperature, and the rigidity of the number}
\phantomsection\label{c:note.singularity.setting-one-temperature-rigidity-number}
The exact equilibrium condition is stated next, and it is stated in order to be examined rather than assumed: Corollary \ref{c:singularity.universality-temperature-under-g} records what it gives, Theorem \ref{c:singularity.exact-equilibrium-over-determines-boundary-modes} shows that it gives too much, and the premise the chain actually adopts, (R) below, is imposed mode by mode instead.

\begin{definition}[The geometric modular condition (G)]\label{c:obj.singularity.geometric-modular-condition-g}
\changedmark{BurntOrange}{changed since the deposit}
Let $K_{\mathrm{geo}} := \sum_k \omega_k K_k$ generate the Killing boost flow of the boundary modes, the flow whose bivector at each point of $E$ is the curvature-selected $\sigma_r$ of Theorem \ref{c:singularity.curvature-gauge-field-exact-shape-decoupling} and Corollary \ref{c:singularity.matching-steps-4b-4c-join-equation}, with $\omega_k$ the mode's weight. A \emph{faithful} state satisfies the \emph{geometric modular condition} (G) if its modular automorphism group is the phase flow $\operatorname{Ad}(e^{itK_{\mathrm{geo}}})$ generated by $K_{\mathrm{geo}}$, not its real rotor flow, up to a single constant reparametrisation. Faithfulness is part of the condition, not a consequence of it: a state has a modular automorphism group only if it is faithful, so (G) presupposes hypothesis (iii) of Theorem \ref{c:singularity.modular-descent} and cannot be used to derive it.
\end{definition}

\begin{theorem}[KMS uniqueness at a factor]\label{c:singularity.kms-uniqueness-factor}
\changedmark{BrickRed}{corrected since the deposit}
Let $\Aa$ be a matrix factor, $\alpha_t = \operatorname{Ad}(e^{itK})$ a flow and $\beta \neq 0$. A state $\omega = \tr(\varrho\,\cdot\,)$ is $(\alpha,\beta)$-KMS if and only if $\varrho = e^{-\beta K}/\tr e^{-\beta K}$. In particular there is exactly one KMS state for each $\beta$, and, for $K$ not a multiple of $1$, the KMS states of the flow form a line parametrised by $\beta$.
\end{theorem}
\begin{proof}
In finite dimensions the KMS condition reads $\omega(ab) = \omega\bigl(b\,\alpha_{i\beta}(a)\bigr) = \omega\bigl(b\,e^{-\beta K}a\,e^{\beta K}\bigr)$ for all $a, b$. Writing both sides as traces and using cyclicity, the left side is $\tr(b\varrho a)$ and the right side is $\tr\bigl(e^{\beta K}\varrho\, b\, e^{-\beta K} a\bigr)$, so for all $b$,
\[
b\,\bigl(\varrho\, e^{\beta K}\bigr) \;=\; \bigl(e^{\beta K}\varrho\bigr)\, b .
\]
Taking $b = 1$ gives $\varrho e^{\beta K} = e^{\beta K}\varrho$, so $\varrho$ commutes with $K$; the displayed identity then says that $\varrho\,e^{\beta K}$ commutes with everything, hence is central, hence scalar because $\Aa$ is a factor. So $\varrho \propto e^{-\beta K}$, and normalisation fixes the constant. The converse is the same computation read backwards. The script checks the identity on all sixteen matrix-unit pairs for $\varrho \propto \operatorname{diag}(1, \tfrac{16}{9})$ and exhibits a pair at which the tracial state fails it.

\end{proof}

\begin{fact}[Imported: the general form]\label{c:singularity.imported-general-form}
For a $C^{*}$-dynamical system the $(\alpha,\beta)$-KMS states at fixed $\beta$ form a weak$^{*}$-compact convex simplex whose extreme points are exactly those whose GNS von Neumann algebra is a factor, distinct extremal KMS states being disjoint; a factor state is therefore extremal, and uniqueness at fixed $\beta$ is the statement that the simplex is a point. \emph{Provenance:} Bratteli and Robinson \cite[Thm.~5.3.30]{BR2}, the KMS condition itself being \cite{HHW1967}. The companion statement that two faithful normal states with the same modular automorphism group differ by an element affiliated to the centre, hence coincide on a factor, is Connes' Radon-Nikodym cocycle theorem \cite[Thm.~1.2.1]{Connes1973}.
\emph{Source:} O. Bratteli and D. W. Robinson, Operator Algebras and Quantum Statistical Mechanics 2: Equilibrium States. Models in Quantum Statistical Mechanics, 2nd ed., Springer (1997)., Thm.~5.3.30; R. Haag, N. M. Hugenholtz and M. Winnink, On the equilibrium states in quantum statistical mechanics, Commun. Math. Phys. 5 (1967), 215--236.; A. Connes, Une classification des facteurs de type III, Ann. Sci. 'Ec. Norm. Sup'er. (4) 6 (1973), 133--252., Thm.~1.2.1.
\end{fact}

\begin{corollary}[Universality of the temperature under (G)]\label{c:singularity.universality-temperature-under-g}
A state satisfying (G) has $\beta_k = \beta$ independent of $k$ over the modes of nonzero weight: one temperature for all the boundary modes. Conversely a state with per-mode marginals $\mathrm{Gibbs}(\beta_k\omega_k)$ and no inter-mode correlation satisfies (G) only if the $\beta_k$ of the modes of nonzero weight coincide.
\end{corollary}
\begin{proof}
By Theorem \ref{c:singularity.kms-uniqueness-factor} applied to $K_{\mathrm{geo}}$, a state satisfying (G) is $\varrho \propto e^{-\beta K_{\mathrm{geo}}} = \bigotimes_k e^{-\beta\omega_k K_k}$, whose $k$-th marginal is $\mathrm{Gibbs}(\beta\omega_k)$: the same $\beta$ at every mode. Conversely the modular generator of a product state is $\log\varrho = -\sum_k \beta_k\omega_k N_k + \mathrm{const}$, proportional to $K_{\mathrm{geo}} = \sum_k \omega_k N_k + \mathrm{const}$ if and only if $\beta_k\omega_k = \lambda\omega_k$ for a single $\lambda$ and every $k$ of nonzero weight (under Premise \ref{c:singularity.q-joint-quasi-freeness-s-wave}'s linear assignment the $k = 0$ mode is excluded from this clause and from the converse: at $\omega_0 = 0$ the marginal is $\tfrac12\,\mathrm{id}$ for every $\beta_0$ and constrains nothing).

\end{proof}

\begin{theorem}[Rate rigidity: the temperature is not algebra intrinsic]\label{c:singularity.rate-rigidity-temperature-not-algebra-intrinsic}
Let $\alpha_t = \operatorname{Ad}(e^{itK})$ be the boost flow of the single-mode algebra $\Aa$, let $\theta$ be a $*$-automorphism of $\Aa$, and let $\omega_\beta$ denote the $\beta$-KMS state of $\alpha$.
\begin{enumerate}
\item If $\theta$ commutes with the flow then it fixes every Gibbs state of $K$: $\omega_\beta \circ \theta = \omega_\beta$ for every $\beta$.
\item If $\theta$ reverses the flow, $\theta\circ\alpha_t = \alpha_{-t}\circ\theta$, then $\omega_\beta\circ\theta = \omega_{-\beta}$.
\item Under a constant reparametrisation of the flow, $\alpha_t \mapsto \alpha_{ct}$ with $c > 0$, the $\beta$-KMS state of $\alpha$ is the $(\beta/c)$-KMS state of the reparametrised flow: $\omega_\beta$ is unchanged and its inverse temperature is rescaled.
\item Hence no automorphism intertwining the flow carries $\omega_\beta$ to $\omega_{\beta'}$ for $|\beta'| \neq |\beta|$: $|\beta|$ is an invariant of the pair (state, parametrised flow), and the automorphism group acts on the line of KMS states with orbits $\{\beta, -\beta\}$, so that line carries no automorphism-theoretic structure beyond the orientation flip of (2). A condition that singles out one point of the line is therefore either a stipulation of the state itself or a reference to data outside the pair $(\Aa, \alpha)$; by (3), the value of $|\beta|$ is meaningful only relative to a parametrisation of the flow, and it is that parametrisation, the conversion of the flow parameter into geometric time, that the algebra does not supply.
\end{enumerate}
\end{theorem}
\begin{proof}
Every automorphism of a matrix factor is inner, $\theta = \operatorname{Ad}(U)$. (1) Commutation gives $Ue^{itK}U^{*} = c(t)e^{itK}$ with $|c(t)| = 1$; taking determinants, $c(t)^{2} = 1$, and $c$ is continuous with $c(0) = 1$, so $c \equiv 1$ and $U$ commutes with $K$. As $K = \operatorname{diag}(-\tfrac12, \tfrac12)$ is nondegenerate, $U$ is diagonal and $\operatorname{Ad}(U)$ fixes every diagonal matrix, in particular $e^{-\beta K}/Z$. (2) Flow reversal gives $UKU^{*} = -K$, so $U$ exchanges the spectral projectors of $K$ and $Ue^{-\beta K}U^{*} = e^{+\beta K}$. (3) The reparametrised flow is $\operatorname{Ad}(e^{it(cK)})$, generated by $cK$, and by Theorem \ref{c:singularity.kms-uniqueness-factor} its $\beta'$-KMS state is $e^{-\beta' cK}/Z$, which equals $e^{-\beta K}/Z$ exactly when $\beta' = \beta/c$. (4) By Theorem \ref{c:singularity.kms-uniqueness-factor} the KMS states are exactly the Gibbs states, and an automorphism intertwining the flow either commutes with it or reverses it (it conjugates $K$ to $\pm K$ up to a constant, the spectrum $\{-\tfrac12, \tfrac12\}$ of $K$ leaving no other option), so by (1) and (2) it acts on the parameter through $\{\pm 1\}$ only. That is the whole mathematical content of the clause; the sentence about what a selecting condition must refer to is its reading, and clause (3) is the reason the reading is exact: $\beta$ is the ratio of the state's modular parameter to the flow's parameter, and the second of these is a choice the algebra does not make. The script verifies clauses (1) to (3) on rational data.

\end{proof}

\begin{remark}[Two clocks, one calibration]\label{c:note.singularity.two-clocks-one-calibration}
Theorem \ref{c:singularity.rate-rigidity-temperature-not-algebra-intrinsic} is a no-go, and a constructive one: it says exactly what kind of input is needed and how much of it. The inverse temperature is a ratio of two clocks, the modular clock of the state and the geometric clock of the Killing boost, and a ratio of clocks is meaningless until one is calibrated. The chain calibrates exactly once, in Proposition \ref{c:singularity.plane-rate}: an interval of $2\pi$ in rapidity is $2\pi/\kappa = 8\pi M$ of Schwarzschild time. Everything numerical in this step passes through that conversion, as the step already said, and the rigidity theorem now proves that no rearrangement of the algebra can remove the need for it. The residual sign in clause (2) is not new either: it is the orientation of the boost, the labelling of outgoing against ingoing, flagged in Remark \ref{c:note.singularity.frame-made-explicit} as the $\mathbb{Z}_2$ attached to the weights of $v$ and $v^{*}$. The rigidity theorem identifies that $\mathbb{Z}_2$ as the entire automorphism-theoretic freedom in the temperature.
\end{remark}

\subsection*{Fixing the rate, and where the premise has to lie}
\phantomsection\label{c:note.singularity.setting-fixing-rate-where-premise-has}
Theorem \ref{c:singularity.rate-rigidity-temperature-not-algebra-intrinsic}(3) says what the missing datum is: a conversion of the flow parameter into geometric time. The geometry supplies one, and supplies it at a definite place. Wick rotation of the exterior region turns the Killing boost into a rotation about the bifurcation sphere, and the section is a smooth manifold there for exactly one period of imaginary Killing time. That period is computed first; the premise then says that the state's correlation functions lie on the section.

\begin{proposition}[The Euclidean section, and its period]\label{c:singularity.euclidean-section-period}
\changedmark{BurntOrange}{changed since the deposit}
Let $t_s = -i\tau$ in the static gauge of the Schwarzschild gravitational field on flat spacetime (the section ``The metric of the gauge field, and its static form''; Definition \emph{Gauge Theory Gravity: the gravitational field on flat spacetime}). The Wick-rotated field lives on the flat Euclidean background, and its derived metric is the Riemannian metric $f\,d\tau^{2} + f^{-1}dr^{2} + r^{2}d\Omega^{2}$ on $r > 2M$; let $\varrho(r) := \int_{2M}^{r} f^{-1/2}\,dr'$ be proper radial distance from the horizon locus. Then, with $F := \sqrt f$ regarded as a function of $\varrho$,
\[
f\,d\tau^{2} + f^{-1}dr^{2} \;=\; d\varrho^{2} + F(\varrho)^{2}\,d\tau^{2}, \qquad F(0) = 0, \qquad \frac{dF}{d\varrho}\Big|_{\varrho = 0} \;=\; \frac{f'(2M)}{2} \;=\; \kappa \;=\; \frac{1}{4M} ,
\]
so that near $\varrho = 0$ the $(\varrho, \tau)$ plane carries the metric $d\varrho^{2} + \kappa^{2}\varrho^{2}d\tau^{2} + O(\varrho^{4})\,d\tau^{2}$. If $\tau$ is periodic with period $p$, the section is a smooth Riemannian manifold at $\varrho = 0$ (the tip of the section, the two-sphere where the Euclidean Killing flow has its fixed points, the Euclidean counterpart of the bifurcation sphere) if and only if
\[
\kappa\,p \;=\; 2\pi, \qquad \text{i.e.} \qquad p \;=\; \frac{2\pi}{\kappa} \;=\; 8\pi M ;
\]
for any other period the origin of the $(\varrho, \tau)$ plane is a conical point of total angle $\kappa p$. The period of imaginary Killing time is therefore fixed by the gravitational field, at the horizon locus, and it is the conversion factor of Proposition \ref{c:singularity.plane-rate} read in imaginary time. The regularity condition is local: it lives in a polar chart about the tip of the Wick-rotated flat background; the Lorentzian bifurcation surface itself lies beyond one flat chart and is not used.
\end{proposition}
\begin{proof}
$d\varrho = f^{-1/2}dr$ gives $d\varrho^{2} = f^{-1}dr^{2}$, and $F^{2} = f$, which is the first display. $F(0) = \sqrt{f(2M)} = 0$. By the chain rule $dF/d\varrho = (d\sqrt f/dr)(dr/d\varrho) = \bigl(f'/2\sqrt f\bigr)\sqrt f = f'/2$, an exact identity of rational functions in $r$, and $f'(2M)/2 = \kappa$ is Proposition \ref{c:singularity.computed}. Since $F$ is smooth in $r$ and $\varrho$ is a smooth function of $r$ with $d\varrho/dr = f^{-1/2}$, one has $\varrho = 2\sqrt{(r - 2M)/f'(2M)}\,(1 + O(r - 2M))$, so $F = \kappa\varrho + O(\varrho^{3})$ ($F$ is odd in $\varrho$ under the reflection $\varrho \mapsto -\varrho$ that extends the metric across the origin), which is the expansion. The rest is the elementary geometry of a surface of revolution: with $\phi := \kappa\tau$ the leading metric is $d\varrho^{2} + \varrho^{2}d\phi^{2}$, the flat plane in polar coordinates, which is smooth at the origin exactly when $\phi$ has period $2\pi$; the circumference of the geodesic circle of radius $\varrho$ is $pF(\varrho) = \kappa p\,\varrho + O(\varrho^{3})$, so the total angle at the origin is $\kappa p$, and a cone of total angle other than $2\pi$ is not a smooth manifold. The sphere factor $r^{2}d\Omega^{2}$ is smooth and nondegenerate at $r = 2M$ and plays no role. The script verifies $\sqrt f\,(d\sqrt f/dr) = f'/2$ as an identity of rational functions, the vanishing of $f$ at $2M$, and $f'(2M)/2 = 1/4M$ exactly.

\end{proof}

\phantomsection\label{c:note.singularity.setting-fixing-rate-where-premise-has-2}
The Euclidean-regularity reading of the period is Gibbons-Hawking \cite{GibbonsHawking1977, HartleHawking1976}; the proposition's contribution is the exact rational-function form $dF/d\varrho = f'/2$ in $F = \sqrt f$ and the period computed rather than chosen.

\phantomsection\label{c:note.singularity.setting-fixing-rate-where-premise-has-3}
(R) is the KMS condition at a stated period, imposed on the marginals and on nothing else; it is not a condition on the joint state of the boundary modes, and Theorem \ref{c:singularity.exact-equilibrium-over-determines-boundary-modes} is the reason it must not be. What it adds to ``KMS at some period'' is exactly the value $2\pi/\kappa$, and Theorem \ref{c:singularity.rate-rigidity-temperature-not-algebra-intrinsic} says that this is the one datum no condition inside the algebra can supply. The condition makes sense for every state, faithful or not: $F_{a,b}$ is defined by evaluation, not through a modular group. That is what makes the next theorem's last clause a derivation rather than a restatement.

\begin{theorem}[The rate, and the discharge of hypothesis (iii)]\label{c:singularity.rate-discharge-hypothesis-iii}
Assume (P) and (R). Then at every mode $k$ the marginal is the Gibbs state of $K_k$ at the same inverse temperature,
\[
\beta_k \;=\; \beta \;=\; \frac{2\pi}{\kappa} \;=\; 8\pi M , \qquad T \;=\; \frac{\kappa}{2\pi} \;=\; \frac{1}{8\pi M} \;=\; \TH ,
\]
with occupancy $\omega(v_k^{*}v_k) = (e^{\beta\omega_k} + 1)^{-1} \in (0,1)$. In particular every marginal is faithful, hence cyclic and separating for its mode factor: hypothesis (iii) of Theorem \ref{c:singularity.modular-descent} is a consequence of (R), not an assumption; and the temperature is the same at every mode without any condition on the joint state.
\end{theorem}
\begin{proof}
Fix a mode and drop its index. By (P, \ref{c:singularity.p-univalence}) and Theorem \ref{c:singularity.ascended-marginals-are-diagonal} the marginal is $\operatorname{diag}(1 - n, n)$ in the number basis with $n := \omega(v^{*}v) \in [0,1]$. The flow acts by $\alpha_t(v) = e^{-it\omega}v$ and $\alpha_t(v^{*}) = e^{it\omega}v^{*}$ (Theorem \ref{c:singularity.modular-structure-completed-pair}(2) at imaginary rapidity, or directly from $Kv = v(K - 1)$). Take $a = v^{*}$, $b = v$: $F(t) = \omega(v^{*}\alpha_t(v)) = e^{-it\omega}n$, and $\omega(\alpha_t(v)v^{*}) = e^{-it\omega}\,\omega(vv^{*}) = e^{-it\omega}(1 - n)$. The gluing condition of (R, \ref{c:singularity.r-regularity-fixed-locus}) reads $e^{p\omega}e^{-it\omega}n = e^{-it\omega}(1 - n)$, that is
\[
e^{p\omega}\,n \;=\; 1 - n \qquad\Longleftrightarrow\qquad n \;=\; \frac{1}{e^{p\omega} + 1} ,
\]
which is the Fermi-Dirac occupancy at inverse temperature $p$, strictly between $0$ and $1$ for every finite $p$ and $\omega > 0$. The pair $a = v$, $b = v^{*}$ gives the same equation with $n$ and $1 - n$ exchanged and $p\omega$ negated, hence the same solution; the pairs $a = b$ give $0 = 0$ on both sides since $v^{2} = 0$. So the marginal is $\operatorname{diag}(1, e^{-p\omega})/(1 + e^{-p\omega}) = e^{-p\omega K}/Z$, the Gibbs state of $K$ at $\beta = p = 2\pi/\kappa$ by Theorem \ref{c:singularity.descended-form-state-independent}, and $T = 1/\beta = \kappa/2\pi$ is Proposition \ref{c:singularity.plane-rate}'s conversion. Faithfulness is $0 < n < 1$: the Gibbs matrix has full rank, so the state is faithful on the mode factor, which for a matrix factor is cyclic and separating. The boundary cases $\beta = \infty$ (unoccupied, $n = 0$) and $\beta = 0$ (tracial, $n = \tfrac12$) fail the gluing condition at every finite nonzero $p$, for $\omega \neq 0$ (at a zero-weight mode the gluing itself degenerates to $n = \tfrac12$, Premise \ref{c:singularity.q-joint-quasi-freeness-s-wave}), so on the modes of nonzero weight (R) excludes them rather than presupposing their exclusion. Universality across modes needs no further input: each mode meets the same section with the same $\kappa$, which is constant over the bifurcation sphere (Proposition \ref{c:singularity.computed} is radius-only, by spherical symmetry). The script verifies the gluing equation, its unique solution, and the failure of both boundary cases, on rational data.

\end{proof}

\phantomsection\label{c:note.singularity.setting-fixing-rate-where-premise-has-4}
The role of the Euclidean rotation in this argument is stated exactly, because it is easy to claim more for it than is true.

\begin{proposition}[The Euclidean rotor, and where it lies]\label{c:singularity.euclidean-rotor-where-it-lies}
Let $\sigma = e_0e_1$ be the boost bivector of the hyperbolic atom in the radial plane, with $\sigma^{2} = +1$, and let $\Aa \cong \Mat_2(\C)$ be the mode algebra, the complexification of the atom.
\begin{enumerate}
\item The boost rotor $R_\lambda = \exp(\lambda\sigma/2)$ continues to imaginary rapidity $\lambda = i\theta$ as $R_{i\theta} = \exp(i\theta\sigma/2) = \cos(\theta/2) + i\sigma\sin(\theta/2)$, a unitary of $\Aa$ whose generator $i\sigma$ squares to $-1$; $R_{2\pi i} = -1$ and $R_{4\pi i} = +1$, so the rotor has period $4\pi$ and its conjugation action period $2\pi$.
\item On the two-dimensional module on which $v, v^{*}$ act as ladder operators, $R_{i\theta} = \operatorname{diag}(e^{-i\theta/2}, e^{i\theta/2})$: a rotation through $2\pi$ multiplies every spinor by $-1$, while its conjugation action on $\Aa$ is the identity.
\item The generator $i\sigma$ is not an element of the real algebra $\Cl(1,3)$. The dual member of the curvature-selected pair, $I\sigma_r = e_0e_1e_2e_3\,e_0e_1 = e_2e_3$, is the \emph{angular} plane: its rotor $\exp(\theta I\sigma_r/2)$ fixes $e_0$, $e_1$ and $\sigma_r$ and rotates $e_2$ into $e_3$, a rotation along the bifurcation sphere, not around it. The Euclidean rotation around the bifurcation sphere is therefore realised by no real rotor of $\Cl(1,3)$, and in particular not by the curvature-selected pair $\{\sigma_r, I\sigma_r\}$; it is realised by the unitary $R_{i\theta}$ of the complexified radial plane.
\end{enumerate}
\end{proposition}
\begin{proof}
(1) With $(i\sigma)^{2} = -\sigma^{2} = -1$ the exponential series splits into cosine and sine, giving the displayed form; it is unitary because $i\sigma$ is skew-adjoint ($\sigma$ is self-adjoint, being $2N - 1$); $R_{2\pi i} = \cos\pi + i\sigma\sin\pi = -1$ and $R_{4\pi i} = +1$. (2) $\sigma = \operatorname{diag}(-1, 1)$ on the module (Theorem \ref{c:singularity.descent-parity-average}'s proof), so $e^{i\theta\sigma/2} = \operatorname{diag}(e^{-i\theta/2}, e^{i\theta/2})$, which is $-1$ at $\theta = 2\pi$; conjugation by a scalar is trivial. (3) The even part of $\Cl(1,3)$ is spanned by $1$, the six frame bivectors and $I$, all with real coefficients, and $i\sigma_r$ is none of these: the bivectors of $\Cl(1,3)$ with square $-1$ are the spatial planes. $I\sigma_r = e_2e_3$ is one line of blade multiplication, and a rotor of the plane $e_2e_3$ fixes every vector orthogonal to that plane, $e_0$ and $e_1$ included, hence fixes $\sigma_r = e_0e_1$; the script verifies $I\sigma_r = e_2e_3$, the action of its rotor on all four frame vectors, and the unitary rotor's algebra on the module, blade-exactly.

\end{proof}

\begin{remark}[What the curvature-selected pair does and does not select]\label{c:note.singularity.what-curvature-selected-pair-does-does}
Corollary \ref{c:singularity.matching-steps-4b-4c-join-equation} uses the pair $\{\sigma_r, I\sigma_r\}$ to identify the modular bivector by equation, and that use stands: $\pm\sigma_r$ are the members with square $+1$. What the pair does \emph{not} contain is the Euclidean partner of the boost. The rotation whose period fixes the temperature lies in the complexified radial plane, and its $2\pi$ is the periodicity of the Euclidean section (Proposition \ref{c:singularity.euclidean-section-period}), a fact about the smooth structure at the bifurcation sphere, not about the eigenplanes of the shape operator. The chain's contribution to the premise is therefore not a location for it inside the curvature but three more modest things: the value of the period, computed from the section (Proposition \ref{c:singularity.euclidean-section-period}); the sign of the statistics, located in the algebra (Proposition \ref{c:singularity.two-kms-functionals-sign}); and the proof that the premise cannot be replaced by anything internal (Theorem \ref{c:singularity.rate-rigidity-temperature-not-algebra-intrinsic}, Proposition \ref{c:singularity.why-calibration-cannot-be-internal}).
\end{remark}

\begin{proposition}[The two KMS functionals, and the sign in $Z$]\label{c:singularity.two-kms-functionals-sign}
On the mode factor with flow $\alpha_t = \operatorname{Ad}(e^{it\omega K})$, let $\Gamma := 1 - 2N = -\sigma$ be the mode's fermion parity, and for homogeneous $b$ write $\epsilon_b = \pm 1$ for even and odd $b$. Fix $\beta > 0$.
\begin{enumerate}
\item The functionals $\phi = \tr(\varrho\,\cdot\,)$ satisfying the \emph{ungraded} KMS identity $\phi(ab) = \phi\bigl(b\,\alpha_{i\beta}(a)\bigr)$ for all $a, b$ are the multiples of $\tr(e^{-\beta\omega K}\,\cdot\,)$; normalised, this is the unique KMS state, with partition function $Z = 1 + e^{-\beta\omega}$ and occupancy $(e^{\beta\omega} + 1)^{-1}$.
\item The functionals satisfying the \emph{graded} identity $\phi(ab) = \epsilon_b\,\phi\bigl(b\,\alpha_{i\beta}(a)\bigr)$ for all $a$ and homogeneous $b$ are the multiples of $\tr(\Gamma e^{-\beta\omega K}\,\cdot\,)$; normalised, this is the parity-twisted functional with $Z_\Gamma = 1 - e^{-\beta\omega}$ and ``occupancy'' $-e^{-\beta\omega}/(1 - e^{-\beta\omega}) < 0$: it is not a state.
\item The $+1$ of the Fermi-Dirac denominator is $Z = 1 + e^{-\beta\omega}$: the two terms are the two states of the module, one per power of $v^{*}$ before $(v^{*})^{2} = 0$, and the sign between them is $+$ for the KMS state and $-$ for the twisted functional. The gluing condition of (R) (Premise \ref{c:singularity.r-regularity-fixed-locus}), which reverses the order of $a$ and $b$ without a sign, is the ungraded identity, and selects the state.
\end{enumerate}
\end{proposition}
\begin{proof}
(1) is Theorem \ref{c:singularity.kms-uniqueness-factor} with $K$ replaced by $\omega K$, and $Z = \tr e^{-\beta\omega K} = e^{\beta\omega/2}(1 + e^{-\beta\omega})$, the prefactor cancelling in the normalised state. (2) Write $\phi = \tr(\varrho\,\cdot\,)$ and $H := \omega K$. The graded identity with $b$ homogeneous is $\tr(b\varrho a) = \epsilon_b\tr(e^{\beta H}\varrho\,b\,e^{-\beta H}a)$ for all $a$, i.e.\ $b\,(\varrho e^{\beta H}) = \epsilon_b\,(e^{\beta H}\varrho)\,b$. With $b = 1$: $\varrho$ commutes with $e^{\beta H}$, hence with $H$, hence is diagonal. Then $c := \varrho e^{\beta H}$ is diagonal, commutes with the even elements and anticommutes with $v$ and $v^{*}$; for $c = \operatorname{diag}(c_0, c_1)$ the relation $cv = -vc$ reads $c_0 = -c_1$, so $c \propto \Gamma$ and $\varrho \propto \Gamma e^{-\beta H}$. Conversely $\Gamma e^{-\beta H}$ satisfies the graded identity because $\Gamma$ commutes with $H$ and $b\Gamma = \epsilon_b\Gamma b$. Its trace is $e^{\beta\omega/2}(1 - e^{-\beta\omega})$ and $\tr(\Gamma e^{-\beta H}N) = -e^{-\beta\omega}e^{\beta\omega/2}$, giving the displayed ratio. (3) is read off (1) and (2): $\tr(e^{-\beta H}) $ and $\tr(\Gamma e^{-\beta H})$ differ exactly in the sign of the occupied term. The script verifies both identities on all sixteen matrix-unit pairs, the failure of each functional to satisfy the other identity, and the negativity of the twisted occupancy, at $e^{-\beta\omega} = \tfrac{9}{16}$.

\end{proof}

\begin{remark}[The Euclidean reading of the sign, flagged]\label{c:note.singularity.euclidean-reading-sign-flagged}
Under continuation to imaginary time the ungraded identity is the antiperiodic boundary condition for the mode's Euclidean two-point function on a circle of circumference $\beta$, and the graded identity the periodic one: this is the Matsubara dichotomy, and it is quoted, not proved. On the smooth Euclidean section a spinor field expressed in the Killing frame, which is the polar frame of the $(\varrho, \tau)$ plane, is multiplied by $R_{i\theta}$ of Proposition \ref{c:singularity.euclidean-rotor-where-it-lies}(2) relative to a frame regular at the origin, and so reverses sign under one full turn: single-valuedness on the section, in the frame in which the modes are defined, is the antiperiodic condition, and hence the ungraded identity of (R). This is the standard spin-structure argument for the thermal, rather than parity-twisted, character of the Hartle-Hawking and Unruh states of fermions, and it is where $R_{2\pi i} = -1$ enters. The chain records it as the reason (R) is stated with the ungraded gluing; it does not derive (R) from it, and the dependence of the boundary condition on the choice of frame is exactly why the statement has to be made in the Killing frame.
\end{remark}

\begin{proposition}[Why the calibration cannot be internal]\label{c:singularity.why-calibration-cannot-be-internal}
\changedmark{BrickRed}{corrected since the deposit}
The only finite rate the Schwarzschild gravitational field on flat spacetime (Definition \emph{Gauge Theory Gravity: the gravitational field on flat spacetime}) supplies to the boost flow is the surface gravity $\kappa = 1/4M$, a horizon quantity: on the singular side the connection rate diverges as $\sqrt{M/2r^{3}}$ (Proposition \ref{c:singularity.connection-bivector-computed}(3)) and offers no finite period to be identified with $2\pi$. (R) (Premise \ref{c:singularity.r-regularity-fixed-locus}) applies across the horizon locus in the chart: the Euclidean section is regular at its tip, which fixes the period $2\pi/\kappa$; for a state that is KMS for the Killing flow, this is also the period at which it is regular across the horizon locus in the Painlev\'e-Gullstrand chart. The horizon locus is met in a cross-section $\Sigma_{t_0} = \{t = t_0,\ r = 2M\}$ of the horizon locus in the Painlev\'e-Gullstrand chart, of area $16\pi M^{2}$, on which the Killing field $K = \hb^{-1}(e_t)$ is null and normal to the horizon, with $(K \cdot D)K = \kappa K$, $\kappa = 1/4M$; the boost flow is the flow of $K$, which carries cross-sections to cross-sections. There, by Propositions \ref{c:singularity.unipotent-ray} and \ref{c:singularity.solution-where-divergence-lies}(3), the gauge field is invertible and unimodular and division does not fail. Hence the transport principle, being a theory of the loci where division fails, makes no statement at the locus where the temperature is calibrated, and no strengthening of the framework's statements at the singular boundary can supply (R).

The bifurcation surface itself, where the boost flow has its fixed points, lies beyond one flat chart (it is the meeting of the four regions of the maximal extension); the chart contains the future horizon, on whose cross-sections the area, the surface gravity and the boost flow are those stated.
\end{proposition}
\begin{proof}
The bifurcation surface named below lies beyond one flat chart, and the statements made there hold on the cross-sections of the horizon locus $r = 2M$ in the Painlev\'e-Gullstrand chart, where Propositions \ref{c:singularity.unipotent-ray} and \ref{c:singularity.solution-where-divergence-lies} are stated (read through Fact \emph{Imported: the dictionary between Gauge Theory Gravity and curved-spacetime language}). The rate statements are the cited propositions (the divergent rate is \ref{c:singularity.connection-bivector-computed}(3)). For the locus: the fixed point set of the Killing boost is where the boost Killing field vanishes, which for Schwarzschild is the bifurcation surface $r = 2M$; the Euclidean section of Proposition \ref{c:singularity.euclidean-section-period}, whose period Theorem \ref{c:singularity.rate-discharge-hypothesis-iii} uses, is the region $r \geq 2M$ and does not contain the singular boundary at all. Proposition \ref{c:singularity.unipotent-ray} shows $\hb$ invertible in closed form there and Proposition \ref{c:singularity.solution-where-divergence-lies}(3) shows $\det\hb = 1$, so the chain's own Step 1 (locate the failure, Step 1 of \ref{c:transport.transport-algorithm}) classifies the horizon as a locus where nothing algebraic happens. A premise imposed there is therefore outside the principle's scope by the principle's own diagnosis.

\end{proof}

\begin{remark}[The cost of (R), and the dictionary]\label{c:note.singularity.cost-r-dictionary}
\changedmark{gray}{reworded}
(R) is not derived from the transport principle, and Proposition \ref{c:singularity.why-calibration-cannot-be-internal} says why it cannot be, so the correct reading is that the chain has localised the last import of its state selection rather than removed it (Step 4e's (Q) is a further import, separately recorded, for the Planck sector alone). It is the same import every derivation of the Hawking temperature makes: Euclidean regularity at the bifurcation point \cite{GibbonsHawking1977, HartleHawking1976}, regularity on the future horizon \cite{Unruh1976}, nonsingularity of the quasi-free state \cite{KayWald1991}, the short-distance condition on the state near the horizon \cite{FredenhagenHaag1990}. What the chain contributes is smaller and exact: the period is computed from the section rather than chosen (Proposition \ref{c:singularity.euclidean-section-period}); the premise is imposed on the marginals, mode by mode, because Theorem \ref{c:singularity.exact-equilibrium-over-determines-boundary-modes} shows that its joint form destroys the correlations the chain needs; the sign of the statistics is located in the algebra (Proposition \ref{c:singularity.two-kms-functionals-sign}); and the premise's irreducibility is a theorem (Theorem \ref{c:singularity.rate-rigidity-temperature-not-algebra-intrinsic}). It does not supply the premise with a location inside the curvature: Proposition \ref{c:singularity.euclidean-rotor-where-it-lies}(3) shows that the Euclidean rotation is not a member of the curvature-selected pair.

The dictionary against the standard alternatives, now as points of an explicitly parametrised line rather than incommensurable constructions: $\beta_k = \infty$ at every mode is the Boulware analogue \cite{Boulware1975}, unoccupied boundary modes with a pure descended state, zero deficit by Theorem \ref{c:singularity.deficit-theorem-cost-descent-exactly} and no flux, excluded by (R); $\beta_k = 2\pi/\kappa$ at every mode is the Hartle-Hawking and Unruh analogue, selected by (R), the product member of that class being the (G) state; $\beta_k$ finite but varying in $k$ is excluded by (R) at the offending modes, and has no geometric modular flow by Corollary \ref{c:singularity.universality-temperature-under-g}; $\beta_k = 0$ is the tracial state, excluded with the other boundary case.

One route to a better-placed premise is worth recording, because it is not the one the singular side offers. The Bisognano-Wichmann theorem does not assume the factor $2\pi$: it derives it, from the spectrum condition together with covariance and locality \cite{BW1976, Sewell1982}. The ascent of Step 3 lands in $\Cl(2,4)$, the conformal algebra, whose parabolic stabiliser of $\nInf$ contains exactly the translations such a spectrum condition needs, so the premise could be restated as positivity of the energy of the boundary modes with respect to a generator the ascent itself installed, rather than as smoothness of a Euclidean section. That would still be an import, positivity of energy being no Clifford fact, but it would be stated in the algebra the chain produces. Transcribing the theorem needs an action of the relevant subgroup on the boundary modes, a wedge and locality structure, and analyticity of the state, each of which is genuine work and any of which could fail; the item belongs with the correspondence paper T9 \cite{ButlerT9}, on the conformal spinor space of $\Cl(2,4)$, rather than here.
\end{remark}

\begin{theorem}[Purity is an output: the partner is the commutant]\label{c:singularity.purity-output-partner-commutant}
Let $\omega$ be a faithful state on the mode factor with $\varrho = \operatorname{diag}(\lambda_0^{2}, \lambda_1^{2})$, $\lambda_i > 0$. In its GNS representation: the space is $\C^{2}\otimes\C^{2}$ and the cyclic vector is $\Omega = \sum_i \lambda_i|ii\rangle$, cyclic and separating for $\pi(\Aa) = \Aa\otimes 1$; the commutant $\pi(\Aa)' = 1 \otimes \Aa$ is the partner copy and equals $J\pi(\Aa)J$; any two purifications of $\omega$ differ by a unitary on the purifying factor alone; and when $\omega$ is the Gibbs state at $\beta$, $\Omega$ is the thermofield double with $\tan\theta = e^{-\omega_k/2T}$, which is exactly the vector Lemma \ref{c:singularity.schmidt-forcing-proved-outright} produces. Hypothesis (i) of Theorem \ref{c:singularity.modular-descent}, the existence of the partner, and the thermofield form are therefore consequences of the one-sided state rather than independent choices.
\end{theorem}
\begin{proof}
The first two clauses are Lemma \ref{c:singularity.modular-theory-matrix-factor-proved-outright} read in the direction of construction rather than of analysis: that lemma computes $\Delta = \rho\otimes\bar\rho^{-1}$ and $J|ij\rangle = |ji\rangle$ for this very $\Omega$ and establishes that $\Omega$ is cyclic and separating, the identification of the commutant being Tomita's theorem, whose finite-dimensional content is that lemma's polar decomposition. For the third: a purification of $\varrho$ has the $\lambda_i$ as its Schmidt coefficients, so it is $(1\otimes W)\Omega$ for an isometry $W$ between purifying spaces. For the fourth, substituting $\varrho = \operatorname{diag}(1, e^{-\omega_k/T})/Z$ gives $\lambda_1/\lambda_0 = e^{-\omega_k/2T} = \tan\theta$. The script instantiates the purification, the modular eigenvalues and the angle at the rational triple $(\tfrac35, \tfrac45)$.

\end{proof}

\begin{theorem}[Exact equilibrium over-determines the boundary modes]\label{c:singularity.exact-equilibrium-over-determines-boundary-modes}
Let a state satisfy the geometric modular condition exactly. Then $\varrho \propto e^{-\beta K_{\mathrm{geo}}} = \bigotimes_k e^{-\beta\omega_k K_k}$ is a product across modes, and consequently: every inter-mode correlation vanishes, $C(\varrho) = \sum_k S(\varrho_k) - S(\varrho) = 0$; the descent deficit vanishes, $\rel{\varrho}{\EN\varrho} = 0$ by Theorem \ref{c:singularity.deficit-theorem-cost-descent-exactly}, the state being already a fixed point of $\EN$; and no information whatever is carried by the inter-quanta correlations of Corollary \ref{c:singularity.deficit-carried-quanta-correlations}, the two levels agreeing exactly. The exact geometric modular condition is therefore available only in the stationary idealisation, and the correlations on which Steps 5 and 6b depend are precisely the deviation from it.
\end{theorem}
\begin{proof}
The product form is Theorem \ref{c:singularity.kms-uniqueness-factor} applied to $K_{\mathrm{geo}}$, whose exponential factorises because the $K_k$ commute. A product state has $S(\varrho) = \sum_k S(\varrho_k)$, so $C(\varrho) = 0$. A product of diagonal matrices is diagonal in the joint number basis, hence a fixed point of $\EN$, and the deficit theorem (\ref{c:singularity.deficit-theorem-cost-descent-exactly}) gives $S(\EN\varrho) - S(\varrho) = 0$. The last clause is Corollary \ref{c:singularity.deficit-carried-quanta-correlations}, which identifies the deficit with $C(\varrho) - C(\EN\varrho)$. The script verifies the fixed-point and factorisation clauses in exact arithmetic against an entangled even control failing both.

\end{proof}

\begin{remark}[Why this rescues rather than damages the chain]\label{c:note.singularity.why-rescues-rather-than-damages-chain}
Read carelessly, Theorem \ref{c:singularity.exact-equilibrium-over-determines-boundary-modes} is an obstruction: the state that exact thermality requires is the state with no information in it. Read carefully, it is the reason the order of argument in this subsection matters. Had per-mode thermality been derived from the exact equilibrium condition, as hypothesis (ii) suggested, then thermality and correlation would have been in direct competition and the chain could not have had both. Theorem \ref{c:singularity.ascended-marginals-are-diagonal} shows that per-mode thermality does not rest on equilibrium at all: it rests on parity, which survives an arbitrary non-stationary history. The state class the chain actually needs is
\[
\bigl\{\ \text{even states, per-mode populations at } \beta = 2\pi/\kappa, \ \text{correlations free}\ \bigr\},
\]
in which every single quantum is exactly thermal and the joint state of boundary modes and radiation together is free to carry everything the evaporation writes. Where it is carried is sharpened by the rate model and should be recorded here: under (I) and (L$_E$) the window marginal itself is fixed to be the product member at every stage (Theorem \ref{c:singularity.eviction-window-thermal-product-radiation-entropy}(1)), its internal correlation and descent deficit identically zero, and the writing lies across the window/radiation cut and among the quanta of the radiation, not inside the boundary-mode marginal. That this class is where (P) and (R) put the state, and that the dynamics of the boundary modes keeps it there, is Corollary \ref{c:singularity.chains-state-class-stability} below; the finite-mode model of the section ``The finite-mode model: the dissolution step by step'' lies in it, and Lemma \ref{c:singularity.parity-protection} is the model's proof that it does. The theorem also settles a question the chain could otherwise have been asked indefinitely: no sharpening of Step 4d will ever supply Step 6b's correlations, because no equilibrium principle can.
\end{remark}

\begin{proposition}[Number-diagonal dynamics; the model's scrambling derived]\label{c:singularity.number-diagonal-dynamics-models-scrambling-derived}
Under (L$_E$) every generator available to the internal dynamics of the boundary modes lies in the abelian algebra generated by $\{\sigma_k\}$, hence is diagonal in the joint number basis and its exponential is a diagonal unitary. Consequently each mode's occupation probabilities are constants of the motion of the boundary modes, so populations set at injection by the local geometry do not drift; and the diagonal scrambling of the finite-mode model, the controlled-$Z$ circuit of the section ``The model'', is not a modelling convenience but the general form of what (L$_E$) permits.
\end{proposition}
\begin{proof}
The $\sigma_k = 2N_k - 1$ commute and are simultaneously diagonal in the joint number basis; a polynomial in commuting diagonal operators is diagonal and its exponential is a diagonal unitary, which multiplies each joint-number amplitude by a phase and so leaves every joint-number probability, hence every marginal population, unchanged. With Theorem \ref{c:singularity.ascended-marginals-are-diagonal} fixing the off-diagonal entries at zero, each $\varrho_k$ is constant in time. The script verifies the invariance for a controlled-$Z$ circuit on the model's mode space and exhibits an orthogonal, parity-preserving hopping rotation that does change the marginal populations, which is exactly the generator (L$_E$, \ref{c:singularity.l-boundary-mode-locality}) excludes.

\end{proof}

\begin{corollary}[The chain's state class, and its stability]\label{c:singularity.chains-state-class-stability}
Let $\mathcal{C}$ be the set of even states of the boundary modes every single-mode marginal of which is the Gibbs state of $K_k$ at $\beta = 2\pi/\kappa$. Then:
\begin{enumerate}
\item under (P) and (R) the state of the boundary modes lies in $\mathcal{C}$, and every member of $\mathcal{C}$ satisfies (P) and (R);
\item $\mathcal{C}$ is invariant under every internal dynamics that (L$_E$) permits: a state placed in $\mathcal{C}$ at injection stays in $\mathcal{C}$ thereafter, whatever the boundary modes write;
\item $\mathcal{C}$ contains the product state $e^{-\beta K_{\mathrm{geo}}}/Z$, which is its only member satisfying (G), and it contains states of arbitrarily large total correlation $C(\varrho)$ and descent deficit; (R) constrains none of the correlations.
\end{enumerate}
So the premises fix the marginals once and for all, and leave the joint state exactly as free as Steps 5 and 6b need it; the rate model then exercises that freedom across the window/radiation cut, its window marginal being itself the product member at every stage (Theorem \ref{c:singularity.eviction-window-thermal-product-radiation-entropy}(1)).
\end{corollary}
\begin{proof}
(1) The first clause is Theorem \ref{c:singularity.rate-discharge-hypothesis-iii}. For the second: a member of $\mathcal{C}$ is even, hence satisfies (P, \ref{c:singularity.p-univalence}); and its marginals are Gibbs at $\beta = 2\pi/\kappa$, so by Theorem \ref{c:singularity.kms-uniqueness-factor} each satisfies the ungraded KMS identity at that $\beta$, which for the ladder pairs is the gluing condition of (R, \ref{c:singularity.r-regularity-fixed-locus}) (the computation in the proof of Theorem \ref{c:singularity.rate-discharge-hypothesis-iii}, read backwards). (2) A dynamics permitted by (L$_E$, \ref{c:singularity.l-boundary-mode-locality}) is conjugation by a diagonal unitary, which commutes with the total parity $\Gamma$ (both are diagonal), so it preserves evenness; and by Proposition \ref{c:singularity.number-diagonal-dynamics-models-scrambling-derived} it preserves every marginal. (3) The product state is even and has the right marginals, and by Theorem \ref{c:singularity.exact-equilibrium-over-determines-boundary-modes} it is the only (G) state with them. For the correlated members: take any even state with the right marginals, for instance the model's mode-space state after an arbitrary controlled-$Z$ pattern with all pairs at $\beta = 2\pi/\kappa$, whose marginals Lemma \ref{c:singularity.parity-protection} fixes and whose total correlation grows with the number of pairs. (R), being a condition on the marginals alone, is insensitive to which member of $\mathcal{C}$ the state is. The script verifies (1) to (3) on a two-mode entangled even state with the prescribed marginals: it is in $\mathcal{C}$, it is not a product, its marginals and evenness survive a controlled-$Z$, and the product member is recovered by $\EN$.

\end{proof}

\begin{remark}[The step's status, and what is still open]\label{c:note.singularity.steps-status-what-still-open}
Collecting: hypothesis (i) is Theorem \ref{c:singularity.purity-output-partner-commutant}; hypothesis (ii) is not needed for the form at all (Theorem \ref{c:singularity.descended-form-state-independent}), and its per-mode content follows from (P) (Theorem \ref{c:singularity.ascended-marginals-are-diagonal}); hypothesis (iii) follows from (R) (Theorem \ref{c:singularity.rate-discharge-hypothesis-iii}, whose gluing condition is stated by evaluation and does not presuppose a modular group); and the value of $\TH$ follows from (R) with the period computed in Proposition \ref{c:singularity.euclidean-section-period}, and provably from nothing internal (Theorem \ref{c:singularity.rate-rigidity-temperature-not-algebra-intrinsic}, Proposition \ref{c:singularity.why-calibration-cannot-be-internal}). The exact equilibrium condition (G) is not a premise of this step: it is examined and found too strong for the evaporating boundary modes (Theorem \ref{c:singularity.exact-equilibrium-over-determines-boundary-modes}), the class the chain uses is Corollary \ref{c:singularity.chains-state-class-stability}'s, and the one place the chain later adopts the (G) member is Step 4e's (Q), for the equilibrium reading of the Planck sector alone. The status line in the table accompanying \ref{c:singularity.imported-equivalence-einstein-equations} is accordingly \emph{proved from (P), (R)} rather than \emph{up to a state}, where (R) is the KMS condition at the geometric period, imposed mode by mode; the premises of the later steps, (Q) in Step 4e, (C1) to (C3) in Step 6a and (L), (U), (C), (I), (L$_E$) in Step 6b, are named where those steps are proved.

Four things remain open and are recorded as such. (a) (R) is an import, and Proposition \ref{c:singularity.why-calibration-cannot-be-internal} shows it cannot be removed by strengthening the framework's statements at the singular boundary; the open route is the spectrum-condition restatement of Remark \ref{c:note.singularity.cost-r-dictionary}, which belongs with T9 \cite{ButlerT9}. (b) (L$_E$) is a statement about which generators the localisation theorems supply; the class-relative no-go is Theorem \ref{c:singularity.class-relative-no-go-no-supplied} (the section ``The metric-level sheaf, and the hopping exclusion''): no supplied generator hops, the supplied class being the second quantisations of function-linear maps of the mode line bundle and the products of the parity sections, an ultralocal class of multiplication operators. What would settle the premise outright is Remark \ref{c:note.singularity.declared-residual-size}'s residual question, whether the physical dynamics on $E$ carries a differential or a pairing term; the necessity witness of Corollary \ref{c:singularity.envelope-endpoint-why-boundary-mode-locality} is itself a pairing rotation, which the theorem's criterion does not exclude; Proposition \ref{c:singularity.l-stabiliser-tangent} reads the two residual types in $\mathfrak{so}(2\nu)$. (c) Theorems \ref{c:singularity.kms-uniqueness-factor}, \ref{c:singularity.rate-rigidity-temperature-not-algebra-intrinsic} and \ref{c:singularity.purity-output-partner-commutant} and Proposition \ref{c:singularity.two-kms-functionals-sign} are proved in the matrix factor, which is where this step works, and their multi-mode applications, Corollary \ref{c:singularity.universality-temperature-under-g} and Theorem \ref{c:singularity.exact-equilibrium-over-determines-boundary-modes}, extend them to finite collections of modes factor by factor, and where the section ``Step 4e: compute in the completion; from fermions to photons'' invokes them on the $\mathbb{Z}$-array the reading is through finite mode subsets, on which every quantity involved is defined, no global entropy being claimed on the infinite array; in the field-theoretic setting the horizon algebra is a type III$_1$ factor, the Gibbs form is unavailable, and the KMS condition must be used directly, in the simplex form of Fact \ref{c:singularity.imported-general-form}, with \cite{KayWald1991} as the template. (d) Theorem \ref{c:singularity.rate-discharge-hypothesis-iii} treats the sphere directions of $E$ as inert, correct for Schwarzschild by symmetry; the corresponding statement over a rotating horizon, where $\kappa$ is still constant but the Killing flow is not the static one, is untouched, and belongs with the Kerr programme of \cite{ButlerT4, ButlerT6}.
\end{remark}

\section{Step 4e: compute in the completion; from fermions to photons}

\phantomsection\label{c:note.singularity.setting-step-4e-compute-completion-fermions}
The chain's pairs are fermionic; the directly derived spectrum is Fermi-Dirac. The observed Hawking flux is dominated by photons, whose occupancy is Bose-Einstein. The bridge is bosonisation, and its setting is not imported: the boundary modes form a one-dimensional fermion array with the boost as its clock, which is exactly the setting of the boson-fermion correspondence \cite{PressleySegal1986}, met from the physics side by the standard s-wave reduction of near-horizon field theory \cite{RW2005}. \emph{Scope declaration, once for the section:} every statement here concerns the s-wave-reduced sector; the four-dimensional lift is a point of contact at the end, flagged as assembly. One matter of notation belongs in the declaration. The paper carries three index sets for the boundary modes: the points of $E \cong \R \times S^{2}$ (Fact \ref{c:singularity.imported-one-sided-completion-singularity}), the $\mathbb{Z}$-array of this section with weights $\omega_k = \kappa k$, and the Planck cells of the stretched singularity (Step 6a). Axiom (C1) names the identification of the first with the third; the identification of the first with the second, the aggregation of $E$'s sphere factor into a single s-wave array with a filled sea, is likewise an identification and not a theorem, made here as part of the s-wave reduction and flagged as such.

\phantomsection\label{c:note.singularity.setting-step-4e-compute-completion-fermions-2}
(Q) is a genuine fourth premise, and it must be recorded that it stands in explicit tension with the principal content of Steps 5 and 6b: for the given marginals the quasi-free state is the product member of Corollary \ref{c:singularity.chains-state-class-stability}'s class, the unique member satisfying (G), and Theorem \ref{c:singularity.exact-equilibrium-over-determines-boundary-modes} rejects exactly that member as carrying none of the correlations those steps need. What (Q) yields is therefore the \emph{equilibrium reading} of the flux: the Planck sector below is a statement about the stationary idealisation of the boundary modes, not about the evaporating hole, for whose state the exact formula fails, the sum acquiring connected inter-mode correlators that (P) and (R) leave free (Corollary \ref{c:singularity.chains-state-class-stability}(3)). One consequence of the rate model belongs here rather than being discovered against it: under (I) and (L$_E$) the window marginal is exactly this product at every stage (Theorem \ref{c:singularity.eviction-window-thermal-product-radiation-entropy}(1)), so for the model's mode space the product clause of (Q) holds as a theorem (the weight assignment, the stationarity and the $\mathbb{Z}$-array remain (Q)'s own), and the tension is relocated rather than dissolved: the product form empties the \emph{internal} correlations of the boundary modes, and the correlations Steps 5 and 6b depend on are carried across the window/radiation cut, to which Theorem \ref{c:singularity.exact-equilibrium-over-determines-boundary-modes}'s rejection does not apply. the table accompanying \ref{c:singularity.imported-equivalence-einstein-equations}'s status line for the step carries the premise.

\begin{theorem}[The bosonisation bridge, s-wave sector]\label{c:singularity.bosonisation-bridge-s-wave-sector}
\changedmark{BrickRed}{corrected since the deposit}
Let $\{\psi_k, \psi^\dagger_k\}_{k \in \mathbb{Z}}$ be the boundary modes, as fermion operators on the filled-sea Fock space (the infinite filled-sea array whose algebraic model is the completed tower limit of \cite[Thm.~5.6]{ButlerT3}, cited for the structure; its provenance here is the s-wave aggregation declared above, not a tower, the mod-8 clock having stopped after one ascent), and define the currents $a_m := \sum_k :\psi^\dagger_k\,\psi_{k+m}:$ with $:\;:$ normal ordering at the sea level; with the sea filled below the level, $a_m$ annihilates the sea for $m > 0$ and $a_{-m} = a_m^{\dagger}$ creates particle-hole pairs. (The opposite index convention flips the sign of the central term; the verification script fixes the convention by computation.) Then:
\begin{enumerate}
\item (Current algebra.) $[a_m, a_n] = m\,\delta_{m+n,0}\,\id$ on the infinite array; on a finite window the commutator equals $m\,\delta_{m+n,0}\,\id$ plus an \emph{exactly identified} edge operator ($B_m - T_m - m$ in the proof's notation), supported on the bottom and top $m$ sites and vanishing on all bulk states: the number-conserving bilinears satisfy the canonical commutation relations of a free boson. A boson is a pair of the pair.
\item (KMS restriction.) Each bilinear $\psi^\dagger_k\psi_{k+m}$ carries a definite boost weight, fixed by the difference $\omega_{k+m} - \omega_k$ of its two modes' weights. Under (Q)'s weight clause, $\omega_k = \kappa k$, that difference is $\kappa m$ for every $k$, so each current $a_m$ is itself a weight vector and the current subalgebra is globally invariant under the modular flow of Theorem \ref{c:singularity.modular-structure-completed-pair}; for weights that are not affine in $k$ the summands of $a_m$ rotate at different rates and the span of the currents need not be invariant, so this invariance uses (Q). Under (Q) the joint state is KMS for that flow, and a KMS state restricted to a flow-invariant subalgebra is KMS at the same temperature for the restricted flow. Without (Q) no joint KMS state is available: (R) is imposed on the marginals alone, and Theorem \ref{c:singularity.exact-equilibrium-over-determines-boundary-modes} is the reason its joint form is rejected elsewhere in the chain.
\item (The Planck factor.) Under (Q), with the marginal occupancies $n_k = (e^{\beta\omega_k} + 1)^{-1} = (e^{\beta\kappa k} + 1)^{-1}$ at the premise's linear weights, the current occupancy is, exactly,
\[
\langle a_{-m} a_m \rangle \;=\; \sum_{k \in \mathbb{Z}} n_{k+m}\,(1 - n_k) \;=\; \frac{m}{e^{\beta\kappa m} - 1} \qquad (m > 0) ,
\]
with no vacuum piece in this convention ($a_m$ annihilates the sea): Bose-Einstein occupancy at exactly $\TH$. The Planck spectrum of the bosonic Hawking flux is derived, not conjectured, for the s-wave sector, in the equilibrium reading (Q) fixes; for a general member of the chain's state class the left side acquires the connected correlators (Q) sets to zero, and no exact spectrum is claimed.
\end{enumerate}
The single assembly-level ingredient in the section is the charge-zero Fock isomorphism (the charge-zero sector of the fermionic Fock space is canonically the bosonic Fock space, with fermions recovered as vertex operators), cited to \cite{PressleySegal1986} and used only for the completeness statement that the two descriptions are one algebra; parts (1)-(3) are proved outright.
\end{theorem}
\begin{proof}
(1) \emph{Finite window first, exactly.} On the window $W = \{-K, \dots, K-1\}$ set $a_m = \sum_{k \in S_m} \psi^\dagger_k \psi_{k+m}$ with $S_m = \{k : k, k+m \in W\}$; normal-ordering constants are scalars and drop out of every commutator, so they play no role here. The CAR give, for single bilinears,
\[
[\psi^\dagger_k \psi_{k+m},\; \psi^\dagger_l \psi_{l+n}] = \delta_{k+m,\,l}\,\psi^\dagger_k\psi_{l+n} - \delta_{l+n,\,k}\,\psi^\dagger_l\psi_{k+m} .
\]
For $n = -m$ ($m > 0$), summing over $k \in S_m$ and $l \in S_{-m}$ and writing $N_j = \psi^\dagger_j\psi_j$:
\[
[a_m, a_{-m}] \;=\; \sum_{k \in S_m} \bigl(N_k - N_{k+m}\bigr)
\;=\; \underbrace{\sum_{j=-K}^{-K+m-1} N_j}_{=:\,B_m} \;-\; \underbrace{\sum_{j=K-m}^{K-1} N_j}_{=:\,T_m} ,
\]
an \emph{exact operator identity}: the interior terms cancel in the telescoping sum and only the bottom and top $m$ sites survive. On any state that agrees with the filled sea on those edge sites, $B_m = m$ and $T_m = 0$, so $[a_m, a_{-m}] = m\,\id$ there; the edge operators are the entire finite-size correction, and they vanish on bulk states at every window size, not merely asymptotically. For $m + n \neq 0$ the same computation cancels all interior terms against each other and leaves an edge-supported remainder, zero on bulk states. The infinite array (modelled on the tower limit of \cite[Thm.~5.6]{ButlerT3}) has no edges, and the identities hold outright.

(2) A bilinear $\psi^\dagger_k\psi_{k+m}$ is an eigenvector of the modular flow with definite weight (both factors carry definite weight by Theorem \ref{c:singularity.modular-structure-completed-pair}(2)), and its weight is fixed by the difference $\omega_{k+m} - \omega_k$. The current $a_m$ sums these bilinears over $k$, so it is an eigenvector exactly when that difference does not depend on $k$, that is when the weights are affine in $k$. This is where the weight clause of (Q, \ref{c:singularity.q-joint-quasi-freeness-s-wave}) enters: with $\omega_k = \kappa k$ the difference is $\kappa m$ for every $k$, each $a_m$ is a weight vector, and the linear span of the currents is invariant. For weights not affine in $k$ the summands of $a_m$ rotate at different rates and the span need not be invariant. The restriction statement is definitional: if $(\omega, \sigma_t)$ is KMS at $\beta$ on an algebra and $A$ is a globally invariant subalgebra, the correlation functions of elements of $A$, their analytic continuations, and the KMS boundary identity are computed in $A$ exactly as in the ambient algebra.

(3) Under (Q, \ref{c:singularity.q-joint-quasi-freeness-s-wave}), Wick's theorem gives $\langle a_{-m} a_m\rangle = \sum_{k} n_{k+m}(1 - n_k)$: the four-point functions factorise because (Q) sets every connected part to zero, which is exactly what a general member of the chain's state class denies. With the present convention the sum needs no vacuum subtraction: at $\beta \to \infty$ the summand $\mathbf{1}[k + m < 0]\,\mathbf{1}[k \geq 0]$ is empty for $m > 0$, consistent with $a_m$ annihilating the sea. Write $x := e^{-\beta\kappa} \in (0,1)$, $n_k = x^k/(1 + x^k)$, and
\[
g(k) := n_{k+m}(1 - n_k) = \frac{x^{k+m}}{(1 + x^{k+m})(1 + x^k)} \;\geq 0.
\]
The two-sided series $S := \sum_{k \in \mathbb{Z}} g(k)$ converges absolutely: $g(k) \leq x^{k+m}$ for $k \geq 0$ and $g(k) \leq x^{-k}$ for $k \leq -m$ (bound each Fermi factor by $1$ and the other by its numerator). Two exact pointwise identities finish the computation. \emph{Detailed balance:} $n_k(1 - n_{k+m}) = x^{-m}\, g(k)$, read off from the displayed formula by exchanging which factor keeps its numerator. \emph{Difference identity:} $g(k) - n_k(1 - n_{k+m}) = n_{k+m} - n_k$, by expanding both sides over the common denominator. Summing the difference identity over $\mathbb{Z}$ (all three series converge absolutely) and evaluating the telescoping middle sum by its partial sums,
\[
\sum_{k=-K}^{K} (n_{k+m} - n_k) \;=\; \sum_{j=K+1}^{K+m} n_j \;-\; \sum_{j=-K}^{-K+m-1} n_j \;\xrightarrow[K \to \infty]{}\; 0 - m = -m ,
\]
since $n_j \to 0$ as $j \to +\infty$ and $n_j \to 1$ as $j \to -\infty$. Therefore $S - x^{-m} S = -m$, i.e.
\[
S \;=\; \frac{m}{x^{-m} - 1} \;=\; \frac{m\,x^m}{1 - x^m} \;=\; \frac{m}{e^{\beta\kappa m} - 1} ,
\]
which is the assertion: Bose-Einstein occupancy, exactly, with the detailed-balance relation supplying the physical content. The deposited script verifies the two pointwise identities, the re-indexing, and the value against exact rational tail bounds; and, for (Q)'s role, that the Planck summand is the exact expectation $\langle n_{k+m}(1 - n_k)\rangle$ on the product member of the class while a correlated member with the same marginals shifts it by exactly the connected correlator.

\end{proof}

\begin{remark}[The four-dimensional lift, flagged]\label{c:note.singularity.four-dimensional-lift-flagged}
What the theorem delivers is the horizon-sector spectrum in the s-wave reduction. The lift to four dimensions is partial-wave reconstruction, and the greybody filtering is inherited verbatim from the exterior geometry, which coincides with classical general relativity on the invertible region; the emitted flux is transmission times occupancy, Fermi-Dirac for fermions and Bose-Einstein for bosons, at exactly $\TH$. Both statements are assembly-level and are points of contact here, not claims of this paper; the interface computation (horizon-sector state feeding exterior modes in the standard way) is the same assumption every derivation of Hawking flux makes, and is flagged as such.
\end{remark}

\section{Step 5: descend, and state what descent costs}

\phantomsection\label{c:note.singularity.setting-step-5-descend-state-what}
One mode is not the boundary mode space. The chain's principal content, the two-level entropy computation of Step 6b and the correlation transfer of the model, concerns the descent of many modes at once, and the multi-mode form of the pinching is where that computation becomes an equation. This step defines it, fixes the scope of a coincidence that Step 4d's identification anticipates, and quantifies the descent exactly: the algorithm's Step 5, ``descend, and state what descent costs'' \cite[\S 8.1, Step 5]{ButlerT1}, executed at the level of states. The single-mode descent it composes is the operator $E$ installed in Step 4a, whose projectors Step 4c identified with the boost's spectral pair.

\begin{definition}[The multi-mode descent]\label{c:obj.singularity.multi-mode-descent}
For $N$ completed pairs $(v_k, v_k^*)$ with commuting occupation projectors $P_k := v_k^* v_k$, and for each $\varepsilon \in \{0,1\}^N$ the joint occupation projector $\Pi_\varepsilon := \prod_k \bigl(\varepsilon_k P_k + (1 - \varepsilon_k)(1 - P_k)\bigr)$ (orthogonal, complete: $\sum_\varepsilon \Pi_\varepsilon = 1$), the \emph{multi-mode descent} is the pinching
\[
\EN(x) \;:=\; \sum_{\varepsilon \in \{0,1\}^N} \Pi_\varepsilon\, x\, \Pi_\varepsilon ,
\]
the composition of the mode-wise maps of Proposition \ref{c:singularity.descent-what-it-erases}. In the joint number representation, $\EN$ retains every matrix element diagonal in the joint occupation basis and erases every off-diagonal one: complete dephasing in the joint number basis. Its kernel is spanned by the operators carrying at least one off-diagonal mode factor, which is precisely the sector in which inter-mode coherence lies; its range is the algebra of joint-number-diagonal operators, the classical algebra of occupation records. That this range, and not the larger total-parity-even algebra, is the metric level on many modes is not a consequence of the grading: it is the composite argument of Proposition \ref{c:singularity.p-reduced-metric-level-reading-many}, stated at the end of this step, where its effect is visible. One word of disambiguation, on which everything depends: ``pair'' here means one completed ladder pair, one \emph{mode}, in Rule 1's sense; the ``pairs'' of the section ``The finite-mode model: the dissolution step by step'''s model are Hawking pairs, each two modes, one radiation and one window, so the model's $n$ pairs are $2n$ modes and $\EN$ there pinches all $2n$, one $\sigma_k$ per mode and never one per Hawking pair (on that wrong reading $\sigma$ would be $Z_{R_k}Z_{L_k}$, the strictness witness would fail, and the deficit would be exactly zero).
\end{definition}

\begin{proposition}[Pinching against partial trace: the scope of a coincidence]\label{c:singularity.pinching-against-partial-trace-scope-coincidence}
\leavevmode
\begin{enumerate}
\item[(i)] (One mode: they agree.) For any global state whose Schmidt basis across the pair/partner cut is the number basis (as boost invariance forces in Lemma \ref{c:singularity.schmidt-forcing-proved-outright}, and as the support constraint forces in the model, Lemma \ref{c:singularity.parity-protection}), the mode marginal $\tr_{\mathrm{partner}}\,\rho$ is diagonal in the number basis and hence a fixed point of $E$: partial trace and pinching coincide, mode by mode.
\item[(ii)] (Many modes: they differ, by exactly the inter-mode coherence.) On $N \geq 2$ modes, $\tr_L$ and $\EN$ are different maps: $\tr_L$ retains the radiation's off-diagonal joint-number matrix elements, $\EN$ annihilates them, and the two agree on every single-mode marginal. The two maps answer different questions: $\tr_L$ asks what the radiation alone contains; $\EN$ asks what the metric level can carry. Conflating them is harmless for single-mode spectra and fatal for the computation of correlations, which is why the conflation is scoped once, in this step: Step 4d's identification uses clause (i), and Step 6b's computation uses clause (ii).
\end{enumerate}
\end{proposition}
\begin{proof}
(i) Diagonal matrices are fixed points of $E$ by the display of Proposition \ref{c:singularity.descent-what-it-erases}. (ii) Definition \ref{c:obj.singularity.multi-mode-descent}: any state with a nonzero joint-number coherence separates the two maps; the minimal instance, a two-pair state whose pair marginal carries the coherence $c_1 s_1 \gamma$, is computed exactly in Proposition \ref{c:singularity.purity-correlation-transfer-recovery}, and its descent, which annihilates that coherence while fixing both populations and both single-mode marginals, is computed exactly in Proposition \ref{c:singularity.two-levels-finite-model}. The script verifies both clauses.

\end{proof}

\begin{theorem}[The deficit theorem: the cost of descent, exactly]\label{c:singularity.deficit-theorem-cost-descent-exactly}
Let $E$ be the pinching by any complete family of orthogonal projectors ($\EN$ included) and $\rho$ a density matrix. Then $E$ is idempotent, trace-preserving, and self-adjoint for the Hilbert-Schmidt pairing, and
\[
S\bigl(E(\rho)\bigr) - S(\rho) \;=\; \rel{\rho}{E(\rho)} \;\geq\; 0 ,
\]
with $S$ the von Neumann entropy and $D$ the relative entropy, and with equality if and only if $\rho = E(\rho)$. The level discrepancy of Step 6b is this quantity: not a slogan but a relative entropy, computed once and for all.
\end{theorem}
\begin{proof}
Idempotency and trace preservation are the orthogonality and completeness of the $\Pi_\varepsilon$. Self-adjointness: $\tr(E(a)\,b) = \sum_\varepsilon \tr(\Pi_\varepsilon a \Pi_\varepsilon b) = \sum_\varepsilon \tr(a\, \Pi_\varepsilon b \Pi_\varepsilon) = \tr(a\,E(b))$ by cyclicity of the trace. The key step: $E(\rho)$ commutes with every $\Pi_\varepsilon$, hence so does every function of it, hence $E\bigl(\log E(\rho)\bigr) = \log E(\rho)$; self-adjointness with $b = \log E(\rho)$ then gives
\[
\tr\bigl(\rho \log E(\rho)\bigr) \;=\; \tr\bigl(E(\rho) \log E(\rho)\bigr) ,
\]
whence $\rel{\rho}{E\rho} = \tr\rho\log\rho - \tr\rho\log E\rho = -S(\rho) + S(E\rho)$. Nonnegativity and the equality case are Klein's inequality, the single classical input of the theorem \cite[\S 11.3]{NielsenChuang}. The script verifies the structural clauses on a generic rational matrix and the key identity, in moment form $\tr\bigl(\rho\,(E\rho)^m\bigr) = \tr\bigl(E\rho\,(E\rho)^m\bigr)$, on the model's scrambled two-pair state, in exact arithmetic.

\end{proof}

\begin{corollary}[The deficit is carried by the quanta correlations]\label{c:singularity.deficit-carried-quanta-correlations}
Write $\rho_k$ for the single-mode marginals of $\rho$ and $C(\rho) := \sum_k S(\rho_k) - S(\rho)$ for the total correlation. If every $\rho_k$ is diagonal in its number basis (as Theorem \ref{c:singularity.modular-descent} guarantees for thermofield injections and Lemma \ref{c:singularity.parity-protection} for the model), then the single-mode marginals of $\EN(\rho)$ equal those of $\rho$, and
\[
\rel{\rho}{\EN\rho} \;=\; C(\rho) - C(\EN\rho) :
\]
the deficit is the total correlation of the ascended state minus the classical, number-basis correlation that survives descent. No single-quantum quantity contributes anything: the cost of the descent consists entirely of inter-quanta correlations. In particular, for a globally pure $\rho$, $S(\rho) = 0$ and the deficit is the whole descended entropy, $\rel{\rho}{\EN\rho} = S(\EN\rho)$.
\end{corollary}
\begin{proof}
Partial trace over the modes other than $k$ commutes with the pinchings on those modes, so $\tr_{\neq k}\,\EN\rho = E_k(\rho_k) = \rho_k$ when $\rho_k$ is diagonal. Then $C(\rho) - C(\EN\rho) = \bigl[\sum_k S(\rho_k) - S(\rho)\bigr] - \bigl[\sum_k S(\rho_k) - S(\EN\rho)\bigr] = S(\EN\rho) - S(\rho)$, which is the deficit by Theorem \ref{c:singularity.deficit-theorem-cost-descent-exactly}.

\end{proof}

\begin{remark}[Where the descent acts, and Hawking's loss derived]\label{c:note.singularity.where-descent-acts-hawkings-loss-derived}
The chain must still supply the reader one sentence of locality, and here it is. The localisation theorem \cite[Thm.~4.4, Cor.~4.5]{ButlerT4}, in the form the Schwarzschild composite gives it at the black hole \cite[Thms 8.7, 8.12]{ButlerT4}, attaches the completed modes, and with them the descent, to the degenerate locus: $\EN$ is the map from the ascended boundary mode space of the singularity to its metric-level shadow, and nothing in the physical closure applies $\EN$ to the exterior radiation, whose observed single-quantum thermality is the coincidence of Proposition \ref{c:singularity.pinching-against-partial-trace-scope-coincidence}(i), entanglement seen one mode at a time. Applying $\EN$ to the radiation regardless is, however, a perfectly well-defined operation, and it is precisely what it means to read the radiation at the $\Cl(1,3)$ level: keep the occupation records, erase the coherences, every mode at once. Theorem \ref{c:singularity.deficit-theorem-cost-descent-exactly} quantifies that reading exactly: it costs $\rel{\rho}{\EN\rho}$, which by Corollary \ref{c:singularity.deficit-carried-quanta-correlations} is exactly the inter-quanta correlation content. One conflation must be avoided here, and the model is what prevents it: the deficit measures how \emph{much} coherent content the metric-level reading destroys, and it is identical in the model's scrambled and unscrambled histories, the accumulated phases squaring away in the diagonal (Proposition \ref{c:singularity.two-levels-finite-model}); the \emph{message} is written in which coherences carry it, the dressed correlators of Proposition \ref{c:singularity.purity-correlation-transfer-recovery}(3), a pattern that moves phase, not entropy, and to which the deficit is provably insensitive. What descent destroys is the sector where the information lies; what the deficit measures is the size of that destruction, never its content. Hawking's information loss is thereby derived rather than merely undergone, as a \emph{quantity}: it is the deficit of the wrong level, in closed form, the content itself residing in the destroyed sector's pattern, which no function of the deficit recovers. Two more locations must be stated plainly. First, the warrant: reading $\EN$, rather than the single average over $\operatorname{Ad}(\Gamma)$, as the metric-level map is Proposition \ref{c:singularity.p-reduced-metric-level-reading-many} below, the composite of the pointwise descents, its definitional antecedent quantified there; under the total-parity grading alone the average fixes every even-total state outright, the model's radiation state included, and the two-level discrepancy would be identically zero. Second, the cost's location: in the rate model the window marginal is joint-diagonal at every stage (Theorem \ref{c:singularity.eviction-window-thermal-product-radiation-entropy}(1)), hence a fixed point of $\EN$, so the reading the closure does perform, of the boundary modes, costs nothing; the entire nonzero deficit of the section ``The descent, applied: the two levels in numbers'' is the cost of the counterfactual reading, of the radiation, at the metric level. the section ``The descent, applied: the two levels in numbers'' exhibits both in exact numbers.
\end{remark}

\begin{proposition}[(P$_{\mathrm{loc}}$), reduced: the metric-level reading on many modes]\label{c:singularity.p-reduced-metric-level-reading-many}
\changedmark{gray}{reworded}
The mode-wise descents commute, each is the unique conditional expectation of the boundary mode space onto the $\operatorname{Ad}(\sigma_k)$-fixed subalgebra (the elements even in mode $k$, of dimension $4^{N}/2$ on $N$ modes, not the two-dimensional even part of the mode factor alone), and their composite, in any order, is $\EN$: a conditional expectation onto the intersection of the mode-wise ranges, the joint-number-diagonal algebra. Consequently, if the metric-level reading of the full boundary mode space is taken to be the simultaneous performance of the pointwise descents that the localisation attaches, one per point of $E$ (Fact \ref{c:singularity.imported-one-sided-completion-singularity}), then that reading is $\EN$ and the metric level is the joint-diagonal algebra, strictly smaller than the total-parity-even algebra on two or more modes. The antecedent is a \emph{choice}, not a consequence: it adopts the sheaf assembly of the pointwise levels over the net assembly, the two quantified, and not decided between, in the section ``The metric-level sheaf, and the hopping exclusion''.
\end{proposition}
\begin{proof}
Each $E_k(x) = \tfrac12(x + \sigma_k x \sigma_k)$ is a conditional expectation onto the $\operatorname{Ad}(\sigma_k)$-fixed subalgebra. The automorphism is inner and local, which is what removes any graded-tensor subtlety from the argument: in a Jordan-Wigner representation the string factors of $v_k$ and $v_k^{*}$ cancel in the product $\sigma_k = v_k^{*}v_k - v_k v_k^{*}$, which is single-site, so the fermionic mode-$k$ grading coincides with the spin-picture one and $\operatorname{Ad}(\sigma_k)$ is an automorphism of the full mode algebra. The uniqueness argument of Theorem \ref{c:singularity.descent-parity-average} runs verbatim on the full mode algebra: a conditional expectation $F$ onto that subalgebra satisfies $\sigma_k F(x)\sigma_k = F(\sigma_k x \sigma_k)$ by bimodularity ($\sigma_k$ lying in the range), so $F$ annihilates the $\operatorname{Ad}(\sigma_k)$-anti-fixed part and equals $E_k$ on the direct sum. The $\sigma_k$ commute, so the $E_k$ commute; their composite in any order is the group average over $\mathbb{Z}_2^{N}$, which is the pinching $\EN$ of Definition \ref{c:obj.singularity.multi-mode-descent}; and the range of a composite of commuting conditional expectations is the intersection of the ranges: an element fixed by every $\operatorname{Ad}(\sigma_k)$ is diagonal in the joint number basis, and conversely. Strict containment for two or more modes is witnessed by elements of the kind $X_{R_1}X_{L_1}Z_{L_2}$, total-parity even and not joint-diagonal. The script verifies the commutation, the composite-equals-pinching identity, and the idempotence, on generic rational data. The consequence clause takes the modes $k$ to be the ladder pairs that \ref{c:singularity.imported-one-sided-completion-singularity} localises, one per point of $E$.

\end{proof}

\phantomsection\label{c:note.singularity.setting-step-5-descend-state-what-2}
What remains of the former premise is the proposition's antecedent, and it is thin. The metric level was defined pointwise: at each mode it is the image of that mode's descent (Proposition \ref{c:singularity.descent-what-it-erases}, Theorem \ref{c:singularity.descent-parity-average}), and the localisation attaches exactly one descent to each point of $E$; reading a region of the boundary mode space at the metric level means surviving every one of its points' readings, which is the composite. The total-parity average is not an alternative candidate: it is the composite of nothing the closure constructs, there being no ``global pair'' for it to descend. Nor does the reduction touch Fact \ref{c:singularity.imported-univalence-superselection-rule}: the superselection rule continues to constrain observability by total parity, and cross-mode even elements such as $X_{R_1}X_{L_1}Z_{L_2}$ remain perfectly good observables of the ascended description; they are simply not metric-level records, being annihilated by the composite, which is exactly the recoverability of Proposition \ref{c:singularity.purity-correlation-transfer-recovery}(3): readable in the completion, invisible in the original space. Had the antecedent been dropped and the total-parity average used in $\EN$'s place, the two-level discrepancy of Steps 5 and 6b would be identically zero, the model's radiation state being even; the script verifies this too. The antecedent enters the table accompanying \ref{c:singularity.imported-equivalence-einstein-equations} as the residue of (P$_{\mathrm{loc}}$), definitional rather than physical. the section ``The metric-level sheaf, and the hopping exclusion'' states the residue in structural terms without reducing it: ``the pointwise metric level'' names two algebras, the even part of the mode factor and the even-in-that-mode part of the boundary mode space, and a region can be read by assembling either contravariantly (a sheaf, a limit: survive every point's descent) or covariantly (a net, a colimit: the algebra the points' retained parts generate). On a finite truncation the sheaf assembly gives $\operatorname{range}(\EN)$, at once the meet of the big fibres and the join of the small ones, and the net assembly gives the whole mode algebra with a discrepancy of zero (Theorems \ref{c:singularity.finite-truncation-meet-big-fibres-equals}, \ref{c:singularity.net-assembly-net-it-everything}). The chain adopts the sheaf assembly; that adoption is the antecedent, restated, and Remark \ref{c:note.singularity.what-antecedent-stated-plainly} records that the appendix quantifies both and decides neither. The sheaf exists whenever the completed mode bundle does, that is under the standing hypothesis $w_1(\shL) = 0$, and needs nothing further.

\section{Step 6a: iterate, or stop; the clock stops, and the boundary-mode capacity}

\phantomsection\label{c:note.singularity.setting-step-6a-iterate-or-stop}
The algorithm's last instruction is Step 6: iterate, or stop \cite[\S 8.1, Step 6]{ButlerT1}. It is discharged before any counting begins: the ascent $\Cl(1,3) \to \Cl(2,4)$ fixes the transport charge, the mod-8 clock stops at $d \equiv 6$, and by \cite[Thm.~3.6]{ButlerT3} no floor of the tower over the completed algebra ever develops a further central obstruction: the iteration is Morita-stationary and nothing structural remains to repair. What the boundary mode space at which the iteration stops \emph{holds} is another matter, and it is the task of this step and the next: Step 6a counts the capacity of the boundary modes, and Step 6b runs the two-level entropy computation on it.

This step counts the modes available for entanglement, and its physical content is not derivable from the algebra: it enters through three postulates, stated here as named axioms so that the count itself is a theorem and every premise is visible. Of their derivation programme, the corank half is discharged: the forced weighted completion of \cite[\S 8]{ButlerT4} fixes one ladder pair per point of the blown-up singular boundary (Remark \ref{c:note.singularity.corank-limit-weighted-answer}). What remains is the quantitative half, the cell calibration, the aggregation over the sphere factor of $E$, and the epoch from the modular period, assigned to the series' counting paper T8 \cite{ButlerT8}, with the connection-level geometry in the interior paper T6 \cite{ButlerT6}. Nothing in Steps 1 through 5 depends on this section.

\begin{lemma}[The causal epoch, computed exactly]\label{c:singularity.causal-epoch-computed-exactly}
\changedmark{gray}{reworded}
For the Schwarzschild gravitational field on flat spacetime (Definition \emph{Gauge Theory Gravity: the gravitational field on flat spacetime}), along the rain worldlines of the interior, the radial infall worldlines with unit energy at infinity whose velocity $\hb^{-1}(x')$ is $D$-parallel and which are the worldlines of the Painlev\'e-Gullstrand frame of the section ``Step 1: locate the failure'' (velocity $e_t$ in the gauge $\hb = 1 - \beta P$), $|dt/dr| = \sqrt{r/2M}\;r/(2M - r)$ in static Killing time $t$, so the span of Killing time between crossing $r = 2M(1 - \varepsilon)$ and reaching the singular locus $r = 0$ is
\[
\Delta t(\varepsilon) \;=\; \int_0^{2M(1-\varepsilon)} \sqrt{\frac{r}{2M}}\,\frac{r\,dr}{2M - r}
\;=\; 4M\Bigl(\tfrac12\ln\frac{1+u}{1-u} - u - \tfrac{u^{3}}{3}\Bigr)\Big|_{u = \sqrt{1-\varepsilon}}
\;=\; 2M\Bigl(\ln\tfrac{1}{\varepsilon} - \tfrac83 + 2\ln 2\Bigr) + O(\varepsilon) :
\]
of order $M$ for any fixed regulator $\varepsilon$, with only logarithmic sensitivity to the horizon-crossing cutoff, the coefficient $2M$ of the logarithm being independent of the worldline's energy. The informal estimate ``$\Delta t \sim M$ up to a logarithm'' is exactly this statement, with the logarithm identified.
\end{lemma}
\begin{proof}
Read in Gauge Theory Gravity, a geodesic of the static metric below is a worldline whose velocity $\hb^{-1}(x')$ is $D$-parallel (Fact \emph{Imported: the dictionary between Gauge Theory Gravity and curved-spacetime language}, clause (iii)). For a radial geodesic of the static metric with conserved energy $\mathcal{E} = f\,dt/d\lambda$ per unit mass, the normalisation gives $(dr/d\lambda)^{2} = \mathcal{E}^{2} - f$, so $dt/dr = \mathcal{E}/\bigl(f\sqrt{\mathcal{E}^{2} - f}\bigr)$; at $\mathcal{E} = 1$, $\sqrt{1 - f} = \sqrt{2M/r}$ and $|dt/dr| = \sqrt{r/2M}\,\bigl|r/(r - 2M)\bigr|$, which is the stated integrand for $r < 2M$. Substituting $r = 2Mu^{2}$ turns the integral into $4M\int_0^{\sqrt{1-\varepsilon}} u^{4}/(1 - u^{2})\,du$, and the integrand identity $u^{4}/(1 - u^{2}) = -u^{2} - 1 + 1/(1 - u^{2})$ gives the antiderivative $-u^{3}/3 - u + \tfrac12\ln\frac{1+u}{1-u}$, checked by differentiation; the closed form follows. For the expansion, $u = \sqrt{1 - \varepsilon} = 1 - \varepsilon/2 + O(\varepsilon^{2})$ gives $\tfrac12\ln\frac{1+u}{1-u} = \tfrac12\ln\frac{4}{\varepsilon} + O(\varepsilon) = \tfrac12\ln\tfrac1\varepsilon + \ln 2 + O(\varepsilon)$ and $u + u^{3}/3 = \tfrac43 + O(\varepsilon)$, so $\Delta t = 2M\ln\tfrac1\varepsilon + 4M(\ln 2 - \tfrac43) + O(\varepsilon)$. Near the horizon $dt/dr \approx 1/f$ for every $\mathcal{E}$, which is why the logarithm's coefficient is universal. Both the integrand identity and the substitution are verified exactly in the deposited script, as identities of rational functions, and the closed form is checked against the integral by exact rational Riemann bounds at the rational endpoint $\varepsilon = 1999 \cdot 10^{-6}$.

\end{proof}

\phantomsection\label{c:note.singularity.setting-step-6a-iterate-or-stop-2}
One identification has to be made explicit before the count, because two different surfaces are in play. The ladder lies on the blown-up boundary $E \cong \R \times S^{2}$ of \cite[\S 8]{ButlerT4}, which is at $r = 0$; the count below is taken on the stretched singularity $r = r_c$, at finite radius. The stretched singularity is a regulator surface for $E$: a slice of the interior, at the radius where the framework's own curvature scalar $2M/r^{3}$ reaches the Planck value, on which the modes of $E$ are counted, in the same way that a stretched horizon regulates a horizon. Axiom (C1) \emph{is} that identification; it is not derived from the completion, which places the ladder at $r = 0$ and says nothing about a cutoff.

\begin{fact}[Imported: the induced geometry of an interior slice]\label{c:singularity.imported-induced-geometry-interior-slice}
\changedmark{gray}{reworded}
Inside the horizon, for the Schwarzschild field in the Painlev\'e-Gullstrand gauge field $\hb = 1 - \beta P$, $\beta = \sqrt{2M/r}$, of Fact \emph{The gauge field and its cometric} [T4, Fact 8.1], the background surface $r = r_c$ ($0 < r_c < 2M$) is spacelike: its tangent vectors $e_t$, $e_\theta$, $e_\phi$ have physical squares $Q(\hb^{-1}(e_t)) = 1 - 2M/r_c < 0$ and $-r_c^{2}$, $-r_c^{2}\sin^{2}\theta$ in the coordinate normalisation, so its induced metric is $ds_3^{2} = |f(r_c)|\,dt^{2} + r_c^{2}\,d\Omega^{2}$ with $|f(r_c)| = 2M/r_c - 1$: Killing time $t$ is a spatial coordinate along it. On a surface of constant $r$ the static time and the Painlev\'e-Gullstrand time differ by a function of $r$ alone, so this is also the induced metric read in the static chart. \emph{Provenance:} this is read off the static form of the metric, the section ``The metric of the gauge field, and its static form'', with $f < 0$ for $r < 2M$.
\emph{Source:} none given: the source says the result is read off the static form of the metric in its own appendix.
\end{fact}

\begin{proposition}[The stretched singularity, and the count]\label{c:singularity.stretched-singularity-count}
\changedmark{gray}{reworded}
Assume (C1) and (C2). For the Schwarzschild field in the Painlev\'e-Gullstrand gauge field $\hb = 1 - \beta P$, $\beta = \sqrt{2M/r}$, of Fact \emph{The gauge field and its cometric} [T4, Fact 8.1], the three-volume of the background slice $r = r_c$ per unit Killing time, from the metric induced by $\hb^{-1}$ (Fact \ref{c:singularity.imported-induced-geometry-interior-slice}), is
\[
4\pi\,r_c^{2}\sqrt{|f(r_c)|} \;=\; 4\pi\sqrt{2M}\,r_c^{3/2}\sqrt{1 - \tfrac{r_c}{2M}} \;=\; 8\pi M\sqrt{1 - \tfrac{r_c}{2M}} \;=\; 8\pi M\,\bigl(1 + O(M^{-2/3})\bigr) ,
\]
using $r_c^{3} = 2M$: the cutoff radius conspires with the anisotropic stretch to leave \emph{linear} mass dependence at leading order, special to this cutoff (a mass-independent cutoff radius gives $M^{1/2}$ and can never reach the area law). With $\nu$ fermionic modes per Planck cell, and writing $N_{\mathrm{6a}}$ for the number of boundary modes (the notation Steps 6b and Fact \ref{c:singularity.assembled-page-curve-up-typicality} use for this section's count),
\[
\frac{dN_{\mathrm{6a}}}{dt} \;=\; 8\pi\nu M\sqrt{1 - \tfrac{r_c}{2M}} \;=\; \frac{\nu}{\TH}\sqrt{1 - \tfrac{r_c}{2M}} :
\]
the stretched singularity registers $\nu$ modes per Hawking period up to the correction factor $\sqrt{1 - r_c/2M}$, which tends to $1$ as $M \to \infty$ in Planck units ($r_c/2M = (2M)^{-2/3}$). It keeps time with the horizon's own clock.
\end{proposition}
\begin{proof}
Read in Gauge Theory Gravity: the metric below is the derived metric of a gravitational field $(\hb, \Omega)$ on flat spacetime (Definition \emph{Gauge Theory Gravity: the gravitational field on flat spacetime}), and the manifold is its invertible region.

By Fact \ref{c:singularity.imported-induced-geometry-interior-slice} the volume element of the slice is $\sqrt{|f(r_c)|}\,r_c^{2}\,dt\,d\Omega$, and integrating over the sphere gives $4\pi r_c^{2}\sqrt{2M/r_c - 1} = 4\pi r_c^{2}\sqrt{2M/r_c}\,\sqrt{1 - r_c/2M}$. With (C1, \ref{c:singularity.c1-curvature-cutoff}), $r_c^{3} = 2M$, so $r_c^{2}\sqrt{2M/r_c} = \sqrt{2M}\,r_c^{3/2} = \sqrt{2M}\sqrt{2M} = 2M$, which is the first display. Multiplying by $\nu$ per Planck cell (C2, \ref{c:singularity.c2-cell-density}) gives the second, and $8\pi M = 1/\TH$ is Theorem \ref{c:singularity.modular-descent}. The script verifies the slice-volume identity along the whole (C1) locus, parametrised by rational $r_c$ with $2M := r_c^{3}$, as an identity of rational functions in $r_c$, together with the rate identity with $\pi$ cancelled; a mass-independent cutoff radius is run as a control and fails the linear-in-$M$ form, the $M^{1/2}$ scaling that can never reach the area law.

\end{proof}

\begin{proposition}[The area-law test]\label{c:singularity.area-law-test}
\changedmark{gray}{reworded}
Assume (C1)-(C3), write $N_{\mathrm{6a}} = (dN_{\mathrm{6a}}/dt)\,\Delta t$ for the mode count of the counted epoch, and let $C := N_{\mathrm{6a}}\ln 2$ be its boundary-mode capacity in nats, $\ln 2$ per fermionic mode. Then, at leading order in $M$, $C = S_{\mathrm{BH}} = 4\pi M^2$ if and only if $\nu\,\Delta t = M/(2\ln 2)$: the capacity matches the Bekenstein-Hawking entropy \cite{Bekenstein1973, Hawking1975}, the coefficient calibrated by $\nu\,\Delta t$ rather than predicted (the closing paragraph of this section quantifies the conventions), exactly when the counted epoch is of order one Schwarzschild time, which is what (C3) supplies via Lemma \ref{c:singularity.causal-epoch-computed-exactly}, the $\ln 2$ being absorbed into the order-one calibration of $\nu\,\Delta t$. It is the capacity, not the bare count, that premise (C) and Fact \ref{c:singularity.assembled-page-curve-up-typicality} compare with $S_{\mathrm{BH}}$; the count itself then reads $N_{\mathrm{6a}} = S_{\mathrm{BH}}/\ln 2$. Counting over the hole's whole lifetime instead gives $N_{\mathrm{6a}} \sim M^4$, which shares with the interior-volume growth of Christodoulou-Rovelli \cite{CR2015} the linearity in advanced time and the unresolved tension with the constancy of the horizon entropy, though not the mass scaling: their $dV/dv \propto M^{2}$ gives $V \sim M^{5}$ over the lifetime, while $dN_{\mathrm{6a}}/dt \propto M$ gives $M^{4}$ here.
\end{proposition}
\begin{proof}
$N_{\mathrm{6a}} = (dN_{\mathrm{6a}}/dt)\,\Delta t = 8\pi\nu M\,\Delta t$ at leading order by Proposition \ref{c:singularity.stretched-singularity-count}, on the stretched singularity of \ref{c:singularity.c1-curvature-cutoff} with $\nu$ modes per Planck cell (\ref{c:singularity.c2-cell-density}), so $C = 8\pi\nu M\,\Delta t\,\ln 2$, and $C = 4\pi M^{2}$ if and only if $\nu\,\Delta t\,\ln 2 = M/2$. The lifetime count uses the evaporation time $\sim M^{3}$ in place of $\Delta t$. The script verifies the calibration with $\ln 2$ carried as a formal indeterminate: under $\nu\,\Delta t = M/(2\ln 2)$ the capacity equals $4\pi M^{2}$ identically, and the alternative 1-nat convention $\nu\,\Delta t = M/2$ misses it by the factor $\ln 2$ exactly, a $31$ per cent discrepancy no ``$\approx$'' should conceal. The reading of the condition as a counted epoch (\ref{c:singularity.c3-counted-epoch}) of order one Schwarzschild time rests on \ref{c:singularity.causal-epoch-computed-exactly}, which computes the causal epoch $\Delta t$ as of order $M$ for any fixed regulator.

\end{proof}

\phantomsection\label{c:note.singularity.setting-step-6a-iterate-or-stop-3}
\emph{The status, quantified.} Given (C1)-(C3), everything in this section is a theorem, and the postulates carry all of the physical content: (C1) is a UV regulator motivated by the framework's own scale; (C2)'s $\nu$ has its corank factor fixed at one by the forced weighted completion (Remark \ref{c:note.singularity.corank-limit-weighted-answer}, \cite[\S 8.9]{ButlerT4}), its remaining content being the cell calibration; (C3) is proved to be of order $M$ (Lemma \ref{c:singularity.causal-epoch-computed-exactly}) but its selection among possible spans is a physical identification. What survives all conventions, independent of (C2) and (C3): $dN_{\mathrm{6a}}/dt \propto 1/\TH$, and the linearity in $M$ that makes the area law reachable at all.

\section{Step 6b: iterate, or stop; two levels, the Page envelope, and eviction}

\subsection*{The paradox restated}
\phantomsection\label{c:note.singularity.setting-paradox-restated}
Hawking's argument: general relativity drives information into the singularity and destroys it. Quantum mechanics: information cannot be destroyed. The chain's finding: both statements are true, at different levels of description, and the contradiction is a category error.

\begin{center}
\emph{Information is lost in $\Cl(1,3)$. Information is preserved in $\Cl(2,4)$. These are the same physics read at two levels, and the discrepancy between the levels is not destruction: it is the fixed projection of Proposition \ref{c:singularity.descent-what-it-erases}, extended to the boundary mode space by Definition \ref{c:obj.singularity.multi-mode-descent}, and its size is an equation: the relative entropy of Theorem \ref{c:singularity.deficit-theorem-cost-descent-exactly}, carried entirely by the correlations of the quanta (Corollary \ref{c:singularity.deficit-carried-quanta-correlations}).}
\end{center}

General relativity is right that at the metric level the information is not conserved. Quantum mechanics is right that at the full level it is. The deficit between the two levels is not unaccounted for: traced through Theorem \ref{c:singularity.modular-descent}, it is precisely the Hawking flux, thermal at $\TH$, the shadow of the entanglement with the adjoined partner; and its magnitude is not estimated but computed, $\rel{\rho}{\EN\rho} = C(\rho) - C(\EN\rho)$, the correlations the descent cannot carry (Corollary \ref{c:singularity.deficit-carried-quanta-correlations}; the section ``The descent, applied: the two levels in numbers'' evaluates it for the model in exact arithmetic). The paradox does not need resolving; it needs the two levels computed separately. What remains a genuine question is the modern sharp form: whether the exterior radiation is unitary \emph{by itself} and follows the Page curve, i.e.\ whether the projection's deficit is returned to the radiation before the end. That question this section takes as far as the two levels allow.

\begin{assembled}[Assembled: the Page curve, up to typicality]\label{c:singularity.assembled-page-curve-up-typicality}
\changedmark{BurntOrange}{changed since the deposit}
Assume (i) the ascended description is globally pure, (ii) at each stage the joint state on radiation $\otimes$ partner-mode space is Page-typical, and (iii) the inputs of Theorem \ref{c:singularity.eviction-window-thermal-product-radiation-entropy}, whose clause (2) supplies the thermal arm. Let $\Sth(t)$ be the thermal accumulation of the emitted radiation and $C(t) = N_{\mathrm{6a}}(M(t))\ln 2 \approx S_{\mathrm{BH}}(M(t))$ the boundary-mode capacity of Step 6a, $N_{\mathrm{6a}}$ its mode count. Then
\[
S_{\mathrm{rad}}(t) \;=\; \min\bigl(\Sth(t),\, C(t)\bigr)
\]
up to subleading corrections: rising along Hawking's curve, turning over at the crossing (about half the initial entropy), falling along the shrinking capacity, vanishing at complete evaporation. \emph{Assembly and provenance:} the capacity arm is the Schmidt inequality, self-contained (a pure bipartite state has equal marginal spectra, so $S_R = S_L \leq \ln\dim L = C$); the thermal arm is Theorem \ref{c:singularity.eviction-window-thermal-product-radiation-entropy}(2) and needs boundary-mode locality, which is why (iii) is in the hypothesis line; saturation for typical states is Page's theorem \cite{Page1993} (\ref{c:singularity.imported-page-average-entropy-subsystem}), the imported ingredient and the reason this is a Fact rather than a theorem of this paper; hypothesis (i) supplies purity, the emission count supplies $\dim R$, Proposition \ref{c:singularity.stretched-singularity-count} with the epoch supplies $\dim L$, Proposition \ref{c:singularity.area-law-test} identifies the capacity with the area law, and monotone mass loss makes $C$ decreasing.
\end{assembled}
\begin{proof}
The capacity arm is the Schmidt inequality. By hypothesis (i) the joint state on radiation $R$ and partner modes $L$ is pure, so its two marginals have the same nonzero spectrum and
\[
S_{\mathrm{rad}} \;=\; S_R \;=\; S_L \;\leq\; \ln\dim L .
\]
The partner-mode space is the boundary-mode space of the counted epoch: Proposition \ref{c:singularity.stretched-singularity-count} with the epoch gives $\dim L = 2^{N_{\mathrm{6a}}}$, so $\ln\dim L = N_{\mathrm{6a}}\ln 2 = C(t)$, and Proposition \ref{c:singularity.area-law-test} identifies this capacity with $S_{\mathrm{BH}}(M(t))$ at leading order. Monotone mass loss makes $C$ decreasing.

The thermal arm is clause (2) of Theorem \ref{c:singularity.eviction-window-thermal-product-radiation-entropy}, whose inputs (hypothesis (iii)) include boundary-mode locality: the radiation's entropy is its thermal accumulation $\Sth(t)$, and the emission count supplies $\dim R$, with $\ln\dim R = \Sth(t)$ at leading order. Hence $S_{\mathrm{rad}}(t) \leq \min(\Sth(t), C(t))$.

Saturation is Page's theorem (\ref{c:singularity.imported-page-average-entropy-subsystem}): by hypothesis (ii) the joint state is Page-typical, so the smaller marginal's entropy is within $m/2n \leq 1/2$ nat of $\ln\min(\dim R, \dim L)$, a subleading correction. Therefore $S_{\mathrm{rad}}(t) = \min(\Sth(t), C(t))$ up to subleading corrections. The curve rises with $\Sth$ along Hawking's curve, turns over where the two arms cross (about half the initial entropy, as the paper's figure of the envelope shows), falls with the shrinking capacity, and vanishes at complete evaporation, where $C = 0$.

\end{proof}

\begin{remark}[The island dictionary, and one dissonance]\label{c:note.singularity.island-dictionary-one-dissonance}
In the quantum-extremal-surface derivations \cite{Penington2020, AEMM2019} the falling branch arises when an interior island joins the computation of the radiation's entropy, contributing its area over $4G$. Here the same role is played by the singularity's boundary mode space, which by Proposition \ref{c:singularity.area-law-test} has exactly the right size. The dictionary: island $\leftrightarrow$ partner modes; extremal-surface area term $\leftrightarrow$ stretched-singularity capacity; saddle exchange $\leftrightarrow$ the crossing of the two counts. One dissonance is recorded: in the islands derivation the exterior semiclassical theory computes its own correction, whereas here the correction is carried by the ascended boundary modes. Whether these are two descriptions of one mechanism or genuinely different claims is the sharpest open question the chain faces; the eviction theorem below converts it from an unresolved tension into a programme.
\end{remark}

\subsection*{The eviction theorem: the falling branch without typicality}
\phantomsection\label{c:note.singularity.setting-eviction-theorem-falling-branch-without}
The rate model replaces typicality with mechanism. Modes come in units, entropies in $\ln 2$. Four of its five inputs are the model's own and are collected here as a numbered premise, each clause citable; the fifth, (L$_E$), is Premise \ref{c:singularity.l-boundary-mode-locality}, restated after the environment in the form the rate model needs.

\phantomsection\label{c:note.singularity.setting-eviction-theorem-falling-branch-without-2}
To these add (L$_E$) \emph{boundary-mode locality}, in the form the rate model needs: the window's internal dynamics is diagonal in the window's joint number basis (Premise \ref{c:singularity.l-boundary-mode-locality} with Proposition \ref{c:singularity.number-diagonal-dynamics-models-scrambling-derived}; diagonal, not merely number conserving, since a parity-preserving hopping changes the marginals), and, a clause the premise does not contain and which is stated here as part of the rate model, the emitted quanta, once in the radiation, are free and never return to the window. Both halves are used in the proof: a returning quantum or a non-diagonal window unitary breaks the identity below, the second by the example in Corollary \ref{c:singularity.envelope-endpoint-why-boundary-mode-locality}. No further assumption about where the window places correlations is needed: (L$_E$) already fixes every window marginal, and that is what the theorem below depends on. An additional ``no attic'' axiom is sometimes imposed at this point; in the model as just set up it is not an axiom but a proposition:

\begin{proposition}[No attic, derived in the model]\label{c:singularity.no-attic-derived-model}
\changedmark{gray}{reworded}
Under (L), (U) and (C): modes exiting the window carry no q-number content out of the two-sector decomposition; their correlations are placed, at eviction, in $L(t) \cup R(t)$. The no-attic axiom (A) is a theorem of the model's kinematics.
\end{proposition}
\begin{proof}
By (L) of \ref{c:singularity.rate-models-kinematic-inputs-l-u} the state space at every stage is $H_{L(t)} \otimes H_{R(t)}$, with the exact dimensions of (C): there is no third tensor factor. An eviction is the transfer of one mode's factor from the $L$ side of the decomposition to the $R$ side (in physical terms, the mode leaves the shrinking boundary mode space and its content is carried with it into the radiation sector; in the finite-mode model of the section ``The finite-mode model: the dissolution step by step'' this is implemented literally). A unitary step (U) preserves the global state, hence all its correlations; since the factor carrying them lands in $R$, the correlations of the evicted mode remain in $L \cup R$, which is the assertion. What this proof does \emph{not} establish, and what carries the physical content, is the exhaustiveness of (L)'s two sectors for real evaporation: that no remnant or third sector exists dynamically. That claim is exactly the closure-completeness programme, argued from the tower's floor-one completeness \cite[Cor.~6.2]{ButlerT3} and flagged as the conjectural step; here it lies inside (L), where it is visible, rather than inside an unexamined axiom. One candidate third sector is excluded by theorem rather than by hypothesis: the completion at the black hole forms no shadow wall anywhere in the physical closure, the conserved quantity exiting the manifold through the blown-up boundary \cite[Thm.~8.12, Rem.~8.13]{ButlerT4}, so no degeneracy-generated sector hides inside the two-sector decomposition. The dynamical no-remnant clause of (L) remains the programme's.

\end{proof}

\begin{theorem}[Eviction: the window is a thermal product, and the radiation entropy is its content]\label{c:singularity.eviction-window-thermal-product-radiation-entropy}
\changedmark{gray}{reworded}
Under (L), (U), (C), (I) and (L$_E$), the window dynamics diagonal in the joint number basis and the radiation quanta free and never returning, with the no-attic property supplied by Proposition \ref{c:singularity.no-attic-derived-model}, write $W(t)$ for the set of pairs whose window half is in the window at stage $t$ and $C(t) = |W(t)|\ln 2$ for the capacity. Then, in the rate model:
\begin{enumerate}
\item (The identity.) At every stage the window state is the product of its members' thermal marginals, $\rho_L = \bigotimes_{k \in W(t)}\operatorname{diag}(c_k^{2}, s_k^{2})$, and consequently
\[
S_R(t) \;=\; S_L(t) \;=\; \sum_{k \in W(t)} h_k \qquad \text{exactly.}
\]
\item (Envelope and endpoint.) $S_R(t) \leq \min(\Sth(t), C(t))$ at every stage, and $S_R = 0$ at complete evaporation.
\item (Rising branch, exact.) While no eviction has occurred, $W(t)$ is the set of all injected pairs and $S_R(t) = \Sth(t)$ exactly.
\item (Falling branch, quantified.) $S_R(t) = C(t) - \sum_{k \in W(t)}(\ln 2 - h_k)$: the radiation entropy lies below the capacity by exactly the sum of the window members' per-quantum shortfalls, with equality $S_R = C$ if and only if every window member is maximally entangled. In the maximally entangled idealisation, every $h_k = \ln 2$, this gives $S_R(t) = \min(\Sth(t), C(t))$ exactly at every stage.
\end{enumerate}
\end{theorem}
\begin{proof}
By (L) and (U) of \ref{c:singularity.rate-models-kinematic-inputs-l-u} the global state is pure on $H_{L(t)} \otimes H_{R(t)}$ at every stage, so $S_R = S_L$ (equal marginal spectra of a pure state), and by (C) $\dim H_{L(t)} = 2^{|W(t)|}$, so $S_L \leq C(t)$ is the dimension bound.

For (1), consider the joint number basis of all the qubits ever injected, labelled $R_k, L_k$. By (I) the state at injection of pair $k$ is the previous state tensored with $c_k|0_{R_k}0_{L_k}\rangle + s_k|1_{R_k}1_{L_k}\rangle$, so its support lies in the configurations with $R_k = L_k$ for every injected $k$ and its amplitude on such a configuration is $\pm\prod_k c^{(k)}_{b_k}$, with $c^{(k)}_0 = c_k$, $c^{(k)}_1 = s_k$. By (L$_E$, \ref{c:singularity.l-boundary-mode-locality}) every subsequent operation is either a diagonal unitary on the window, which multiplies each configuration by a phase and preserves both the support and every squared amplitude, or an eviction, which relabels a qubit from $L$ to $R$ and touches no amplitude (no attic, by \ref{c:singularity.no-attic-derived-model}); the radiation quanta are free. So at every stage the support constraint $R_k = L_k$ holds and the squared amplitude of each configuration is $\prod_k (c^{(k)}_{b_k})^{2}$. Now trace out the radiation. Every $R_k$ is in the radiation (the emitted quanta never return), so two configurations that agree on the radiation agree on every $L_k$ with $k \in W(t)$ by the support constraint: $\rho_L$ has no off-diagonal entries. Its diagonal entry at the window configuration $(b_k)_{k \in W(t)}$ is the total probability of the global configurations extending it, which is $\prod_{k \in W(t)}(c^{(k)}_{b_k})^{2}$ times the sum over the other pairs' configurations, and that sum is $1$. This is the product state displayed, and its entropy is $\sum_{k \in W(t)} H(c_k^{2}) = \sum_{k \in W(t)} h_k$.

(2) $W(t)$ is a subset of the injected pairs and every $h_k \geq 0$, so $\sum_{W(t)} h_k \leq \sum_{k \leq t} h_k = \Sth(t)$; and $\sum_{W(t)} h_k \leq |W(t)|\ln 2 = C(t)$ since each $h_k \leq \ln 2$. At complete evaporation $W = \varnothing$ and the sum is empty. (3) is (1) with $W(t)$ the full set of injected pairs. (4) is (1) rewritten as $|W|\ln 2 - \sum_{W}(\ln 2 - h_k)$, with equality in $S_R \leq C$ exactly when every term of the shortfall vanishes; in the idealisation the shortfall is identically zero, so $S_R = C$ whenever an eviction has occurred and $S_R = \Sth$ before, which is $\min(\Sth, C)$ in both regimes. The script verifies (1) on the thermal model at every stage, from the exact spectrum of $\rho_L$, as an identity of formal logarithms; the strict inequality of (4) at every stage with a nonempty window, by rigorous rational bounds; and (4)'s exact saturation on an idealised run simulated on exact density matrices along a decreasing capacity schedule.

\end{proof}

\phantomsection\label{c:note.singularity.setting-eviction-theorem-falling-branch-without-3}
The theorem's content is that under boundary-mode locality the falling branch is not a statistical effect and not the outcome of a scrambling mechanism: it is the window's contents, read off. Typicality is not used anywhere, and neither is any assumption about where correlations are deposited, because under (L$_E$) nothing the window does can move a marginal. What the theorem does not deliver, for any mode the chain admits, is exact saturation of the envelope: maximal entanglement is $c_k = s_k$, which by the thermofield angle $\tan^{2}\theta_k = e^{-\omega_k/\TH}$ is the zero-frequency case $\omega_k = 0$; the linear assignment of Premise \ref{c:singularity.q-joint-quasi-freeness-s-wave} supplies exactly one such mode, the sea level, so saturation, which needs every window member at zero frequency, is unavailable to any window of two or more members; at every nonzero weight $h_k < \ln 2$ strictly, and the falling branch runs strictly inside the capacity. What this paper claims without typicality is therefore the envelope, the exact rising branch, the quantified falling branch $S_R = \sum_{W} h_k$ and the endpoint, and not the saturated Page curve, which remains available on exactly two routes: typicality (Fact \ref{c:singularity.assembled-page-curve-up-typicality}), or the zero-frequency idealisation. the section ``The finite-mode model: the dissolution step by step'' exhibits the shortfall in numbers, and Proposition \ref{c:singularity.models-entropy-sequence-exactly} states the conclusion.

\begin{corollary}[The envelope and the endpoint, and why boundary-mode locality is needed for them]\label{c:singularity.envelope-endpoint-why-boundary-mode-locality}
From (L), (U), (C), (I) and (L$_E$), with no typicality and no assumption about the placement of correlations:
\[
S_R(t) \;\leq\; \min\bigl(\Sth(t),\, C(t)\bigr) \qquad \text{and} \qquad S_R(\text{complete evaporation}) = 0 .
\]
The premise (L$_E$) cannot be dropped: a window unitary that is not number diagonal can leave the remaining window qubit far more mixed than its thermal marginal, and then $S_R$ exceeds $\Sth$. Two pairs with $(c, s) = (\tfrac{24}{25}, \tfrac{7}{25})$, the window rotation $|00\rangle \mapsto \tfrac35|00\rangle + \tfrac45|11\rangle$, $|11\rangle \mapsto -\tfrac45|00\rangle + \tfrac35|11\rangle$ on $(L_1, L_2)$, and then the eviction of $L_1$ leave $\rho_{L_2} = \operatorname{diag}(\tfrac{5968}{15625}, \tfrac{9657}{15625})$, so $S_R = S_L \approx 0.96$ bits against $\Sth = 2H(\tfrac{576}{625}) \approx 0.55$ nats, that is $0.79$ bits.
\end{corollary}
\begin{proof}
The bound and the endpoint are Theorem \ref{c:singularity.eviction-window-thermal-product-radiation-entropy}(2). For the necessity of (L$_E$, \ref{c:singularity.l-boundary-mode-locality}): the rotation is a unitary on the window, so it satisfies (U, \ref{c:singularity.rate-models-kinematic-inputs-l-u}) and does not change $S_L$, but it is not number diagonal and it moves the marginal of $L_2$ off $\operatorname{diag}(c^{2}, s^{2})$; after the eviction the window is $L_2$ alone, whose reduced state is the displayed diagonal matrix (the script computes it exactly), and its entropy exceeds $\Sth$ by rigorous rational bounds. Every input other than (L$_E$) holds in the example. The rotation is a pairing rotation, $\Delta N = \pm 2$, not a hop: it lies outside the class within which Theorem \ref{c:singularity.class-relative-no-go-no-supplied} proves no-hopping, which is why (L$_E$)'s exclusion of it is a premise and not a theorem (Remark \ref{c:note.singularity.declared-residual-size}).

\end{proof}

\begin{corollary}[What an event returns to the radiation]\label{c:singularity.what-event-returns-radiation}
\changedmark{gray}{reworded}
At an emission event let $A$ be the block the event adds to the radiation, the new quantum $R_k$ together with the qubits $L_j$ evicted at that event, and $R_{<}$ the radiation already present. Then, under the hypotheses of Theorem \ref{c:singularity.eviction-window-thermal-product-radiation-entropy},
\[
S(\rho_A) \;=\; h_k + \sum_{j\ \mathrm{evicted}} h_j , \qquad I(A{:}R_{<}) \;=\; 2\sum_{j\ \mathrm{evicted}} h_j :
\]
the new quantum arrives uncorrelated with the earlier radiation, and every evicted qubit arrives fully correlated with the earlier radiation as a whole, $I(L_j{:}R_{<}) = 2h_j$, the maximum a qubit of entropy $h_j$ can carry. That correlation is not, in general, with its own partner $R_j$ alone: the pair coherence is damped by the scrambling (Proposition \ref{c:singularity.purity-correlation-transfer-recovery}(2)), so $I(L_j{:}R_j) < 2h_j$ once a gate has acted, and the remainder lies in the multi-qubit correlations of Proposition \ref{c:singularity.purity-correlation-transfer-recovery}(3). In the maximally entangled idealisation with one eviction per event, $I(A{:}R_{<}) = 2\ln 2 = S(\rho_A)$: the block is returned to the radiation fully correlated, at the rate the capacity falls. This is the sense in which the rate model reproduces the mirror behaviour of \cite{HP2007}: the window returns information at the rate it loses capacity, by eviction, not by scrambling.
\end{corollary}
\begin{proof}
Every member of $A$ has its partner outside $A$ ($L_k \in W$ for the new quantum, $R_j \in R_{<}$ for an evicted $L_j$), so by the support constraint $\rho_A$ is diagonal with the product distribution of its members, and $S(\rho_A) = h_k + \sum_j h_j$. By purity, $S(\rho_{R_<})$ before the event equals $S_L$ before it, and it is unchanged by the event (a unitary on the complement of $R_<$ and the adjunction of a product factor leave $\rho_{R_<}$ fixed); $S(\rho_{AR_<})$ after the event is $S_L$ after it. Theorem \ref{c:singularity.eviction-window-thermal-product-radiation-entropy}(1) gives $S_L^{\mathrm{after}} - S_L^{\mathrm{before}} = h_k - \sum_j h_j$, so $I(A{:}R_<) = S(\rho_A) + S_L^{\mathrm{before}} - S_L^{\mathrm{after}} = 2\sum_j h_j$. The per-qubit statement is the same computation with $A$ replaced by the single $L_j$. The script computes $S(\rho_A)$, $S(\rho_{R_<})$ and $S(\rho_{AR_<})$ from exact spectra at every event of the model's history and verifies the identity. The clause on the partner correlation rests on \ref{c:singularity.purity-correlation-transfer-recovery}, whose (2) gives the damped pair coherence and whose (3) places the remainder in multi-qubit correlations.

\end{proof}

\begin{remark}[What was obtained, and what it cost]\label{c:note.singularity.what-was-obtained-what-it-cost}
Typicality is absent from what the model claims, and the axiom count is correspondingly small: the envelope, the exact rising branch, the quantified falling branch, the turnover and the endpoint follow mechanically from the window identity, saturation of the falling branch remaining typicality's alone (Fact \ref{c:singularity.assembled-page-curve-up-typicality}; the one-mode zero-frequency idealisation aside), the envelope and endpoint needing no typicality and no placement rule (Corollary \ref{c:singularity.envelope-endpoint-why-boundary-mode-locality}), and the no-attic axiom (A) is a proposition of the model rather than an assumption, its physical force consolidated into the exhaustiveness of (L), which is the no-remnant hypothesis in its correct location: exactly the islands-side choice in the dissonance of Remark \ref{c:note.singularity.island-dictionary-one-dissonance}, with its derivation programme aimed at the closure-completeness of the tower \cite[Cor.~6.2]{ButlerT3}. What remains assumed: boundary-mode locality (L$_E$) with its never-return clause (its no-hopping half a theorem within the class of supplied generators, the section ``The metric-level sheaf, and the hopping exclusion''), which is what makes the falling branch mechanical (the capacity crossing is the rate-model shadow of the extremal-surface exchange); (I)'s freshness clause, the rate model's own; the capacity postulates (C1)-(C3) of Step 6a feeding (C); the standing hypothesis (U); and, for the Planck sector alone, the joint quasi-freeness (Q) of Step 4e, in its recorded tension with Steps 5 and 6b; the state of Theorem \ref{c:singularity.modular-descent} is not a standing hypothesis but is fixed by (P) and (R) (Corollary \ref{c:singularity.chains-state-class-stability}), and (L)'s exhaustiveness clause remains the conjectural input. Every premise in the chain is stated, named and quantified; and the two-level discrepancy itself is quantified separately, once and exactly, by Theorem \ref{c:singularity.deficit-theorem-cost-descent-exactly}. The division of tasks with the sheaf paper is likewise exact: its Schwarzschild composite proves the kinematics (the carrier module, its uniqueness, the rain-finiteness of the chain's dyad, the localised ladder pair) and poses back to this paper the operator and state level \cite[\S 8.9]{ButlerT4}; that residue is discharged in the section ``What selects the state'', down to the single import (R).
\end{remark}

\begin{remark}[Monogamy, firewalls, and the small-corrections theorem]\label{c:note.singularity.monogamy-firewalls-small-corrections-theorem}
Two classical objections must be located relative to the eviction mechanism, because the chain's state class invites both. First, monogamy: the firewall argument of \cite{AMPS2013} forbids a late quantum from being maximally entangled with its interior partner and with the early radiation at once. The rate model never asks for that: a window qubit is entangled with its own partner while it is in the window, and its correlation with the early radiation as a whole is acquired at eviction, when the qubit itself is transferred to the radiation sector (Corollary \ref{c:singularity.what-event-returns-radiation}), so at no stage does one quantum carry both entanglements. The cost is stated plainly: eviction is a physical transfer of an interior degree of freedom into the radiation, and the smoothness of the horizon it presupposes is asserted by the model, not derived, which is exactly the point \cite{AMPS2013} press. Second, the small-corrections theorem of \cite{Mathur2009} shows that no accumulation of order-$\epsilon$ corrections to the pairwise Hawking process can return the radiation entropy to zero. Eviction does not evade that theorem by small corrections; it violates its locality hypothesis outright, the evicted qubit being an order-one, nonlocal transfer out of the interior. That is a legitimate alternative and a costly one: the falling branch of Theorem \ref{c:singularity.eviction-window-thermal-product-radiation-entropy} stands or falls with the physical reality of eviction, which is the same commitment the islands programme makes in its own language (Remark \ref{c:note.singularity.island-dictionary-one-dissonance}).
\end{remark}

\section{The finite-mode model: the dissolution step by step}

\phantomsection\label{c:note.singularity.setting-finite-mode-model-dissolution-step}
This section contains the paper's new computation: an exactly solvable model, small enough that every claim is a finite calculation in exact arithmetic, in which the chain's conceptual endpoint can be followed step by step, at both levels. The model is a toy and is offered as one; its role is to exhibit, not to derive. What it exhibits, exactly rather than approximately: every single radiated quantum is thermal; the joint radiation state ends pure; the information is carried by inter-quanta correlations invisible to any single-quantum measurement; the entropy of the radiation rises exactly along $\Sth$, falls under the eviction mechanism of Theorem \ref{c:singularity.eviction-window-thermal-product-radiation-entropy}, and returns to zero, with the falling branch strictly inside the envelope by exactly the per-quantum shortfall the theorem quantifies; and, in the section ``The descent, applied: the two levels in numbers'', the descent of Definition \ref{c:obj.singularity.multi-mode-descent} is applied to the radiation outright, so that the metric-level reading of the same history is also computed rather than described.

\subsection*{The model}
\phantomsection\label{c:note.singularity.setting-model}
Qubits stand for modes. At emission event $k$ ($k = 1, \dots, n$) a thermofield pair is created between a radiation qubit $R_k$ and a window qubit $L_k$:
\[
|\psi_k\rangle \;=\; c_k\,|0_{R_k} 0_{L_k}\rangle + s_k\,|1_{R_k} 1_{L_k}\rangle, \qquad c_k = \cos\theta_k, \quad s_k = \sin\theta_k, \quad \tan^2\theta_k = e^{-\omega_k/\TH} ,
\]
the qubit image of Theorem \ref{c:singularity.modular-descent}'s state, with rational data (the script uses $(c_k, s_k) = (\tfrac45, \tfrac35), (\tfrac45, \tfrac35), (\tfrac45, \tfrac35), (\tfrac{12}{13}, \tfrac{5}{13})$ for $k = 1, \dots, 4$; every pair has $c_k > s_k$, so every $\omega_k > 0$). The window holds at most $N_{\max} = 2$ qubits; while $|L| \leq N_{\max}$ nothing else happens (Phase I). When an emission would exceed capacity, the model applies a \emph{scrambling unitary} $U$ to the window, diagonal in the number basis (a circuit of controlled-$Z$ gates, entries $\pm 1$, encoding whatever pattern the dynamics carries, including any infalling message), and then \emph{evicts} the oldest window qubit: the evicted qubit is transferred to the radiation as a late quantum (Phase II). This implements the inputs of the section ``The eviction theorem: the falling branch without typicality'' with a window that fills to $N_{\max} = 2$ over the first two events and then stays full until it is finally emptied: the no-attic property of Proposition \ref{c:singularity.no-attic-derived-model} holds because the evicted qubit itself joins $R$, (I) because each pair arrives fresh, and (L$_E$) because the scrambling is diagonal. At the end all window qubits have been evicted and the radiation holds all $2n$ qubits.

\begin{lemma}[Parity protection]\label{c:singularity.parity-protection}
At every stage of the model, and for \emph{any} sequence of number-basis-diagonal window unitaries, the reduced density matrix of each individual qubit, radiation or window, early or evicted, is exactly its thermal matrix $\operatorname{diag}(c_k^2, s_k^2)$: diagonal, with the original populations, regardless of the scrambling.
\end{lemma}
\begin{proof}
The initial global state is supported on computational configurations satisfying $R_k = L_k$ for every $k$ (each pair contributes $|00\rangle$ or $|11\rangle$). A diagonal unitary maps each configuration to itself times a phase: the support constraint is preserved forever, and so is each configuration's squared amplitude. Off-diagonals: the coherence $\langle 0|\rho_{R_k}|1\rangle$ is a sum over pairs of global configurations that differ in the bit $R_k$ and agree everywhere else; but agreement at $L_k$ contradicts the support constraint for one of the two, so the sum is empty and the coherence is exactly zero. (For $L_k$, exchange the roles.) Diagonals: the populations are sums of squared amplitudes, unchanged by diagonal unitaries; they equal $c_k^2$ and $s_k^2$ from the start. Eviction relabels a qubit from $L$ to $R$ and changes no amplitudes.

\end{proof}

\begin{remark}[Thermality records what the projection cannot see]\label{c:note.singularity.thermality-records-what-projection-cannot-see}
The lemma is now also a special case: by Theorem \ref{c:singularity.ascended-marginals-are-diagonal} every even state of the boundary modes has diagonal single-mode marginals, and the model's states are even, so the protection is not a property of the phase-type scrambling but of fermion parity, and Proposition \ref{c:singularity.number-diagonal-dynamics-models-scrambling-derived} shows that under boundary-mode locality the scrambling could not have been anything but phase type. Within the model the lemma is the model's instance of Theorem \ref{c:singularity.descended-form-state-independent}: no encoding performed by the (phase-type) internal dynamics can be visible to any single-quantum measurement, exactly. One quantum at a time, the radiation carries nothing, and this is a theorem of the encoding, not an approximation: thermality is what the projection records in place of what it cannot see. The choice of \emph{diagonal} scrambling is what makes the statement exact rather than approximate, and it is the exact match to the chain, where single-mode thermality is likewise exact (Theorem \ref{c:singularity.modular-descent}) rather than emergent.
\end{remark}

\begin{proposition}[Purity, correlation transfer, and recovery]\label{c:singularity.purity-correlation-transfer-recovery}
In the model with two pairs and one controlled-$Z$ between $L_1$ and $L_2$ (the minimal scrambled instance; the script runs larger ones):
\begin{enumerate}
\item the global state is pure at every stage, and at completion the radiation alone carries it: $S(R_{\mathrm{all}}) = 0$ exactly;
\item the pair coherence is damped by exactly the visibility $\gamma = c_2^2 - s_2^2 = \cos 2\theta_2$: the two-qubit marginal of $(R_1, L_1)$ has coherence $c_1 s_1 \gamma$ in place of $c_1 s_1$, so its R\'enyi-2 purity drops from $1$ to $c_1^4 + s_1^4 + 2 c_1^2 s_1^2 \gamma^2$: two-body correlations alone no longer hold the whole pair;
\item the deficit is not lost but relocated to three-body correlations, and it is recoverable: the correlator $\langle X_{R_1} X_{L_1} \rangle$ equals $2 c_1 s_1 \gamma$ while the dressed correlator $\langle X_{R_1} X_{L_1} Z_{L_2} \rangle$ equals $2 c_1 s_1$ exactly, the visibility $\gamma$ having moved from the undressed correlator to the dressed one, so the pattern of applied gates (the message) is determined by multi-qubit correlators although every single-qubit state, by Lemma \ref{c:singularity.parity-protection}, is identically thermal.
\end{enumerate}
\end{proposition}
\begin{proof}
(1) Unitarity and relabelling. (2) Write the four-qubit state $\sum_{b_1 b_2} c^{(1)}_{b_1} c^{(2)}_{b_2} (-1)^{b_1 b_2} |b_1 b_1 b_2 b_2\rangle$ with $c^{(k)}_0 = c_k$, $c^{(k)}_1 = s_k$. Tracing out $(R_2, L_2)$: the coherence of $(R_1, L_1)$ between $|00\rangle$ and $|11\rangle$ picks up $\sum_{b_2} |c^{(2)}_{b_2}|^2 (-1)^{b_2} = c_2^2 - s_2^2 = \gamma$; the populations are untouched. The stated purity is $\tr \rho^2$ for the two-by-two supported block $\begin{psmallmatrix} c_1^2 & c_1 s_1 \gamma \\ c_1 s_1 \gamma & s_1^2 \end{psmallmatrix}$. (3) The computations are short enough to print. Write the scrambled state as $|\Psi\rangle = \sum_{b_1, b_2} c^{(1)}_{b_1} c^{(2)}_{b_2} (-1)^{b_1 b_2} |b_1 b_1 b_2 b_2\rangle$ in the order $(R_1, L_1, R_2, L_2)$. The operator $X_{R_1}X_{L_1}$ flips both bits of the first pair, so
\[
\langle X_{R_1}X_{L_1}\rangle
= \sum_{b_1, b_2} c^{(1)}_{b_1} c^{(1)}_{\bar b_1}\, \bigl(c^{(2)}_{b_2}\bigr)^2\, (-1)^{b_1 b_2}(-1)^{\bar b_1 b_2}
= 2 c_1 s_1 \sum_{b_2} \bigl(c^{(2)}_{b_2}\bigr)^2 (-1)^{b_2}
= 2 c_1 s_1 \bigl(c_2^2 - s_2^2\bigr) = 2 c_1 s_1 \gamma ,
\]
using $(-1)^{b_1 b_2}(-1)^{\bar b_1 b_2} = (-1)^{b_2}$ since $b_1 + \bar b_1 = 1$. Dressing with $Z_{L_2}$ inserts one more $(-1)^{b_2}$, which cancels the previous one:
\[
\langle X_{R_1}X_{L_1} Z_{L_2}\rangle
= 2 c_1 s_1 \sum_{b_2} \bigl(c^{(2)}_{b_2}\bigr)^2 (-1)^{b_2}(-1)^{b_2}
= 2 c_1 s_1 .
\]
For the unscrambled state the phases are absent and the same two sums give $2 c_1 s_1$ and $2 c_1 s_1 \gamma$ respectively: the pair of correlators $(\langle XX\rangle, \langle XXZ\rangle)$ takes the value $(2c_1s_1,\, 2c_1s_1\gamma)$ without the gate and $(2c_1s_1\gamma,\, 2c_1s_1)$ with it, so the gate pattern is read off from three-qubit correlations, while Lemma \ref{c:singularity.parity-protection} fixes every single-qubit marginal at the identical thermal matrix in both histories. With the script's numbers ($c_1, s_1 = \tfrac45, \tfrac35$; $c_2, s_2 = \tfrac45, \tfrac35$; $\gamma = \tfrac{7}{25}$) the two signatures are $(\tfrac{24}{25}, \tfrac{168}{625})$ and $(\tfrac{168}{625}, \tfrac{24}{25})$, confirmed exactly.

\end{proof}

\subsection*{The model, run exactly}
\phantomsection\label{c:note.singularity.setting-model-run-exactly}
The script runs the full model ($n = 4$ pairs, $N_{\max} = 2$, thermal angles from the rational triples $(c_k, s_k) = (\tfrac45, \tfrac35), (\tfrac45, \tfrac35), (\tfrac45, \tfrac35), (\tfrac{12}{13}, \tfrac{5}{13})$, a controlled-$Z$ between the two oldest window qubits before each eviction) and tabulates, after every event, the exact purity $\tr\rho_{R(t)}^{2}$ and the exact von Neumann entropy $S(\rho_{R(t)})$ of the radiation so far, together with the thermal accumulation and the capacity. Every entry of the table accompanying \ref{c:singularity.purity-correlation-transfer-recovery} is a rational number or an exact expression in the binary entropy $H(p)$ of a rational argument; the entropy columns are reported in bits, that is $H/\ln 2$ and $C/\ln 2 = |W|$, and the identities between the columns are verified as identities of formal logarithms; the decimals are for orientation only.

 \begin{tabular}{cccccccc}
\toprule
event & $|R|$ & window & $\tr\rho_R^{2}$ & $S_R = S_L$ (bits) & $\approx$ & $\Sth$ (bits) & $C$ (bits) \\
\midrule
1 & 1 & $\{L_1\}$ & $337/625$ & $h$ & $0.943$ & $h$ & 1 \\
2 & 2 & $\{L_1, L_2\}$ & $113569/390625$ & $2h$ & $1.885$ & $2h$ & 2 \\
3 & 4 & $\{L_2, L_3\}$ & $113569/390625$ & $2h$ & $1.885$ & $3h$ & 2 \\
4 & 6 & $\{L_3, L_4\}$ & $7198657/17850625$ & $h + h'$ & $1.547$ & $3h + h'$ & 2 \\
5 & 7 & $\{L_4\}$ & $21361/28561$ & $h'$ & $0.605$ & $3h + h'$ & 1 \\
6 & 8 & $\varnothing$ & $1$ & $0$ & $0$ & $3h + h'$ & 0 \\
\bottomrule
\end{tabular}

\emph{The finite-mode model, event by event, entropies in bits: $h = H(\tfrac{9}{25})/\ln 2 = H(\tfrac{16}{25})/\ln 2$ and $h' = H(\tfrac{25}{169})/\ln 2$ are the per-quantum entropies of the model's two kinds of pair, and $C$ is $|W|$. Events 1 and 2 are Phase I ($S_R = \Sth$ exactly); events 3 to 6 are Phase II, where $S_R = \sum_{k \in W}h_k$ lies strictly below the capacity whenever the window is nonempty, because no pair is maximally entangled, by exactly the shortfall Theorem \ref{c:singularity.eviction-window-thermal-product-radiation-entropy} quantifies. The descended reading of the same history, monotone with no falling branch, is Proposition \ref{c:singularity.two-levels-finite-model}(2).}

\begin{proposition}[The model's entropy sequence, exactly]\label{c:singularity.models-entropy-sequence-exactly}
In the run of the table accompanying \ref{c:singularity.purity-correlation-transfer-recovery}, Theorem \ref{c:singularity.eviction-window-thermal-product-radiation-entropy}(1) reads as follows: at every stage the window state $\rho_L$ is diagonal in the number basis, with the product distribution of its current members, so that, with $W(t)$ the set of pairs whose window half is still in the window,
\[
S_R(t) \;=\; S_L(t) \;=\; \sum_{k \in W(t)} H(c_k^{2}) \qquad \text{exactly, and} \qquad \tr\rho_R(t)^{2} \;=\; \prod_{k \in W(t)}\bigl(c_k^{4} + s_k^{4}\bigr) .
\]
Consequently $S_R = \Sth$ exactly through Phase I; in Phase II, $S_R \leq C$ with equality if and only if every window member is maximally entangled, so that with the model's pairs the inequality is strict at every Phase II stage with a nonempty window; and $S_R = 0$ at completion. The model therefore realises the envelope, the exact rising branch and the endpoint of Theorem \ref{c:singularity.eviction-window-thermal-product-radiation-entropy}, and exhibits the falling branch's departure from exact saturation, not its saturation.
\end{proposition}
\begin{proof}
The radiation contains $R_k$ for every injected pair $k$, and the window contains $L_k$ for $k \in W(t)$ only. By the support constraint $R_k = L_k$ (Lemma \ref{c:singularity.parity-protection}), two global configurations that agree on every radiation qubit agree on every window qubit, so tracing out the radiation leaves no off-diagonal entries: $\rho_L$ is diagonal, with entries the probabilities of the window configurations. A window configuration is a choice of $b_k \in \{0, 1\}$ for $k \in W(t)$, and its probability is $\prod_{k \in W(t)}(c^{(k)}_{b_k})^{2}$ times the sum over the remaining pairs' configurations, which is $1$; the controlled-$Z$ phases square away. So the distribution is the product of the members' thermal distributions, its Shannon entropy is $\sum_{k \in W(t)}H(c_k^{2})$, its collision probability is $\prod_{k \in W(t)}(c_k^{4} + s_k^{4})$, and $S_R = S_L$, $\tr\rho_R^{2} = \tr\rho_L^{2}$ by purity of the global state. The Phase I identity is the first two rows of the table; $C = |W(t)|\ln 2$ and $H(c_k^{2}) \leq \ln 2$ with equality only at $c_k = s_k$ give the Phase II clause; and $W = \varnothing$ at completion gives $0$. The script computes the exact spectrum of $\rho_L$ at every stage, verifies the entropy identity as an identity of formal logarithms, and decides the strict inequalities by exact rational bounds on the logarithms of the primes involved.

\end{proof}

\subsection*{The descent, applied: the two levels in numbers}
\phantomsection\label{c:note.singularity.setting-descent-applied-two-levels-numbers}
Everything so far in this section lies at the ascended level: the states, marginals and correlators of the completed description. The chain's conclusion, however, is a comparison of two levels, so the model must also be read in the second one. This subsection does exactly that: it applies the multi-mode descent $\EN$ of Definition \ref{c:obj.singularity.multi-mode-descent} to the radiation itself and tabulates both readings side by side; per Remark \ref{c:note.singularity.where-descent-acts-hawkings-loss-derived}, this is the counterfactual reading, the one the closure does not perform, and it is where the entire nonzero deficit lies, the window marginal being an $\EN$ fixed point at every stage. Every number below is a rational number, computed and verified exactly.

\begin{proposition}[The two levels of the finite model]\label{c:singularity.two-levels-finite-model}
\changedmark{gray}{reworded}
In the full model run of the section ``The model, run exactly'' ($n = 4$ pairs, $N_{\max} = 2$), write $\rho_R$ for the state of the complete radiation at evaporation's end (all eight qubits) and $E := \EN$ for the joint-number dephasing. Then:
\begin{enumerate}
\item (Ascended level.) $\tr \rho_R^2 = 1$: $S_2(\rho_R) = 0$. The radiation is pure, and all information is internal to it (Proposition \ref{c:singularity.purity-correlation-transfer-recovery}).
\item (Descended level.) Exactly,
\[
\tr\bigl(E\rho_R\bigr)^2 \;=\; \prod_{k=1}^{4}\bigl(c_k^4 + s_k^4\bigr) \;=\; \frac{817544276833}{6972900390625},
\qquad
S_2(E\rho_R) = \log_2\frac{6972900390625}{817544276833} \approx 3.09 \ \text{bits},
\]
and the von Neumann entropy of $E\rho_R$ is $\sum_k H(c_k^2) = \Sth$ exactly, the full thermal accumulation. Stage by stage, the descended entropy is monotone non-decreasing: read at the $\Cl(1,3)$ level, the radiation's entropy only rises, with no falling branch and no return to zero. This is Hawking's curve, recovered as the wrong-level reading of a history whose ascended state is pure.
\item (The deficit, located.) $\rel{\rho_R}{E\rho_R} = S(E\rho_R) - S(\rho_R) = \Sth$: at completion the radiation's entire information content is the deficit, and the deficit is inter-quanta. Every single-qubit marginal is identical at the two levels (Corollary \ref{c:singularity.deficit-carried-quanta-correlations} with Lemma \ref{c:singularity.parity-protection}); the decoding correlators of Proposition \ref{c:singularity.purity-correlation-transfer-recovery}, whose signatures $(\tfrac{24}{25}, \tfrac{168}{625})$ against $(\tfrac{168}{625}, \tfrac{24}{25})$ distinguish the scrambled from the unscrambled history in the completion, are annihilated exactly by $E$: both become $0$, in both histories, so the two descended histories are indistinguishable by any metric-level measurement, that is, by any joint-number-diagonal observable (Proposition \ref{c:singularity.p-reduced-metric-level-reading-many}; under the total-parity grading alone the even correlator $X_{R_1}X_{L_1}Z_{L_2}$ would descend, which is exactly why the pointwise reading is the one the chain needs), and the message does not descend.
\item (What survives, and what is lost.) Descent preserves each intra-pair classical record exactly: the $(R_k, L_k)$ marginal of $E\rho_R$ is the perfectly correlated $\operatorname{diag}(c_k^2, 0, 0, s_k^2)$. And descent destroys every inter-pair correlation exactly: the joint distributions of $(R_k, R_l)$, $k \neq l$, factorise. What the metric level retains is the thermal populations together with the pairwise emission records; everything the scrambling wrote \emph{across} the pairs, which is where Proposition \ref{c:singularity.purity-correlation-transfer-recovery} located the information, is erased.
\end{enumerate}
\end{proposition}
\begin{proof}
Computation in exact rational arithmetic, deposited with the script. (1) is Proposition \ref{c:singularity.purity-correlation-transfer-recovery}(1). (2) The diagonal of $E\rho_R$ is the configuration distribution of the pure state; by the support constraint ($R_k = L_k$) a configuration is a choice of $b_k \in \{0,1\}$ per pair with probability $\prod_k (c^{(k)}_{b_k})^2$, the accumulated CZ phases squaring away, so the distribution is a product over pairs, its collision probability is $\prod_k (c_k^4 + s_k^4)$, and its Shannon entropy is $\sum_k H(c_k^2) = \Sth$; the monotonicity claim is checked stage by stage (the script also confirms that pinching never increases purity at any stage). (3) is Theorem \ref{c:singularity.deficit-theorem-cost-descent-exactly} with $S(\rho_R) = 0$ and Corollary \ref{c:singularity.deficit-carried-quanta-correlations}, with Lemma \ref{c:singularity.parity-protection} for the single-qubit marginals; the correlator annihilations are checked exactly, in both histories ($X_{R_1}X_{L_1}$ and $X_{R_1}X_{L_1}Z_{L_2}$ are off-diagonal in the joint number basis, so the general vanishing is also visible by inspection of Definition \ref{c:obj.singularity.multi-mode-descent}). That the joint-number-diagonal observables are the metric-level measurements is \ref{c:singularity.p-reduced-metric-level-reading-many}. (4) Both clauses are read off the diagonal of $E\rho_R$: the intra-pair record is the support constraint surviving as classical correlation, and the inter-pair factorisation is the product structure exhibited in (2); both are checked exactly, for every pair and for the sampled inter-pair joints.

\end{proof}

\begin{remark}[The dissolution, now at both levels]\label{c:note.singularity.dissolution-now-both-levels}
the table accompanying \ref{c:singularity.purity-correlation-transfer-recovery} and Proposition \ref{c:singularity.two-levels-finite-model} are the paradox in miniature, every entry exact. In the completion: purity, the Page shape, a recoverable message. In the original space: exact thermality, monotone entropy, no message, and no trace that a message was ever present. Same history, two levels, and the map between them is the explicit dephasing $\EN$, whose cost is the theorem $\rel{\rho}{\EN\rho} = C(\rho) - C(\EN\rho)$. What the metric level retains is exactly what survives $\EN$, populations and per-pair records; what it cannot see is exactly what carries the information. Thermality is what the projection records in place of what it cannot carry (Remark \ref{c:note.singularity.thermality-records-what-projection-cannot-see}); the deficit theorem gives the total, carried by inter-quanta correlations and insensitive to their pattern; and the model exhibits it in exact arithmetic.
\end{remark}

\section{The modular flow of the boundary modes inside, and the entangling half of (L$_E$)}

\phantomsection\label{c:note.singularity.setting-modular-flow-boundary-modes-inside}
Two results that the chain's later readers supply are recorded here at their own level. The first follows the modular flow of Theorem \ref{c:singularity.modular-structure-completed-pair} from the horizon, where Proposition \ref{c:singularity.plane-rate} fixes its rate, into the interior, where the geometry paper's weighted limit gives the collar an exact frame: the flow is the holonomy along the $t$-lines of the exceptional set, and the inverse temperature its Hubble-normalised rate would fix, for any KMS state of that flow on the collar, is a universal number, $6\pi$; whether such a state exists is left open. The second reads Premise \ref{c:singularity.l-boundary-mode-locality} against the Lie algebra of the boundary mode space: the half of (L$_E$) that Theorem \ref{c:singularity.class-relative-no-go-no-supplied} proves within the supplied class lies in the stabiliser of the vacuum, and of the half it premises, the pairing part is exactly the tangent space of the vacuum's orbit, the only part that can move a Witt decomposition, while the differential part of Remark \ref{c:note.singularity.declared-residual-size} stays in the stabiliser.

\begin{assembled}[Imported: the interior collar in Kasner form]\label{c:singularity.imported-interior-collar-kasner-form}
\changedmark{BurntOrange}{changed since the deposit}
In units $2M = 1$ the Schwarzschild interior gravitational field $(\hb, \Omega)$ on flat spacetime (Definition \emph{Gauge Theory Gravity: the gravitational field on flat spacetime}), in synchronous gauge, has derived metric $d\tau^2 - A\,dt^2 - r^2 d\Omega^2$ with $A = 1/r - 1$ and $\tau = \tfrac23 r^{3/2}(1 + \tfrac{3}{10}r + O(r^2))$; the orthonormal physical frame $(e_0, e_1, e_2, e_3)$, $e_0 = g_\tau$, adapted to $(t, \theta, \varphi)$ is exactly $D$-parallel along the free-fall worldlines of velocity $g_\tau$; the rotation-gauge field on the unit direction $e_1$ along a $t$-line is $-\sigma_1 B_1$ with $\sigma_1 = \partial_\tau \log\sqrt{A}$ and $B_1 = e_0 \wedge e_1$ the boost of the $(t, r)$ plane; the Hubble-normalised rates are $\tau\sigma_1 = -\tfrac13 - \tfrac{4}{15}R + O(R^2)$ and $\tau\sigma_2 = \tau\sigma_3 = \tfrac23 - \tfrac{2}{15}R + O(R^2)$ with $R = (3\tau/2)^{2/3}$, whose limits are the Kasner exponents $(p_1, p_2, p_3) = (-\tfrac13, \tfrac23, \tfrac23)$; the three limiting boosts $B_i$ generate $\mathfrak{so}(1,3)$ and do not commute; and the exceptional set of the weighted blow-up is $\{\tau = 0\} \times \R_t \times S^2$. \emph{Provenance:} \cite[Lemmas 5.1 to 5.3, Def.~5.5, Thm.~5.6]{ButlerT6}; the exponents are classical.
\end{assembled}
\begin{proof}
The collected statements are read in Gauge Theory Gravity through Fact \emph{Imported: the dictionary between Gauge Theory Gravity and curved-spacetime language}: the metric is the derived metric, parallelism is $D$-parallelism, and the Levi-Civita connection is the rotation-gauge field $\Omega$. In units $2M = 1$ the Schwarzschild interior takes the synchronous form $d\tau^2 - A\,dt^2 - r^2 d\Omega^2$ with $A = 1/r - 1$ and $\tau = \tfrac23 r^{3/2}(1 + \tfrac{3}{10}r + O(r^2))$; the orthonormal frame $(\partial_\tau, e_1, e_2, e_3)$ adapted to $(t, \theta, \varphi)$ is exactly parallel along $\partial_\tau$; the Levi-Civita connection on the unit direction $e_1$ along a $t$-line is $-\sigma_1 B_1$ with $\sigma_1 = \partial_\tau \log\sqrt{A}$ and $B_1 = e_0 \wedge e_1$ the boost of the $(t, r)$ plane; the Hubble-normalised rates are $\tau\sigma_1 = -\tfrac13 - \tfrac{4}{15}R + O(R^2)$ and $\tau\sigma_2 = \tau\sigma_3 = \tfrac23 - \tfrac{2}{15}R + O(R^2)$ with $R = (3\tau/2)^{2/3}$, whose limits are the Kasner exponents $(p_1, p_2, p_3) = (-\tfrac13, \tfrac23, \tfrac23)$; the three limiting boosts $B_i$ generate $\mathfrak{so}(1,3)$ and do not commute; and the exceptional set of the weighted blow-up is $\{\tau = 0\} \times \R_t \times S^2$. \emph{Provenance:} \cite[Lemmas 5.1 to 5.3, Def.~5.5, Thm.~5.6]{ButlerT6} (\ref{c:interior.interior-connection-exactly}, \ref{c:interior.kasner-form-interior}, \ref{c:interior.weighted-limit-theorem}); the exponents are classical.

This restates \ref{c:interior.exact-parallelism-convergence-rate}.

\end{proof}

\begin{proposition}[The modular flow of the boundary modes as a holonomy, and its universal temperature]\label{c:singularity.modular-flow-boundary-modes-holonomy-universal}
\changedmark{gray}{reworded}
In the collar of Fact \ref{c:singularity.imported-interior-collar-kasner-form}:
\begin{enumerate}
\item[(i)] the transport along a $t$-line at fixed $r$ is a boost in the plane $B_1$ of the completed pair, with rapidity per unit Killing time $|A'(r)|/2 = 1/(2r^2)$ (in general units $M/r^2$), an identity of rational functions; at the horizon $r = 1$ it equals $\tfrac12 = \kappa$, the rate of Proposition \ref{c:singularity.plane-rate}, so the modular flow of Theorem \ref{c:singularity.modular-structure-completed-pair} with the period of Premise \ref{c:singularity.r-regularity-fixed-locus} is the $t$-line holonomy at the horizon; inside, the rate grows as $r^{-2}$ and the Killing-time KMS period $2\pi/(|A'|/2) = 4\pi r^2$ shrinks with it;
\item[(ii)] Hubble-normalised, the same rate is $\tau\,|A'|/(2\sqrt{A}) = \tfrac13 S(r)/\sqrt{1 - r} = \tfrac13 + \tfrac{4}{15}r + \cdots$ with $\tau = \tfrac23 r^{3/2} S(r)$, whose limit at the exceptional set is $|p_1| = \tfrac13$; the period $2\pi$ of the boost's rapidity divided by this limit is
\[
\beta_{\mathrm{H}} \;=\; \frac{2\pi}{|p_1|} \;=\; 6\pi ,
\]
independent of $M$ and of the position on the exceptional set. If the KMS state of the boundary modes on the completed pair extends to a KMS state of the $t$-line holonomy flow on the collar, then its inverse temperature in Hubble-normalised time at the exceptional set is this $\beta_{\mathrm{H}}$; whether it so extends is open (the section ``What is established, what is premised, what is open'');
\item[(iii)] for the sphere directions the corresponding number would be $2\pi/p_2 = 3\pi$. The limiting boosts $B_1$, $B_2$, $B_3$ pairwise fail to commute, $B_1 \times B_2 = -R_{12} \neq 0$ (commutator product, $A \times B = (AB - BA)/2$), so no two of the three limiting flows have proportional generators. Consequently, on any von Neumann algebra carrying the three flows as automorphism groups, a faithful normal state that is KMS at some $\beta \neq 0$ for one of them is KMS at no $\beta' \neq 0$ for either other, by the uniqueness of the modular group; at $\beta = 0$ a trace is KMS for every flow, and is excluded. This series constructs no such algebra for the collar; on the one algebra this paper does construct, the matrix factor of Theorem \ref{c:singularity.modular-structure-completed-pair}, only $B_1$ acts. The direction the modular flow of the boundary modes selects is $B_1$, the $t$-direction, by (i).
\end{enumerate}
\end{proposition}
\begin{proof}
(i) In the synchronous form the $t$-line at fixed $r$ has unit tangent $e_1$ and the connection on it is $\sigma_1 B_1$ per unit proper length $\sqrt{A}\,dt$ of the line, with $\sigma_1 = \partial_\tau\log\sqrt A = \tfrac12 (A'/A)\,dr/d\tau = \tfrac12 A'/\sqrt{A}$, since $dr/d\tau = +\sqrt A$ in the orientation of Fact \ref{c:singularity.imported-interior-collar-kasner-form}, where $\tau$ increases with $r$ (and $\sigma_1 < 0$ there, $A$ decreasing in $r$, matching the sign of $\tau\sigma_1 \to -\tfrac13$); per unit Killing time the rapidity rate is therefore $\sigma_1\sqrt A = \tfrac12 A'(r)$, of magnitude $-\tfrac12 A'(r)$, which for $A = 1/r - 1$ is $1/(2r^2)$: the script verifies $-\tfrac12 A' = 1/(2r^2)$ as an identity in $\mathbb{Q}(r)$, its value $\tfrac12$ at $r = 1$, and the values $2, 8, 32$ at $r = \tfrac12, \tfrac14, \tfrac18$. That $\kappa = 1/4M = \tfrac12$ in these units is Proposition \ref{c:singularity.plane-rate}, and that the modular flow is the boost of this plane at the geometric period is Theorem \ref{c:singularity.modular-structure-completed-pair} with Premise \ref{c:singularity.r-regularity-fixed-locus}. (ii) With $\tau = \tfrac23 r^{3/2} S(r)$, $S = 1 + \tfrac{3}{10}r + \tfrac{9}{56}r^2 + \cdots$ the series of Fact \ref{c:singularity.imported-interior-collar-kasner-form}, $\tau |A'|/(2\sqrt A) = \tfrac23 r^{3/2} S \cdot \tfrac{1}{2r^2} \cdot \sqrt{r/(1-r)} = \tfrac13 S(r)(1 - r)^{-1/2}$; the script forms the product of exact truncated series and reads off $\tfrac13 + \tfrac{4}{15}r$, whose constant term is $|p_1|$ (Fact \ref{c:singularity.imported-interior-collar-kasner-form}: $\tau\sigma_1 \to -\tfrac13$, the sign being that of $A' < 0$), and likewise the sphere rate $\tau\sigma_2 = \tau\sqrt{A}/r = \tfrac23 S(r)(1 - r)^{1/2} = \tfrac23 - \tfrac{2}{15}r + \cdots$, whose constant term is $p_2$; the script then divides $2\pi$ by each constant term, obtaining $6\pi$ and $3\pi$ from the series themselves. The conditional clause is the definition of the KMS condition: a KMS state of the flow generated by a boost of rapidity rate $\rho$ per unit parameter has inverse temperature $2\pi/\rho$ in that parameter (Theorem \ref{c:singularity.modular-structure-completed-pair}, Premise \ref{c:singularity.r-regularity-fixed-locus}, applied along the $t$-line), and the rate at the exceptional set in Hubble-normalised time is $|p_1|$. (iii) The commutator is computed blade-exactly in $\Cl(1,3)$. If a faithful normal state is KMS at $\beta \neq 0$ for a flow $\alpha_t$ then its modular group is $\alpha_{-\beta t}$, and the modular group is unique (Tomita-Takesaki theory, \cite[Sec.~5.3]{BR2}; Fact \ref{c:singularity.imported-general-form} imports the state-side consequences of the same theory), so a second flow $\alpha'_t$ with the state KMS at $\beta' \neq 0$ would satisfy $\alpha_{-\beta t} = \alpha'_{-\beta' t}$ for all $t$, making the generators proportional, which the commutator excludes for the pairs among $B_1$, $B_2$, $B_3$.

\end{proof}

\begin{remark}[Where a spectrum would lie]\label{c:note.singularity.where-spectrum-would-lie}
The proposition fixes the thermal direction and does not fix the sphere: the KMS condition for the $B_1$ flow determines no state on the holonomies along the $\theta$- and $\varphi$-lines of the exceptional set, but it is not silent about them either, since the $B_1$ flow moves the angular subalgebra ($[B_1, B_2] = -R_{12}$ lies outside the span of $B_2$, $B_3$, $R_{23}$), so a state on the algebra the three boosts generate that is KMS for $B_1$ is constrained on the sphere holonomies through that action, and which $SO(3)$-invariant states meet the constraint is not settled here. That is where a spectrum, in the sense of a two-point function on the exceptional set, would have to come from, and it is recorded as a successor problem in the section ``What is established, what is premised, what is open'', beside the question whether the KMS state of the boundary modes on the completed pair extends to a KMS state at $\beta_{\mathrm{H}} = 6\pi$ on the algebra of $t$-line holonomies of the whole collar, whose obstruction, if any, is the $O(\tau^{2/3})$ running of the rate before its limit.
\end{remark}

\subsection*{The entangling half of (L$_E$)}
\phantomsection\label{c:note.singularity.setting-entangling-half-l-e}
Theorem \ref{c:singularity.class-relative-no-go-no-supplied} proves the no-hopping half of Premise \ref{c:singularity.l-boundary-mode-locality} within the supplied class and leaves the no-pairing half premised; Remark \ref{c:note.singularity.declared-residual-size} names two residual generator types, differential (kinetic) terms and pairing terms. The following proposition says which of the two can touch entanglement, and it is exactly one.

\begin{proposition}[(L$_E$) as the stabiliser and the tangent]\label{c:singularity.l-stabiliser-tangent}
On the Fock space of $\nu \geq 2$ boundary modes with creation operators $v^*_i$, the bivectors of $\mathfrak{so}(2\nu)$ are the quadratic elements $v^*_i v_j - \tfrac12\delta_{ij}$, $v_i v_j$ and $v^*_i v^*_j$, and split as hops (the $\nu(\nu - 1)$ elements $v^*_i v_j$ with $i \neq j$ together with the $\nu$ Cartan elements $v^*_i v_i - \tfrac12$), pair annihilations $v_i v_j$ and pair creations $v^*_i v^*_j$ ($\nu(\nu - 1)/2$ each), $\nu^2 + \nu(\nu - 1) = \nu(2\nu - 1) = \dim\mathfrak{so}(2\nu)$; the same elements without the $\tfrac12$ span, with the identity, $\mathfrak{so}(2\nu) \oplus \C$, and are not themselves closed under the bracket. On the vacuum $|0\rangle$ every hop with $i \neq j$ and every pair annihilation vanishes and every Cartan element acts as the scalar $-\tfrac12$, so the hops and the pair annihilations together are the stabiliser of the vacuum ray; the pair creations applied to $|0\rangle$ are linearly independent and span, modulo $|0\rangle$, the tangent space at $|0\rangle$ of the orbit of Witt vacua, of complex dimension $\nu(\nu - 1)/2$, while the tangent space of the orbit of the vector $|0\rangle$ itself, the span of $X|0\rangle$ over $X \in \mathfrak{so}(2\nu)$, has dimension $1 + \nu(\nu - 1)/2$, the extra dimension being the scalar action of the Cartan elements. Consequently every hop rotation $\exp\theta(v^*_i v_j - v^*_j v_i)$ fixes the vacuum and every pairing rotation $\exp\theta(v^*_i v^*_j - v_j v_i)$ moves it; the necessity witness of Corollary \ref{c:singularity.envelope-endpoint-why-boundary-mode-locality}, $|00\rangle \mapsto \tfrac35|00\rangle + \tfrac45|11\rangle$, is the pairing rotation on modes $(1,2)$ at $(\cos\theta, \sin\theta) = (\tfrac35, \tfrac45)$. The half of Premise \ref{c:singularity.l-boundary-mode-locality} that Theorem \ref{c:singularity.class-relative-no-go-no-supplied} proves, within the supplied class, therefore lies in the stabiliser of the vacuum; of the half it premises, the exclusion of unsupplied generators, the pairing terms are the tangent space of the vacuum's orbit, the part that changes the Witt decomposition and the only part that can move the vacuum, and so the only part that can create entanglement from it, at one dimension per pair of modes, while the differential terms of Remark \ref{c:note.singularity.declared-residual-size} are hops and lie in the stabiliser.
\end{proposition}
\begin{proof}
The count is the dimension of $\mathfrak{so}(2\nu)$ on $2\nu$ Majorana generators (its elements are the bivectors, \ref{c:std-algebra.bivector-lie-algebra}): the bivectors $\tfrac12[\gamma_a, \gamma_b]$ in the Majorana basis $\gamma_{2i-1} = v_i + v^*_i$, $\gamma_{2i} = i(v_i - v^*_i)$ are, up to scalars, the elements listed, the diagonal ones being $v^*_i v_i - v_i v^*_i = 2(v^*_i v_i - \tfrac12)$ by the anticommutation relations (\ref{c:std-algebra.witt-basis-car}); without the $\tfrac12$ the bracket $[v^*_i v^*_j, v_j v_i] = v^*_i v_i + v^*_j v_j - 1$ leaves the span of the listed elements and lands in their span with the identity. On the vacuum $v_j|0\rangle = 0$ annihilates every off-diagonal hop and every annihilation, and $v^*_i v_i|0\rangle = 0$ gives the scalar $-\tfrac12$ for the Cartan elements; $v^*_i v^*_j|0\rangle = |ij\rangle$ are distinct basis states, hence independent and independent of $|0\rangle$, so the span of $X|0\rangle$ over $X \in \mathfrak{so}(2\nu)$ is that of $|0\rangle$ and the $|ij\rangle$, of dimension $1 + \nu(\nu - 1)/2$. The orbit of the ray of $|0\rangle$ under $\mathrm{Spin}(2\nu, \C)$ is the projective variety of pure spinors, of dimension $\nu(\nu - 1)/2$ (Cartan's theorem on pure spinors \cite{Chevalley1954}, \ref{c:std-algebra.pure-spinors}), and the orbit of the vector is the cone over it, of dimension one more, so the creations span the whole tangent space of the ray orbit and, with the scalar direction, the whole tangent space of the vector orbit. The rotations are the exponentials of the corresponding generators, and the hop rotations fix $|0\rangle$ because their generators annihilate it. The script performs the count and the closure of the corrected generators under the bracket at $\nu = 2, 3, 4$, the annihilation and scalar statements, the tangent ranks $1 + \nu(\nu - 1)/2$ and $\nu(\nu - 1)/2$ by exact rank computations on the images of all $\nu(2\nu - 1)$ generators, the fixing and moving at $\nu = 4$ for every pair by the exponentials themselves, formed in closed form and verified orthogonal, the identification of the witness, and, as the converse the statement needs, that a hop rotation applied off the vacuum does entangle. A differential term of Dirac type, being first order in the modes, is a hop in this decomposition and lies in the stabiliser: of Remark \ref{c:note.singularity.declared-residual-size}'s two residual types, only the pairing term is in the tangent.

\end{proof}

\begin{remark}[What this does to the residual]\label{c:note.singularity.what-does-residual}
Remark \ref{c:note.singularity.declared-residual-size}'s question, whether the physical dynamics on the singular boundary is generated within the supplied class or carries a differential or pairing term, splits into two halves of different kinds. A differential term would hop occupation between points of $E$ and would break the joint diagonality that Step 6b uses, and off the vacuum it entangles, as any hop does (the script exhibits one on $|10\rangle$); but it cannot move the vacuum, and so cannot create entanglement from it. A pairing term moves the vacuum, and only a pairing term does. The premised content of (L$_E$) is thereby separated into its two parts, both class-relative in the sense of Theorem \ref{c:singularity.class-relative-no-go-no-supplied}: the exclusion of pairing terms, which is the statement that the closure supplies no generator in the $\nu(\nu - 1)/2$-dimensional tangent of the vacuum ray, and the exclusion of the unsupplied hops, the differential dynamics, which lie in the stabiliser and are excluded by the premise for the diagonality Step 6b uses rather than for what they do to the vacuum. Its necessity, Corollary \ref{c:singularity.envelope-endpoint-why-boundary-mode-locality}, is the statement that one generator in the tangent, acting once on one pair, is what the return of the boundary modes needs. The framework compares two Witt decompositions and does not derive them; nothing here changes that, and the derivation of the one pairing angle the chain uses, the angle of Premise \ref{c:singularity.r-regularity-fixed-locus}, remains outside its reach (Proposition \ref{c:singularity.why-calibration-cannot-be-internal}).
\end{remark}

\section{Machine verification}

\phantomsection\label{c:note.singularity.setting-machine-verification}
Every computational identity in this paper, chain, gauge sector, state selection and model alike, has been checked by computer in exact rational arithmetic, in the manner of the verifications deposited with the earlier papers of the series \cite{ButlerT1, ButlerT3, ButlerT4}. The script \texttt{verify\_T05.py}, deposited with this paper alongside the series' shared exact-arithmetic library \texttt{GA.py} \cite{ButlerVerifyT05}, performs two hundred and thirty-three checks in nine kinds of exact arithmetic: the basis-blade Clifford product over the rationals, which also carries the point-level transport of the section ``The point-level transport, instantiated'' in $\Cl(2,4)$ with the vector slot represented; rational-function arithmetic over $\mathbb{Q}(r)$ for the Schwarzschild fibre, in which the scalar identities of the point-level transport (the null direction, the dual norms, T04's Gram and its dual) are verified as identities rather than at samples; arithmetic for the field $\mathbb{Q}(r)[\beta]$ with $\beta^2 = 2M/r$, in which the gauge sector is verified as identities rather than at samples; fermionic Fock-space arithmetic on a finite site array for the current algebra; a multi-qubit simulator with rational amplitudes (controlled-$Z$ scrambling contributes only signs) for the finite-mode model; rational-function arithmetic in the radial variable for the appendix, so that the curvature identities are verified as identities of rational functions; formal-logarithm arithmetic, in which the von Neumann entropy of a rational spectrum is an exact vector over the primes, identities between entropies are dictionary equalities, and inequalities are decided by rigorous rational bounds on the logarithms of the primes involved; exact truncated power series over the rationals, for the Hubble-normalised rates of the section ``The modular flow of the boundary modes inside, and the entangling half of (L$_E$)''; and Gaussian-rational arithmetic on the Jordan-Wigner Fock module, for the $\mathfrak{so}(2\nu)$ computations of the same section.

The chain coverage. Step 1 in full: $P^2 = 0$ on every basis vector, the closed-form inverse at rational $\beta$ and at the horizon, the null locus of the metric exactly at $\lambda = -\beta \pm 1$, computed through $\hb^{-1}$, with a non-null control value and with the $\hb$-images of the flat null vectors confirmed $g$-null, and the monotone folding of the normalised coefficients; the gauge sector: $\beta' = -\beta/2r$, the Newtonian-field identity $\beta\beta' = -M/r^2$ and the divergence rate $(\beta')^2 = M/2r^3$ as identities in $\mathbb{Q}(r)[\beta]$, the torsion equations of the PG coframe on the $\theta^0$, $\theta^1$ and $\theta^2$ legs as exact two-form identities, the radial curvature $\mathcal{R} = (\beta\beta')'\,\sigma_r = (2M/r^3)\sigma_r$ with the abelian vanishing of the wedge term, unimodularity of $\hb$ at sample radii including the horizon, the determinant computed from the map $a \mapsto a - \beta(a \cdot e_t)e_r$ itself, the rescaled-limit identities of Proposition \ref{c:singularity.degeneration-exactly} on all sixteen basis pairs with the limit form's exact signature $(0,1,3)$ by congruence diagonalisation, the rank, kernel and image of $P$, the nullity computation of the folded direction in both forms, and the Kretschmann identity $48M^2/r^6$. Step 2: the mod-8 arithmetic, and the centre of $\Cl(1,3)$ computed blade-wise and exhaustively. Steps 3, 4a and 4c in the hyperbolic atom with rational normalisation: CAR, pseudo-inverse identities, the exact-rank proof that $v$ has no inverse, the complementary idempotents, $\sigma = e_0 e_1$ with $\sigma^2 = 1$, the rotor flow with the rational pair $(\cosh, \sinh) = (\tfrac53, \tfrac43)$ scaling $v$ and $v^*$ by exactly $\tfrac19$ and $9$, and the spectral-equals-descent projector identity. Step 4b: the shape operator $S = -\tfrac12(B + 3\sigma_r B\sigma_r) = 1 - 3\Pi$ verified blade-exactly on all six frame planes, $S^2 + S = 2$ as an operator identity, the spectrum $(-2,-2,1,1,1,1)$ with vanishing trace, the eigenvalue ratio $-2 : 1$ against the appendix's rational-function values, the sign $\mathcal{R} = -(M/r^{3})S$ on the radial, dual and mixed planes, the $\hb$-invariance of the radial plane, and the matching characterisation (eigenvalue $-2$ plus square $+1$ single out the radial plane among the six, exhaustively); and Proposition \ref{c:singularity.plane-rate} directly: $\kappa = f'(2M)/2 = 1/4M$ from the derivative, the conversion $2\pi/\kappa = 8\pi M$ with $\pi$ cancelled, and the invariant separation $\sigma_r^{2} = +1$, $(I\sigma_r)^{2} = -1$ blade-exactly. Step 4d: the finite-dimensional KMS identity verified on all sixteen matrix-unit pairs for $\rho \propto \operatorname{diag}(1, x)$ at both $x = e^{-\omega/T} = \tfrac{9}{16}$ (the physical, positive-temperature point) and $x = \tfrac{16}{9}$ (population-inverted, so that both signs of the inverse temperature are exercised), the thermofield reductions $\operatorname{diag}(\tfrac{16}{25}, \tfrac{9}{25})$ and $\operatorname{diag}(\tfrac{9}{25}, \tfrac{16}{25})$ with $\tan^2\theta = e^{-\omega/T}$ at the two points, and the two examples bounding Theorem \ref{c:singularity.modular-descent}'s hypotheses ($c|01\rangle + s|10\rangle$ descending thermally without (ii); $c|{+}0\rangle + s|{-}1\rangle$ descending with nonzero coherence); and Lemma \ref{c:singularity.schmidt-forcing-proved-outright} as an exact eigenspace computation over $\mathbb{Q}$: the generator's $0$-eigenspace computed exactly, a $\pm 1$ eigenvector shown to be a product vector failing separation and cyclicity by exact rank, a $0$-eigenvector with both coefficients nonzero shown cyclic and separating by exact rank, and the invariance-up-to-phase clause discriminating the two. Step 4e: the current commutators on bulk states, which is also how the script fixed the index convention; the exact edge-operator identity $[a_m, a_{-m}] = B_m - T_m$ verified as an operator identity on sea, bulk and edge-excited states for $m = 1, 2$; the detailed-balance and difference identities of the Planck computation verified pointwise, the telescoping re-indexing verified at finite truncation, and the closed-form value confirmed against an explicit exact geometric tail bound; and premise (Q) isolated as an operator computation: the number operators built as matrices, the Planck summand $\langle n_{k+m}(1 - n_k)\rangle$ taken as a trace against the state and verified to factorise into its two-point pieces on a quasi-free product state carrying Theorem \ref{c:singularity.rate-discharge-hypothesis-iii}'s gluing occupancy $n = \tfrac{9}{25}$, and the connected correlator computed outright on a correlated even state with identical marginals, the summand shifting by exactly it. Step 6a: the exact slice volume $r_c^{4}|f(r_c)| = (2M)^{2}(1 - r_c/2M)$ verified as an identity of rational functions along the whole (C1) locus ($2M = r_c^{3}$, $r_c$ an indeterminate) with a mass-independent-cutoff control failing the linear-in-$M$ form, the identity $(dN_{\mathrm{6a}}/dt)^{2} = (\nu/\TH)^{2}(1 - r_c/2M)$ with $\pi$ cancelled, the capacity calibration $C = N_{\mathrm{6a}}\ln 2 = 4\pi M^{2}$ if and only if $\nu\Delta t = M/(2\ln 2)$ verified with $\ln 2$ carried as a formal indeterminate, the 1-nat convention $\nu\Delta t = M/2$ failing by exactly $\ln 2$, and the causal-epoch lemma (the geodesic integrand after the substitution $r = 2Mu^{2}$, the partial-fraction identity and the antiderivative as rational-function identities; the closed form bracketed by exact Riemann sums at the rational endpoint $\varepsilon = 1999 \cdot 10^{-6}$, with the closed form of the un-square-rooted integrand shown to lie outside the bracket, and the expansion checked to $O(\varepsilon)$). Step 6b: Theorem \ref{c:singularity.eviction-window-thermal-product-radiation-entropy}(1) on the thermal model, the window's exact spectrum equal to the product of its members' thermal distributions at every stage and $S_R = \sum_{k \in W}h_k$ as an identity of formal logarithms, with clause (4)'s shortfall identity; the maximally entangled idealisation simulated on exact density matrices along a decreasing capacity schedule, the window's entropy computed from its exact spectrum and equal to $\min(\Sth, C)$ at every event, ending at zero, with the window exactly maximally mixed throughout Phase II; the necessity of (L$_E$) in Corollary \ref{c:singularity.envelope-endpoint-why-boundary-mode-locality}, the two-pair rotation example computed exactly with $S_R > \Sth$ decided by rigorous bounds; and Corollary \ref{c:singularity.what-event-returns-radiation}'s two identities at every event, $S(\rho_A)$, $S(\rho_{R_<})$ and $S(\rho_{AR_<})$ each taken from an exact spectrum; Fact \ref{c:singularity.assembled-page-curve-up-typicality}'s envelope directly on the thermal model, both arms at every stage, the rising branch exact and the crossing of the two arms exhibited; and Proposition \ref{c:singularity.no-attic-derived-model} both kinematically (the two sectors disjoint and exhaustive at every stage) and dynamically (second and third trace moments of $\rho_R$ against $\rho_L$ agreeing exactly, so no q-number content leaves the two-sector decomposition). Five checks verify what the section ``Step 5: descend, and state what descent costs'' through the section ``The finite-mode model: the dissolution step by step'' state about their own limits: the window marginal, verified diagonal entrywise, equals the product of its own single-mode marginals entry by entry and is a fixed point of $\EN$ at every stage, so the reading of the boundary modes costs nothing and the information lies across the window/radiation cut; the saturation boundary is computed through the thermofield angle, $h = \ln 2$ exactly at $x = e^{-\omega/T} = 1$, the one zero-weight mode, strictly below and monotone toward the boundary otherwise; the deficit is insensitive to the message, the scrambled and unscrambled histories carrying identical final configuration distributions and identical deficit $\Sth$ while differing in amplitude signs alone; and Proposition \ref{c:singularity.p-reduced-metric-level-reading-many} takes the last two, one on each front: its composite argument (the mode-wise parity averages idempotent, commuting, and composing in either order to the pinching, on generic rational data), and its antecedent's effect (the total-parity average fixing the model's even radiation state outright, the even correlator $X_{R_1}X_{L_1}Z_{L_2}$ commuting with $\Gamma$ while $\EN$ annihilates it).

The point-level coverage. The instantiated lift of Proposition \ref{c:singularity.transport-schwarzschild-fibre}, parts L1 to L14, with the vector slot represented: the incidence condition checked as the blade-exact nullity $X_p(x)^{2} = 0$ in $\Cl(2,4)$ on tangent vectors with generic rational components, the coordinate slots extracted from the multivector itself, with a wrong-coefficient control that squares to $-Q$, not $0$; the preservation of the quadric value of arbitrary six-component vectors under $\iota$, a non-incident control included; the involution property on the same vectors; the recovery $\dsc(\iota(X_p(x))) = x/Q_p(x)$ on the invertible locus, componentwise on the represented vector; the nullity of the radial null direction $\et + f\er$, symbolically in $\mathbb{Q}(r)$ and blade-exactly at the rational-orthonormal radius $r = 4$, with the vanishing first slot of its transported value exhibiting the ideal point $[\,x + \tfrac12\nInf\,]$ as a multivector identity, the read-out undefined on it, and a perturbed non-null control that descends; the fibre norm $Q_p(\et) = f$; the dual norm of the radical generator $k$, $Q^*_p(k) = -r$, as an identity of rational functions in the static chart and as the contraction $k^{\mathsf T}Gk$ against the Painlev\'e-Gullstrand Gram of \cite[Thm.~8.7]{ButlerT4} built over $\mathbb{Q}(r)$, with the mixed vector $\rho - m\,dt_{\mathrm{PG}}$ as a multi-entry control contraction and $\det G = -r$, and the dual computed through the Gram, $Gk = (-r, 0, 0, 0)$, the PG-time component of $k^{\sharp}$ vanishing identically so that $k^{\sharp}$ is radial in Step 1's own frame; the horizon pole of the static norm $1/f$ of $dt_s$ probed by evaluation; $Q^*_p(k)$ rebuilt pointwise from its two unreduced terms $2M/f$ and $-r/f$ at off-horizon radii, each term separately singular at the horizon where the value $-2M$ is certified by the rational-function identity whose content is exactly that cancellation, and $0$ at $r = 0$, so that $k$ descends at every $r > 0$, horizon included, its transported value tends to the ideal stratum, and the failure locus is $E$ and not a shadow wall.

The descent coverage. The pinching's structural clauses (idempotent, trace-preserving, Hilbert-Schmidt self-adjoint) on a generic rational matrix; the single-mode coincidence of pinching with partial trace and its two-mode failure, quantified by the coherence $c_1 s_1 \gamma$; the deficit theorem's key identity in moment form by full matrix products on a state with nonzero coherences, with a control showing that it fails when the diagonal factor is replaced by the state itself; the invariance of every single-qubit marginal under descent; the two-level values $\tr\rho_R^2 = 1$ against $\tr(E\rho_R)^2 = \prod_k(c_k^4 + s_k^4) = \tfrac{817544276833}{6972900390625}$ and $S(E\rho_R) = \sum_k H(c_k^{2}) = \Sth$ as formal logarithms; the stagewise monotonicity of the descended entropy and the purity-decrease of pinching at every stage; the exact annihilation of both decoding correlators in both histories; the survival of the intra-pair classical records as $\operatorname{diag}(c^2, 0, 0, s^2)$; and the exact factorisation of every sampled inter-pair joint distribution.

The state-selection coverage. Theorem \ref{c:singularity.descent-parity-average}: the identity $E(x) = \tfrac12(x + \sigma x\sigma)$ on a generic rational matrix, idempotence, the conditional-expectation property $E(axb) = aE(x)b$ for even $a, b$, the range through $E(1) = 1$ and $E(\sigma) = \sigma$, and the kernel through $E(v) = E(v^{*}) = 0$ together with the exhaustive statement that $E$ annihilates exactly the off-diagonal part, alongside the atom's CAR with $\sigma = e_0e_1$ blade-exactly. Theorem \ref{c:singularity.descended-form-state-independent}: for a density matrix with nonzero coherences, $E(\varrho)$ is diagonal, is a state, and equals $\operatorname{diag}(1,x)/(1+x)$ at the $x$ the theorem names, with the injectivity of $x \mapsto p$ on a rational witness pair. Theorem \ref{c:singularity.ascended-marginals-are-diagonal}: an even three-mode state, built as a rational mixture of two even pure states, verified even by $\Gamma\varrho\Gamma = \varrho$, with all three single-mode marginals diagonal and $\omega(v_k) = \omega(v_k^{*}) = 0$ at every mode; a parity-mixed control with marginal coherence $\tfrac{12}{25}$, so the premise demonstrably has an effect; and the model's pair state, controlled-$Z$ sign included, verified even with all marginals diagonal. Theorem \ref{c:singularity.kms-uniqueness-factor}: the KMS identity on all sixteen matrix-unit pairs for $\varrho \propto \operatorname{diag}(1, x)$ at $x = \tfrac{9}{16}$ and at $x = \tfrac{16}{9}$, a matrix-unit pair at which the tracial state fails it, and the commutation the identity forces. Corollary \ref{c:singularity.universality-temperature-under-g}: with weights $(1,2)$, equal inverse temperatures as the exact rational condition $x_2 = x_1^{2}$, the joint populations of the two-mode product, a witness pair for which the condition fails, and the factorisation of the (G) state. Theorem \ref{c:singularity.rate-rigidity-temperature-not-algebra-intrinsic}, at both sample points: the diagonal unitary commuting with $K$ and fixing the Gibbs state; the exchange unitary with $XKX = -K$ carrying $\beta$ to $-\beta$ (the two sample points to each other) and failing to commute with $K$; and the reparametrisation clause, the same state being KMS at $\beta$ for the flow of $H$, at $\beta/2$ for the flow of $2H$, and not at $\beta$ for the flow of $2H$. Proposition \ref{c:singularity.euclidean-section-period}: $dF/d\varrho = f'/2$ as an identity of rational functions in $F = \sqrt f$ (both sides $(1 - F^{2})^{2}/4M$), its horizon value $1/4M$, and $f(2M) = 0$. Theorem \ref{c:singularity.rate-discharge-hypothesis-iii}: the gluing equation $e^{p\omega}n = 1 - n$ with its unique solution $n = \tfrac{9}{25}$ at $e^{-p\omega} = \tfrac{9}{16}$ (linearity in $n$ with nonzero coefficient, a probe list of failing candidates), the same equation from the other ladder ordering, the failure of both boundary cases, and the full rank of the fixed marginal. Proposition \ref{c:singularity.why-calibration-cannot-be-internal}: the connection rate $(M/2)/r^{3}$ with its pole at $r = 0$ against the finite horizon rate $\kappa = 1/4M$, and $f(2M) = 0$ marking the fixed locus. Proposition \ref{c:singularity.euclidean-rotor-where-it-lies}: $I\sigma_r = e_2e_3$ blade-exactly, the action of its rotor on all four frame vectors (fixing $e_0$, $e_1$ and $\sigma_r$, rotating $e_2$ into $e_3$), the identity $(a\sigma_r + bI\sigma_r)^{2} = (a^{2} - b^{2}) + 2abI$ that confines the real square-$(-1)$ members of the pair's span to $\pm I\sigma_r$; and, over the Gaussian rationals, the complexified rotor's unitarity, its composition law $R(\tfrac35,\tfrac45)^{2} = R(-\tfrac{7}{25},\tfrac{24}{25})$, $R_{2\pi i} = -1$ on the module and $R_{4\pi i} = +1$ by squaring the quarter rotor, and the triviality of conjugation by $-1$ on a generic complex matrix. Proposition \ref{c:singularity.two-kms-functionals-sign}: the Gibbs state satisfying the ungraded identity on all sixteen matrix-unit pairs and failing the graded one; the parity-twisted functional satisfying the graded identity and failing the ungraded one; its occupancy $-\tfrac{9}{7}$, negative; and $Z = 1 + x$ against $Z_\Gamma = 1 - x$. Theorem \ref{c:singularity.exact-equilibrium-over-determines-boundary-modes}: the two-mode Gibbs product as a fixed point of $\EN$ and as the tensor product of its own marginals, against an entangled even control failing both and losing purity under descent. Theorem \ref{c:singularity.purity-output-partner-commutant}: the thermofield vector purifying $\operatorname{diag}(\tfrac{9}{25},\tfrac{16}{25})$, the angle $\tan^{2}\theta = \tfrac{16}{9}$, the modular operator recomputed as $\rho\otimes\bar\rho^{-1}$, and the invariance of the marginal under a unitary on the purifying factor. Proposition \ref{c:singularity.number-diagonal-dynamics-models-scrambling-derived}: a controlled-$Z$ circuit on the model's four-mode space preserving every joint-number probability and every marginal, commuting with the total parity, and still changing the joint coherence; against an orthogonal parity-preserving hopping rotation that changes the marginal populations, exhibiting exactly what (L$_E$) excludes. Corollary \ref{c:singularity.chains-state-class-stability}, at both sample points: an entangled even two-mode state with both marginals Gibbs at the same $\beta$ ($\operatorname{diag}(\tfrac{16}{25}, \tfrac{9}{25})$ at $x = \tfrac{9}{16}$, $\operatorname{diag}(\tfrac{9}{25}, \tfrac{16}{25})$ at $x = \tfrac{16}{9}$) is in the class and is not a product, the product member with the same marginals is a fixed point of $\EN$, and a controlled-$Z$ keeps the state in the class while changing it.

The model coverage: the support constraint and normalisation at every stage; the parity-protection lemma verified literally (every single-qubit marginal equal to its thermal matrix, at every stage, scrambling included); the Schmidt purity equality $\tr\rho_R^2 = \tr\rho_L^2$ at every stage; final purity of the radiation; the R\'enyi-2 Page shape; the entropy sequence of Proposition \ref{c:singularity.models-entropy-sequence-exactly}, with $\rho_L$'s exact spectrum equal to the product distribution of its members at every stage, $S_R = S_L = \sum_{k \in W}H(c_k^{2})$ as an identity of formal logarithms, every purity of the table accompanying \ref{c:singularity.purity-correlation-transfer-recovery}, the Phase I identity $S_R = \Sth$, the strict Phase II inequality $S_R < C$, and $S_R = 0$ at the end; the coherence damping with $\gamma = c_2^2 - s_2^2$ and the pair-purity formula; and the decoding table, with $(\langle XX\rangle, \langle XXZ\rangle)$ exchanging values between the scrambled and unscrambled histories while all single-qubit data are fixed, including the printed numerical signatures $(\tfrac{24}{25}, \tfrac{168}{625})$.

The appendix coverage: the Painlev\'e-Gullstrand metric read off from $\hb^{-1}$ blade-exactly; the cross-term elimination to static form as an exact congruence $T^{\mathsf T}G_{\mathrm{PG}}T = \operatorname{diag}(f, -1/f)$, the matrices formed and multiplied in $\mathbb{Q}(r)[\beta]$; the Cartan derivation itself, parametrised by $F = \sqrt f$ so that every coefficient is rational: the first structure equations on the $\theta^{0}$, $\theta^{1}$, $\theta^{2}$ legs with the stated connection forms, and the second structure equations producing $R^{0}{}_{1}$, $R^{0}{}_{2}$, $R^{1}{}_{2}$ with the stated coefficients, one perturbed connection coefficient shown to break both structure equations as the block's negative control; the three curvature eigenvalues $2M/r^3$, $-M/r^3$, $2M/r^3$ as exact rational-function identities, the radial and angular ones equal and $-2$ times the mixed one; and $\kappa = f'(2M)/2 = 1/4M$ exactly; the lift of Definition \ref{c:obj.singularity.fibrewise-conformal-lift} is null in $\Cl(2,4)$ with $\nO \cdot \nInf = -1$ and not with $-\tfrac12$, blade-exactly.

The the section ``The modular flow of the boundary modes inside, and the entangling half of (L$_E$)'' coverage, twenty-one checks: in $\mathbb{Q}(r)$, the rational-function identity $-A'/2 = 1/(2r^2)$ with its horizon value and three interior values; in exact truncated series, $\tfrac13 S(r)(1 - r)^{-1/2} = \tfrac13 + \tfrac{4}{15}r + \cdots$ and $\tfrac23 S(r)(1 - r)^{1/2} = \tfrac23 - \tfrac{2}{15}r + \cdots$, with the two inverse temperatures $6\pi$ and $3\pi$ obtained by dividing $2\pi$ by the constant terms of the series so formed; blade-exactly in $\Cl(1,3)$, the commutator product of $B_1$ and $B_2$ equal to $-e_1 e_2$ (Proposition \ref{c:singularity.modular-flow-boundary-modes-holonomy-universal}); and, in Gaussian-rational arithmetic on the Jordan-Wigner Fock module at $\nu = 2, 3, 4$, the linear independence of the $\nu(2\nu - 1)$ generators, their closure under the bracket once the Cartan elements carry the $-\tfrac12$ and the failure of closure without it, the annihilation of the vacuum by every off-diagonal hop and pair annihilation with the scalar $-\tfrac12$ from every Cartan element, the tangent ranks $1 + \nu(\nu - 1)/2$ for the vector orbit and $\nu(\nu - 1)/2$ for the ray orbit as exact ranks of the images of all generators, and at $\nu = 4$ the rotations formed as closed-form exponentials of their generators at a rational point of the unit circle, verified orthogonal, every hop rotation fixing the vacuum and every pairing rotation moving it, a hop rotation on $|10\rangle$ producing the mixed marginal $\operatorname{diag}(\tfrac{9}{25}, \tfrac{16}{25})$, and the identification of the necessity witness (Proposition \ref{c:singularity.l-stabiliser-tangent}).

The Appendix C coverage, on the two-mode Jordan-Wigner algebra and a three-point finite model, all matrices exact rational: the CAR across modes with the strings cancelling in $\sigma_k = v_k^{*}v_k - v_kv_k^{*}$, single-site both, the Jordan-Wigner locality used in the proof of Proposition \ref{c:singularity.p-reduced-metric-level-reading-many}; $\operatorname{Ad}(\sigma_x)$ fixing $v_y$ while $\operatorname{Ad}(\sigma_y)$ negates it (Theorem \ref{c:singularity.net-assembly-net-it-everything}); the composite of the two parity averages, in either order, equal to the joint-diagonal pinching on a generic rational matrix; the two fibres computed as exact kernels and ranks, the big fibres of dimension eight containing the small fibres of dimension two, and the meet of the big fibres, the join of the small fibres and the joint-diagonal algebra coinciding as four-dimensional subspaces (Theorem \ref{c:singularity.finite-truncation-meet-big-fibres-equals}); words in the generators spanning the full sixteen-dimensional two-mode algebra by exact rank, and the isotony of the net assembly under the one-mode inclusion (Theorem \ref{c:singularity.net-assembly-net-it-everything}); the tensoriality lemma's finite model, a map of the three-point module commuting with every point multiplier exactly when it is diagonal, all nine matrix units exhaustively (Lemma \ref{c:singularity.function-linear-endomorphisms-are-bundle-maps}); $d\Gamma$ of a bundle map built from $\psi^{\dagger}a\psi$, equal to $a_xN_x + a_yN_y$, diagonal and conserving both occupations, while $d\Gamma$ of a hop fails to commute with $N_x$ and carries $|0_x 1_y\rangle$ into the $(1,0)$ sector (Proposition \ref{c:singularity.bundle-maps-quantise-without-hopping}); the controlled-$Z$ generator $\sigma_x\sigma_y$ conserving both occupations; the pairing generator $v^{*}_xv^{*}_y + v_yv_x$ commuting with $\Gamma$, off-diagonal in the joint number basis, and failing to commute with $N_x$ and with $N_x + N_y$, the witness that fixes Theorem \ref{c:singularity.class-relative-no-go-no-supplied}'s scope; and Corollary \ref{c:singularity.envelope-endpoint-why-boundary-mode-locality}'s rotation shown to be a pairing rotation, orthogonal, parity-preserving, with $\Delta N = 2$ on $|00\rangle$ and no hopping component (Remark \ref{c:note.singularity.declared-residual-size}). Facts \ref{c:singularity.imported-no-equivariant-value-wall} and \ref{c:sheaf.function-linearity-tame-constructions} are imports: the computational content of the former, the radical-gauge solve, the intertwiner identity $\Phi^{\mathsf T}h_u\Phi = h_{u+v}$, its cocycle, the derived contragredient and the translation action on candidates, is machine-verified in \texttt{verify\_T04.py} on the Schwarzschild boundary data at two masses, and is not re-verified here. No floating point is used anywhere; the script exits nonzero on any failure, and all two hundred and thirty-three checks pass.

\section{Position in the series}

\phantomsection\label{c:note.singularity.setting-position-series}
This paper is T5 of the series. This paper relies on: T1 \cite{ButlerT1} for the no-go theorem, the minimal completion, the zero-divisor taxonomy, the pinching and one-way theorems, and the six-step transport algorithm that the table accompanying \ref{c:singularity.imported-equivalence-einstein-equations} instantiates; T3 \cite{ButlerT3} for the tower and its CAR limit, the frozen clock at $d \equiv 6$ that stops the iteration at Step 6a, and the closure-completeness corollary that (L)'s exhaustiveness clause aims at; T4 \cite{ButlerT4} for the localisation theorems reproduced in the section ``The imported localisation theorems, verbatim'', which fix the weight of Proposition \ref{c:singularity.degeneration-exactly}, place the ladder pair on $E$, and discharge the boundary case of the no-repair import (Facts \ref{c:sheaf.schwarzschild-degeneration-boundary-tame}, \ref{c:singularity.imported-one-sided-completion-singularity} and \ref{c:singularity.imported-no-equivariant-value-wall}); and T6 \cite{ButlerT6} for the interior collar and its $t$-line holonomy, which the section ``The modular flow of the boundary modes inside, and the entangling half of (L$_E$)'' identifies with the modular flow of the boundary modes inside. The sheaf paper T4 in turn reads this paper's Schwarzschild data (Propositions \ref{c:singularity.unipotent-ray} to \ref{c:singularity.degeneration-exactly} and the section ``The Schwarzschild curvature and the surface gravity: the metric-picture cross-check'') as its worked composite; the two papers are cited in both directions and neither proof depends on the other's conclusion.

Papers that rely on this one: the interior paper T6 \cite{ButlerT6} imports the curvature eigenvalues and the Cartan conventions of the section ``The Schwarzschild curvature and the surface gravity: the metric-picture cross-check'', the rescaled-limit collapse of Proposition \ref{c:singularity.degeneration-exactly}, and the type-N reading of Proposition \ref{c:singularity.type-n-only-type-n}; the Hawking paper T7 \cite{ButlerT7} imports the single-mode thermality of the section ``Step 4d: compute in the completion; thermality, and the state that carries it'' (Theorems \ref{c:singularity.modular-descent} and \ref{c:singularity.descended-form-state-independent}, premises (P) and (R), the gluing occupancy of Theorem \ref{c:singularity.rate-discharge-hypothesis-iii}), the rate-rigidity Theorem \ref{c:singularity.rate-rigidity-temperature-not-algebra-intrinsic}, and the state class of Corollary \ref{c:singularity.chains-state-class-stability}; the entropy paper T8 \cite{ButlerT8} imports the counting axioms and eviction dynamics of the section ``Step 6a: iterate, or stop; the clock stops, and the boundary-mode capacity'' and the section ``Step 6b: iterate, or stop; two levels, the Page envelope, and eviction'', the geometric normalisation of the section ``The Schwarzschild curvature and the surface gravity: the metric-picture cross-check'', the atom's Euclidean period from Proposition \ref{c:singularity.euclidean-rotor-where-it-lies}, and the capacity comparison of Proposition \ref{c:singularity.area-law-test}; the twistor paper T9 \cite{ButlerT9} takes the chain end to end as the instance whose completion it develops, and receives the spectrum-condition restatement of Remark \ref{c:note.singularity.cost-r-dictionary}; the dark-energy paper T10 \cite{ButlerT10} imports the s-wave current algebra of the section ``Step 4e: compute in the completion; from fermions to photons'' under premise (Q) (Premise \ref{c:singularity.q-joint-quasi-freeness-s-wave}); the dark-matter paper T14 \cite{ButlerT14} imports the deficit theorem (Theorem \ref{c:singularity.deficit-theorem-cost-descent-exactly}) as an entropy statement, the exhaustiveness and purity clauses (L) and (U) of Premise \ref{c:singularity.rate-models-kinematic-inputs-l-u}, and the ascend-then-descend structure run tower by tower; the laboratory paper T16 \cite{ButlerT16} imports the modular period $\beta = 8\pi M$ with the descent as a one-way conditional expectation (Theorem \ref{c:singularity.modular-descent}, Remark \ref{c:note.singularity.steps-status-what-still-open}), the thermofield angle of Lemma \ref{c:singularity.schmidt-forcing-proved-outright}, and the verbatim imports of the section ``The imported localisation theorems, verbatim''; and the programme paper T17 \cite{ButlerT17} carries this paper's status row, its check count and its premises in the series ledger.

\section{What is established, what is premised, what is open}

\phantomsection\label{c:note.singularity.setting-what-established-what-premised-what}
The chain's concessions are argued where they arise, which scatters them; this section is the one place they stand together, so that the front matter has a single statement to agree with.

\emph{Established, unconditionally or from the imported Facts alone.} The gauge sector of the section ``Step 1: locate the failure'': the unipotent ray, the connection bivector in closed form, the exact degeneration, non-removability by the Kretschmann invariant. The shape decoupling $\mathcal{R} = -(M/r^{3})S$ and the matching's two equations (Theorem \ref{c:singularity.curvature-gauge-field-exact-shape-decoupling}, Corollary \ref{c:singularity.matching-steps-4b-4c-join-equation}), with the matching's residual transport step stated in the corollary. The parity-average identity and the outright uniqueness of the conditional expectation (Theorem \ref{c:singularity.descent-parity-average}). Single-mode modular descent (Theorem \ref{c:singularity.modular-descent}), the over-determination theorem (Theorem \ref{c:singularity.exact-equilibrium-over-determines-boundary-modes}), rate rigidity (Theorem \ref{c:singularity.rate-rigidity-temperature-not-algebra-intrinsic}), and the canonical purification (Theorem \ref{c:singularity.purity-output-partner-commutant}), all in the matrix factor and its finite tensor powers. The deficit theorem and its correlation identity (Theorem \ref{c:singularity.deficit-theorem-cost-descent-exactly}, Corollary \ref{c:singularity.deficit-carried-quanta-correlations}), as statements about the map $\EN$. The exact edge-operator current algebra (Theorem \ref{c:singularity.bosonisation-bridge-s-wave-sector}(1)). The causal-epoch integral (Lemma \ref{c:singularity.causal-epoch-computed-exactly}). From the section ``The metric-level sheaf, and the hopping exclusion'': that the sheaf assembly of the pointwise metric levels is a sheaf whose finite truncations are $\operatorname{range}(\EN)$ and the net assembly a net whose finite truncations are the full mode algebras (Theorems \ref{c:singularity.finite-truncation-meet-big-fibres-equals}, \ref{c:singularity.net-assembly-net-it-everything}, given the completed mode bundle of Fact \ref{c:singularity.imported-one-sided-completion-singularity}, hence under its standing hypothesis $w_1(\shL) = 0$), and that no generator in the supplied class hops (Theorem \ref{c:singularity.class-relative-no-go-no-supplied}, from the imported Fact \ref{c:sheaf.function-linearity-tame-constructions}); and the boundary no-value theorem, imported verbatim as Fact \ref{c:singularity.imported-no-equivariant-value-wall}. And the finite-mode model of the section ``The finite-mode model: the dissolution step by step'', every number exact. Two much-used results are placed deliberately between the lists: Corollary \ref{c:singularity.chains-state-class-stability}, a theorem from (P), (R) and, for its stability clause, (L$_E$), and Theorem \ref{c:singularity.eviction-window-thermal-product-radiation-entropy}, whose window identity with the envelope, the exact rising branch, the quantified falling branch and the endpoint is a theorem from its five named inputs. The dynamical input the section ``Step 6b: iterate, or stop; two levels, the Page envelope, and eviction'' uses is (L$_E$)'s \emph{full} diagonality: the half of (L$_E$) that the section ``The metric-level sheaf, and the hopping exclusion'' proves, no hopping within the supplied class, does not carry the envelope alone, since Corollary \ref{c:singularity.envelope-endpoint-why-boundary-mode-locality}'s necessity witness is a pairing rotation that hops nothing. Each result is exactly as strong as its premises, and both belong to what the premises yield rather than to the unconditional list. Two further results belong to this list. The arithmetic of Proposition \ref{c:singularity.modular-flow-boundary-modes-holonomy-universal}, from the geometry paper's collar alone: the $t$-line holonomy rate $M/r^2$ per Killing time, its Hubble-normalised limit $|p_1| = \tfrac13$ and the number $2\pi/|p_1| = 6\pi$, and the non-proportionality of the three limiting boosts; its identification of that holonomy with the modular flow of the boundary modes rests on Premise \ref{c:singularity.r-regularity-fixed-locus}, as clause (i) states, and whether any state carries $6\pi$ as a temperature is in the Open list below. And the stabiliser-tangent split of Proposition \ref{c:singularity.l-stabiliser-tangent}, an exact statement about $\mathfrak{so}(2\nu)$ on the Fock space, which places the pairing part of the premised half of (L$_E$) in the tangent of the vacuum orbit and its differential part in the stabiliser.

\emph{Premised, named where each is used.} (P) univalence (Premise \ref{c:singularity.p-univalence}) and (R) the KMS condition at the geometric period (Premise \ref{c:singularity.r-regularity-fixed-locus}), in Step 4d; the metric-level reading $\EN$ on many modes, licensed by Proposition \ref{c:singularity.p-reduced-metric-level-reading-many} with its definitional antecedent, the sheaf assembly of the section ``The metric-level sheaf, and the hopping exclusion'' (Step 5); (Q) joint quasi-freeness with the s-wave weights (Premise \ref{c:singularity.q-joint-quasi-freeness-s-wave}), the equilibrium reading of the Planck sector, in Step 4e, in its recorded tension with Steps 5 and 6b; (C1) to (C3) in Step 6a (Axioms \ref{c:singularity.c1-curvature-cutoff}, \ref{c:singularity.c2-cell-density}, \ref{c:singularity.c3-counted-epoch}); and, in Step 6b, (L) with its exhaustiveness postulate, (U), (C), (I) with its freshness clause (Premise \ref{c:singularity.rate-models-kinematic-inputs-l-u}), and (L$_E$), stated as Premise \ref{c:singularity.l-boundary-mode-locality} in Step 4d and used in Step 6b as full diagonality in the joint number basis with the never-return clause: its no-hopping half is a theorem within the supplied class (Theorem \ref{c:singularity.class-relative-no-go-no-supplied}), its exclusion of pairing terms and of unsupplied generators is a premise, and it is the premised part, not the proved part, that Step 6b's envelope turns on (Corollary \ref{c:singularity.envelope-endpoint-why-boundary-mode-locality}). The imported localisation theorems that carry Steps 2 through 4a are reproduced verbatim as Facts \ref{c:sheaf.schwarzschild-degeneration-boundary-tame} and \ref{c:singularity.imported-one-sided-completion-singularity}.

\emph{Open, and stated as such.} Whether the KMS state of the boundary modes on the completed pair extends to a KMS state on the algebra of $t$-line holonomies of the whole collar, which would carry the universal Hubble-normalised inverse temperature $6\pi$ of Proposition \ref{c:singularity.modular-flow-boundary-modes-holonomy-universal}(ii) (the obstruction, if any, is the running of the rate before its limit), so that the universal $6\pi$ is a temperature of something only if this is settled; and the $SO(3)$-invariant state on the sphere holonomies of the exceptional set with its two-point function, the one place a spectrum could enter (Remark \ref{c:note.singularity.where-spectrum-would-lie}). The derivation of (R), shown to lie outside the framework's reach (Proposition \ref{c:singularity.why-calibration-cannot-be-internal}); of (P$_{\mathrm{loc}}$) only the definitional antecedent remains, quantified at Proposition \ref{c:singularity.p-reduced-metric-level-reading-many} and restated by the section ``The metric-level sheaf, and the hopping exclusion'' as the choice of the sheaf assembly over the net assembly, which the appendix quantifies on both sides and does not decide. (L)'s exhaustiveness, the no-remnant hypothesis in its correct location. The identifications joining the three mode index sets (the section ``Step 4e: compute in the completion; from fermions to photons'''s declaration; Axiom (C1)). Saturation of the falling branch, available only through typicality or the one-mode zero-frequency idealisation. The separation the model enforces between the deficit, a size, and the message, a pattern the deficit cannot see. The type III$_1$ setting, and Kerr. The islands dissonance of Remark \ref{c:note.singularity.island-dictionary-one-dissonance}, which the body calls the sharpest open question the chain faces: whether the extremal-surface derivations and the eviction mechanism are two descriptions of one thing. The matching's residual transport step (Corollary \ref{c:singularity.matching-steps-4b-4c-join-equation}). Of the hopping no-go (Remark \ref{c:note.singularity.steps-status-what-still-open}(b)) the class-relative form is proved (the section ``The metric-level sheaf, and the hopping exclusion''), for a class of multiplication operators; what remains is whether the dynamics on $E$ carries a differential or a pairing term. The four-dimensional lift of the s-wave sector (Remark \ref{c:note.singularity.four-dimensional-lift-flagged}). And the horizon smoothness that eviction presupposes (Remark \ref{c:note.singularity.monogamy-firewalls-small-corrections-theorem}). The chain's conclusion survives at exactly this scope: the division transported rather than forbidden, the temperature from named premises, the two-level discrepancy an equation; what the premises yield is quantified above, and what they do not yield, this section records.

\section{The Schwarzschild curvature and the surface gravity: the metric-picture cross-check}

\phantomsection\label{c:note.singularity.setting-schwarzschild-curvature-surface-gravity-metric}
This appendix computes the curvature spectrum and the surface gravity in the metric picture, by the Cartan structure equations in the static chart, where the computation is shortest. It plays two roles. It supplies the mixed- and angular-plane eigenvalues that the abelian radial computation of Proposition \ref{c:singularity.solution-where-divergence-lies}(1) does not reach, completing the input to Theorem \ref{c:singularity.curvature-gauge-field-exact-shape-decoupling}; and it is the independent cross-check, in a second formalism, of the gauge-theoretic the section ``Step 1: locate the failure'': on the invertible region the two pictures agree by Fact \ref{c:singularity.imported-equivalence-einstein-equations}, so any disagreement anywhere would falsify one of them. The computations are classical differential geometry, included in full per the series' standard of proof; the verification script checks every derivative identity and eigenvalue exactly, with $f$ handled as a rational function.

\subsection*{The metric of the gauge field, and its static form}
\phantomsection\label{c:note.singularity.setting-metric-gauge-field-static-form}
The metric is $g(a,b) = \hb^{-1}(a) \cdot \hb^{-1}(b)$ with $\hb^{-1} = 1 + \beta P$ (Proposition \ref{c:singularity.unipotent-ray}). On the $(t,r)$ frame vectors:
\[
g(\et, \et) = (\et + \beta\er)^2 = 1 - \beta^2, \qquad
g(\et, \er) = (\et + \beta\er)\cdot\er = -\beta, \qquad
g(\er, \er) = \er^2 = -1 ,
\]
and the angular directions are untouched. This is the Painlev\'e-Gullstrand line element
\[
ds^2 = \Bigl(1 - \tfrac{2M}{r}\Bigr)\,dt^2 - 2\beta\,dt\,dr - dr^2 - r^2 d\Omega^2 , \qquad \beta^2 = \tfrac{2M}{r} .
\]
Substituting $dt = dt_s + (\beta/f)\,dr$ with $f := 1 - 2M/r = 1 - \beta^2$ eliminates the cross term exactly:
\[
f\,dt^2 - 2\beta\,dt\,dr - dr^2
\;=\; f\,dt_s^2 + 2\beta\,dt_s dr + \tfrac{\beta^2}{f}dr^2 - 2\beta\,dt_s dr - \tfrac{2\beta^2}{f}dr^2 - dr^2
\;=\; f\,dt_s^2 - \tfrac{\beta^2 + f}{f}\,dr^2
\;=\; f\,dt_s^2 - \tfrac{dr^2}{f} ,
\]
since $\beta^2 + f = 1$: the static Schwarzschild form, with the same areal radius $r$. Curvature is a pointwise invariant, so it may be computed in the static chart.

\subsection*{The curvature operator, by the Cartan structure equations}
\phantomsection\label{c:note.singularity.setting-curvature-operator-cartan-structure-equations}
Take the orthonormal coframe $\theta^{\hat 0} = \sqrt f\,dt_s$, $\theta^{\hat 1} = dr/\sqrt f$, $\theta^{\hat 2} = r\,d\vartheta$, $\theta^{\hat 3} = r\sin\vartheta\,d\varphi$. The first structure equations $d\theta^{\hat a} = -\,\omega^{\hat a}{}_{\hat b} \wedge \theta^{\hat b}$ are solved by
\[
\omega^{\hat 0}{}_{\hat 1} = \tfrac{f'}{2\sqrt f}\,\theta^{\hat 0}, \qquad
\omega^{\hat 2}{}_{\hat 1} = \tfrac{\sqrt f}{r}\,\theta^{\hat 2}, \qquad
\omega^{\hat 3}{}_{\hat 1} = \tfrac{\sqrt f}{r}\,\theta^{\hat 3}, \qquad
\omega^{\hat 3}{}_{\hat 2} = \tfrac{\cot\vartheta}{r}\,\theta^{\hat 3},
\]
all other independent components zero; each is verified by substitution (for instance $d\theta^{\hat 0} = \tfrac{f'}{2\sqrt f}\,dr \wedge dt_s = \tfrac{f'}{2\sqrt f}\,\theta^{\hat 1}\wedge\theta^{\hat 0}$, matched by $-\,\omega^{\hat 0}{}_{\hat 1}\wedge\theta^{\hat 1}$). The second structure equations $R^{\hat a}{}_{\hat b} = d\omega^{\hat a}{}_{\hat b} + \omega^{\hat a}{}_{\hat c}\wedge\omega^{\hat c}{}_{\hat b}$ then give, writing each curvature two-form on its own plane:
\[
R^{\hat 0}{}_{\hat 1} = -\tfrac{f''}{2}\,\theta^{\hat 0}\wedge\theta^{\hat 1}, \qquad
R^{\hat 0}{}_{\hat 2} = -\tfrac{f'}{2r}\,\theta^{\hat 0}\wedge\theta^{\hat 2}, \qquad
R^{\hat 1}{}_{\hat 2} = -\tfrac{f'}{2r}\,\theta^{\hat 1}\wedge\theta^{\hat 2},
\]
with $R^{\hat 0}{}_{\hat 3}, R^{\hat 1}{}_{\hat 3}$ mirroring the $\hat 2$ entries, and
\[
R^{\hat 2}{}_{\hat 3} = \tfrac{1 - f}{r^2}\,\theta^{\hat 2}\wedge\theta^{\hat 3} .
\]
As a sample of the derivations: $R^{\hat 0}{}_{\hat 1} = d\omega^{\hat 0}{}_{\hat 1} = \bigl(\sqrt f\,(\tfrac{f'}{2\sqrt f})' + \tfrac{f'^2}{4f}\bigr)\theta^{\hat 1}\wedge\theta^{\hat 0} = \tfrac{f''}{2}\,\theta^{\hat 1}\wedge\theta^{\hat 0}$, the $f'^2$ terms cancelling; $R^{\hat 1}{}_{\hat 2} = d(-\sqrt f\,d\vartheta) = -\tfrac{f'}{2\sqrt f}dr\wedge d\vartheta = -\tfrac{f'}{2r}\theta^{\hat 1}\wedge\theta^{\hat 2}$; and $R^{\hat 2}{}_{\hat 3} = d(-\cos\vartheta\,d\varphi) + \omega^{\hat 2}{}_{\hat 1}\wedge\omega^{\hat 1}{}_{\hat 3} = \tfrac{1}{r^2}\theta^{\hat 2}\wedge\theta^{\hat 3} - \tfrac{f}{r^2}\theta^{\hat 2}\wedge\theta^{\hat 3}$. The curvature operator on the six-dimensional bivector space is therefore diagonal in the coordinate planes, with eigenvalues proportional to
\[
-\tfrac{f''}{2} \ \ (\text{radial boost plane } \hat 0\hat 1), \qquad
-\tfrac{f'}{2r} \ \ (\text{the four mixed planes}), \qquad
\tfrac{1-f}{r^2} \ \ (\text{angular plane } \hat 2\hat 3).
\]
For $f = 1 - 2M/r$: $f' = 2M/r^2$, $f'' = -4M/r^3$, so the three values are
\[
\tfrac{2M}{r^3}, \qquad -\tfrac{M}{r^3}, \qquad \tfrac{2M}{r^3} :
\]
the doubled-magnitude eigenvalue $2M/r^3$ is attained exactly on the radial plane $\hat 0\hat 1$ and its Hodge dual $\hat 2\hat 3$, at every radius, with the four mixed planes at half the magnitude. The extreme eigenvalue therefore distinguishes the pair $\{\sigma_r, I\sigma_r\}$ invariantly (the dual written $I\sigma_r$ here as in the body), in agreement with the gauge-sector value of Proposition \ref{c:singularity.solution-where-divergence-lies}(1) on the radial plane, and the shape these eigenvalues assemble into is Theorem \ref{c:singularity.curvature-gauge-field-exact-shape-decoupling}'s $S = 1 - 3\Pi$, with $\mathcal{R} = -(M/r^{3})S$ in this convention exactly as in the gauge-sector computation (only the ratio $2 : 1$ and the identity of the distinguished plane are used in the chain).

\begin{fact}[Imported: surface gravity as redshifted acceleration]\label{c:singularity.imported-surface-gravity-redshifted-acceleration}
\changedmark{gray}{reworded}
Let $(\hb, \Omega)$ be a static gravitational field on flat spacetime (Definition \emph{Gauge Theory Gravity: the gravitational field on flat spacetime}) with a static Killing horizon: a Killing field $K$, a physical vector field whose covariant derivative $A \mapsto A \cdot D K$ is antisymmetric in the flat product, and a locus of the background on which $Q(K) = 0$ and $K$ is normal to it. The surface gravity defined on the horizon locus by $(K \cdot D)K = \kappa\,K$ equals the limit, at the horizon locus, of the proper acceleration of static observers (velocity $v = K/\sqrt{Q(K)}$, acceleration $(v \cdot D)v$) multiplied by the redshift factor $\sqrt{Q(K)}$; $Q(K)$ is the derived-metric norm $g(\xi,\xi)$ of the Killing vector $\xi$ over which $K$ sits (Fact \emph{Imported: the dictionary between Gauge Theory Gravity and curved-spacetime language}, clause (vi)). \emph{Provenance:} R.~M.~Wald, General Relativity, University of Chicago Press (1984), \S 12.5, read in Gauge Theory Gravity through the dictionary.
\emph{Source:} R. M. Wald, General Relativity, University of Chicago Press (1984), §12.5.
\end{fact}

\begin{proposition}[$\kappa = 1/4M$, computed]\label{c:singularity.computed}
\changedmark{gray}{reworded}
For the Schwarzschild gravitational field of the section ``The metric of the gauge field, and its static form'' (the stored field $\hb = 1 - \beta P$ on flat spacetime, Definition \emph{Gauge Theory Gravity: the gravitational field on flat spacetime}), the static Killing field $K = \hb^{-1}(e_t)$ has $Q(K) = f = 1 - 2M/r$, and the static observers, of velocity $v = K/\sqrt f$, have redshifted proper acceleration $\sqrt f\,\sqrt{-Q((v \cdot D)v)} = f'(r)/2$ at every $r > 2M$, so the surface gravity of the horizon locus $r = 2M$ is
\[
\kappa \;=\; \lim_{r \to 2M} \frac{f'(r)}{2} \;=\; \frac{1}{2}\cdot\frac{2M}{(2M)^2} \;=\; \frac{1}{4M} .
\]
\end{proposition}
\begin{proof}
The components below are those of the covariant derivative $D$ in the static frame, read through Fact \emph{Imported: the dictionary between Gauge Theory Gravity and curved-spacetime language} (clauses (ii) and (vi)). The static observer has four-velocity $u = \xi/\sqrt f$ with $\xi = \partial_{t_s}$. Its acceleration $A^\mu = u^\nu \nabla_\nu u^\mu$ has the single component $A^r = \Gamma^r_{tt}\,(u^t)^2$, with $\Gamma^r_{tt} = -\tfrac12 g^{rr}\partial_r g_{tt} = -\tfrac12(-f)(f') = \tfrac12 f f'$ and $(u^t)^2 = 1/f$, so $A^r = f'/2$. Its proper magnitude is $\sqrt{-g_{rr}}\,|A^r| = f'/(2\sqrt f)$, and the redshifted magnitude is $\sqrt f \cdot f'/(2\sqrt f) = f'/2$, finite and radius-independent in its limit. By Fact \ref{c:singularity.imported-surface-gravity-redshifted-acceleration}, $\kappa = f'(2M)/2 = 1/4M$. The conversion quoted in Proposition \ref{c:singularity.plane-rate} follows: one period $2\pi$ of rapidity is $2\pi/\kappa = 8\pi M$ of Schwarzschild time.

\end{proof}

\section{The imported localisation theorems, verbatim}

\phantomsection\label{c:note.singularity.setting-imported-localisation-theorems-verbatim}
The two theorems of T4 \cite{ButlerT4} that carry the localisation weight at Steps 2, 3, 4a, 6a and 6b, the no-value theorem at the wall used at Step 3, and the function-linearity principle used by the section ``The metric-level sheaf, and the hopping exclusion'', are reproduced here in full, hypotheses included, so that this paper can be refereed without the sheaf paper on the desk. The notation is that of T4 \cite{ButlerT4}, aligned with this paper's where the two differ: $\shW$ is the weighted coframe module on the real oriented blow-up $\Mt$ of the singular worldline, $E \cong \R \times S^2$ its boundary, $m = \sqrt{2M}$, $\rho$ the weighted radial coframe element, $(\nO, \nInf)$ the adjoined hyperbolic plane, and T4's $dt$ and $dt_{\mathrm{S}}$ are written $dt_{\mathrm{PG}}$ and $dt_s$ here. The proofs of Facts \ref{c:sheaf.schwarzschild-degeneration-boundary-tame} and \ref{c:singularity.imported-one-sided-completion-singularity}, and the machine verification of every computational identity in them (seventy checks, exact rational functions in $u = \sqrt r$), are Section 8 of T4 \cite{ButlerT4}; Fact \ref{c:singularity.imported-no-equivariant-value-wall} is proved in its Section 10 (thirty-four checks on the same boundary data), and Fact \ref{c:sheaf.function-linearity-tame-constructions} in its Section 9. Two further alignments concern Fact \ref{c:singularity.imported-no-equivariant-value-wall}: that paper's wall $\Sigma$ is the boundary $E$ here, and its coframe module $T^*M$ at the wall point is the weighted module $\shW$ at a point of $E$, on whose frame Gram, the matrix of Fact \ref{c:sheaf.schwarzschild-degeneration-boundary-tame}, the machine checks in fact run. The coupling constant $\nu$ of Fact \ref{c:singularity.imported-one-sided-completion-singularity} is T4's and is unrelated to the cell density $\nu$ of Axiom (C2).

\begin{assembled}[Imported: one-sided completion at the singularity]\label{c:singularity.imported-one-sided-completion-singularity}
Adjoin the trivial hyperbolic plane $(\nO, \nInf)$ to $\shW$ and couple it to the dying direction by the coupling construction of \cite[Def.~4.1]{ButlerT4}, whose standing hypothesis $w_1(\shL) = 0$ (the dying line is orientable) is met here, the radical line being trivialised globally by $k = d\Phi$ (clause (5) of Fact \ref{c:sheaf.schwarzschild-degeneration-boundary-tame}), taking the normalised transversal $u = \partial_t/m$ (so that $k(u) = 1$, as that construction requires), coupling constants satisfying $\mu\nu < 0$, and bump supported inside the horizon. Then the ambient form $h$ is nondegenerate of constant signature $(2,4)$ at every point of $\Mt$, the boundary $E$ included; the subalgebra generated by $\shW$ is canonically the degeneration sheaf; and the ladder pair
\[
k^* = \frac{\nO}{2\nu}, \qquad k k^* + k^* k = 1, \qquad k^2 = (k^*)^2 = 0
\]
exists exactly over $E$. The ascent $\Cl(1,3) \to \Cl(2,4)$, i.e.\ $\Mat_2(\HH) \to \Mat_4(\HH)$, is thereby performed at the Schwarzschild singularity, localised on $E \cong \R \times S^2$. Moreover no shadow wall forms anywhere in the manifold.
\emph{Provenance:} \cite[Thm.~8.12]{ButlerT4}, verbatim up to the recorded notational alignment, with one clause added here: the parenthetical $w_1(\shL) = 0$ and its discharge by $k = d\Phi$ are not in the source's statement, which invokes \cite[Def.~4.1]{ButlerT4} with that hypothesis standing; the clause records that the hypothesis holds on the Schwarzschild data, by clause (5) of Fact \ref{c:sheaf.schwarzschild-degeneration-boundary-tame}, so that the import is applied within its hypotheses.
\end{assembled}
\begin{proof}
Adjoin the trivial hyperbolic plane $(\nO, \nInf)$ to $\shW$ and couple it to the dying direction by the coupling construction of \cite[Def.~4.1]{ButlerT4}, whose standing hypothesis $w_1(\shL) = 0$ (the dying line is orientable) is met here, the radical line being trivialised globally by $k = d\Phi$ (clause (5) of Fact \ref{c:sheaf.schwarzschild-degeneration-boundary-tame}), taking the normalised transversal $u = \partial_t/m$ (so that $k(u) = 1$, as that construction requires), coupling constants satisfying $\mu\nu < 0$, and bump supported inside the horizon. Then the ambient form $h$ is nondegenerate of constant signature $(2,4)$ at every point of $\Mt$, the boundary $E$ included; the subalgebra generated by $\shW$ is canonically the degeneration sheaf; and the ladder pair
\[
k^* = \frac{\nO}{2\nu}, \qquad k k^* + k^* k = 1, \qquad k^2 = (k^*)^2 = 0
\]
exists exactly over $E$. The ascent $\Cl(1,3) \to \Cl(2,4)$, i.e.\ $\Mat_2(\HH) \to \Mat_4(\HH)$, is thereby performed at the Schwarzschild singularity, localised on $E \cong \R \times S^2$. Moreover no shadow wall forms anywhere in the manifold.
\emph{Provenance:} \cite[Thm.~8.12]{ButlerT4} (\ref{c:sheaf.one-sided-completion-singularity}), verbatim up to the recorded notational alignment, with one clause added here: the parenthetical $w_1(\shL) = 0$ and its discharge by $k = d\Phi$ are not in the source's statement, which invokes \cite[Def.~4.1]{ButlerT4} with that hypothesis standing; the clause records that the hypothesis holds on the Schwarzschild data, by clause (5) of Fact \ref{c:sheaf.schwarzschild-degeneration-boundary-tame}, so that the import is applied within its hypotheses.

\end{proof}

\begin{assembled}[Imported: no equivariant value at the wall]\label{c:singularity.imported-no-equivariant-value-wall}
At a point of $E$ with the collar constants $\mu, \nu$ and the bump held fixed, call a vector $u'$ with $k(u') = 1$ a \emph{candidate value} for the dying direction $k$; the candidates form an affine space over $\ker k$. For every $\lambda \in \ker k$ the composite $\theta_\lambda := \hat g_\lambda \circ \Phi_\lambda^{-1}$, with $\hat g_\lambda := g_\lambda \oplus \id$ on $T = \shW \oplus \operatorname{span}\{\nO, \nInf\}$, is an automorphism of the completed structure $(T, h_u, k)$: an $h_u$-isometry restricting to $g_\lambda$ on $\shW$ and fixing $k$, hence an automorphism of the completed Clifford algebra preserving its grading and its degeneration subalgebra. Its contragredient action on vectors, $(\theta_\lambda|_{\shW})^{-\mathsf{T}} = \id - \lambda \otimes k$, translates every candidate by $-\lambda$. A natural assignment of a candidate value, one per completed structure and commuting with automorphisms, would therefore be a fixed point of a free translation action of $(\ker k, +)$: none exists. Nor is a ladder partner for $k$ any more canonical than a value: for $\lambda \neq 0$, $h_u(k, x) = 0$ for every $x \in \operatorname{Fix}(\theta_\lambda)$, whereas a ladder partner has $h_u(k, x) = \tfrac12$; and the construction attaches to each candidate $u' = u + v$ the partner $\Phi_v(k^*) = k^* + \tfrac12\tilde v$ of the transversal-independence theorem, injectively in $v$. An analogue of the no-repair theorem of \cite[Thm.~4.2]{ButlerT1} thus holds at the wall itself, where the base form carries no hyperbolic partner for $k$ and that theorem's hypotheses fail.
\emph{Provenance:} \cite[Thm.~10.4]{ButlerT4}, verbatim up to the recorded notational alignment ($E$ for that paper's $\Sigma$, $\shW$ for its $T^*M$). The objects in the statement are that paper's: the \emph{completed structure} $(T, h_u, k)$ is the ambient module $T$ with the ambient form $h_u$ of the coupling construction built on the transversal $u$ (Fact \ref{c:singularity.imported-one-sided-completion-singularity}) and the dying direction $k$, carrying the $\mathbb{Z}_2$-grading of its Clifford algebra and the subalgebra generated by $\shW$; the named frame $(\nO, \nInf)$ and the partner $k^* = \nO/2\nu$ are choices inside it, not part of it. Its supporting results are \cite[Lems 10.1, 10.2, Thm.~10.3]{ButlerT4}: the base unipotents $g_\lambda(\alpha) = \alpha + \alpha(\lambda)k$, $G$-isometries fixing $k$; for $v \in \ker k$ the representative $\tilde v \in \shW$ with $G(\alpha, \tilde v) = \alpha(v)$, normalised by $G(\tilde v, \tilde v) + 2\tilde v(u) = 0$; and the unipotent intertwiner $\Phi_v$, the identity on $\shW$ with $\Phi_v(\nO) = \nO + \nu\tilde v$, $\Phi_v(\nInf) = \nInf + \mu\tilde v$, an isometry $(T, h_{u+v}) \to (T, h_u)$ satisfying the cocycle $\Phi^u_v \circ \Phi^{u+v}_w = \Phi^u_{v+w}$, so that the completions built on all transversals are canonically identified, at the wall, by a coherent system of isomorphisms preserving the form, the grading and the degeneration subalgebra but neither the frame nor the ladder pair. Off the wall the two ambient forms need not be isometric (the shadow wall moves with the transversal); nothing off the wall is used. All of it is machine-verified there on the Schwarzschild boundary data at two masses.
\end{assembled}
\begin{proof}
At a point of $E$ with the collar constants $\mu, \nu$ and the bump held fixed, call a vector $u'$ with $k(u') = 1$ a \emph{candidate value} for the dying direction $k$; the candidates form an affine space over $\ker k$. For every $\lambda \in \ker k$ the composite $\theta_\lambda := \hat g_\lambda \circ \Phi_\lambda^{-1}$, with $\hat g_\lambda := g_\lambda \oplus \id$ on $T = \shW \oplus \operatorname{span}\{\nO, \nInf\}$, is an automorphism of the completed structure $(T, h_u, k)$: an $h_u$-isometry restricting to $g_\lambda$ on $\shW$ and fixing $k$, hence an automorphism of the completed Clifford algebra preserving its grading and its degeneration subalgebra. Its contragredient action on vectors, $(\theta_\lambda|_{\shW})^{-\mathsf{T}} = \id - \lambda \otimes k$, translates every candidate by $-\lambda$. A natural assignment of a candidate value, one per completed structure and commuting with automorphisms, would therefore be a fixed point of a free translation action of $(\ker k, +)$: none exists. Nor is a ladder partner for $k$ any more canonical than a value: for $\lambda \neq 0$, $h_u(k, x) = 0$ for every $x \in \operatorname{Fix}(\theta_\lambda)$, whereas a ladder partner has $h_u(k, x) = \tfrac12$; and the construction attaches to each candidate $u' = u + v$ the partner $\Phi_v(k^*) = k^* + \tfrac12\tilde v$ of the transversal-independence theorem, injectively in $v$. An analogue of the no-repair theorem of \cite[Thm.~4.2]{ButlerT1} thus holds at the wall itself, where the base form carries no hyperbolic partner for $k$ and that theorem's hypotheses fail.
\emph{Provenance:} \cite[Thm.~10.4]{ButlerT4} (\ref{c:sheaf.no-equivariant-value-wall}), verbatim up to the recorded notational alignment ($E$ for that paper's $\Sigma$, $\shW$ for its $T^*M$). The objects in the statement are that paper's: the \emph{completed structure} $(T, h_u, k)$ is the ambient module $T$ with the ambient form $h_u$ of the coupling construction built on the transversal $u$ (Fact \ref{c:singularity.imported-one-sided-completion-singularity}) and the dying direction $k$, carrying the $\mathbb{Z}_2$-grading of its Clifford algebra and the subalgebra generated by $\shW$; the named frame $(\nO, \nInf)$ and the partner $k^* = \nO/2\nu$ are choices inside it, not part of it. Its supporting results are \cite[Lems 10.1, 10.2, Thm.~10.3]{ButlerT4} (\ref{c:sheaf.base-unipotents-wall}, \ref{c:sheaf.transversal-independence-wall}, \ref{c:sheaf.solve-closed-radical-gauge}): the base unipotents $g_\lambda(\alpha) = \alpha + \alpha(\lambda)k$, $G$-isometries fixing $k$; for $v \in \ker k$ the representative $\tilde v \in \shW$ with $G(\alpha, \tilde v) = \alpha(v)$, normalised by $G(\tilde v, \tilde v) + 2\tilde v(u) = 0$; and the unipotent intertwiner $\Phi_v$, the identity on $\shW$ with $\Phi_v(\nO) = \nO + \nu\tilde v$, $\Phi_v(\nInf) = \nInf + \mu\tilde v$, an isometry $(T, h_{u+v}) \to (T, h_u)$ satisfying the cocycle $\Phi^u_v \circ \Phi^{u+v}_w = \Phi^u_{v+w}$, so that the completions built on all transversals are canonically identified, at the wall, by a coherent system of isomorphisms preserving the form, the grading and the degeneration subalgebra but neither the frame nor the ladder pair. Off the wall the two ambient forms need not be isometric (the shadow wall moves with the transversal); nothing off the wall is used. All of it is machine-verified there on the Schwarzschild boundary data at two masses.

\end{proof}

\section{The metric-level sheaf, and the hopping exclusion}

\phantomsection\label{c:note.singularity.setting-metric-level-sheaf-hopping-exclusion}
Two further results are proved here rather than imported; both are used in the body (Proposition \ref{c:singularity.p-reduced-metric-level-reading-many}'s antecedent paragraph and Premise \ref{c:singularity.l-boundary-mode-locality}) and both were developed alongside Fact \ref{c:singularity.imported-no-equivariant-value-wall}. Throughout, the collar constants and bump are fixed and the completed mode bundle $\Aa$ of Fact \ref{c:singularity.imported-one-sided-completion-singularity} is in place: its existence as an untwisted bundle over $E$ is exactly the standing hypothesis $w_1(\shL) = 0$, met here by $k = d\Phi$. Given $\Aa$, the ladder $(k, k^{*})$ and the parity $\sigma := k^{*}k - kk^{*}$ are global sections over $E$, and $\Aa_x \cong \Mat_2$ denotes the mode factor at $x \in E$ with parity automorphism $\alpha_x := \operatorname{Ad}(\sigma(x))$. Since $\sigma$ is quadratic in the ladder, $\sigma$, $\alpha_x$, the idempotent $\tfrac12(\id + \alpha_x)$ and everything built from them are insensitive to the sign of the ladder, which is what $w_1$ measures; nothing in this appendix depends on the standing hypothesis beyond the existence of $\Aa$ itself.

\subsection*{The metric-level sheaf: two assemblies of the pointwise levels}
\phantomsection\label{c:note.singularity.setting-metric-level-sheaf-two-assemblies}
The body defines the metric level pointwise (Proposition \ref{c:singularity.descent-what-it-erases}, Theorem \ref{c:singularity.descent-parity-average}) and reads a region of the boundary mode space through $\EN$, the composite of the pointwise descents (Proposition \ref{c:singularity.p-reduced-metric-level-reading-many}). This subsection separates what that reading assumes. ``The pointwise metric level at $x$'' names two algebras, and the choice between them, made without comment in the body, is the whole of Proposition \ref{c:singularity.p-reduced-metric-level-reading-many}'s antecedent.

\begin{definition}[The two fibres, and the two assemblies]\label{c:obj.singularity.two-fibres-two-assemblies}
For $x \in E$ the \emph{small fibre} is the even part of the mode factor, $\Aa_x^{+} := \{a \in \Aa_x : \alpha_x(a) = a\} = \operatorname{span}\{1, \sigma(x)\}$, and the \emph{even subbundle} $\Aa^{+} \subseteq \Aa$ has these fibres. On a finite truncation $F = \{x_1, \dots, x_N\} \subseteq E$ with its $F$-mode algebra, the \emph{big fibre} at $x_j$ is $B_j := \operatorname{Fix}(\operatorname{Ad}\sigma(x_j))$, the elements of the whole $F$-mode algebra even in mode $j$, of dimension $4^N/2$; the small fibre lies inside it as the elements of mode $j$'s own factor. The \emph{sheaf assembly} is $\mathfrak{M}(U) := \Gamma(U, \Aa^{+})$ for open $U \subseteq E$, families $s$ with $s(x) \in \Aa_x^{+}$ at every $x \in U$, with the restriction maps of sections. The \emph{net assembly} is $J(F) :=$ the algebra generated in the $F$-mode algebra by $\bigcup_j B_j$, for finite $F$, with the inclusion maps $J(F) \hookrightarrow J(F')$ for $F \subseteq F'$.
\end{definition}

\begin{proposition}[The sheaf assembly is a sheaf]\label{c:singularity.sheaf-assembly-sheaf}
$\mathfrak{M}$ is a sheaf of algebras on $E$, with stalk at $x$ the germs valued in $\Aa_x^{+}$. It exists whenever the completed mode bundle $\Aa$ does, that is under the standing hypothesis $w_1(\shL) = 0$, and needs nothing further: the fibrewise projection $\tfrac12(\id + \alpha_x)$ is a smooth idempotent bundle map, so $\Aa^{+}$ is a locally trivial subbundle, and it is defined by $\sigma$, which does not see the sign of the ladder. For $w_1(\shL) \neq 0$ the mode bundle itself, and with it $\mathfrak{M}$, exists only in the $\shL$-twisted variant.
\end{proposition}
\begin{proof}
Sections of a vector bundle form a sheaf: locality and gluing are inherited from the sections of $\Aa$, membership being a fibrewise condition, and restriction is restriction of sections. Multiplication is fibrewise and preserves evenness. The subbundle statement is the idempotent, smooth because $\sigma$ is a smooth section; the stalk is then as stated. The $w_1$ clause is Fact \ref{c:singularity.imported-one-sided-completion-singularity}'s standing hypothesis, which is what puts $\Aa$ in place; that the rest is sign-insensitive is the quadratic dependence of $\sigma$ on the ladder, $\sigma \mapsto \sigma$ under $(k, k^{*}) \mapsto (-k, -k^{*})$.

\end{proof}

\begin{theorem}[On a finite truncation, meet of the big fibres equals join of the small fibres equals $\operatorname{range}(E_F)$]\label{c:singularity.finite-truncation-meet-big-fibres-equals}
For finite $F = \{x_1, \dots, x_N\} \subseteq E$,
\[
\bigcap_{j} B_j \;=\; \Bigl\langle \bigcup_j \Aa_{x_j}^{+} \Bigr\rangle \;=\; \operatorname{range}(E_F),
\]
the joint-number-diagonal algebra, of dimension $2^N$: the meet of the big fibres and the join of the small fibres coincide, and both are the algebra of compatible evaluations of $\mathfrak{M}$ at the points of $F$. Here $\langle\,\cdot\,\rangle$ denotes the generated subalgebra.
\end{theorem}
\begin{proof}
The meet: an element fixed by every $\operatorname{Ad}\sigma(x_j)$ is diagonal in the joint number basis, and conversely (Proposition \ref{c:singularity.p-reduced-metric-level-reading-many}), and $\operatorname{range}(E_F)$ is that algebra, of dimension $2^N$. The join: the $\sigma(x_j)$ commute, are single-site in the Jordan-Wigner representation (the strings cancel in $v^{*}v - vv^{*}$), and generate the $2^N$ independent monomials $\prod_{j \in S}\sigma(x_j)$, $S \subseteq F$, each joint-diagonal; so the join is a subalgebra of the meet of the same dimension. Compatible evaluations of $\mathfrak{M}$ at the points of $F$ are exactly tuples in $\prod_j \Aa_{x_j}^{+}$ (every tuple occurs, by bump functions), and the algebra they generate is the join.

\end{proof}

\begin{theorem}[The net assembly is a net, and it is everything]\label{c:singularity.net-assembly-net-it-everything}
For $F$ with at least two points, $J(F)$ is the whole $F$-mode algebra, so the two-level discrepancy computed against it vanishes identically. The assignment $F \mapsto J(F)$ is covariant (isotone) under inclusion: it is a net of algebras, the structure of an algebra of regions, and not a presheaf; it has no restriction maps, and the question whether it is a sheaf does not arise. Its ``stalk'' at a point, the algebra generated by the big fibre at $x$ alone, contains elements odd at every other point of $F$.
\end{theorem}
\begin{proof}
In the Jordan-Wigner representation $\operatorname{Ad}(\sigma_x)$ fixes $v_y$ for $y \neq x$ and negates it at $y = x$, $\sigma_x$ being single-site; hence every generator of every mode lies in some other point's big fibre, $J(F)$ contains them all, and words in them span the full $F$-mode algebra (verified by exact rank on two modes, the section ``Machine verification''). Isotony is the inclusion of the $F$-mode algebra in the $F'$-mode algebra, which carries big fibres to big fibres. The last clause is $v_y \in B_x$ for $x \neq y$ with $\operatorname{Ad}(\sigma_y)(v_y) = -v_y$.

\end{proof}

\begin{remark}[What the antecedent is, stated plainly]\label{c:note.singularity.what-antecedent-stated-plainly}
Theorems \ref{c:singularity.finite-truncation-meet-big-fibres-equals} and \ref{c:singularity.net-assembly-net-it-everything} quantify the two readings of ``the metric level over a region'' and do not decide between them. Both are determined by readings near the region's points; they differ in how the pointwise levels are assembled. The contravariant assembly (the sheaf, a limit: an element must survive every point's descent) gives $\operatorname{range}(\EN)$, the joint-diagonal algebra, and with it the nonzero two-level discrepancy of Steps 5 and 6b. The covariant assembly (the net, a colimit: the algebra generated by what each point's descent retains of the whole mode algebra) gives the whole mode algebra and a discrepancy of zero; it is the assembly an algebraic quantum field theory would adopt for its algebras of regions. The chain adopts the former, and that adoption is Proposition \ref{c:singularity.p-reduced-metric-level-reading-many}'s antecedent: the sheaf assembly with the small fibres as stalks is the metric level if and only if the metric level over a region is the meet, so the antecedent is restated by this appendix, not reduced. What the appendix adds is the structural content of the choice: on the sheaf side the reading is a genuine sheaf whose finite truncations are $\operatorname{range}(E_F)$, on the net side it is a net whose finite truncations are the full mode algebras, and the discrepancy of Theorem \ref{c:singularity.deficit-theorem-cost-descent-exactly} is precisely the cost of reading a region as a limit rather than a colimit. The choice is made in Definition \ref{c:obj.singularity.two-fibres-two-assemblies}, before any sheaf theory runs, and it is a choice about which algebra ``the pointwise metric level'' names: the even part of the mode factor, or the even-in-that-mode part of the boundary mode space.
\end{remark}

\begin{lemma}[Function-linear endomorphisms are bundle maps]\label{c:singularity.function-linear-endomorphisms-are-bundle-maps}
Let $\mathcal{F}$ be the sheaf of sections of a vector bundle over $E$ and $D$ a $C^\infty(E)$-linear endomorphism of $\mathcal{F}(E)$. Then $D$ is tensorial: $(Ds)(x)$ depends only on $s(x)$, so $D$ is a bundle map $\{a_x\}$; for a line bundle, multiplication by a function.
\end{lemma}
\begin{proof}
If $s(x) = 0$, take a local frame $(t_i)$ near $x$ and a bump $\chi$ equal to $1$ near $x$ and supported in the frame's domain; then $\chi^2 s = \sum_i f_i\, \chi t_i$ with $f_i := \chi \cdot (\text{coefficient of } s \text{ on } t_i)$ global and $f_i(x) = 0$, while $(1 - \chi^2)s$ is $(1 - \chi^2)$ times $s$. Linearity gives $(Ds)(x) = \sum_i f_i(x)\,(D\chi t_i)(x) + (1 - \chi^2)(x)\,(Ds)(x) = 0$. So $(Ds)(x)$ is determined by $s(x)$, and $a_x$ is the induced fibre map.

\end{proof}

\begin{proposition}[Bundle maps quantise without hopping]\label{c:singularity.bundle-maps-quantise-without-hopping}
On a finite truncation the second quantisation of a bundle map $\{a_x\}$ of the mode line bundle is the sum of single-point bilinears $d\Gamma(a) = \sum_j \psi^{\dagger}_{x_j} a_{x_j} \psi_{x_j} = \sum_j a_{x_j} N_{x_j}$, a multiplication operator in the joint number basis, commuting with every occupation number. A one-particle map with a cross component $a_{xy} \neq 0$, $x \neq y$, quantises instead to a generator containing $\psi^{\dagger}_x\psi_y$, which carries the sector $(N_x, N_y) = (0,1)$ onto $(1,0)$.
\end{proposition}
\begin{proof}
The block form is the definition of second quantisation applied to a bundle map, each summand a polynomial in one point's ladder, and for a line bundle $\psi^{\dagger}_x a_x \psi_x = a_x N_x$; the cross-component computation is $\psi^{\dagger}_x\psi_y|0_x 1_y\rangle = |1_x 0_y\rangle$. Both are verified as exact matrix identities, together with the controlled-$Z$ generator $\sigma_x\sigma_y$ conserving both occupations: cross-point diagonal products are not hopping and were never excluded.

\end{proof}

\begin{theorem}[The class-relative no-go: no supplied generator hops]\label{c:singularity.class-relative-no-go-no-supplied}
Assume Fact \ref{c:sheaf.function-linearity-tame-constructions}. Let the \emph{supplied class} consist of the second quantisations $d\Gamma(a)$ of the $C^\infty(E)$-linear endomorphisms $a$ of the mode line bundle (the line bundle over $E$ whose fibre at $x$ is the one-particle line of the mode there), together with the products of the parity sections $\sigma(x)$. Then no generator in the supplied class moves occupation from the ladder at one point of $E$ to the ladder at another: every one of them commutes with every $N_x$. Indeed the supplied class is \emph{ultralocal}: each of its generators is a sum of single-point terms, a multiplication operator in the joint number basis, and propagates nothing between points, so the theorem excludes hopping by generators that could not move anything in any case. Within the supplied class, the no-hopping half of Premise \ref{c:singularity.l-boundary-mode-locality} is a theorem. The class does not contain, and the theorem does not exclude, pairing generators of the form $v^{*}_x v^{*}_y + v_y v_x$: these are $C^\infty$-linear in their coefficient functions, commute with the total parity $\Gamma$, are not joint-diagonal, and do not commute with $N_x$; their exclusion remains part of Premise \ref{c:singularity.l-boundary-mode-locality}.
\end{theorem}
\begin{proof}
By Fact \ref{c:sheaf.function-linearity-tame-constructions} and Lemma \ref{c:singularity.function-linear-endomorphisms-are-bundle-maps}, each supplied $a$ is a bundle map, and by Proposition \ref{c:singularity.bundle-maps-quantise-without-hopping} its $d\Gamma$ is a multiplication operator; products of the $\sigma(x)$ are joint-diagonal by construction. Both kinds commute with every $N_x$. The pairing clause is an exact matrix computation on two modes (the section ``Machine verification''): $v^{*}_x v^{*}_y + v_y v_x$ commutes with $\Gamma = \sigma_x\sigma_y$, has off-diagonal entries in the joint number basis, and fails to commute with $N_x$; it is not of the form $d\Gamma(a)$ for a bundle map $a$, since it changes the total occupation by two.

\end{proof}

\begin{remark}[The declared residual, and its size]\label{c:note.singularity.declared-residual-size}
The excluded class is small, as the theorem itself records: for a line bundle the $C^\infty(E)$-linear endomorphisms are multiplication operators, and the supplied class is ultralocal. Every dynamics with a kinetic term lies in the residual: a first-order differential operator of Dirac type has $[D, f] = c(df) \neq 0$, lies outside the $C^\infty$-linear class, and would hop; and so, at the coefficient-linear level, would the pairing terms named in the theorem. The closure constructs none of these, and nothing here excludes one in nature. Two remarks on method. First, support preservation is not the right criterion: a $C^\infty$-linear map preserves supports, but so does every differential operator (Peetre's theorem identifies the support-preserving linear maps of sections with the differential operators), so support preservation alone would not have given Proposition \ref{c:singularity.bundle-maps-quantise-without-hopping}; tensoriality, Lemma \ref{c:singularity.function-linear-endomorphisms-are-bundle-maps}, is what is needed. Second, the necessity witness of Corollary \ref{c:singularity.envelope-endpoint-why-boundary-mode-locality} is a pairing rotation, $|00\rangle \mapsto \tfrac35|00\rangle + \tfrac45|11\rangle$ with $\Delta N = \pm 2$, not a hop: the generator that shows (L$_E$) cannot be dropped is one this theorem's criterion does not exclude, which is consistent, and is why the exclusion of pairing terms stays with the premise. The sharp remaining form of Remark \ref{c:note.singularity.steps-status-what-still-open}(b)'s question is therefore: is the physical dynamics on the singular boundary generated within the supplied class, or does it carry a differential or pairing term? That is a successor question, recorded beside the closure-completeness programme; Proposition \ref{c:singularity.l-stabiliser-tangent} separates the two residual types by what they do to the vacuum, the differential term lying in its stabiliser and the pairing term in the tangent of its orbit.
\end{remark}

\section{Temperature and Hawking radiation}

\begin{remark}[The two boost flows]\label{c:note.decisions.two-boost-flows}
\changedmark{MidnightBlue}{added since the deposit}
The series calls two different one-parameter groups ``the boost flow''. Both are generated by the same element, $K = \sigma/2$ for the completed pair with $\sigma = v^{*}v - vv^{*}$ (and $\sum_j \omega_j K_j$ for the boundary modes), and they differ in the parameter.

The \emph{rotor flow} is $\operatorname{Ad}(R_\lambda)$, $R_\lambda = \exp(\lambda\sigma/2)$, a real rotor flow of the atom: it acts by $v \mapsto e^{-\lambda}v$, $v^{*} \mapsto e^{+\lambda}v^{*}$, so it rescales the modes, its adjoint eigenvalues are real, and it has no period. This is the flow of Theorem \ref{c:singularity.modular-structure-completed-pair}(2) and (5), the only one-parameter flow implemented by real rotors of the atom, and of the clocks of Lemma \emph{The clock inventory of the boundary modes} [T8, Lem 4.2], intrinsic flows of the real Clifford boundary-mode algebra with real adjoint eigenvalues. The scrambling law of \emph{The spreading law, and the chaos bound met, not saturated} [T8, Asm 6.1] uses this flow: a quadratic generator containing the boost scales the ladder operators by $e^{\pm\kappa t}$, and that real growth, exponentiated across the window, gives $t_{*} = \kappa^{-1}\ln N$.

The \emph{phase flow} is $\operatorname{Ad}(e^{itK})$, which on the modes is $a_j \mapsto e^{-i\omega_j t}a_j$: a unitary group whose adjoint eigenvalues are imaginary, periodic in the parameter for each single mode. This is the flow of Theorem \ref{c:singularity.rate-rigidity-temperature-not-algebra-intrinsic}, $\alpha_t = \operatorname{Ad}(e^{itK})$, and the boost dynamics $\tau_t(a_j) = e^{-i\omega_j t}a_j$ of Definition \emph{The boundary-mode algebra, windows, and the boost dynamics} [T7, Def 2.7]. The KMS analysis uses this flow: the descended state of Theorem \ref{c:singularity.modular-descent} is KMS for it, the modular flow in its proof being $\operatorname{Ad}(\rho^{it})$, and the rigidity theorem, the geometric modular condition and the Hawking paper's I-state are stated for it. Proposition \ref{c:singularity.euclidean-rotor-where-it-lies} identifies it as the Euclidean rotation: the boost rotor continued to imaginary rapidity, $R_{i\theta} = \exp(i\theta\sigma/2)$, whose generator $i\sigma$ is not in the real algebra and which is realised by no real rotor of $\Cl(1,3)$.

The two are related by analytic continuation of the parameter, $\lambda = it$: $\operatorname{Ad}(e^{itK}) = \operatorname{Ad}(R_{it})$. The rotor flow fixes the orbits and, through the geometric normalisation, the rate $\kappa = 1/4M$ that converts rapidity to Schwarzschild time, $\lambda = \kappa t$; the phase flow carries the period, $2\pi$ in rapidity for the conjugation action on each atom, which under premise (R) is the modular period $\beta = 2\pi/\kappa = 8\pi M$ in Schwarzschild time. The chaos-bound statement compares one with the other: the rotor flow's growth rate $\kappa$ against $2\pi T$, the phase flow's KMS temperature, which are equal. Where the boundary-mode algebra is complex (the Hawking paper) the phase flow acts within it; where it is a real Clifford algebra (the scrambling results) its generator, $i$ times a bivector, is excluded and only the rotor flow is intrinsic, which is the dynamical face of the second known limit in the conventions. The Hawking paper's dynamics is the phase flow, while the scrambling assumption runs the window's dynamics on the rotor flow. Which of the two is the physical time evolution of a window, in the one Schwarzschild time $t$ that both papers use, is settled by Proposition \emph{The phase flow is the time evolution; the rotor flow is its KMS continuation}: taking a time evolution to be a group of $*$-automorphisms preserving the states, the phase flow is the window's time evolution and the rotor flow its KMS-strip continuation, so for the quadratic dynamics the scrambling law assumes its growth $e^{\kappa t}$ is an imaginary-time rate on the CAR windows; real-time spreading is left open there.
\end{remark}

\begin{remark}[Two clocks, one calibration]\label{c:note.state-selection.two-clocks-one-calibration}
\changedmark{MidnightBlue}{added since the deposit}
Theorem \ref{c:singularity.rate-rigidity-temperature-not-algebra-intrinsic} is the note's no-go, and it is a constructive one: it says exactly what kind of input is needed and how much. The inverse temperature is the ratio of two clocks, the modular clock of the state and the geometric clock of the Killing boost, and a ratio of clocks is meaningless until one of them is calibrated. The chain calibrates exactly once, in \cite[Prop.~6.2]{ButlerT5}: an interval of $2\pi$ in rapidity is $2\pi/\kappa = 8\pi M$ of Schwarzschild time. Everything numerical in Step 4d passes through that one conversion, as the paper already says, and Theorem \ref{c:singularity.rate-rigidity-temperature-not-algebra-intrinsic} now proves that no rearrangement of the algebra can remove the need for it. The residual sign in (2) is not new either: it is the orientation of the boost, the labelling of outgoing against ingoing, which \cite[Rem.~6.5]{ButlerT5} already flags as the $\mathbb{Z}_2$ attached to the weights of $v$ and $v^{*}$. The rigidity theorem identifies that $\mathbb{Z}_2$ as the entire automorphism-theoretic freedom in the temperature.
\end{remark}

\section{Descent, parity and states}

\begin{remark}[Why this rescues rather than damages the chain]\label{c:note.state-selection.why-rescues-rather-than-damages-chain}
\changedmark{MidnightBlue}{added since the deposit}
Read carelessly, Theorem \ref{c:singularity.exact-equilibrium-over-determines-boundary-modes} looks like an obstruction: the state that thermality wants is the state that has no information in it. Read carefully, it is the reason the architecture of this note matters. Had per-mode thermality been derived from the exact equilibrium condition, as hypothesis (ii) suggests, then thermality and correlation would have been in direct competition and the chain could not have had both. Theorem \ref{c:singularity.ascended-marginals-are-diagonal} shows that per-mode thermality does not rest on equilibrium at all: it rests on parity, which survives an arbitrary non-stationary history. The correct state class for the chain is therefore
\[
\bigl\{\ \text{even states, per-mode populations at } \beta = 2\pi/\kappa, \ \text{correlations free}\ \bigr\},
\]
a class in which every single quantum is exactly thermal and the joint state is free to carry everything the evaporation writes. That is exactly the class the finite-ledger model of \cite[\S 12]{ButlerT5} inhabits, and \cite[Lem.~12.1]{ButlerT5} is the model's proof that it does.
\end{remark}

\begin{remark}[The true reading, and why it sharpens the gap]\label{c:note.state-selection.true-reading-why-it-sharpens-gap}
\changedmark{MidnightBlue}{added since the deposit}
Theorem \ref{c:singularity.descended-form-state-independent} should be read as deflationary, not triumphant. For a single fermionic mode, ``thermal'' is not a discriminating property: the state space of one mode, after the descent that the chain performs anyway, is one-dimensional, and every point of it is a Gibbs state. It follows that no part of the physical content of Step 4d lies in the \emph{shape} of the reduced density matrix. All of it lies in the number $\beta_k$, in the claim that $\beta_k$ is the same for all modes, and in the claim that its value is $2\pi/\kappa$. The gap is therefore not ``which class of states'', as the status line ``up to a state'' suggests; it is ``which point of a line, and why the same point at every mode''. Everything that follows is about that line.

This is a sharpening of the gap and not a closing of it, and it is worth being explicit that the sharpening cuts against a certain kind of rhetorical comfort: one cannot claim the Fermi-Dirac \emph{form} as a prediction of the chain, because the descent manufactures that form from any input whatsoever. What the chain can claim is the temperature, and only under the premises named below.
\end{remark}

\begin{remark}[What this says about the two ledgers]\label{c:note.state-selection.what-says-about-two-ledgers}
\changedmark{MidnightBlue}{added since the deposit}
The chain's slogan is that the metric level keeps the occupation records and loses the coherences. Theorem \ref{c:singularity.descent-parity-average} says something sharper and more structural: \emph{the metric level is the parity even level}, and the ladder sector that Rule 1 adjoined is exactly the odd sector. The one-way character of the descent \cite[Fact 5.9]{ButlerT5} is then not a separate fact to be imported but the statement that a $\mathbb{Z}_2$ average cannot be undone: the even subalgebra retains no trace of the odd part it averaged away. It also explains, at a glance, why the descent is the \emph{unique} trace-compatible conditional expectation onto the commutant of $P$: it is the unique one compatible with the grading, and the grading is not a choice but the Clifford algebra's own.
\end{remark}

\begin{remark}[What was bought]\label{c:note.state-selection.what-was-bought}
\changedmark{MidnightBlue}{added since the deposit}
Theorem \ref{c:singularity.ascended-marginals-are-diagonal} replaces hypothesis (ii), a symmetry assumption about the state, with (P), a statement about which elements of the ledger are observable. It is strictly weaker: (ii) constrains the whole state, (P) constrains one complex number per mode. It is also more robust, and this matters for the chain: (ii) is an equilibrium condition, and the evaporating hole is not in equilibrium, while (P) holds throughout an arbitrary non-stationary history. the section ``What equilibrium cannot supply: correlations'' shows that this robustness is not a luxury but a necessity, because the exact equilibrium condition is provably too strong for what Steps 5 and 6b need.
\end{remark}

\section{Entropy and information}

\begin{remark}[The status of (L$_E$)]\label{c:note.state-selection.status-l-e}
\changedmark{MidnightBlue}{added since the deposit}
(L$_E$) is a premise, not a theorem, and it is the true location of a genuine assumption: it says that the framework supplies no hopping generator such as $v_j^{*}v_k$ between distinct points of $E$. That is a statement about what the localisation theorems provide, and they do provide the ladder pair pointwise; it is not a proof that no such coupling can arise from a dynamics not yet written down. Recorded as a programme item in the section ``What remains open''. What Proposition \ref{c:singularity.number-diagonal-dynamics-models-scrambling-derived} buys, if (L$_E$) holds, is worth the exposure: it converts the model's most questionable-looking choice, that the scrambling is phase-type, into the general case, and with it upgrades \cite[Lem.~12.1]{ButlerT5} from a statement about a particular circuit to a statement about the ledger.
\end{remark}

\section{Results used from elsewhere}
The proofs above use the following results of the series, stated briefly; each is proved in the paper named.

\begin{fact}[Exact parallelism, and the convergence rate]\label{c:interior.exact-parallelism-convergence-rate}
In the Schwarzschild interior the adapted physical frame is exactly $D$-parallel along the infalling free-fall worldlines, and the Hubble-normalised expansion rates reach the Kasner exponents $(-\tfrac13, \tfrac23, \tfrac23)$ at the single, non-oscillating rate $O(\tau^{2/3})$.
\emph{Source:} T6, Lem 5.2 (\href{https://doi.org/10.5281/zenodo.22721990}{doi:10.5281/zenodo.22721990}).
\end{fact}

\begin{fact}[The interior connection, exactly]\label{c:interior.interior-connection-exactly}
In the parallel frame of the Schwarzschild interior the rotation-gauge field $\Omega$ is exactly $-\sigma_i B_i$ on each unit tangential direction, plus the round-sphere rotation $(\cot\theta/r)R_{23}$ on $e_3$, with $\Omega(e_0) = 0$.
\emph{Source:} T6, Lem 5.3 (\href{https://doi.org/10.5281/zenodo.22721990}{doi:10.5281/zenodo.22721990}).
\end{fact}

\begin{fact}[Kasner form of the interior]\label{c:interior.kasner-form-interior}
In proper time $\tau$ the derived metric of the Schwarzschild interior field is synchronous and diagonal, with Kasner exponents $(-\tfrac13, \tfrac23, \tfrac23)$ to leading order and explicit series corrections.
\emph{Source:} T6, Lem 5.1 (\href{https://doi.org/10.5281/zenodo.22721990}{doi:10.5281/zenodo.22721990}).
\end{fact}

\begin{fact}[Weighted limit theorem]\label{c:interior.weighted-limit-theorem}
For Schwarzschild and every admissible metric the Hubble-normalised connection tends to $-\sum_i p_i B_i \otimes \hat\theta^i$, the scale tractor extends as a parallel section, and the boundary holonomy algebra is exactly $\Lambda^2 W \cong \mathfrak{so}(1,3)$.
\emph{Source:} T6, Thm 5.6 (\href{https://doi.org/10.5281/zenodo.22721990}{doi:10.5281/zenodo.22721990}).
\end{fact}

\begin{fact}[Classification of the real Clifford algebras]\label{c:primer.classification-real-clifford-algebras}
$\Cl(p,q)$ is a matrix algebra over $\R$, $\C$ or $\HH$, or a sum of two such, according to $p - q \bmod 8$, and it is simple except when $p - q \equiv 1 \pmod 4$.
\emph{Source:} T0, Thm 9.3 (\href{https://doi.org/10.5281/zenodo.22714501}{doi:10.5281/zenodo.22714501}).
\end{fact}

\begin{fact}[Base unipotents at the wall]\label{c:sheaf.base-unipotents-wall}
At a wall point the shears $g_\lambda(\alpha) = \alpha + \alpha(\lambda)\,k$, $\lambda \in \ker k$, are $G$-isometries fixing $k$ that act on transversals as a free and transitive group of translations.
\emph{Source:} T4, Lem 10.1 (\href{https://doi.org/10.5281/zenodo.22721669}{doi:10.5281/zenodo.22721669}).
\end{fact}

\begin{fact}[Function-linearity of the tame constructions]\label{c:sheaf.function-linearity-tame-constructions}
All the tame constructions are linear over smooth functions and stay so under composition and the standard operations, so every endomorphism they supply is tensorial.
\emph{Source:} T4, Prop 9.17 (\href{https://doi.org/10.5281/zenodo.22721669}{doi:10.5281/zenodo.22721669}).
\end{fact}

\begin{fact}[Localised transport]\label{c:sheaf.localised-transport}
The degenerate Clifford sheaf embeds canonically in the completed sheaf, and on the wall the radical covector acquires a partner with which it forms a ladder pair that exists exactly over the wall.
\emph{Source:} T4, Cor 4.5 (\href{https://doi.org/10.5281/zenodo.22721669}{doi:10.5281/zenodo.22721669}).
\end{fact}

\begin{fact}[No equivariant value at the wall]\label{c:sheaf.no-equivariant-value-wall}
The completed structure at the wall has automorphisms that translate candidate values of the dying direction freely, so no equivariant value, and no canonical ladder partner, exists.
\emph{Source:} T4, Thm 10.4 (\href{https://doi.org/10.5281/zenodo.22721669}{doi:10.5281/zenodo.22721669}).
\end{fact}

\begin{fact}[One-sided completion]\label{c:sheaf.one-sided-completion}
Under tameness, with signature $(p,q)$, $p \geq 1$, on the physical side and a trivial radical line bundle, coupling an adjoined hyperbolic plane to the dying direction gives an ambient form of constant signature $(p{+}1, q{+}1)$ on the whole physical side, wall included.
\emph{Source:} T4, Thm 4.4 (\href{https://doi.org/10.5281/zenodo.22721669}{doi:10.5281/zenodo.22721669}).
\end{fact}

\begin{fact}[One-sided completion at the singularity]\label{c:sheaf.one-sided-completion-singularity}
At the blown-up Schwarzschild singularity, coupling an adjoined hyperbolic plane to the dying direction with $\mu\nu < 0$ and bump inside the horizon gives an ambient form of constant signature $(2,4)$ everywhere, boundary included, with no shadow wall and a ladder pair existing exactly over the boundary.
\emph{Source:} T4, Thm 8.12 (\href{https://doi.org/10.5281/zenodo.22721669}{doi:10.5281/zenodo.22721669}).
\end{fact}

\begin{fact}[The Schwarzschild degeneration is boundary-tame in $\shW$]\label{c:sheaf.schwarzschild-degeneration-boundary-tame}
\changedmark{gray}{reworded}
In the weighted module at the blown-up Schwarzschild singularity, the derived cometric of the Painlev\'e-Gullstrand gauge field extends smoothly to the boundary and degenerates there with corank exactly one and to exactly first order, with radical $k = \rho + m\,\dtPG$ annihilating the infalling congruence.
\emph{Source:} T4, Thm 8.7 (\href{https://doi.org/10.5281/zenodo.22721669}{doi:10.5281/zenodo.22721669}).
\end{fact}

\begin{fact}[The solve, closed by the radical gauge]\label{c:sheaf.solve-closed-radical-gauge}
For each $v \in \ker k$ the equation $G(\alpha, \tilde v) = \alpha(v)$ has a solution unique modulo $\R k$, and the normalisation $G(\tilde v, \tilde v) + 2\tilde v(u) = 0$ selects exactly one.
\emph{Source:} T4, Lem 10.2 (\href{https://doi.org/10.5281/zenodo.22721669}{doi:10.5281/zenodo.22721669}).
\end{fact}

\begin{fact}[Transversal independence at the wall]\label{c:sheaf.transversal-independence-wall}
Explicit unipotent isometries, composing coherently, identify the completions built on all collar transversals, preserving form, grading and degeneration subalgebra but not the ladder partner.
\emph{Source:} T4, Thm 10.3 (\href{https://doi.org/10.5281/zenodo.22721669}{doi:10.5281/zenodo.22721669}).
\end{fact}

\begin{fact}[Imported: the average entropy of a subsystem (Page)]\label{c:singularity.imported-page-average-entropy-subsystem}
\changedmark{BurntOrange}{changed since the deposit}
A random pure state of a bipartite system leaves the smaller part within half a nat of maximal entropy: $\langle S_A\rangle \simeq \ln m - m/2n$ (Page).
\emph{Source:} D. N. Page, Average entropy of a subsystem, Phys. Rev. Lett. 71 (1993), 1291--1294.; S. K. Foong and S. Kanno, Proof of Page's conjecture on the average entropy of a subsystem, Phys. Rev. Lett. 72 (1994), 1148--1151.; J. Sanchez-Ruiz, Simple proof of Page's conjecture on the average entropy of a subsystem, Phys. Rev. E 52 (1995), 5653--5655.; S. Sen, Average entropy of a quantum subsystem, Phys. Rev. Lett. 77 (1996), 1--3..
\end{fact}

\begin{fact}[Bivectors form the Lie algebra of the orthogonal group]\label{c:std-algebra.bivector-lie-algebra}
$\Cl^{2}$ is closed under $A \times C = \langle AC \rangle_2$, acts skew-adjointly on $V$ by $A \times x = \langle Ax \rangle_1$, and is isomorphic to $\so(V,Q)$ when $Q$ is nondegenerate.
\emph{Source:} the algebra core of the Transport Lab, proved there with an exact check; standard, see C. Doran and A. Lasenby, Geometric Algebra for Physicists, Cambridge University Press (2003), Ch. 4 (foundations of geometric algebra: rotors, versors).
\end{fact}

\begin{fact}[Cl(1,1) and Cl(2,0) are the real 2x2 matrices]\label{c:std-algebra.cl11-matrix}
$\Cl(H) \cong \Mat_2(\R)$ for a hyperbolic plane, with a hyperbolic pair going to the off-diagonal matrix units; also $\Cl(2,0) \cong \Mat_2(\R)$.
\emph{Source:} the algebra core of the Transport Lab, proved there with an exact check; standard, see P. Lounesto, Clifford Algebras and Spinors, 2nd ed., London Mathematical Society Lecture Note Series 286, Cambridge University Press (2001), Ch. 16 (matrix representations and periodicity of 8).
\end{fact}

\begin{fact}[Orthogonal vectors anticommute]\label{c:std-algebra.orthogonal-anticommute}
$uv = -vu$ iff $B(u,v) = 0$; in an orthogonal basis $e_I e_J = \pm \prod_{I \cap J} Q(e_i)\; e_{I \triangle J}$ with the sorting sign.
\emph{Source:} the algebra core of the Transport Lab, proved there with an exact check; standard, see C. Doran and A. Lasenby, Geometric Algebra for Physicists, Cambridge University Press (2003), Ch. 4 (foundations of geometric algebra).
\end{fact}

\begin{fact}[How the pseudoscalar commutes]\label{c:std-algebra.pseudoscalar-commutation}
$I X_k = (-1)^{k(n-1)} X_k I$; $I\tilde{I} = \prod_i Q(e_i)$, and $I^{2} = 0$ in a degenerate algebra.
\emph{Source:} the algebra core of the Transport Lab, proved there with an exact check; standard, see C. Doran and A. Lasenby, Geometric Algebra for Physicists, Cambridge University Press (2003), Ch. 4 (foundations of geometric algebra).
\end{fact}

\begin{fact}[Pure spinors, Cartan's purity theorem and the Thouless form]\label{c:std-algebra.pure-spinors}
The annihilator of a spinor is totally null; pure spinors (annihilator of dimension $k$) correspond to maximal null subspaces, all definite-parity spinors are pure for $k \le 3$, and $\exp(\tfrac12 Z_{ij} w_i w_j)$ is pure.
\emph{Source:} E. Cartan, Lecons sur la theorie des spineurs, Hermann (1938)., pure spinors and their dimension count; C. Chevalley, The Algebraic Theory of Spinors, Columbia University Press (1954)., Ch. III (pure spinors); D. J. Thouless, Stability conditions and nuclear rotations in the Hartree-Fock theory, Nucl. Phys. 21 (1960) 225., the Thouless form.
\end{fact}

\begin{fact}[The spacetime algebra as a matrix algebra, and the spacetime split]\label{c:std-algebra.sta-matrix}
$\Cl(1,3) \cong \Mat_2(\HH)$, the Dirac matrices represent $\C \otimes \Cl(1,3) \cong \Mat_4(\C)$, and $\sigma_k = \gamma_k\gamma_0$ give $\Cl^{+}(1,3) \cong \Cl(3,0)$ with $\sigma_1\sigma_2\sigma_3 = I$.
\emph{Source:} the algebra core of the Transport Lab, proved there with an exact check; standard, see C. Doran and A. Lasenby, Geometric Algebra for Physicists, Cambridge University Press (2003), Ch. 5 (relativity and spacetime).
\end{fact}

\begin{fact}[The product of two vectors]\label{c:std-algebra.vector-squares}
$v^{2} = Q(v)$ and $uv + vu = 2B(u, v)$, so $uv = B(u, v) + \tfrac12(uv - vu)$.
\emph{Source:} the algebra core of the Transport Lab, proved there with an exact check; standard, see C. Doran and A. Lasenby, Geometric Algebra for Physicists, Cambridge University Press (2003), Ch. 4 (foundations of geometric algebra).
\end{fact}

\begin{fact}[Null pairs satisfy the canonical anticommutation relations]\label{c:std-algebra.witt-basis-car}
For a Witt basis ($B(v_i,v_j) = B(w_i,w_j) = 0$, $B(v_i,w_j) = \tfrac12\delta_{ij}$), $\{v_i,v_j\} = \{w_i,w_j\} = 0$ and $\{v_i,w_j\} = \delta_{ij}$ in $\Cl(V,Q)$.
\emph{Source:} the algebra core of the Transport Lab, proved there with an exact check; standard, see C. Doran and A. Lasenby, Geometric Algebra for Physicists, Cambridge University Press (2003), Ch. 9 (multiparticle states).
\end{fact}

\begin{fact}[Rigidity; the clock is frozen]\label{c:tower.rigidity-clock-frozen}
Every floor of the tower shares the ground floor's charge, division-ring type, centre, simplicity class, Morita class and obstruction taxonomy, so the physical tower $\Cl(1{+}k, 3{+}k) \cong \Mat_{2^{k+1}}(\HH)$ is frozen at clock position $6$.
\emph{Source:} T3, Thm 3.6 (\href{https://doi.org/10.5281/zenodo.22721556}{doi:10.5281/zenodo.22721556}).
\end{fact}

\begin{fact}[Cohomological form of the no-go]\label{c:transport.cohomological-form-no-go}
The obstruction to an invariant pseudo-inverse of a null vector is the tautological class $\theta \in \mathrm{Hom}(N_v, U)$, which vanishes exactly when $U = 0$ and persists unchanged through every ascent and every larger symmetry group containing $N_v$.
\emph{Source:} T1, Thm 4.8 (\href{https://doi.org/10.5281/zenodo.21743863}{doi:10.5281/zenodo.21743863}).
\end{fact}

\begin{fact}[Descent: the read-out map]\label{c:transport.descent-read-out-map}
The read-out $\dsc[s\,\nO + x + t\,\nInf] = x/s$ inverts the lift off the ideal points, returns $x/Q(x)$ after the transported reciprocal, and is undefined exactly where $Q(x) = 0$.
\emph{Source:} T1, Lem 4.23 (\href{https://doi.org/10.5281/zenodo.21743863}{doi:10.5281/zenodo.21743863}).
\end{fact}

\begin{fact}[The door is one-way]\label{c:transport.door-one-way}
There is no unital algebra homomorphism from either minimal completion of a nondegenerate $\Cl(p,q)$ back to $\Cl(p,q)$ (the Rule 2 one existing only when $p - q \equiv 1 \pmod 4$), so every descent is non-multiplicative.
\emph{Source:} T1, Thm 7.6 (\href{https://doi.org/10.5281/zenodo.21743863}{doi:10.5281/zenodo.21743863}).
\end{fact}

\begin{fact}[Existence and classification of type C]\label{c:transport.existence-classification-type-c}
$\Cl(p,q)$, $n \geq 1$, has a nontrivial central idempotent exactly when $p - q \equiv 1 \pmod 4$, and then the only ones are $\varepsilon_\pm = \tfrac12(1 \pm \omega)$, which split the algebra into two simple ideals and are exchanged by the grade involution but by no inner automorphism.
\emph{Source:} T1, Thm 3.6 (\href{https://doi.org/10.5281/zenodo.21743863}{doi:10.5281/zenodo.21743863}).
\end{fact}

\begin{fact}[Existence locus of type N]\label{c:transport.existence-locus-type-n}
$\Cl(p,q)$ contains null vectors exactly when $pq \neq 0$; they then form a cone of dimension $n - 1$, and each is a two-sided zero divisor with no left or right inverse in any unital algebra containing $\Cl(p,q)$.
\emph{Source:} T1, Prop 3.3 (\href{https://doi.org/10.5281/zenodo.21743863}{doi:10.5281/zenodo.21743863}).
\end{fact}

\begin{fact}[No-go at the level of elements]\label{c:transport.no-go-level-elements}
In mixed signature with $n \geq 3$, no assignment of von Neumann pseudo-inverses to null vectors, valued in the Clifford algebra of any extension, is equivariant even under $\Spin^0(p,q)$.
\emph{Source:} T1, Thm 4.2 (\href{https://doi.org/10.5281/zenodo.21743863}{doi:10.5281/zenodo.21743863}).
\end{fact}

\begin{fact}[The pinching, and its kernel]\label{c:transport.pinching-kernel}
The pinching by the descent projector is a conditional expectation onto its commutant whose kernel is exactly the adjoined ladder sector, and it returns transported elements unchanged.
\emph{Source:} T1, Prop 7.5 (\href{https://doi.org/10.5281/zenodo.21743863}{doi:10.5281/zenodo.21743863}).
\end{fact}

\begin{fact}[Swap-and-rescale normal form]\label{c:transport.swap-rescale-normal-form}
In homogeneous coordinates $(s, x, t)$ with $Q(x) = 2st$, the transported reciprocal is the linear involution $(s, x, t) \mapsto (2t, x, \frac12 s)$, and the ideal points $s = 0$ fall into three strata.
\emph{Source:} T1, Prop 4.21 (\href{https://doi.org/10.5281/zenodo.21743863}{doi:10.5281/zenodo.21743863}).
\end{fact}

\begin{fact}[Totality]\label{c:transport.totality}
On the quadric of the completion, $\iota[N] = [e_+ N e_+]$ is a total continuous involution that equals the reciprocal $x \mapsto x/Q(x)$ on invertible vectors and sends each $x_0$ with $Q(x_0) = 0$ to $[x_0 + \frac12 \nInf]$, and it is the only such continuous extension.
\emph{Source:} T1, Thm 4.10 (\href{https://doi.org/10.5281/zenodo.21743863}{doi:10.5281/zenodo.21743863}).
\end{fact}

\begin{fact}[The transport algorithm]\label{c:transport.transport-algorithm}
The transport results assemble into a six-step procedure (locate, diagnose, apply the rule, compute, descend, iterate) in which every step is licensed by a proved result and no case is left uncovered.
\emph{Source:} T1, Sec 8.1 (\href{https://doi.org/10.5281/zenodo.21743863}{doi:10.5281/zenodo.21743863}).
\end{fact}

\begin{fact}[Type P is self-repairing]\label{c:transport.type-p-self-repairing}
In any unital ring the identity assignment $\sigma(e) = e$ on idempotents satisfies both von Neumann identities and is equivariant under every automorphism of every extension preserving the ring, so type P obstructions are repaired at signature increment $(0,0)$.
\emph{Source:} T1, Prop 6.1 (\href{https://doi.org/10.5281/zenodo.21743863}{doi:10.5281/zenodo.21743863}).
\end{fact}

\begin{refutation}[The original conjecture, refuted below]\label{c:transport.original-conjecture-refuted-below}
Element-by-element repair of a null vector by an equivariant partner one ascent up: refuted.
\end{refutation}

\section*{Changes since the deposited version}
\addcontentsline{toc}{section}{Changes since the deposited version}
Since the deposited version 1.0: 6 results corrected, 10 changed, 19 reworded (wording, notation or citations only), 7 results added and 33 cited results added. The comparison is with the source of version 1.0 on Zenodo, \href{https://doi.org/10.5281/zenodo.22721790}{doi:10.5281/zenodo.22721790}. A result counts as revised when the words or mathematics of its statement or proof differ from that source; a citation added in a parenthesis does not count. A result the paper cites from elsewhere is compared with the deposit of the paper that proves it. Each revision is classed as corrected, changed or reworded; the classes were assigned by a review of each difference and the store's record of why it was made, and a revision not yet reviewed is counted as changed.
\subsection*{Corrected}
The deposited version was wrong or claimed too much here.
\begin{itemize}\setlength{\itemsep}{2pt}
\item Postulate~\ref{c:singularity.c2-cell-density} (\emph{C2: cell density}), Axiom 10.2 in the deposit: statement. Corrected: nu is not of order one; the area law needs nu about 0.029. Also scoped to Schwarzschild calibration. \emph{Store notes.} \emph{2026-10-22:} Version 3 adds a note, at Leon's decision to keep (C1) and (C2) rather than replace them by (C2'), that their forms are forced by an algebraic curvature criterion plus isometry invariance leaving one constant (synthesis.c1-c2-forms-forced-one-number) and that (C2') on the exceptional set is equivalent at leading order [...] \emph{2026-10-23:} Version 4 adds a scope sentence, at Leon's decision of 23 October (cutoff premises as a Schwarzschild calibration only), that (C2) with (C1) and (IC) is read as a calibration on the Schwarzschild family, since read literally the count fails the area law on Schwarzschild-de Sitter (u/(3u-2) times A\_b/4 over the hole's modular period [...]
\item Corollary~\ref{c:singularity.matching-steps-4b-4c-join-equation} (\emph{The matching: Steps 4b and 4c join by equation}), Cor 6.4 in the deposit: statement and proof. Corrected: Step 4c's modular flow is the phase flow Ad(exp(it sigma/2)), not the real rotor flow, which is its continuation. \emph{Store notes.} \emph{2026-10-21:} The proof covers only the equation pair; a closing sentence states the role of the three premises the statement's location clause uses (the localisation theorem cited there as T4 Thm 4.4 and Cor 4.5). \emph{2026-10-21:} The two boost flows: flow\_note recorded; formal now name the flow (original formal in source\_formal, proof in orig-proof-*.tex where changed; version bumped).
\item Theorem~\ref{c:singularity.kms-uniqueness-factor} (\emph{KMS uniqueness at a factor}), Thm 7.13 in the deposit: statement. Corrected: the clause on the family of KMS states needs K not a multiple of 1, where the flow is trivial. \emph{Store notes.} \emph{2026-10-22:} Merged the state-selection note's Thm 4.2 (state-selection.kms-uniqueness-factor, retired to tombstones.yaml) into this entity: the same theorem. Kept this statement, levels and checks; the note's entity is recorded above and its proof files are kept in this folder. References to the note's id were repointed here. Extraction-only change, so the version and [...]
\item Proposition~\ref{c:singularity.why-calibration-cannot-be-internal} (\emph{Why the calibration cannot be internal}), Prop 7.25 in the deposit: statement and proof. Corrected: regularity across the horizon is no longer stated as equivalent in general, only for states KMS for the Killing flow; GTG restatement added. \emph{Store notes.} \emph{2026-10-21:} The two boost flows: flow\_note recorded. \emph{2026-10-23:} Restated in Gauge Theory Gravity per Leon 2026-10-23 (DICTIONARY, Bifurcation surfaces): the locus where (R) applies is the horizon locus r = 2M in the Painleve-Gullstrand chart, on a cross-section of area 16 pi M\^{}2 with (K.D)K = kappa K; the bifurcation surface itself noted as beyond one flat chart. Same conclusion; proof given a one-sentence pointer. Fix [...]
\item Theorem~\ref{c:singularity.bosonisation-bridge-s-wave-sector} (\emph{The bosonisation bridge, s-wave sector}), Thm 8.2 in the deposit: statement and proof. Corrected: invariance of the current span under the modular flow holds only for weights affine in k, so it uses Premise (Q). \emph{Store notes.} \emph{2026-10-22:} Clause (2) and its proof now say that the invariance of the span of the currents under the modular flow uses the weight clause of Premise (Q): each current is a weight vector only when the mode weights are affine in k (omega\_k = kappa k under (Q)); the earlier text stated the invariance as if it held for any weights. (Q) was already among the assumes [...] \emph{2026-10-23:} , note only, statement and proof unchanged. Clauses (2) and (3) display the weights omega\_k = kappa k of Premise (Q) version 2 (T05's Premise 8.1). Premise (Q) version 3 (Leon's decision of 23 October) takes kappa(k + 1/2). The weight difference of each bilinear is still kappa m, so (2) holds unchanged; the Planck occupancy of (3), m/(e\^{}\{beta kappa m\} - 1) [...]
\end{itemize}
\subsection*{Changed}
The content changed by a later decision, not because the deposited version was wrong.
\begin{itemize}\setlength{\itemsep}{2pt}
\item Postulate~\ref{c:singularity.c1-curvature-cutoff} (\emph{C1: curvature cutoff}), Axiom 10.1 in the deposit: statement. By decision, scope limited to a Schwarzschild calibration, with a note that the premise's form is forced up to one constant. \emph{Store notes.} \emph{2026-10-23:} Version 3 adds a scope sentence, at Leon's decision of 23 October (cutoff premises as a Schwarzschild calibration only), that (C1) with (C2) and (IC) is read as a calibration on the Schwarzschild family, since read literally the count fails the area law on Schwarzschild-de Sitter (u/(3u-2) times A\_b/4 over the hole's modular period [...] \emph{2026-10-24:} Referee fixes to the GTG wording pass (2026-10-24): 'a gravitational field of nonvanishing field strength' (gauge-independent) in place of 'a nonvanishing gravitational field'; global content read on several charts of flat spacetime, each with its own gravitational field, glued by displacement and rotation gauge transformations (complexified charts for the [...]
\item Postulate~\ref{c:singularity.l-boundary-mode-locality} (\emph{(L$_E$): boundary-mode locality, as descent covariance on finite windows}), Premise 7.30 in the deposit: statement. Premise restated as descent covariance on finite windows (equivalent), with a note that it forbids fast real-time scrambling. \emph{Store notes.} \emph{2026-10-22:} Version 2 restates the premise as descent covariance on finite windows (the dynamics commutes with the descent at all times), equivalent to the version 1 wording by synthesis.boundary-mode-locality-is-descent-covariance, and records that it forbids exponential real-time commutator growth (synthesis.number-diagonal-dynamics-spreads-nothing), so the real-time [...]
\item Postulate~\ref{c:singularity.q-joint-quasi-freeness-s-wave} (\emph{(Q): joint quasi-freeness, with the s-wave weights at the half-integer offset}), Premise 8.1 in the deposit: statement. By decision, mode weights become kappa(k + 1/2) instead of kappa k, so no zero-weight mode; results under kappa k refer to version 2. \emph{Store notes.} \emph{2026-10-23:} Version 3 (Leon's decision of 23 October on the array offset): the weight clause's offset becomes kappa(k + 1/2), the sea level midway between two sites, adopted as a consequence of reading the array's modes as the discretised modes of the field theory's spinor fields on the horizon's light rays (boost weight 1/2) in the state that (N+) [...]
\item Postulate~\ref{c:singularity.r-regularity-fixed-locus} (\emph{(R): regularity at the fixed locus}), Premise 7.19 in the deposit: statement. A note added: for the eternal hole (R) follows from the null ground-state premise; restated in Gauge Theory Gravity. \emph{Store notes.} \emph{2026-10-23:} Version 2 (Leon's decisions 1 and 13 of 23 October): adds a note that the continuum restatement of (R), for the eternal hole, follows from the new postulate synthesis.n-plus-null-ground-state by synthesis.positive-null-energy-gives-r-premises (iii), and that this point-spectrum wording is its discretisation; the premise keeps postulate status and its [...] \emph{2026-10-23:} Restated in Gauge Theory Gravity: the Euclidean section is that of the Schwarzschild field on flat spacetime, and its regularity is named as regularity at the tip, a local condition (DICTIONARY 14 item 5, Bifurcation surfaces section); a parenthetical points to the flag on (N+) for the eternal hole. Premise content, scope and postulate status unchanged [...]
\item Theorem~\ref{c:singularity.modular-descent} (\emph{Modular descent}), Thm 7.3 in the deposit: statement and proof. The flow is now named precisely as the phase flow of sigma, and premise (R) is cited at the normalisation step. \emph{Store notes.} \emph{2026-10-21:} Cited (R) at the normalisation step, which the statement identifies as premise (R): the KMS period 2 pi / kappa of each single-mode marginal. \emph{2026-10-21:} The two boost flows: flow\_note recorded; formal, one-liner, technical now name the flow (original formal in source\_formal, proof in orig-proof-*.tex where changed; version bumped).
\item Definition~\ref{c:obj.singularity.geometric-modular-condition-g} (\emph{The geometric modular condition (G)}), Def 7.12 in the deposit: statement. The definition now specifies the phase flow generated by K\_geo, not its real rotor flow. \emph{Store notes.} \emph{2026-10-21:} The two boost flows: flow\_note recorded; formal, technical, one-liner now name the flow (original formal in source\_formal, proof in orig-proof-*.tex where changed; version bumped). \emph{2026-10-22:} Merged the state-selection note's Def 4.1 (obj.state-selection.geometric-modular-condition, retired to tombstones.yaml) into this entity: the same set of states (the note's definition now says faithful, as T05's does). Kept this statement, levels and checks; the note's entity is recorded above and its proof files are kept in this folder. References to the [...]
\item Proposition~\ref{c:singularity.euclidean-section-period} (\emph{The Euclidean section, and its period}), Prop 7.18 in the deposit: statement. Restated in Gauge Theory Gravity, with a new remark that the regularity condition is local at the tip; period 8 pi M unchanged. \emph{Store notes.} \emph{2026-10-23:} Restated in Gauge Theory Gravity: the Euclidean section is the Wick-rotated static gauge of the Schwarzschild field on the flat Euclidean background; rho = 0 named as the tip of the section (the Euclidean counterpart of the bifurcation sphere), and the regularity condition noted as local (DICTIONARY 14 item 5). Same period 8 pi M.
\item Assembled~\ref{c:singularity.assembled-page-curve-up-typicality} (\emph{Assembled: the Page curve, up to typicality}), Fact 11.1 in the deposit: proof. Status changed from imported Fact to proved assembly, with a step-by-step proof and Page's theorem split out as a separate import. \emph{Store notes.} \emph{2026-10-21:} Kept as imported, so no derivation: no proof uses it, the rate model being built to replace its typicality; that model's premise (U) is taken from its hypothesis (i), recorded as lineage on the premise. \emph{2026-10-21:} Was status imported with no derivation; it is an assembly of three series results (eviction (2), the count, the area law) and Page's theorem, so it is now proved, with Page's theorem the separate import singularity.imported-page-average-entropy-subsystem. The proof is the paper's own 'Assembly and provenance' paragraph set out as steps.
\item Assembled~\ref{c:singularity.imported-interior-collar-kasner-form} (\emph{Imported: the interior collar in Kasner form}), Fact 13.1 in the deposit: statement and proof. Connection values re-signed to the store's GTG sign convention (minus sigma\_1 B\_1); GTG restatement otherwise, content unchanged. \emph{Store notes.} \emph{2026-10-23:} GTG restatement: the Schwarzschild interior is the gravitational field (hb, Omega) on flat spacetime with synchronous derived metric; the frame is the physical frame, D-parallel along the free-fall worldlines; the Levi-Civita connection is the rotation-gauge field. Series, limits, algebra and exceptional set unchanged. Proof: one opening sentence pointing [...]
\item Fact~\ref{c:singularity.imported-page-average-entropy-subsystem} (\emph{Imported: the average entropy of a subsystem (Page)}), Fact 11.1 in the deposit: statement. New separate import stating Page's average-entropy theorem, split out of the assembled Page curve Fact. \emph{Store notes.} \emph{2026-10-21:} Created so that the assembled Page curve (T05 Fact 11.1) can be a proved assembly whose one outside ingredient is this import; T05 cites Page's theorem inside that Fact and in its introduction ("Page's curve and typicality theorem").
\end{itemize}
\subsection*{Reworded}
Wording, notation, terminology or citations only; the mathematics and the claim are as deposited.
\begin{itemize}\setlength{\itemsep}{2pt}
\item Proposition~\ref{c:singularity.connection-bivector-computed} (\emph{The connection bivector, computed}), Prop 2.3 in the deposit: statement and proof. Gauge Theory Gravity phrasing: the connection one-forms are components of the rotation-gauge field; values unchanged.
\item Lemma~\ref{c:singularity.pseudo-inverse-identities-hyperbolic-atom} (\emph{Pseudo-inverse identities, and the hyperbolic atom}), Lem 4.1 in the deposit: proof. Citations only: broken equation references repaired and standard facts cited.
\item Proposition~\ref{c:singularity.transport-schwarzschild-fibre} (\emph{The transport, on the Schwarzschild fibre}), Prop 5.2 in the deposit: statement and proof. Gauge Theory Gravity phrasing: Q\_p is the square of the physical vector; clauses (1) to (4) unchanged.
\item Assembled~\ref{c:singularity.imported-descent-one-way-conditional-expectation} (\emph{Imported: the descent is a one-way conditional expectation}), Fact 5.9 in the deposit: statement and proof. Presentation only: the imported statement and provenance are repeated, with citations formatted.
\item Theorem~\ref{c:singularity.curvature-gauge-field-exact-shape-decoupling} (\emph{The curvature of the gauge field: exact shape decoupling}), Thm 6.1 in the deposit: proof. Citations only: standard facts cited, and the classical type-D spectrum attributed.
\item Remark~\ref{c:note.singularity.cost-r-dictionary} (\emph{The cost of (R), and the dictionary}), Rem 7.26 in the deposit: statement. Wording only: T9 named as the correspondence paper on the conformal spinor space rather than the twistor correspondence.
\item Proposition~\ref{c:singularity.p-reduced-metric-level-reading-many} (\emph{(P$_{\mathrm{loc}}$), reduced: the metric-level reading on many modes}), Prop 9.6 in the deposit: proof. Cross-references updated and a sentence makes explicit that the modes are the localised ladder pairs.
\item Lemma~\ref{c:singularity.causal-epoch-computed-exactly} (\emph{The causal epoch, computed exactly}), Lem 10.4 in the deposit: statement and proof. Gauge Theory Gravity wording: rain geodesics as worldlines of D-parallel velocity; same integral and expansion.
\item Fact~\ref{c:singularity.imported-induced-geometry-interior-slice} (\emph{Imported: the induced geometry of an interior slice}), Fact 10.5 in the deposit: statement. Gauge Theory Gravity wording of the same induced metric, with a note that static and PG times agree on the slice.
\item Proposition~\ref{c:singularity.stretched-singularity-count} (\emph{The stretched singularity, and the count}), Prop 10.6 in the deposit: statement and proof. Gauge Theory Gravity restatement on flat spacetime with cross-references; same claim and mathematics.
\item Proposition~\ref{c:singularity.area-law-test} (\emph{The area-law test}), Prop 10.7 in the deposit: proof. Citations and cross-references added, including the source of the epoch reading; argument unchanged.
\item Proposition~\ref{c:singularity.no-attic-derived-model} (\emph{No attic, derived in the model}), Prop 11.4 in the deposit: proof. Citations and cross-references only.
\item Theorem~\ref{c:singularity.eviction-window-thermal-product-radiation-entropy} (\emph{Eviction: the window is a thermal product, and the radiation entropy is its content}), Thm 11.5 in the deposit: proof. Citations and cross-references only.
\item Corollary~\ref{c:singularity.what-event-returns-radiation} (\emph{What an event returns to the radiation}), Cor 11.7 in the deposit: proof. Cross-references added, with a sentence naming the result the partner-correlation clause rests on.
\item Proposition~\ref{c:singularity.two-levels-finite-model} (\emph{The two levels of the finite model}), Prop 12.5 in the deposit: proof. Cross-references and attributions added; content unchanged.
\item Proposition~\ref{c:singularity.modular-flow-boundary-modes-holonomy-universal} (\emph{The modular flow of the boundary modes as a holonomy, and its universal temperature}), Prop 13.2 in the deposit: statement. Notation only: the bracket written as the commutator product with its normalisation stated; cross-references updated.
\item Fact~\ref{c:singularity.imported-surface-gravity-redshifted-acceleration} (\emph{Imported: surface gravity as redshifted acceleration}), Fact A.1 in the deposit: statement. Gauge Theory Gravity wording of Wald's result with the covariant derivative D and Q(K); content unchanged.
\item Proposition~\ref{c:singularity.computed} (\emph{$\kappa = 1/4M$, computed}), Prop A.2 in the deposit: statement and proof. Gauge Theory Gravity wording; same value kappa = 1/4M.
\item Fact~\ref{c:sheaf.schwarzschild-degeneration-boundary-tame} (\emph{The Schwarzschild degeneration is boundary-tame in $\shW$}), Fact B.1 in the deposit: statement and proof. Gauge Theory Gravity phrasing: the PG cometric is the derived cometric of the gauge field; Gram and identities unchanged.
\end{itemize}
\subsection*{Added}
Results stated here that the deposited paper did not state.
\begin{itemize}\setlength{\itemsep}{2pt}
\item Remark~\ref{c:note.decisions.two-boost-flows} (\emph{The two boost flows}). New remark recording the maintainer's decision that the two flows the series calls "the boost flow" be named distinctly (seed/CONTRADICTIONS\_STRUCTURE.md H2). Each result that speaks of the boost flow records which one it means under extraction.flow\_note.
\item Remark~\ref{c:note.state-selection.two-clocks-one-calibration} (\emph{Two clocks, one calibration}). The state-selection note's duplicates of T05 entities were retired to tombstones.yaml (merged into the T05 entities). References moved from the retired state-selection.rate-rigidity-temperature-not-algebra-intrinsic to singularity.rate-rigidity-temperature-not-algebra-intrinsic; version 1 -> 2.
\item Remark~\ref{c:note.state-selection.why-rescues-rather-than-damages-chain} (\emph{Why this rescues rather than damages the chain}). The state-selection note's duplicates of T05 entities were retired to tombstones.yaml (merged into the T05 entities). References moved from the retired state-selection.ascended-marginals-are-diagonal, state-selection.exact-equilibrium-over-determines-ledger to singularity.ascended-marginals-are-diagonal [...]
\item Remark~\ref{c:note.state-selection.true-reading-why-it-sharpens-gap} (\emph{The true reading, and why it sharpens the gap}). The state-selection note's duplicates of T05 entities were retired to tombstones.yaml (merged into the T05 entities). References moved from the retired state-selection.descended-form-state-independent to singularity.descended-form-state-independent; version 1 -> 2.
\item Remark~\ref{c:note.state-selection.what-says-about-two-ledgers} (\emph{What this says about the two ledgers}). The state-selection note's duplicates of T05 entities were retired to tombstones.yaml (merged into the T05 entities). References moved from the retired state-selection.descent-parity-average to singularity.descent-parity-average; version 1 -> 2.
\item Remark~\ref{c:note.state-selection.what-was-bought} (\emph{What was bought}). The state-selection note's duplicates of T05 entities were retired to tombstones.yaml (merged into the T05 entities). References moved from the retired state-selection.ascended-marginals-are-diagonal to singularity.ascended-marginals-are-diagonal; version 1 -> 2.
\item Remark~\ref{c:note.state-selection.status-l-e} (\emph{The status of (L$_E$)}). The state-selection note's duplicates of T05 entities were retired to tombstones.yaml (merged into the T05 entities). References moved from the retired state-selection.number-diagonal-dynamics-models-scrambling-derived to singularity.number-diagonal-dynamics-models-scrambling-derived; version 1 -> 2.
\end{itemize}
\noindent In addition, 33 results from elsewhere in the series or the algebra core are stated here that the deposited paper cited or used without stating; 4 results gained citations only. These are not marked.
\begin{thebibliography}{99}
\bibitem{LDG1998} A.~Lasenby, C.~Doran and S.~Gull,
Gravity, gauge theories and geometric algebra,
\emph{Phil.\ Trans.\ R.\ Soc.\ Lond.\ A} \textbf{356} (1998), 487--582.
\bibitem{ButlerT1} L.~Butler,
\emph{The Transport Principle: Total Division in Clifford Algebras via Minimal Completion},
Transport Papers T1 (2026).
\bibitem{ButlerT3} L.~Butler,
\emph{The Tower and the CAR Limit: Iterated Transport, the Frozen Clock, and Second Quantisation},
Transport Papers T3 (2026), doi:\href{https://doi.org/10.5281/zenodo.22721556}{10.5281/zenodo.22721556}.
\bibitem{Hawking1975} S.~W.~Hawking,
Particle creation by black holes,
\emph{Commun.\ Math.\ Phys.} \textbf{43} (1975), 199--220.
\bibitem{Israel1976} W.~Israel,
Thermo-field dynamics of black holes,
\emph{Phys.\ Lett.\ A} \textbf{57} (1976), 107--110.
\bibitem{BW1976} J.~J.~Bisognano and E.~H.~Wichmann,
On the duality condition for quantum fields,
\emph{J.\ Math.\ Phys.} \textbf{17} (1976), 303--321.
\bibitem{Sewell1982} G.~L.~Sewell,
Quantum fields on manifolds: PCT and gravitationally induced thermal states,
\emph{Ann.\ Phys.} \textbf{141} (1982), 201--224.
\bibitem{Page1993} D.~N.~Page,
Average entropy of a subsystem,
\emph{Phys.\ Rev.\ Lett.} \textbf{71} (1993), 1291--1294.
\bibitem{Penington2020} G.~Penington,
Entanglement wedge reconstruction and the information paradox,
\emph{J.\ High Energy Phys.} \textbf{2020}, 002 (2020).
\bibitem{AEMM2019} A.~Almheiri, N.~Engelhardt, D.~Marolf and H.~Maxfield,
The entropy of bulk quantum fields and the entanglement wedge of an evaporating black hole,
\emph{J.\ High Energy Phys.} \textbf{2019}, 063 (2019).
\bibitem{HP2007} P.~Hayden and J.~Preskill,
Black holes as mirrors: quantum information in random subsystems,
\emph{J.\ High Energy Phys.} \textbf{2007}, 120 (2007).
\bibitem{CR2015} M.~Christodoulou and C.~Rovelli,
How big is a black hole?,
\emph{Phys.\ Rev.\ D} \textbf{91} (2015), 064046.
\bibitem{ButlerT4} L.~Butler,
\emph{Clifford Sheaves over Degenerate Metrics: Nilpotent Ideals, One-Sided Completion, and the Exceptional Algebra},
Transport Papers T4 (2026), doi:\href{https://doi.org/10.5281/zenodo.22721669}{10.5281/zenodo.22721669}.
\bibitem{DoranLasenby2003} C.~Doran and A.~Lasenby,
\emph{Geometric Algebra for Physicists},
Cambridge University Press (2003).
\bibitem{PressleySegal1986} A.~Pressley and G.~Segal,
\emph{Loop Groups},
Oxford University Press (1986).
\bibitem{RW2005} S.~P.~Robinson and F.~Wilczek,
Relationship between Hawking radiation and gravitational anomalies,
\emph{Phys.\ Rev.\ Lett.} \textbf{95} (2005), 011303.
\bibitem{Penrose1967} R.~Penrose,
Twistor algebra,
\emph{J.\ Math.\ Phys.} \textbf{8} (1967), 345--366.
\bibitem{NielsenChuang} M.~A.~Nielsen and I.~L.~Chuang,
\emph{Quantum Computation and Quantum Information},
Cambridge University Press (2000).
\bibitem{ButlerT6} L.~Butler,
\emph{The Geometry of the Black-Hole Interior: Where Division Fails, the Weighted Limit, and the Kerr Tests},
Transport Papers T6 (2026), doi:\href{https://doi.org/10.5281/zenodo.22721990}{10.5281/zenodo.22721990}.
\bibitem{ButlerVerifyT05} L.~Butler,
\emph{verify\_T05.py: exact-arithmetic verification of the computational identities of ``Division by Zero at the Gravitational Singularity''},
Python script, deposited with this paper alongside the shared library GA.py (2026).
\bibitem{Porteous1995} I.~R.~Porteous,
\emph{Clifford Algebras and the Classical Groups},
Cambridge University Press (1995).
\bibitem{Takesaki1972} M.~Takesaki,
Conditional expectations in von Neumann algebras,
\emph{J.\ Funct.\ Anal.} \textbf{9} (1972), 306--321.
\bibitem{Chandrasekhar1983} S.~Chandrasekhar,
\emph{The Mathematical Theory of Black Holes},
Oxford University Press, Oxford, 1983.
\bibitem{WWW1952} G.~C.~Wick, A.~S.~Wightman and E.~P.~Wigner,
The intrinsic parity of elementary particles,
\emph{Phys.\ Rev.} \textbf{88} (1952), 101--105.
\bibitem{BR2} O.~Bratteli and D.~W.~Robinson,
\emph{Operator Algebras and Quantum Statistical Mechanics 2: Equilibrium States. Models in Quantum Statistical Mechanics},
2nd ed., Springer (1997).
\bibitem{HHW1967} R.~Haag, N.~M.~Hugenholtz and M.~Winnink,
On the equilibrium states in quantum statistical mechanics,
\emph{Commun.\ Math.\ Phys.} \textbf{5} (1967), 215--236.
\bibitem{Connes1973} A.~Connes,
Une classification des facteurs de type III,
\emph{Ann.\ Sci.\ \'Ec.\ Norm.\ Sup\'er.} (4) \textbf{6} (1973), 133--252.
\bibitem{GibbonsHawking1977} G.~W.~Gibbons and S.~W.~Hawking,
Action integrals and partition functions in quantum gravity,
\emph{Phys.\ Rev.\ D} \textbf{15} (1977), 2752--2756.
\bibitem{HartleHawking1976} J.~B.~Hartle and S.~W.~Hawking,
Path-integral derivation of black-hole radiance,
\emph{Phys.\ Rev.\ D} \textbf{13} (1976), 2188--2203.
\bibitem{Unruh1976} W.~G.~Unruh,
Notes on black-hole evaporation,
\emph{Phys.\ Rev.\ D} \textbf{14} (1976), 870--892.
\bibitem{KayWald1991} B.~S.~Kay and R.~M.~Wald,
Theorems on the uniqueness and thermal properties of stationary, nonsingular, quasifree states on spacetimes with a bifurcate Killing horizon,
\emph{Phys.\ Rep.} \textbf{207} (1991), 49--136.
\bibitem{FredenhagenHaag1990} K.~Fredenhagen and R.~Haag,
On the derivation of Hawking radiation associated with the formation of a black hole,
\emph{Commun.\ Math.\ Phys.} \textbf{127} (1990), 273--284.
\bibitem{Boulware1975} D.~G.~Boulware,
Quantum field theory in Schwarzschild and Rindler spaces,
\emph{Phys.\ Rev.\ D} \textbf{11} (1975), 1404--1423.
\bibitem{ButlerT9} L.~Butler,
\emph{The Twistor Correspondence: Ascended Spinors, Incidence, and the Palatial Bridge},
Transport Papers T9 (2026).
\bibitem{ButlerT8} L.~Butler,
\emph{Black-Hole Entropy and Scrambling: The Singularity Ledger and the No-Echo Prediction},
Transport Papers T8 (2026), doi:\href{https://doi.org/10.5281/zenodo.22722176}{10.5281/zenodo.22722176}.
\bibitem{Bekenstein1973} J.~D.~Bekenstein,
Black holes and entropy,
\emph{Phys.\ Rev.\ D} \textbf{7} (1973), 2333--2346.
\bibitem{AMPS2013} A.~Almheiri, D.~Marolf, J.~Polchinski and J.~Sully,
Black holes: complementarity or firewalls?,
\emph{J.\ High Energy Phys.} \textbf{2013}, 062 (2013).
\bibitem{Mathur2009} S.~D.~Mathur,
The information paradox: a pedagogical introduction,
\emph{Class.\ Quantum Grav.} \textbf{26} (2009), 224001.
\bibitem{Chevalley1954} C.~Chevalley,
\emph{The Algebraic Theory of Spinors},
Columbia University Press (1954).
\bibitem{ButlerT7} L.~Butler,
\emph{Hawking Radiation from the Ledger: State, Continuum, Lift, and the Finite-Ledger Seam},
Transport Papers T7 (2026), doi:\href{https://doi.org/10.5281/zenodo.22722121}{10.5281/zenodo.22722121}.
\bibitem{ButlerT10} L.~Butler,
\emph{Dark Energy and the Action of Gravity: MacDowell-Mansouri, the Frozen Point, and One Way to Grip},
Transport Papers T10 (2026), doi:\href{https://doi.org/10.5281/zenodo.22722496}{10.5281/zenodo.22722496}.
\bibitem{ButlerT14} L.~Butler,
\emph{Dark Matter and the Arrow of Time: The Five Dark Towers},
Transport Papers T14 (2026).
\bibitem{ButlerT16} L.~Butler,
\emph{Quantum Mechanics and General Relativity from One Relation: The Two Shadows and the Laboratory},
Transport Papers T16 (2026), doi:\href{https://doi.org/10.5281/zenodo.22723212}{10.5281/zenodo.22723212}.
\bibitem{ButlerT17} L.~Butler,
\emph{The Transport Programme: Method, Status, and the Falsifiable Ledger},
Transport Papers T17 (2026).
\bibitem{ButlerT5} L.~Butler,
\emph{Division by Zero at the Gravitational Singularity: The Transport Principle Applied to the Black Hole, and the Information Paradox as a Category Error},
Transport Papers T5 (2026), doi:\href{https://doi.org/10.5281/zenodo.22721790}{10.5281/zenodo.22721790}.
\end{thebibliography}
\end{document}